\documentclass[multi=snippet,border={5bp 20bp 5bp 5bp}]{standalone}
\usepackage{fontspec,amssymb,array,graphicx,xcolor,tikz,enumitem,booktabs,multicol}
\usepackage[fleqn]{amsmath}
\usetikzlibrary{calc,arrows.meta,angles,quotes,decorations.pathreplacing,patterns,positioning,intersections}
\setmainfont[Path=fonts/,AutoFakeBold=1.2]{FandolKai-Regular.otf}
\definecolor{answerred}{RGB}{205,20,35}
\newenvironment{snippet}[3]{\begin{minipage}{#1bp}\color{answerred}\fontsize{#2}{#3}\selectfont\raggedright}{\end{minipage}}
\setlength{\parindent}{0pt}\setlength{\parskip}{2bp}
\setlength{\abovedisplayskip}{3bp}\setlength{\belowdisplayskip}{3bp}
\setlength{\abovedisplayshortskip}{2bp}\setlength{\belowdisplayshortskip}{2bp}
\setlength{\mathindent}{0bp}\setlength\jot{1pt}\setlist{nosep,leftmargin=*}
\setlength{\columnsep}{11bp}\setlength{\columnseprule}{0.18bp}
\def\columnseprulecolor{\color{answerred!45}}
\setlength{\multicolsep}{0bp}\setlength{\premulticols}{0bp}\setlength{\postmulticols}{0bp}
\setcounter{unbalance}{0}\setcounter{minrows}{1}
\XeTeXlinebreaklocale "zh"\XeTeXlinebreakskip=0pt plus 1pt\emergencystretch=3em\sloppy
\newfontfamily\SourceSans[Path=fonts/,AutoFakeBold=1.2]{DroidSansFallback.ttf}
\begin{document}
% p001_s01
\begin{snippet}{505}{11.5}{14}\begin{multicols}{3}\raggedcolumns 题目要求先求和，因此先用小括号把\(671\)\allowbreak 与\(29\)\allowbreak 合成一个整体，不能仍让\(29\)\allowbreak 先乘\(4\)\allowbreak 。括号里的和算出后，括号外只剩同级的乘除法，按从左到右的顺序就会先乘\(4\)\allowbreak ，最后除以\(2\)\allowbreak 。\end{multicols}\end{snippet}
% p001_s02
\begin{snippet}{505}{11.5}{14}\begin{multicols}{3}\raggedcolumns \textbf{一、2（1）：}\ 先看箭头的起点：\(136\)\allowbreak 与\(46\)\allowbreak 相减，所得的差填入第二层右框，并作为\(450\)\allowbreak 的除数。第二层求出的商再与\(22\)\allowbreak 相乘，得到最下框的数。合成一个算式时，要把\(136\)\allowbreak 减\(46\)\allowbreak 的差括起来，才能保留“先减、再除、最后乘”的流程。\end{multicols}\end{snippet}
% p001_s03
\begin{snippet}{505}{11.5}{14}\begin{multicols}{3}\raggedcolumns \textbf{一、2（2）：}\ 第一层先算\(252\)\allowbreak 除以\(42\)\allowbreak ，这个商在第二层作\(219\)\allowbreak 的减数；第二层得到的整个差，又在第三层作\(731\)\allowbreak 的减数。合成时应把\(219\)\allowbreak 减去这个商的部分整体括起来，否则就会把两次减法误写成从左到右连续相减。小括号内已有除法，按先除后减计算即可。\end{multicols}\end{snippet}
% p001_s04
\begin{snippet}{505}{11.5}{14}\begin{multicols}{3}\raggedcolumns 用前一步的算式替换后一步的结果：\(18\)\allowbreak 来自\(53\)\allowbreak 与\(35\)\allowbreak 的差，所以先把这个差括起来，再乘\(4\)\allowbreak 。最后一步中\(72\)\allowbreak 是除数，应由前面求得的整个积来替换，因而乘法外还要加一层括号。若只括住减法、不括住整个积，最后就会变成先除再乘，改变原来的顺序。\end{multicols}\end{snippet}
% p001_s05
\begin{snippet}{505}{11.5}{14}\begin{multicols}{3}\raggedcolumns 把括号内的商看作一个因数，已知它与\(8\)\allowbreak 的积是\(120\)\allowbreak ，先用\(120\div8\)\allowbreak 求出这个商。再根据“被除数\(=\)\allowbreak 除数\(\times\)\allowbreak 商”，将所得的商乘\(5\)\allowbreak ，就能还原方框里的数。代回检查时先除以\(5\)\allowbreak ，再乘\(8\)\allowbreak ，应回到\(120\)\allowbreak 。\end{multicols}\end{snippet}
% p001_s06
\begin{snippet}{505}{11.5}{14}\begin{multicols}{3}\raggedcolumns 看错的是除数，原来的被除数没有改变；错误算式中的\(23\)\allowbreak 组\(34\)\allowbreak 再加余下的\(18\)\allowbreak ，仍应合成原被除数。先用\(\text{除数}\times\text{商}+\text{余数}\)\allowbreak 还原：\[\begin{aligned}&23\times34+18\\={}&782+18\\={}&800\end{aligned}\]\[\begin{aligned}&800\div32\\={}&25\end{aligned}\]这里\(18<23\)\allowbreak ，符合原错误计算中余数应小于除数的要求。还原后改用真正的除数\(32\)\allowbreak 计算；所得商乘\(32\)\allowbreak 正好是\(800\)\allowbreak ，所以正确计算没有余数。\end{multicols}\end{snippet}
% p001_s07
\begin{snippet}{505}{11.5}{14}\begin{multicols}{3}\raggedcolumns \textbf{一、6（1）：}\ 左边没有括号，应先求\(122\)\allowbreak 与\(40\)\allowbreak 的积，再减\(15\)\allowbreak 、加\(328\)\allowbreak ；右边要先求\(40\)\allowbreak 与\(15\)\allowbreak 的差，再乘\(122\)\allowbreak 。左、右两边分别算得\(5193\)\allowbreak 、\(3378\)\allowbreak ，据此比较大小。也可从结构看：右边从\(122\)\allowbreak 个\(40\)\allowbreak 中减去的是\(122\)\allowbreak 个\(15\)\allowbreak ，左边只减一个\(15\)\allowbreak ，而最后都加\(328\)\allowbreak ，所以左边较大。\end{multicols}\end{snippet}
% p001_s08
\begin{snippet}{505}{11.5}{14}\begin{multicols}{3}\raggedcolumns \textbf{一、6（2）：}\ 两边都先算\(14\)\allowbreak 与\(6\)\allowbreak 的差，得到共同的除数\(8\)\allowbreak 。左边按从左到右计算，是将\(560\div8\)\allowbreak 的商再乘\(5\)\allowbreak ；右边的整个乘积作除数，相当于把同一个商再除以\(5\)\allowbreak 。这个共同的商大于\(0\)\allowbreak ，乘\(5\)\allowbreak 会增大、除以\(5\)\allowbreak 会缩小，因此无须重算两条长算式就能比较。\end{multicols}\end{snippet}
% p002_s01
\begin{snippet}{505}{11.5}{14}\begin{multicols}{3}\raggedcolumns “照这样计算”表示每分钟走的路程不变，前\(2\)\allowbreak 分钟的\(130\)\allowbreak 米平均分成\(2\)\allowbreak 份，就能求出每分钟走多少米。原式中再乘\(32\)\allowbreak 求的是后来走的路程，加最初的\(130\)\allowbreak 米才是全程，不能把\(32\)\allowbreak 分钟当作总时间。原式的后续计算为：\[\begin{aligned}={}&130+65\times32\\={}&130+2080\\={}&2210\text{(米)}\end{aligned}\]也可先将前后的\(2\)\allowbreak 分钟与\(32\)\allowbreak 分钟相加，求出总时间，再用相同的速度乘总时间：\[\begin{aligned}&130\div2\times(2+32)\\={}&65\times34\\={}&2210\text{(米)}\end{aligned}\]\par\textbf{简便提示（学完运算律后）}\par 先求速度的主法保留。计算后段路程时，可用整数乘除的关系先算32除以2；若先合并总时间34分钟，也可先算34除以2。只调整数值计算的次序，没有把32分钟改成总时间。\[\begin{aligned}&130+130\div2\times32\\={}&130+130\times(32\div2)\\={}&130+130\times16\\={}&130+2080\\={}&2210\text{(米)}\end{aligned}\]先合并时间的一路可简算为：\[\begin{aligned}&130\div2\times(2+32)\\={}&130\div2\times34\\={}&130\times(34\div2)\\={}&130\times17\\={}&2210\text{(米)}\end{aligned}\]\end{multicols}\end{snippet}
% p002_s02
\begin{snippet}{505}{11.5}{14}\begin{multicols}{3}\raggedcolumns 先用每段\(6\)\allowbreak 米乘\(15\)\allowbreak 段，求做手工用去的长度；备用的\(48\)\allowbreak 米也已经从原绳剪下，应和手工用绳合计后一起扣除。原位综合式继续计算：\[\begin{aligned}={}&320-(90+48)\\={}&320-138\\={}&182\text{(米)}\end{aligned}\]另一种整理方法是从\(320\)\allowbreak 米中先扣手工用绳，再扣备用绳，两次扣除的对象与前法相同：\[\begin{aligned}&320-6\times15-48\\={}&320-90-48\\={}&230-48\\={}&182\text{(米)}\end{aligned}\]可把剩下的长度、手工用去的长度和备用长度加回去，检查是否正好为\(320\)\allowbreak 米。\end{multicols}\end{snippet}
% p002_s03
\begin{snippet}{505}{11.5}{14}\begin{multicols}{3}\raggedcolumns 没有括号的乘除混合运算同级，应从左到右，不能把末尾的\(25\times4\)\allowbreak 擅自合成除数。这里先乘\(25\)\allowbreak 又除以\(25\)\allowbreak ，会回到最初的\(4\)\allowbreak ，最后还要再乘题末的\(4\)\allowbreak 。所以也可先抵消中间这对相反的运算，再求两个\(4\)\allowbreak 的积。\end{multicols}\end{snippet}
% p002_s04
\begin{snippet}{505}{11.5}{14}\begin{multicols}{3}\raggedcolumns 最里层的\(22+28\)\allowbreak 要最先算，它的和是\(950\)\allowbreak 的除数。求出商后，在中括号内减去\(12\)\allowbreak ，整个差才与括号外的\(21\)\allowbreak 相乘。因此应沿“小括号内加法—中括号内除法—减法—括号外乘法”的顺序判断，不能先算\(21\)\allowbreak 乘\(950\)\allowbreak 。\end{multicols}\end{snippet}
% p002_s05
\begin{snippet}{505}{11.5}{14}\begin{multicols}{3}\raggedcolumns 要求先减法，就先用小括号把\(35\)\allowbreak 与\(25\)\allowbreak 的差组成整体。接着要让这个差先除以\(5\)\allowbreak ，再与\(90\)\allowbreak 相乘，因此还须把“差除以\(5\)\allowbreak ”放入中括号。若只加小括号，括号外的乘除会从左到右，变成先乘后除，不符合指定顺序。\end{multicols}\end{snippet}
% p002_s06
\begin{snippet}{505}{11.5}{14}\begin{multicols}{3}\raggedcolumns 第一季度有\(3\)\allowbreak 个月，\(405\div3\)\allowbreak 表示每月扣多少元物业费；\(945-405\)\allowbreak 表示扣完第一季度后账户还剩多少元。用余额除以每月费用，就是看余额里有几个这样的月费用，所以结果的单位是“个月”。这里的\(4\)\allowbreak 表示还能交费的月份数，不是第一季度已经交过的月数。\end{multicols}\end{snippet}
% p002_s07
\begin{snippet}{505}{11.5}{14}\begin{multicols}{3}\raggedcolumns 每分钟走\(70\)\allowbreak 米，\(25\)\allowbreak 分钟的路程应用\(70\times25\)\allowbreak 求，得到\(1750\)\allowbreak 米。全程是\(2000\)\allowbreak 米，说明已经走了大部分，离图书馆还差\(250\)\allowbreak 米，应在靠近图书馆而尚未到达的位置。原图没有等分刻度，应按路程的相对位置判断，不能把各字母间隔当作相等。\end{multicols}\end{snippet}
% p002_s08
\begin{snippet}{505}{11.5}{14}\begin{multicols}{3}\raggedcolumns 单独的整数口算可借助拆数：\(23\times3\)\allowbreak 分成\(20\times3\)\allowbreak 与\(3\times3\)\allowbreak ，\(360\times4\)\allowbreak 先算\(36\times4\)\allowbreak 再在积末尾添一个\(0\)\allowbreak ；\(163-45\)\allowbreak 可先减\(40\)\allowbreak ，再减\(5\)\allowbreak 。\(84\div4\)\allowbreak 、\(98\div7\)\allowbreak 、\(84\div6\)\allowbreak 可用相应的乘法反查商。含加法和除法的一题，要先把\(36\div4\)\allowbreak 算出，再把三个加数相加；含括号的一题先求\(28\)\allowbreak 与\(32\)\allowbreak 的和，再除以\(5\)\allowbreak 。不要因从左往右看题，就在乘除之前先做加法。\end{multicols}\end{snippet}
% p003_s01
\begin{snippet}{505}{11.5}{14}\begin{multicols}{3}\raggedcolumns \textbf{三、2（1）：}\ 先进入小括号，但括号内仍要先算\(432\div36\)\allowbreak ，再加\(45\)\allowbreak 。括号内求得的是一个完整的和，最后用这个和乘\(24\)\allowbreak ；脱式时尚未计算的\(\times24\)\allowbreak 要照写，不能漏掉。\end{multicols}\end{snippet}
% p003_s02
\begin{snippet}{505}{11.5}{14}\begin{multicols}{3}\raggedcolumns \textbf{三、2（2）：}\ 小括号把\(352\)\allowbreak 与\(272\)\allowbreak 的和合成了一个被除数，所以先求和，再用这个和除以\(39\)\allowbreak 。最后才从\(91\)\allowbreak 中减去所得的商，不能把\(91\)\allowbreak 先减去括号内的和。求商时可用\(39\)\allowbreak 乘所得商，检查是否回到\(624\)\allowbreak 。\end{multicols}\end{snippet}
% p003_s03
\begin{snippet}{505}{11.5}{14}\begin{multicols}{3}\raggedcolumns \textbf{三、2（3）：}\ 按本单元常规算法，小括号内先求\(118\)\allowbreak 的\(4\)\allowbreak 倍，再减\(215\)\allowbreak ，最后从\(592\)\allowbreak 中扣除整个差。不能把“减去差”误写成连续减两次。\par 简便提示（学完运算律后）：减去一个差，可先减去前一个数，再加上后一个数；先算\(592-472\)\allowbreak 可得整十数\(120\)\allowbreak 。\[\begin{aligned}&592-(472-215)\\={}&592-472+215\\={}&120+215\\={}&335\end{aligned}\]\end{multicols}\end{snippet}
% p003_s04
\begin{snippet}{505}{11.5}{14}\begin{multicols}{3}\raggedcolumns \textbf{三、2（4）：}\ 先算最里层\(62\)\allowbreak 与\(78\)\allowbreak 的和，再从\(215\)\allowbreak 中减去这个和，求中括号的值。最后以中括号内的差为一个因数，与\(76\)\allowbreak 相乘。每算完一层就去掉这一层括号，外面尚未进行的运算仍须保留。\par 简便提示（学完运算律后）：括号仍按原顺序算；最后把\(76\)\allowbreak 拆成\(19\times4\)\allowbreak ，让\(4\)\allowbreak 与\(75\)\allowbreak 先凑成\(300\)\allowbreak 。\[\begin{aligned}&76\times75\\={}&19\times4\times75\\={}&19\times(4\times75)\\={}&19\times300\\={}&5700\end{aligned}\]\end{multicols}\end{snippet}
% p003_s04_local1
\begin{snippet}{92.05938110351565}{10}{12.8}\SourceSans\setlength{\mathindent}{2bp}\(\displaystyle=76\times[215-140]\)\end{snippet}
% p003_s05
\begin{snippet}{505}{11.5}{14}\begin{multicols}{3}\raggedcolumns \textbf{三、2（5）：}\ \(302\)\allowbreak 与\(291\)\allowbreak 的差是\(396\)\allowbreak 的除数，应先求出这个差，再做除法。中括号求得的商与\(58\)\allowbreak 相乘才是最后一步。乘积可把\(58\)\allowbreak 看作\(60-2\)\allowbreak 进行检查，但不能把\(58\)\allowbreak 并入前面的除数。\par\textbf{简便提示（学完运算律后）}\par 前面的括号和除法仍按原顺序求得36。最后把58看成60少2，先多算两份36再减回；两个乘法在同一步完成。\[\begin{aligned}&36\times58\\={}&36\times(60-2)\\={}&36\times60-36\times2\\={}&2160-72\\={}&2088\end{aligned}\]\end{multicols}\end{snippet}
% p003_s05_local1
\begin{snippet}{87.0875946044922}{10}{12.8}\SourceSans\setlength{\mathindent}{2bp}\(\displaystyle=[396\div11]\times58\)\end{snippet}
% p003_s06
\begin{snippet}{505}{11.5}{14}\begin{multicols}{3}\raggedcolumns \textbf{三、2（6）：}\ 先从\(328\)\allowbreak 中减\(258\)\allowbreak ，所得的差用来除\(560\)\allowbreak ，这样才得到最外层的除数。最后是\(184\)\allowbreak 除以中括号内的商。连续出现两个除号时，要特别看清括号：外层除的是整个商，不能改成先用\(184\)\allowbreak 除以\(560\)\allowbreak 。\end{multicols}\end{snippet}
% p003_s06_local1
\begin{snippet}{92.06008300781252}{10}{12.8}\SourceSans\setlength{\mathindent}{2bp}\(\displaystyle=184\div[560\div70]\)\end{snippet}
% p003_s07
\begin{snippet}{505}{11.5}{14}\begin{multicols}{3}\raggedcolumns 从条件出发分析：先求一共回收多少千克，再求剩下多少千克，最后求平均每批送多少千克。\par 从问题出发分析：求每批质量，需要知道剩余质量和批数；求剩余质量，需要知道回收总量和已处理质量；求回收总量，需要知道每班回收质量和班数。两幅图的分析方向不同，原位的综合算式相同。\par 简便提示（学完运算律后）：将总量和已处理量分别平均分成4份，再求每份之差。利用乘除运算的性质，先算\(4\div4\)\allowbreak ，同时算\(32\div4\)\allowbreak ，避免先求72。\[\begin{aligned}&4\times18\div4-32\div4\\={}&(4\div4)\times18-32\div4\\={}&1\times18-8\\={}&18-8\\={}&10\text{(千克)}\end{aligned}\]\par 检验：\[\begin{aligned}&10\times4+32\\={}&40+32\\={}&72\text{(千克)}\end{aligned}\]剩余量与已处理量之和等于回收总量\(4\times18=72\)\allowbreak 千克，答案正确。乘法中的\(4\)\allowbreak 表示每班回收的千克数，除法中的\(4\)\allowbreak 表示批数。\end{multicols}\end{snippet}
% p003_s07_local1
\begin{snippet}{69.80000000000001}{8.5}{12.8}\SourceSans\setlength{\mathindent}{2bp}每班\(4\)千克\end{snippet}
% p003_s07_local2
\begin{snippet}{69.54999999999998}{8.5}{12.8}\SourceSans\setlength{\mathindent}{2bp}\(18\)个班\end{snippet}
% p003_s07_local3
\begin{snippet}{68.79999999999995}{8.5}{12.8}\SourceSans\setlength{\mathindent}{2bp}平均每批\(10\)千克\end{snippet}
% p003_s07_local4
\begin{snippet}{69.80000000000001}{8.5}{12.8}\SourceSans\setlength{\mathindent}{2bp}回收总量\(72\)千克\end{snippet}
% p003_s07_local5
\begin{snippet}{69.04999999999998}{8.5}{12.8}\SourceSans\setlength{\mathindent}{2bp}已处理\(32\)千克\end{snippet}
% p003_s07_local6
\begin{snippet}{72.54999999999995}{8.5}{12.8}\SourceSans\setlength{\mathindent}{2bp}剩余\(40\)千克\end{snippet}
% p003_s07_local7
\begin{snippet}{69.04999999999995}{8.5}{12.8}\SourceSans\setlength{\mathindent}{2bp}\(4\)批\end{snippet}
% p003_s07_local8
\begin{snippet}{77.05000000000001}{8.5}{12.8}\SourceSans\setlength{\mathindent}{2bp}剩余\(40\)千克\end{snippet}
% p003_s07_local9
\begin{snippet}{69.29999999999998}{8.5}{12.8}\SourceSans\setlength{\mathindent}{2bp}\(4\)批\end{snippet}
% p003_s07_local10
\begin{snippet}{72.29999999999995}{8.5}{12.8}\SourceSans\setlength{\mathindent}{2bp}回收总量\(72\)千克\end{snippet}
% p003_s07_local11
\begin{snippet}{69.29999999999995}{8.5}{12.8}\SourceSans\setlength{\mathindent}{2bp}已处理\(32\)千克\end{snippet}
% p003_s07_local12
\begin{snippet}{69.54999999999998}{8.5}{12.8}\SourceSans\setlength{\mathindent}{2bp}平均每批\(10\)千克\end{snippet}
% p003_s07_local13
\begin{snippet}{72.79999999999995}{8.5}{12.8}\SourceSans\setlength{\mathindent}{2bp}每班\(4\)千克\end{snippet}
% p003_s07_local14
\begin{snippet}{69.04999999999995}{8.5}{12.8}\SourceSans\setlength{\mathindent}{2bp}\(18\)个班\end{snippet}
% p003_s07_local15
\begin{snippet}{80}{10.0}{14}\SourceSans\setlength{\mathindent}{2bp}\(\text{(千克)}\)\end{snippet}
% p004_s01
\begin{snippet}{505}{11.5}{14}\begin{multicols}{3}\raggedcolumns 先把人工每分钟的\(40\)\allowbreak 件看作一份，机器人的数量是这样的\(3\)\allowbreak 份，再多\(20\)\allowbreak 件，所以应先求三倍再加\(20\)\allowbreak 。题目比较的是全自动分拣机与机器人，要用\(220\)\allowbreak 件减机器人的数量，不能直接减人工的\(40\)\allowbreak 件。两者都是“每分钟”的数量，直接相减后单位仍写“件”。\end{multicols}\end{snippet}
% p004_s01_local1
\begin{snippet}{80}{10.0}{14}\SourceSans\setlength{\mathindent}{2bp}\(\text{(件)}\)\end{snippet}
% p004_s02
\begin{snippet}{505}{11.5}{14}\begin{multicols}{3}\raggedcolumns 先由日期选票价：\(10\)\allowbreak 月\(15\)\allowbreak 日仍在旺季，学生每人\(30\)\allowbreak 元、教师每人\(45\)\allowbreak 元。原方法分别用单价乘人数求两类票款，再从\(9000\)\allowbreak 元中减去总票款。另法先将\(270+15\)\allowbreak 人都按学生票计费，这会使每位教师少算\(45-30\)\allowbreak 元，所以须按\(15\)\allowbreak 位教师补足差价：\end{multicols}\[\begin{aligned}&9000-[30\times(270+15)+(45-30)\times15]\\={}&9000-[30\times285+15\times15]\\={}&9000-[8550+225]\\={}&9000-8775\\={}&225\text{(元)}\end{aligned}\]\begin{multicols}{3}\raggedcolumns 补差只针对教师，不能再对全部师生补差；票款与剩款相加应回到\(9000\)\allowbreak 元。\par 简便提示（学完运算律后）：减去两笔票款的和，也可连续减去这两笔钱；先减学生票款得\(900\)\allowbreak 元，更便于再减教师票款。\[\begin{aligned}&9000-(30\times270+45\times15)\\={}&9000-(8100+675)\\={}&9000-8100-675\\={}&900-675\\={}&225\text{(元)}\end{aligned}\]\end{multicols}\end{snippet}
% p004_s02_local1
\begin{snippet}{80}{10.0}{14}\SourceSans\setlength{\mathindent}{2bp}\(\text{(元)}\)\end{snippet}
% p004_s03
\begin{snippet}{505}{11.5}{14}\begin{multicols}{3}\raggedcolumns \textbf{四、4（方案一）：}\ 平均每个班分到多少袋糖果？\par 所问的单位是“袋”，所以要选每箱\(16\)\allowbreak 袋与\(3\)\allowbreak 箱这两条信息求总袋数，每箱\(480\)\allowbreak 块不参与本问。先分给\(6\)\allowbreak 个年级，再将每个年级所得平均分给\(4\)\allowbreak 个班，应用连续两次除法：\[\begin{aligned}&16\times3\div6\div4\\={}&48\div6\div4\\={}&8\div4\\={}&2\text{(袋)}\end{aligned}\]也可先按“每个年级\(4\)\allowbreak 班\(\times6\)\allowbreak 个年级”求全校班数，再将总袋数一次平均分到这些班。简便提示（学完运算律后）：24=3×8，除以24就是依次除以3和8。前面先乘3、这里再除以3，仍是每箱的16袋，只需再除以8。\[\begin{aligned}&16\times3\div(4\times6)\\={}&16\times3\div24\\={}&16\times3\div(3\times8)\\={}&16\times3\div3\div8\\={}&16\div8\\={}&2\text{(袋)}\end{aligned}\]答：平均每班分到2袋糖果。\end{multicols}\end{snippet}
% p004_s04
\begin{snippet}{505}{11.5}{14}\begin{multicols}{3}\raggedcolumns \textbf{四、4（方案二）：}\ 若改问每班分到多少“块”，就要用每箱\(480\)\allowbreak 块和\(3\)\allowbreak 箱求总块数，不能再用袋数代替块数。另一合理提问：选①③④⑤，平均每班分到多少块糖果？\par 先把总块数平均分给\(6\)\allowbreak 个年级，再把每个年级所得分给\(4\)\allowbreak 个班：\[\begin{aligned}&480\times3\div6\div4\\={}&1440\div6\div4\\={}&240\div4\\={}&60\text{(块)}\end{aligned}\]另一思路先求\(6\)\allowbreak 个年级一共有多少班，再用总块数除以总班数；两法平均分的总份数相同。简便提示（学完运算律后）：24=3×8，除以24就是依次除以3和8。先乘3再除以3，还原为每箱480块，再除以8即可。\[\begin{aligned}&480\times3\div(4\times6)\\={}&480\times3\div24\\={}&480\times3\div(3\times8)\\={}&480\times3\div3\div8\\={}&480\div8\\={}&60\text{(块)}\end{aligned}\]答：平均每班分到60块糖果。\end{multicols}\end{snippet}
% p004_s05
\begin{snippet}{505}{11.5}{14}\begin{multicols}{3}\raggedcolumns 四数各用一次、中间结果为非负整数，交换加数或因数次序不另计；原位一式先用\(9\div3\)\allowbreak 与\(1\)\allowbreak 凑出\(4\)\allowbreak ，再与\(6\)\allowbreak 相乘。除它以外，还可寻找下面四种结构。\par 把\(24\)\allowbreak 看成\(30-6\)\allowbreak ，先由\(9+1\)\allowbreak 凑出\(10\)\allowbreak ，再乘\(3\)\allowbreak ：\[\begin{aligned}&3\times(9+1)-6\\={}&3\times10-6\\={}&30-6\\={}&24\end{aligned}\]\par 把\(24\)\allowbreak 看成\(9+15\)\allowbreak ，用\(6-1\)\allowbreak 凑出\(5\)\allowbreak ，再乘\(3\)\allowbreak ：\[\begin{aligned}&9+3\times(6-1)\\={}&9+3\times5\\={}&9+15\\={}&24\end{aligned}\]\par 把\(24\)\allowbreak 看成\(8\times3\)\allowbreak ，分别用两对数的差凑出\(8\)\allowbreak 与\(3\)\allowbreak ：\[\begin{aligned}&(9-1)\times(6-3)\\={}&8\times3\\={}&24\end{aligned}\]\par 把\(24\)\allowbreak 看成\(6+18\)\allowbreak ，用\(3-1\)\allowbreak 凑出\(2\)\allowbreak ，再乘\(9\)\allowbreak ：\[\begin{aligned}&6+9\times(3-1)\\={}&6+9\times2\\={}&6+18\\={}&24\end{aligned}\]\end{multicols}\end{snippet}
% p005_s01
\begin{snippet}{505}{11.5}{14}\begin{multicols}{3}\raggedcolumns 加号把算式分成左边的乘积和右边的商，两部分互不依赖，且都比加法先算，所以乘法、除法可以在同一步完成。具体是先求\(12\)\allowbreak 个\(8\)\allowbreak 的积和\(15\)\allowbreak 除以\(3\)\allowbreak 的商，再相加。不能看见乘法在前，就把乘积先与\(15\)\allowbreak 相加后再除以\(3\)\allowbreak 。\end{multicols}\end{snippet}
% p005_s02
\begin{snippet}{505}{11.5}{14}\begin{multicols}{3}\raggedcolumns \textbf{一、2（1）：}\ 要让\(400\)\allowbreak 减\(200\)\allowbreak 先算，就先给这两个数的差加小括号。又要求除法早于乘法，所以还要把“这个差除以\(25\)\allowbreak ”整个商括起来，作为\(32\)\allowbreak 的另一个因数。若只有内层小括号，括号外乘除同级从左到右，会先乘\(32\)\allowbreak ，不符合本小题顺序。\end{multicols}\end{snippet}
% p005_s03
\begin{snippet}{505}{11.5}{14}\begin{multicols}{3}\raggedcolumns \textbf{一、2（2）：}\ 本小题只需把\(400\)\allowbreak 减\(200\)\allowbreak 的差提前，因此将减法部分用小括号括起来。算完差后，外面剩乘法和除法，按从左到右就会先乘\(32\)\allowbreak ，再除以\(25\)\allowbreak 。与上一小题相比，这里不应再把除法部分整体括起来，否则就会让除法提前。\end{multicols}\end{snippet}
% p005_s04
\begin{snippet}{505}{11.5}{14}\begin{multicols}{3}\raggedcolumns 从条件开始分析：先把两种单价相加，再求这个和的3倍，最后加多出的10元。\par 方法一：\[\begin{aligned}&(12+96)\times3+10\\={}&108\times3+10\\={}&324+10\\={}&334\text{(元)}\end{aligned}\]答：每盒即食海鲜礼盒334元。\par 方法二：用乘法分配律，先分别求两种单价的3倍，再合并并加多出的部分：\[\begin{aligned}&12\times3+96\times3+10\\={}&36+288+10\\={}&324+10\\={}&334\text{(元)}\end{aligned}\]两种方法结果相同。\par 检验：\[\begin{aligned}&(334-10)\div3\\={}&324\div3\\={}&108\end{aligned}\]与\(12+96=108\)\allowbreak 一致。多出的10元不参加3倍计算。\end{multicols}\end{snippet}
% p005_s04_local1
\begin{snippet}{55.549999999999955}{7.5}{12.8}\SourceSans\setlength{\mathindent}{2bp}扇贝单价\(12\)元\end{snippet}
% p005_s04_local2
\begin{snippet}{55.549999999999955}{7.5}{12.8}\SourceSans\setlength{\mathindent}{2bp}梭子蟹单价\(96\)元\end{snippet}
% p005_s04_local3
\begin{snippet}{93.29999999999995}{8.5}{12.8}\SourceSans\setlength{\mathindent}{2bp}两种单价和\(108\)元\end{snippet}
% p005_s04_local4
\begin{snippet}{40.799999999999955}{8.5}{12.8}\SourceSans\setlength{\mathindent}{2bp}\(3\)倍\end{snippet}
% p005_s04_local5
\begin{snippet}{117.29999999999995}{8.5}{12.8}\SourceSans\setlength{\mathindent}{2bp}单价和的\(3\)倍\(324\)元\end{snippet}
% p005_s04_local6
\begin{snippet}{40.299999999999955}{8.5}{12.8}\SourceSans\setlength{\mathindent}{2bp}多\(10\)元\end{snippet}
% p005_s04_local7
\begin{snippet}{111.79999999999995}{8.5}{12.8}\SourceSans\setlength{\mathindent}{2bp}礼盒单价\(334\)元\end{snippet}
% p005_s05
\begin{snippet}{505}{11.5}{14}\begin{multicols}{3}\raggedcolumns “\(630\)\allowbreak 减去\(78\)\allowbreak 的差”要先求出，必须用括号表示这个差作为后面运算的整体。差算出后再乘\(14\)\allowbreak 、除以\(7\)\allowbreak ，按同级运算从左到右计算。也可想到\(14\)\allowbreak 里面有\(2\)\allowbreak 个\(7\)\allowbreak ，乘\(14\)\allowbreak 再除以\(7\)\allowbreak 相当于乘\(2\)\allowbreak ，这样可以避免求较大的中间乘积。\par 按题意列式并逐步计算：\[\begin{aligned}&(630-78)\times14\div7\\={}&552\times14\div7\\={}&7728\div7\\={}&1104\end{aligned}\]若按前述简算思路，中间的\(552\)\allowbreak 直接乘\(2\)\allowbreak ，同样得到\(1104\)\allowbreak 。\par\textbf{简便提示（学完运算律后）}\par 除以7只作用在前面整个乘积上，可先用14除以7得到2；括号内的减法与14除以7互不依赖，可同一步算完。\[\begin{aligned}&(630-78)\times14\div7\\={}&(630-78)\times(14\div7)\\={}&552\times2\\={}&1104\end{aligned}\]\end{multicols}\end{snippet}
% p005_s06
\begin{snippet}{505}{11.5}{14}\begin{multicols}{3}\raggedcolumns \textbf{一、5①：}\ 左式先算\(75\div5\)\allowbreak ，再乘\(3\)\allowbreak ，最后才从\(150\)\allowbreak 中减去这个积，算得\(105\)\allowbreak 。右式因有括号，先从\(150\)\allowbreak 减去\(75\)\allowbreak ，再把所得的差除以\(5\)\allowbreak 、乘\(3\)\allowbreak ，算得\(45\)\allowbreak 。两式数字和运算符相同，括号使\(150\)\allowbreak 参与运算的先后不同，应分别按顺序算出后再比较。\end{multicols}\end{snippet}
% p005_s07
\begin{snippet}{505}{11.5}{14}\begin{multicols}{3}\raggedcolumns \textbf{一、5②及发现：}\ 先求\(6+2\)\allowbreak 的和，两式都可先看成由\(240\div8\)\allowbreak 得到同一个正数。左边再乘\(4\)\allowbreak ，所得数比这个共同的商大；右边把\(8\times4\)\allowbreak 作为整个除数，相当于把共同的商再除以\(4\)\allowbreak ，所得数比它小，所以左边较大。无需计算一般小数除法，也能完成比较。结合前一小题可知，括号通过改变运算顺序，能使相同数字和运算符得到不同结果。\end{multicols}\end{snippet}
% p005_s08
\begin{snippet}{505}{11.5}{14}\begin{multicols}{3}\raggedcolumns \textbf{一、6①：}\ 把前一步得到的\(9\)\allowbreak 用它的来源\(36\div4\)\allowbreak 替换，放回\(12+9\)\allowbreak 中，就保留了先除后加的顺序。再把这个和作为一个因数，与\(5\)\allowbreak 相乘，所以整个加法部分要加小括号。括号内除法本来就先于加法，不必再给除法单独加括号。\end{multicols}\end{snippet}
% p005_s09
\begin{snippet}{505}{11.5}{14}\begin{multicols}{3}\raggedcolumns \textbf{一、6②：}\ \(25\)\allowbreak 来自\(60\)\allowbreak 与\(35\)\allowbreak 的差，先给减法加小括号，再把这个差乘\(4\)\allowbreak ，替换第二步的\(100\)\allowbreak 。最后一步是\(500\)\allowbreak 除以\(100\)\allowbreak ，所以前面整个乘积都要作除数，外面还需加中括号。若只括住差，最后的\(\times4\)\allowbreak 就可能跑到除法后，改变给出的分步关系。\end{multicols}\end{snippet}
% p006_s01
\begin{snippet}{505}{11.5}{14}\begin{multicols}{3}\raggedcolumns 先以每分钟\(80\)\allowbreak 米乘已经走的\(5\)\allowbreak 分钟，求已走路程；再从\(1800\)\allowbreak 米全程中扣除它，得到后面要走的路程。把这段剩余路程平均分给\(15\)\allowbreak 分钟，才是剩下每分钟需要走的米数，不能用全程直接除以\(15\)\allowbreak 。题目只要求列综合算式，写清先求剩余路程、再除以剩余时间的关系即可。\end{multicols}\end{snippet}
% p006_s02
\begin{snippet}{505}{11.5}{14}\begin{multicols}{3}\raggedcolumns 信息①②先给出盒数和每盒个数，应按“每盒\(12\)\allowbreak 个\(\times15\)\allowbreak 盒”求羽毛球总数。信息③说明\(80\)\allowbreak 个已经分走，剩下的才能平均分给\(5\)\allowbreak 个班，所以先扣除再平均分：\[\begin{aligned}&(12\times15-80)\div5\\={}&(180-80)\div5\\={}&100\div5\\={}&20\text{(个)}\end{aligned}\]另一种思路把总数与已分出的数量都平均分成\(5\)\allowbreak 份，再求每份的差；这样相减后恰好是一班所得：简便提示（学完运算律后）：这一路径可先用\(15\div5=3\)\allowbreak ，避免先算\(12\times15\)\allowbreak ：\[\begin{aligned}&12\times15\div5-80\div5\\={}&12\times(15\div5)-80\div5\\={}&12\times3-16\\={}&36-16\\={}&20\text{(个)}\end{aligned}\]检查时用每班所得乘\(5\)\allowbreak ，加回已分的\(80\)\allowbreak 个，应等于原有的\(180\)\allowbreak 个。\end{multicols}\end{snippet}
% p006_s03
\begin{snippet}{505}{11.5}{14}\begin{multicols}{3}\raggedcolumns 第一行的\(4\)\allowbreak 个方框代表同一个数，\(100\)\allowbreak 平均分成\(4\)\allowbreak 份，就能先求出一个方框。第二行用\(100\)\allowbreak 除以方框的值，得到的是两个三角形的和，还要再平均分成\(2\)\allowbreak 份，才是一个三角形。最后按题目要求将一个方框和两个三角形的值相乘；不能把第二行的两个三角形误当成三角形的平方。\par 简便提示（学完运算律后）：求出方框是\(25\)\allowbreak 、三角形是\(2\)\allowbreak 以后，把两个\(2\)\allowbreak 先相乘凑成\(4\)\allowbreak ，再与\(25\)\allowbreak 凑\(100\)\allowbreak 。\[\begin{aligned}&25\times2\times2\\={}&25\times(2\times2)\\={}&25\times4\\={}&100\end{aligned}\]\end{multicols}\end{snippet}
% p006_s04
\begin{snippet}{505}{11.5}{14}\begin{multicols}{3}\raggedcolumns “先求和”要求先把\(24\)\allowbreak 与\(6\)\allowbreak 相加，再用这个和乘\(2\)\allowbreak ，组成一个完整的积。最后的除法是让\(600\)\allowbreak 除以这个积，因此乘法部分也要整体括起来。只给加法加括号时，外层会先除后乘，无法让乘法排在除法之前。\end{multicols}\end{snippet}
% p006_s05
\begin{snippet}{505}{11.5}{14}\begin{multicols}{3}\raggedcolumns 先将\(520\)\allowbreak 与\(360\)\allowbreak 的和作为一个整体，再从\(900\)\allowbreak 中减去它，最后以整个差作\(160\)\allowbreak 的除数，才能逐层保留原来的三个步骤。选项D把“减去和”改写成连续减两个数，与B数值相同，但没有保留先求和的步骤结构。若只要求结果相同，B、D都得\(8\)\allowbreak ；本题按给出的分步算式合成，采用B，不能说D的数值算错。\end{multicols}\end{snippet}
% p006_s06
\begin{snippet}{505}{11.5}{14}\begin{multicols}{3}\raggedcolumns 第一空没有括号，乘除同级从左到右；先乘\(8\)\allowbreak 再除以\(25\)\allowbreak 后，还要继续乘末尾的\(8\)\allowbreak ，不能把它漏掉或当成除数的一部分。第二空的括号把\(25\times8\)\allowbreak 组成整个除数，被除数与除数是两个相同且不为\(0\)\allowbreak 的乘积，所以相除得\(1\)\allowbreak 。两式的关键差别是最后一个\(8\)\allowbreak 在括号里面还是外面。\par 简便提示（学完运算律后）：只对没有括号的第一式，乘以\(25\)\allowbreak 与除以\(25\)\allowbreak 可以抵消，最后一个\(8\)\allowbreak 仍是乘数；有括号的第二式整体相除，答案仍为\(1\)\allowbreak 。\[\begin{aligned}&25\times8\div25\times8\\={}&(25\div25)\times(8\times8)\\={}&1\times64\\={}&64\end{aligned}\]\end{multicols}\end{snippet}
% p006_s07
\begin{snippet}{505}{11.5}{14}\begin{multicols}{3}\raggedcolumns 根据减法各部分的关系，减数加差正好等于被减数。因此题目给出的“被减数、减数与差的和”，可以看成两个相同的被减数相加。把\(120\)\allowbreak 平均分成\(2\)\allowbreak 份即可求出被减数；无需先知道减数与差各是多少。\end{multicols}\end{snippet}
% p006_s08
\begin{snippet}{505}{11.5}{14}\begin{multicols}{3}\raggedcolumns 从最后的\(24\)\allowbreak 逆着找：最后一步除以\(2\)\allowbreak ，倒推时先乘\(2\)\allowbreak ，恢复前一个积。这个积由\(8\)\allowbreak 乘括号内的商得到，再除以\(8\)\allowbreak 恢复该商；括号里原来除以\(3\)\allowbreak ，最后乘\(3\)\allowbreak 恢复方框里的被除数。倒推每一步都用原运算的逆运算，求得后再按先除以\(3\)\allowbreak 、乘\(8\)\allowbreak 、除以\(2\)\allowbreak 的正向顺序检验。\end{multicols}\end{snippet}
% p007_s01
\begin{snippet}{505}{11.5}{14}\begin{multicols}{3}\raggedcolumns “今年比去年多\(1200\)\allowbreak 千克”要在去年的\(6800\)\allowbreak 千克上相加，先得到今年产量。明年的两倍关系以今年为标准，须把这个和整体乘\(2\)\allowbreak ，再减去少的\(500\)\allowbreak 千克。不能只把去年产量乘\(2\)\allowbreak ，也不能将\(500\)\allowbreak 千克放进倍数计算的括号里。\end{multicols}\end{snippet}
% p007_s02
\begin{snippet}{505}{11.5}{14}\begin{multicols}{3}\raggedcolumns \textbf{三、1：}\ 先算最内层\(85\)\allowbreak 与\(47\)\allowbreak 的差，再将这个差乘\(6\)\allowbreak ，求出中括号内的积。最后从\(3744\)\allowbreak 中减去这个积；减法虽写在前面，却要等减数全部求出后才算。书写脱式时，括号外的\(3744\)\allowbreak 和减号应一直保留到最后一步。\end{multicols}\end{snippet}
% p007_s02_local1
\begin{snippet}{79.33829650878907}{10}{12.8}\SourceSans\setlength{\mathindent}{2bp}\(\displaystyle=3744-[38\times6]\)\end{snippet}
% p007_s03
\begin{snippet}{505}{11.5}{14}\begin{multicols}{3}\raggedcolumns \textbf{三、2：}\ 先求\(160\)\allowbreak 与\(512\)\allowbreak 的和，再把这个和除以\(32\)\allowbreak ，得到中括号的值。外面的\(25\)\allowbreak 最后与这个商相乘，不能一开始就与\(160\)\allowbreak 相乘。中间除法可用\(32\)\allowbreak 乘所得的商检查；检查后的商才带入最后的乘法。\par 简便提示（学完运算律后）：前面的括号照常算到\(25\times21\)\allowbreak ；最后把\(21\)\allowbreak 拆成\(20+1\)\allowbreak ，先算\(25\times20=500\)\allowbreak 。\[\begin{aligned}&25\times21\\={}&25\times(20+1)\\={}&500+25\\={}&525\end{aligned}\]\end{multicols}\end{snippet}
% p007_s03_local1
\begin{snippet}{87.07808837890627}{10}{12.8}\SourceSans\setlength{\mathindent}{2bp}\(\displaystyle=25\times[672\div32]\)\end{snippet}
% p007_s04
\begin{snippet}{505}{11.5}{14}\begin{multicols}{3}\raggedcolumns \textbf{三、3：}\ 小括号内先求\(23\)\allowbreak 与\(32\)\allowbreak 的和，它是\(275\)\allowbreak 的除数。中括号内得到的商再与\(45\)\allowbreak 相乘，按“加、除、乘”三步完成。这里的\(45\)\allowbreak 是最后的因数，不能与括号内的和一同作除数。\end{multicols}\end{snippet}
% p007_s04_local1
\begin{snippet}{87.0875946044922}{10}{12.8}\SourceSans\setlength{\mathindent}{2bp}\(\displaystyle=[275\div55]\times45\)\end{snippet}
% p007_s05
\begin{snippet}{505}{11.5}{14}\begin{multicols}{3}\raggedcolumns \textbf{三、4：}\ 先求\(38\)\allowbreak 与\(55\)\allowbreak 的和，再用它与\(73\)\allowbreak 相乘，整个积才是外面减法的被减数。最后减\(2876\)\allowbreak 时，应按数位对齐并注意退位。用最终的差加回\(2876\)\allowbreak ，应能回到前一步求出的乘积，可据此检查减法。\par\textbf{简便提示（学完运算律后）}\par 括号按原顺序算得93后，把93看成100少7。先多算7份73再减回；乘100与乘7在同一步完成。\[\begin{aligned}&73\times93\\={}&73\times(100-7)\\={}&73\times100-73\times7\\={}&7300-511\\={}&6789\end{aligned}\]\end{multicols}\end{snippet}
% p007_s05_local1
\begin{snippet}{84.3201843261719}{10}{12.8}\SourceSans\setlength{\mathindent}{2bp}\(\displaystyle=[73\times93]-2876\)\end{snippet}
% p007_s06
\begin{snippet}{505}{11.5}{14}\begin{multicols}{3}\raggedcolumns \textbf{三、5：}\ 最里层先求\(384\)\allowbreak 与\(576\)\allowbreak 的和，再除以\(12\)\allowbreak ，得到括号内的商。最后把这个商乘\(65\)\allowbreak ；可先算\(65\times8\)\allowbreak ，再在积的末尾添一个\(0\)\allowbreak ，检查与\(80\)\allowbreak 相乘的结果。计算中不能把\(65\)\allowbreak 放进括号参与求和。\end{multicols}\end{snippet}
% p007_s06_local1
\begin{snippet}{87.07808837890626}{10}{12.8}\SourceSans\setlength{\mathindent}{2bp}\(\displaystyle=65\times[960\div12]\)\end{snippet}
% p007_s07
\begin{snippet}{505}{11.5}{14}\begin{multicols}{3}\raggedcolumns \textbf{三、6：}\ 先用\(12\times3\)\allowbreak 求最内层的乘积，再用它除\(252\)\allowbreak ，得到最外层真正的除数。最后才用\(336\)\allowbreak 除以这个商。两个除号之间有括号约束，不能按从左到右改成先算\(336\div252\)\allowbreak ；还可用最后的商乘\(7\)\allowbreak 检查是否回到\(336\)\allowbreak 。\end{multicols}\end{snippet}
% p007_s07_local1
\begin{snippet}{92.06008300781252}{10}{12.8}\SourceSans\setlength{\mathindent}{2bp}\(\displaystyle=336\div[252\div36]\)\end{snippet}
% p007_s08
\begin{snippet}{505}{11.5}{14}\begin{multicols}{3}\raggedcolumns 从问题倒着分析，一把椅子的单价需要“椅子总价\(\div5\)\allowbreak 把”；椅子总价又要从\(230\)\allowbreak 元总金额中扣去桌款。桌款按“每张\(80\)\allowbreak 元\(\times2\)\allowbreak 张”计算，所以分析图从最下层乘法，经过减法，再到最上层除法。另法把总金额与桌款分别平均分成\(5\)\allowbreak 份：\[\begin{aligned}&230\div5-80\times2\div5\\={}&46-160\div5\\={}&46-32\\={}&14\text{(元)}\end{aligned}\]这里每份总费用中既有一把椅子的费用，也有五分之一的桌款，减掉对应的一份桌款才得到椅子单价。不能将\(230\)\allowbreak 元直接除以\(2+5\)\allowbreak ，因为桌子与椅子的单价不同。\par 检验：\[\begin{aligned}&80\times2+14\times5\\={}&160+70\\={}&230\text{(元)}\end{aligned}\]与小票总金额一致。\end{multicols}\end{snippet}
% p007_s08_local1
\begin{snippet}{72.04999999999995}{8.5}{12.8}\SourceSans\setlength{\mathindent}{2bp}椅子单价\(14\)元\end{snippet}
% p007_s08_local2
\begin{snippet}{71.79999999999995}{8.5}{12.8}\SourceSans\setlength{\mathindent}{2bp}椅子总价\(70\)元\end{snippet}
% p007_s08_local3
\begin{snippet}{71.79999999999995}{8.5}{12.8}\SourceSans\setlength{\mathindent}{2bp}椅子数量\(5\)把\end{snippet}
% p007_s08_local4
\begin{snippet}{72.29999999999995}{8.5}{12.8}\SourceSans\setlength{\mathindent}{2bp}总金额\(230\)元\end{snippet}
% p007_s08_local5
\begin{snippet}{72.04999999999995}{8.5}{12.8}\SourceSans\setlength{\mathindent}{2bp}桌子总价\(160\)元\end{snippet}
% p007_s08_local6
\begin{snippet}{72.29999999999995}{8.5}{12.8}\SourceSans\setlength{\mathindent}{2bp}桌子单价\(80\)元\end{snippet}
% p007_s08_local7
\begin{snippet}{71.29999999999995}{8.5}{12.8}\SourceSans\setlength{\mathindent}{2bp}桌子数量\(2\)张\end{snippet}
% p007_s08_local8
\begin{snippet}{80}{10.0}{14}\SourceSans\setlength{\mathindent}{2bp}\(\text{(元)}\)\end{snippet}
% p008_s01
\begin{snippet}{505}{11.5}{14}\begin{multicols}{3}\raggedcolumns 每块地每月\(66\)\allowbreak 元，需把\(5\)\allowbreak 块和\(8\)\allowbreak 个月都计入费用；先求五块地每月租金，再乘月数，或先求一块地八个月租金，再乘块数，意义都清楚。优惠前合计\(2640\)\allowbreak 元，达到\(2000\)\allowbreak 元门槛，才能从总租金中减一次\(300\)\allowbreak 元，不能每块地或每个月各减一次。另一思路先求一块地整个租期的费用：\[\begin{aligned}&66\times8\times5-300\\={}&528\times5-300\\={}&2640-300\\={}&2340\text{(元)}\end{aligned}\]\par 简便提示（学完运算律后）：列式仍是每块每月租金乘块数、月数；计算时先把\(5\)\allowbreak 与\(8\)\allowbreak 相乘凑成\(40\)\allowbreak ，再算总租金。原租金\(2640\)\allowbreak 元满足门槛，优惠只减一次。\[\begin{aligned}&66\times5\times8-300\\={}&66\times(5\times8)-300\\={}&66\times40-300\\={}&2640-300\\={}&2340\text{(元)}\end{aligned}\]\end{multicols}\end{snippet}
% p008_s01_local1
\begin{snippet}{80}{10.0}{14}\SourceSans\setlength{\mathindent}{2bp}\(\text{(元)}\)\end{snippet}
% p008_s02
\begin{snippet}{505}{11.5}{14}\begin{multicols}{3}\raggedcolumns 先求下午人数时，要把上午的\(420\)\allowbreak 人乘\(2\)\allowbreak ，再减\(60\)\allowbreak 人；但这只是下午总人数，题目还要求比上午多多少，需再减去上午人数。另法直接看倍数差：下午的两份上午人数减去上午的一份后，还剩一份，再扣少的\(60\)\allowbreak 人：\[\begin{aligned}&420\times(2-1)-60\\={}&420\times1-60\\={}&420-60\\={}&360\text{(人)}\end{aligned}\]这样直接求的是人数差，不能在这个差上再次减\(420\)\allowbreak 人。\end{multicols}\end{snippet}
% p008_s02_local1
\begin{snippet}{80}{10.0}{14}\SourceSans\setlength{\mathindent}{2bp}\(\text{(人)}\)\end{snippet}
% p008_s03
\begin{snippet}{505}{11.5}{14}\begin{multicols}{3}\raggedcolumns 把未修复长度看作\(1\)\allowbreak 份，已修复的长度是这样的\(3\)\allowbreak 份，\(440\)\allowbreak 米全长共对应\(4\)\allowbreak 份。先将全长平均分成\(4\)\allowbreak 份，求出未修复长度，再将这一份平均分给\(5\)\allowbreak 天，才是每天的任务。\(440\div5\)\allowbreak 会把已经修复的部分也计入任务；而直接除以\(3\)\allowbreak ，又会把总长错当成已修复长度。检查时以每天的长度乘\(5\)\allowbreak 还原未修部分，再乘\(3\)\allowbreak 求已修部分，两部分相加应为\(440\)\allowbreak 米。\par 简便提示（学完运算律后）：连续除以\(4\)\allowbreak 和\(5\)\allowbreak ，等于除以它们的积；先凑\(4\times5=20\)\allowbreak 可简算。\[\begin{aligned}&440\div4\div5\\={}&440\div[4\times5]\\={}&440\div20\\={}&22\text{(米)}\end{aligned}\]\end{multicols}\end{snippet}
% p008_s03_local1
\begin{snippet}{80}{10.0}{14}\SourceSans\setlength{\mathindent}{2bp}\(\text{(米)}\)\end{snippet}
% p008_s04
\begin{snippet}{505}{11.5}{14}\begin{multicols}{3}\raggedcolumns 甲数除以乙数商\(5\)\allowbreak ，表示甲数是乙数的\(5\)\allowbreak 倍，不是比乙数多\(5\)\allowbreak 。把乙数看作\(1\)\allowbreak 份，甲数就是\(5\)\allowbreak 份，和\(300\)\allowbreak 对应\(6\)\allowbreak 份，先求一份，再求五份。填空顺序先甲后乙，因此要把五份的数写在前、一份的数写在后；两数相加及相除都可反查条件。\end{multicols}\end{snippet}
% p008_s05
\begin{snippet}{505}{11.5}{14}\begin{multicols}{3}\raggedcolumns 括号可嵌套时最大240；若只添一对小括号，最大60。\par 比较不同括号结构时，要保持\(1200,\allowbreak 30,\allowbreak 20,\allowbreak 2\)\allowbreak 和原有除、减、除的次序不变，只改变哪一部分先算。题面未限定只添一对小括号，允许嵌套时，可把\(30-20\)\allowbreak 的差再除以\(2\)\allowbreak 组成较小的除数，得到原位所示最大值；若课堂另限一对括号，最大值范围随之改变，不能混用两种条件。最小值对应先算最前面的除法，再减\(20\)\allowbreak ，最后除以\(2\)\allowbreak 。\par 为检查极值不遗漏，四个数保持次序时共有五种完全括号结构；原位已写最大、最小用到的两种，其余三种为：\par \((1200\div(30-20))\div2=60\)\allowbreak ；\par \((1200\div30)-(20\div2)=30\)\allowbreak ；\par \(1200\div(30-(20\div2))=60\)\allowbreak 。把这三种与原位两种的结果比较，便完成了允许嵌套范围内的穷尽检查。\par 原位最大值所用算式继续计算：\[\begin{aligned}={}&1200\div(10\div2)\\={}&1200\div5\\={}&240\end{aligned}\]原位最小值所用算式继续计算：\[\begin{aligned}={}&(40-20)\div2\\={}&20\div2\\={}&10\end{aligned}\]\end{multicols}\end{snippet}
% p009_s01
\begin{snippet}{505}{11.5}{14}\begin{multicols}{3}\raggedcolumns 判断两条直线的关系，应想象把画出的部分向两端延长：①②⑦方向相同，延长后仍不相交；④虽然图上还没交到一起，延长后仍会相交。⑤⑥都已相交，但只有⑥交成直角，因此垂直是相交中的一种特殊情况。③画的是曲线，不能因为两条曲线看起来间隔相近，就用直线平行的定义判断。\end{multicols}\end{snippet}
% p009_s02
\begin{snippet}{505}{11.5}{14}\begin{multicols}{3}\raggedcolumns 本题两个条件必须同时满足：先找延长后仍不相交的平行边，再找相交成直角的边。第②图的上下底互相平行，左边又分别与两底成直角，所以两种关系齐全。第①③⑤图虽有平行线段，却没有所需直角；第④图虽有直角，却找不到平行边，不能只满足一个条件就打勾。\end{multicols}\end{snippet}
% p009_s03
\begin{snippet}{505}{11.5}{14}\begin{multicols}{3}\raggedcolumns 在同一平面内讨论两条不同直线：延长后不相交的归为平行，能够相交的归为相交，先分成这两类。若交角恰好是直角，则这两条直线还具有互相垂直的关系，所以垂直不能与平行、相交并列成三类。写垂直符号时，应把\(m\)\allowbreak 、\(n\)\allowbreak 分别放在符号两侧，读的时候也要说清哪两条直线互相垂直。\end{multicols}\end{snippet}
% p009_s04
\begin{snippet}{505}{11.5}{14}\begin{multicols}{3}\raggedcolumns 点到地面的距离，取的是从该点向地面作出的垂线段长度，而不是任意一条斜着连到地面的线段。题目说三条线段中有一条垂线段，又因垂线段最短，比较\(3\)\allowbreak 、\(4\)\allowbreak 、\(5\)\allowbreak 米就能认出它。这里不需要把三条长度相加，也不能取平均值。\end{multicols}\end{snippet}
% p009_s05
\begin{snippet}{505}{11.5}{14}\begin{multicols}{3}\raggedcolumns 长方形的两组对边方向相同且不相交，所以对边平行；每两条邻边交成直角，所以邻边垂直。沿操场跑一圈要依次经过两条\(100\)\allowbreak 米的长和两条\(60\)\allowbreak 米的宽，可先把一长一宽相加，再乘\(2\)\allowbreak 。不能把长乘宽，那求的是面积，单位也不会是路程所用的米。\end{multicols}\end{snippet}
% p009_s06
\begin{snippet}{505}{11.5}{14}\begin{multicols}{3}\raggedcolumns 沿同一方向对折时，展开后的折痕朝向一致，把它们延长也不相交，因此比较的是平行关系。若一次沿横向、一次沿竖向对折，两道折痕会相交成直角，因此比较的是垂直关系。判断依据是折痕的方向与交角，不是折叠次数相同就得到相同关系。\end{multicols}\end{snippet}
% p010_s01
\begin{snippet}{505}{11.5}{14}\begin{multicols}{3}\raggedcolumns \textbf{一、7（续）：}\ 把步道两边看成一对平行线，路宽要沿垂直于两边的方向量。无论在步道的哪一处量，这样得到的垂线段都一样长，因此能保证整条步道宽度一致；斜着量的线段长不能作为路宽。\end{multicols}\end{snippet}
% p010_s02
\begin{snippet}{505}{11.5}{14}\begin{multicols}{3}\raggedcolumns \textbf{一、8（两个空）：}\ 先换成钟面上的时刻：\par\noindent\(\displaystyle 21\text{时}=\text{晚上}9\text{时}\)\allowbreak \par 整点时分针指向\(12\)\allowbreak ，时针指向\(9\)\allowbreak ，两针之间隔着\(3\)\allowbreak 个大格，每个大格是\(30^\circ\)\allowbreak ，合起来恰好一个直角。判断是否垂直，要看两针所在直线相交成的角。\end{multicols}\end{snippet}
% p010_s03
\begin{snippet}{505}{11.5}{14}\begin{multicols}{3}\raggedcolumns 先数图形中互相平行的对边有几组，再抓住题目中的“只有一组”。长方形和平行四边形都有两组对边平行，不能满足“只有一组”；梯形的定义正好符合这个条件。判断依据是对边的位置关系，不能只看图形有几条边。\end{multicols}\end{snippet}
% p010_s04
\begin{snippet}{505}{11.5}{14}\begin{multicols}{3}\raggedcolumns \textbf{二、2（读图条件）：}\ 点到直线的距离，是从这个点向直线所作垂线段的长度。图中\(AC\perp AB\)\allowbreak ，垂足是\(A\)\allowbreak ，因此应读取线段\(AC\)\allowbreak 的长度。已核定\(AC=3\)\allowbreak 厘米，所以点\(C\)\allowbreak 到直线\(AB\)\allowbreak 的距离是\(3\)\allowbreak 厘米，选\(\mathrm A\)\allowbreak ；线段\(BC\)\allowbreak 的\(5\)\allowbreak 厘米不是所求距离。\end{multicols}\end{snippet}
% p010_s04_local1
\begin{snippet}{7.699999999999989}{11.5}{14}\SourceSans\setlength{\mathindent}{2bp}{\color{black}\(3\)}\end{snippet}
% p010_s04_local2
\begin{snippet}{18}{11.5}{14}\SourceSans\setlength{\mathindent}{2bp}\(\mathrm A\)\end{snippet}
% p010_s05
\begin{snippet}{505}{11.5}{14}\begin{multicols}{3}\raggedcolumns \textbf{二、3：}\ 逐项对照概念：互相垂直必须先相交，并且交成直角，所以A不成立。平行说的是同一平面内两条直线的位置关系，即使无限延长也不相交，因而B成立。线段是否平行与它们是否等长是两回事，C不成立，D也随之不成立。\end{multicols}\end{snippet}
% p010_s06
\begin{snippet}{505}{11.5}{14}\begin{multicols}{3}\raggedcolumns 条件有两个：新直线必须经过给定的线外点，还必须与已知直线保持相同方向。满足这两个条件的位置只有一个；若把方向转动，就会与原直线相交，若只平移到别处，又不经过给定点。因此“过一点可以画无数条直线”不能直接用于本题。\end{multicols}\end{snippet}
% p010_s07
\begin{snippet}{505}{11.5}{14}\begin{multicols}{3}\raggedcolumns 汽车沿直线行驶时，两侧车轮前进的方向相同，而且两条轮印之间的距离保持不变。把两条直轮印所在的直线延长，它们仍不会相交，所以用平行关系描述。这里比较的是两侧轮印的位置关系，与每条轮印有多长无关。\end{multicols}\end{snippet}
% p010_s08
\begin{snippet}{505}{11.5}{14}\begin{multicols}{3}\raggedcolumns 先把\(n\perp l\)\allowbreak 换一个读法理解为\(l\)\allowbreak 也垂直于\(n\)\allowbreak ，这样\(m\)\allowbreak 和\(l\)\allowbreak 就有了同一个参照：直线\(n\)\allowbreak 。在同一平面内，两条不同直线都与\(n\)\allowbreak 成直角，它们的方向相同，互相平行。不能因为题中出现两次“垂直”，就认为\(m\)\allowbreak 与\(l\)\allowbreak 也垂直。\end{multicols}\end{snippet}
% p010_s09
\begin{snippet}{505}{11.5}{14}\begin{multicols}{3}\raggedcolumns 判断两条边是否垂直，要寻找它们相交处的直角。三角尺自带一个直角，组成这个直角的两条直角边正好满足条件；它的斜边并不与两条直角边垂直。观察其余物品时，也应按图示的角判断，不能只凭名称猜测。\end{multicols}\end{snippet}
% p010_s10
\begin{snippet}{505}{11.5}{14}\begin{multicols}{3}\raggedcolumns 先让一条直角边与\(l\)\allowbreak 重合，才能确定另一条直角边与\(l\)\allowbreak 垂直。接着保持方向平移，使直角顶点到达\(A\)\allowbreak ，再沿另一条直角边画线，这样所画直线既经过\(A\)\allowbreak 又垂直于\(l\)\allowbreak 。最后在交点处标直角符号，不能先标符号再随意画线。\end{multicols}\end{snippet}
% p010_s11
\begin{snippet}{505}{11.5}{14}\begin{multicols}{3}\raggedcolumns 把其中一条步道上的出发点看成线外一点，另一条步道看成直线。从这一点到对面步道的所有连接线段中，垂线段最短，因此人行道应与两条平行步道都成直角。沿斜线连接虽也能到达对面，路程却会更长。\end{multicols}\end{snippet}
% p011_s01
\begin{snippet}{505}{11.5}{14}\begin{multicols}{3}\raggedcolumns 改直道路时，比较的道路起点和终点保持不变，只改变它们之间的路线。连接这两个固定端点的线段，比绕弯或折转的路线短，因此能缩短路程。“两点确定一条直线”只说明直线的位置能确定，没有说明哪条路线最短。\end{multicols}\end{snippet}
% p011_s02
\begin{snippet}{505}{11.5}{14}\begin{multicols}{3}\raggedcolumns \textbf{三、1（1）：}\ 把乘数\(23\)\allowbreak 分成个位的\(3\)\allowbreak 和十位的\(20\)\allowbreak ，分别求\(114\)\allowbreak 的\(3\)\allowbreak 倍与\(20\)\allowbreak 倍，最后合并两个部分积。第二次乘的\(2\)\allowbreak 实际表示\(2\)\allowbreak 个十，所以这一行的末位要从十位写起，不能与第一行的个位对齐。相加时再按相同数位对齐，避免把\(2280\)\allowbreak 误当成\(228\)\allowbreak 。\end{multicols}\end{snippet}
% p011_s02_local1
\begin{snippet}{27.0}{10.5}{14}\SourceSans\setlength{\mathindent}{2bp}\(2622\)\end{snippet}
% p011_s03
\begin{snippet}{505}{11.5}{14}\begin{multicols}{3}\raggedcolumns \textbf{三、1（2）：}\ 被除数的前两位\(44\)\allowbreak 比\(48\)\allowbreak 小，要用完整的\(448\)\allowbreak 试商，所以商只有一位。试商时取不超过\(448\)\allowbreak 的最大整倍数：\(48\)\allowbreak 的\(9\)\allowbreak 倍还不超，\(10\)\allowbreak 倍就超过了。相减后剩余\(16\)\allowbreak ，它小于\(48\)\allowbreak ，说明不能再分出一份；检验时应把余数加回去。\end{multicols}\end{snippet}
% p011_s03_local1
\begin{snippet}{47.0}{10.5}{14}\SourceSans\setlength{\mathindent}{2bp}\(9\cdots\cdots16\)\end{snippet}
% p011_s04
\begin{snippet}{505}{11.5}{14}\begin{multicols}{3}\raggedcolumns \textbf{三、1（3）：}\ 前两位\(37\)\allowbreak 不够除以\(62\)\allowbreak ，因此从\(372\)\allowbreak 整体试商，商写在个位。可先用\(60\)\allowbreak 的倍数估计，再用\(62\)\allowbreak 乘试得的商核实；\(6\)\allowbreak 个\(62\)\allowbreak 正好用完\(372\)\allowbreak ，不留下余数。商的位置不能因为被除数有三位就多写一位。\end{multicols}\end{snippet}
% p011_s04_local1
\begin{snippet}{15.0}{10.5}{14}\SourceSans\setlength{\mathindent}{2bp}\(6\)\end{snippet}
% p011_s05
\begin{snippet}{505}{11.5}{14}\begin{multicols}{3}\raggedcolumns \textbf{三、1（4）：}\ 先看前两位\(59\)\allowbreak 里面有几个\(18\)\allowbreak ，这一位的商写在十位。第一次相减剩\(5\)\allowbreak ，落下个位的\(8\)\allowbreak 后组成\(58\)\allowbreak ，再除以\(18\)\allowbreak ，个位继续商\(3\)\allowbreak 。最后剩\(4\)\allowbreak ，小于除数\(18\)\allowbreak 才算分完；验算要用除数乘整个两位数商后，再加余数。\end{multicols}\end{snippet}
% p011_s05_local1
\begin{snippet}{47.0}{10.5}{14}\SourceSans\setlength{\mathindent}{2bp}\(33\cdots\cdots4\)\end{snippet}
% p011_s06
\begin{snippet}{505}{11.5}{14}\begin{multicols}{3}\raggedcolumns \textbf{三、2（1）：}\ 外层是用\(1008\)\allowbreak 除以中括号里的整个数，所以必须先确定这个除数。小括号内\(73\)\allowbreak 与\(32\)\allowbreak 相加得到\(105\)\allowbreak ，再把\(105\)\allowbreak 除以\(5\)\allowbreak ，才得到外层要用的除数\(21\)\allowbreak 。不能先算\(1008\)\allowbreak 除以\(105\)\allowbreak ，否则就改变了中括号把商作为整体的作用。\end{multicols}\end{snippet}
% p011_s06_local1
\begin{snippet}{92.05979309082032}{10}{12.8}\SourceSans\setlength{\mathindent}{2bp}\(\displaystyle=1008\div[105\div5]\)\end{snippet}
% p011_s07
\begin{snippet}{505}{11.5}{14}\begin{multicols}{3}\raggedcolumns \textbf{三、2（2）：}\ 先求小括号里\(27\)\allowbreak 与\(19\)\allowbreak 的和，确定\(966\)\allowbreak 要除以的数。去掉括号后只剩乘除同级运算，应从左往右，先求商，再把这个商扩大到\(45\)\allowbreak 倍。不能把\(46\)\allowbreak 与\(45\)\allowbreak 先相乘作为除数，那会改变原来的运算顺序。\par 简便提示（学完运算律后）：括号与除法照常算到\(21\times45\)\allowbreak ；最后把\(21\)\allowbreak 看作\(20+1\)\allowbreak ，先凑\(900\)\allowbreak 再加\(45\)\allowbreak 。\[\begin{aligned}&21\times45\\={}&(20+1)\times45\\={}&900+45\\={}&945\end{aligned}\]\end{multicols}\end{snippet}
% p011_s08
\begin{snippet}{505}{11.5}{14}\begin{multicols}{3}\raggedcolumns \textbf{三、2（3）：}\ 加号左右是两个互不依赖的部分：左边先求\(28\)\allowbreak 除以\(4\)\allowbreak 的商，右边先求\(63\)\allowbreak 乘\(25\)\allowbreak 的积，因此这两步可以同时完成。两个部分都算好后才相加，不能先把\(4\)\allowbreak 与\(63\)\allowbreak 相加。计算\(63\)\allowbreak 乘\(25\)\allowbreak 时，也可把\(63\)\allowbreak 分成\(60\)\allowbreak 和\(3\)\allowbreak ，分别求积再合并。\par\textbf{运算律口算提示（学完运算律后）}\par 补充乘法部分的口算过程：把63分成60与3，用乘法分配律分别求积。前面的28除以4仍得7，再与1575相加为1582。\[\begin{aligned}&63\times25\\={}&(60+3)\times25\\={}&60\times25+3\times25\\={}&1500+75\\={}&1575\end{aligned}\]\end{multicols}\end{snippet}
% p011_s09
\begin{snippet}{505}{11.5}{14}\begin{multicols}{3}\raggedcolumns 两条新直线都要经过同一个黑点\(A\)\allowbreak ，先用点的位置检查是否画准。画平行线时用直尺固定方向，让三角板平移到作图边经过\(A\)\allowbreak ；画垂线时把一条直角边靠齐\(l\)\allowbreak ，移到另一条直角边经过\(A\)\allowbreak 后描线。垂线与\(l\)\allowbreak 的交点才是垂足，直角符号应标在那里。\end{multicols}\end{snippet}
% p011_s10
\begin{snippet}{505}{11.5}{14}\begin{multicols}{3}\raggedcolumns 已有的两条边是长方形的一组邻边，它们决定了两组对边的方向。分别从两条已有边的外端点画对面边的平行线段，新增的两段相交后就能闭合。完成后检查两组对边分别平行、四个角都是直角；不能只把两端任意连起来。\end{multicols}\end{snippet}
% p012_s01
\begin{snippet}{505}{11.5}{14}\begin{multicols}{3}\raggedcolumns 从\(O\)\allowbreak 到某一条边，先找向这条边作出的垂线段，沿斜线爬行不会更短。题目允许爬到任意一边，还应比较到四条边的垂直距离，图示向上边\(AD\)\allowbreak 的垂线段最短。若\(O\)\allowbreak 恰在上、下两条平行边的中间，到\(AD\)\allowbreak 和\(BC\)\allowbreak 的两条垂线段等长，都可以作为最近路线；不能把未标明的中点位置当作已知条件。\end{multicols}\end{snippet}
% p012_s02
\begin{snippet}{505}{11.5}{14}\begin{multicols}{3}\raggedcolumns 长方形两组对边间的垂直距离，就是它的长和宽，因此已知的一条边长为\(12\)\allowbreak 米。方法一把周长平均分成两份，每份是“一条长加一条宽”，再减去已知的\(12\)\allowbreak 米：\[\begin{aligned}&84\div2-12\\={}&42-12\\={}&30\text{(米)}\end{aligned}\]方法二先从周长中扣去两条\(12\)\allowbreak 米的边，余下的是两条相等的未知边，所以还要除以\(2\)\allowbreak ：\[\begin{aligned}&(84-12\times2)\div2\\={}&(84-24)\div2\\={}&60\div2\\={}&30\text{(米)}\end{aligned}\]答：另一组对边之间的距离是30米。\end{multicols}\end{snippet}
% p012_s03
\begin{snippet}{505}{11.5}{14}\begin{multicols}{3}\raggedcolumns 上一问求出的\(30\)\allowbreak 米与已知的\(12\)\allowbreak 米，是花园相邻两条边的长度。面积表示花园里面的大小，要用长乘宽，相当于每排\(30\)\allowbreak 个\(1\)\allowbreak 平方米的小方格，排出\(12\)\allowbreak 排。结果的单位应为平方米；周长\(84\)\allowbreak 米是四条边的总长，不能直接拿来与\(12\)\allowbreak 相乘求面积。\end{multicols}\end{snippet}
% p012_s03_local1
\begin{snippet}{80}{9.5}{14}\SourceSans\setlength{\mathindent}{2bp}\(\text{(平方米)}\)\end{snippet}
% p012_s04
\begin{snippet}{505}{11.5}{14}\begin{multicols}{3}\raggedcolumns \textbf{五、2（1）：}\ 分别把\(A\)\allowbreak 村、\(B\)\allowbreak 村看成线外一点，把各村近侧的国道边线看成直线。连接村庄与国道的线段中，向该边线作出的垂线段最短，所以两村要各自作一条垂线段。两条路线都从原黑点出发，终点落在对应国道边线上，并在垂足处标直角。\end{multicols}\end{snippet}
% p012_s05
\begin{snippet}{505}{11.5}{14}\begin{multicols}{3}\raggedcolumns \textbf{五、2（2）：}\ 这一问的桥两端已经固定在\(A\)\allowbreak 、\(B\)\allowbreak 两村，直接连接这两个黑点即可。任意绕行路线都会比线段\(AB\)\allowbreak 长，依据是两点之间线段最短。桥线是否垂直国道与本问无关，应以连接两个指定端点为准。\end{multicols}\end{snippet}
% p012_s06
\begin{snippet}{505}{11.5}{14}\begin{multicols}{3}\raggedcolumns 先确定民主路左侧的一段：经过\(A\)\allowbreak ，方向与新华路相同，并延伸到民主路。它与民主路的交点确定后，右段就必须从这个同一交点出发，方向改为与青年路平行。检查两段在民主路上首尾相连，才能成为一条连续的阳光大道，不能分别画成两条不相接的平行线。\end{multicols}\end{snippet}
% p013_s01
\begin{snippet}{505}{11.5}{14}\begin{multicols}{3}\raggedcolumns 图①的两条线同向平放，图③的两条线虽倾斜，方向也相同，延长后都不会相交。图④的两条线相交成直角，所以用垂直关系描述；仅仅相交的②和⑤还不能算垂直。图⑥画的是曲线，不属于本题讨论的两条直线，判断时不能只看它们暂时没有碰到。\end{multicols}\end{snippet}
% p013_s02
\begin{snippet}{505}{11.5}{14}\begin{multicols}{3}\raggedcolumns 判断互相垂直要同时满足“相交”和“成直角”，普通的相交关系还不够。记号\(\perp\)\allowbreak 表示这种直角相交关系，读的时候按两条直线的名称从左往右读。只有这两条垂线的交点才叫垂足，不是所有直线的交点都叫垂足。\end{multicols}\end{snippet}
% p013_s03
\begin{snippet}{505}{11.5}{14}\begin{multicols}{3}\raggedcolumns \textbf{一、3：}\ 整点时分针固定指向\(12\)\allowbreak ，只需找时针与它相隔\(3\)\allowbreak 个大格的位置。除了右边的\(3\)\allowbreak ，左边的\(9\)\allowbreak 也与\(12\)\allowbreak 相隔\(3\)\allowbreak 个大格，所成较小角同样是直角。时针指\(6\)\allowbreak 时两针在同一直线上，不能算互相垂直。\end{multicols}\end{snippet}
% p013_s04
\begin{snippet}{505}{11.5}{14}\begin{multicols}{3}\raggedcolumns \textbf{一、4：}\ 门框上边提供了一个固定方向，调整横批时，让它的下边与这个方向保持一致。还可观察横批下边到门框上边的垂直距离是否各处一样，来检查有没有一端高、一端低。只把横批的一端对齐，另一端仍倾斜，不能保证贴正。\end{multicols}\end{snippet}
% p013_s05
\begin{snippet}{505}{11.5}{14}\begin{multicols}{3}\raggedcolumns \textbf{一、5（两个空）：}\ 先分别看题中指定的两组边，不能把所有窗框混在一起判断。上下两条轨道始终保持同一方向，它们之间的垂直距离不变；竖直边框碰到横向轨道时则形成直角。这两种关系的依据不同：一组看是否相交，另一组还要看相交成的角。\end{multicols}\end{snippet}
% p013_s06
\begin{snippet}{505}{11.5}{14}\begin{multicols}{3}\raggedcolumns 图中直角标记表明\(EF\)\allowbreak 、\(GH\)\allowbreak 、\(AD\)\allowbreak 都是两条平行线之间的垂线段，它们量出的都是同一个间距。既然\(EF\)\allowbreak 长\(5\)\allowbreak 厘米，另外两条垂线段也按这个长度确定；两个数值空后没有预印单位，应带上厘米。再从\(A\)\allowbreak 出发比较到\(FH\)\allowbreak 的路线，\(AD\)\allowbreak 是垂线段而\(AB\)\allowbreak 、\(AC\)\allowbreak 都是斜线段，垂线段最短，所以两条斜线段都比\(5\)\allowbreak 厘米长。\end{multicols}\end{snippet}
% p014_s01
\begin{snippet}{505}{11.5}{14}\begin{multicols}{3}\raggedcolumns 第一空中，\(a\)\allowbreak 与\(b\)\allowbreak 平行，方向相同；\(c\)\allowbreak 既与\(b\)\allowbreak 成直角，也就与\(a\)\allowbreak 成直角。第二空中，\(c\)\allowbreak 、\(d\)\allowbreak 都垂直于同一条直线\(b\)\allowbreak ，在同一平面内这两条不同直线方向相同，因而平行。两空分别借助\(b\)\allowbreak 传递位置关系，不能看到“垂直”两字就都填垂直。\end{multicols}\end{snippet}
% p014_s02
\begin{snippet}{505}{11.5}{14}\begin{multicols}{3}\raggedcolumns 题目要求两种位置关系同时出现，找到一种还不够。字母\(H\)\allowbreak 左右两条竖线方向相同，提供平行关系，中间横线又与两条竖线成直角，提供垂直关系。\(L\)\allowbreak 虽有直角，却缺少另一组平行线段，因此不满足全部条件。\end{multicols}\end{snippet}
% p014_s03
\begin{snippet}{505}{11.5}{14}\begin{multicols}{3}\raggedcolumns 两平行线间的“距离”应沿垂直方向量，不能拿斜线段的长度来比较。无论把垂线段移到两条平行线之间的哪一处，它的长度都不变，这正是“同高”的含义。平行线画得一长一短或位置倾斜，都不会改变这个性质。\end{multicols}\end{snippet}
% p014_s04
\begin{snippet}{505}{11.5}{14}\begin{multicols}{3}\raggedcolumns 三条直线共有三对，每一对至多产生一个交点；把三对的交点分开，就得到最多的情况。若让三条直线共同通过同一个点，三对产生的交点重合，就得到最少的情况。题目已规定两两相交，所以不能用平行线得到零个交点；要分清“相交的次数”和“不同交点的个数”。\end{multicols}\end{snippet}
% p014_s05
\begin{snippet}{505}{11.5}{14}\begin{multicols}{3}\raggedcolumns 把\(b\)\allowbreak 作为共同参照，\(a\)\allowbreak 与它同向，\(c\)\allowbreak 也与它同向，因此\(a\)\allowbreak 与\(c\)\allowbreak 的方向一致。在题设的同一平面内，这两条不同直线延长也不会相交。不能因为\(a\)\allowbreak 和\(c\)\allowbreak 没有直接给出符号，就认为它们的位置关系无法判断。\end{multicols}\end{snippet}
% p014_s06
\begin{snippet}{505}{11.5}{14}\begin{multicols}{3}\raggedcolumns 题目没有限定两次对折的方向，应考虑通常折法中的不同情况。连续沿同一方向对折，展开后的两条折痕可以平行；一次横向、一次竖向对折，展开后的折痕可以垂直。因此不能只凭“对折两次”就认定唯一的位置关系，应保留“可能”的限定。\end{multicols}\end{snippet}
% p014_s07
\begin{snippet}{505}{11.5}{14}\begin{multicols}{3}\raggedcolumns 救援船所在的\(B\)\allowbreak 点和游客所在的\(A\)\allowbreak 点，都是已经固定的端点。直接沿线段\(AB\)\allowbreak 航行，路程比任何绕弯或折转的路线短，因此采用两点之间线段最短的原理。这里没有要求到某条直线上的任意一点，不能套用点到直线的垂线段最短。\end{multicols}\end{snippet}
% p014_s08
\begin{snippet}{505}{11.5}{14}\begin{multicols}{3}\raggedcolumns 把村庄看成线外一点，把笔直河岸看成一条直线，取水点可以选在河岸的不同位置。所有从村庄连接到河岸的线段中，垂线段最短，因此把管道垂直铺到河岸能节省管材。关键是河岸上的终点可以选择，与两个端点已经固定的情形不同。\end{multicols}\end{snippet}
% p015_s01
\begin{snippet}{505}{11.5}{14}\begin{multicols}{3}\raggedcolumns \textbf{三、1（1）：}\ 乘数\(13\)\allowbreak 表示\(3\)\allowbreak 个一和\(1\)\allowbreak 个十，两个部分积分别对应\(258\)\allowbreak 的\(3\)\allowbreak 倍和\(10\)\allowbreak 倍。第二行乘十位上的\(1\)\allowbreak 时，末位从十位写起，这一行实际是\(2580\)\allowbreak 。最后按数位对齐相加，不能把十位乘出的部分积写成只扩大\(1\)\allowbreak 倍。\end{multicols}\end{snippet}
% p015_s01_local1
\begin{snippet}{27.0}{10.5}{14}\SourceSans\setlength{\mathindent}{2bp}\(3354\)\end{snippet}
% p015_s02
\begin{snippet}{505}{11.5}{14}\begin{multicols}{3}\raggedcolumns \textbf{三、1（2）：}\ 先用被除数的前两位\(75\)\allowbreak 除以\(26\)\allowbreak ，因为\(26\)\allowbreak 的\(3\)\allowbreak 倍已超过\(75\)\allowbreak ，十位只能商\(2\)\allowbreak 。相减剩\(23\)\allowbreak ，落下个位的\(4\)\allowbreak 组成\(234\)\allowbreak ，再用它除以\(26\)\allowbreak ，确定个位的商。每次都用商乘除数再相减，最后余数为零，说明恰好分完。\end{multicols}\end{snippet}
% p015_s02_local1
\begin{snippet}{19.0}{10.5}{14}\SourceSans\setlength{\mathindent}{2bp}\(29\)\end{snippet}
% p015_s03
\begin{snippet}{505}{11.5}{14}\begin{multicols}{3}\raggedcolumns \textbf{三、1（3）：}\ 乘数\(25\)\allowbreak 中的个位\(5\)\allowbreak 和十位\(2\)\allowbreak ，分别表示乘\(5\)\allowbreak 和乘\(20\)\allowbreak ，两个部分积要按数位错开后相加。被乘数\(308\)\allowbreak 中间的\(0\)\allowbreak 不能省掉，计算时仍要保留十位的位置。还可用\(25\)\allowbreak 是\(100\)\allowbreak 的四分之一来验算：先把\(308\)\allowbreak 平均分成\(4\)\allowbreak 份，再扩大\(100\)\allowbreak 倍，与竖式结果应相同。\end{multicols}\end{snippet}
% p015_s03_local1
\begin{snippet}{27.0}{10.5}{14}\SourceSans\setlength{\mathindent}{2bp}\(7700\)\end{snippet}
% p015_s04
\begin{snippet}{505}{11.5}{14}\begin{multicols}{3}\raggedcolumns \textbf{三、1（4）：}\ 前两位\(55\)\allowbreak 小于除数\(69\)\allowbreak ，所以商只有一位，要用\(552\)\allowbreak 整体试商。可把\(69\)\allowbreak 看成接近\(70\)\allowbreak 的数估计商，再用准确的除数\(69\)\allowbreak 乘试商核对。试得的\(8\)\allowbreak 倍正好等于被除数，余数为零；估计时用的\(70\)\allowbreak 不能代替最终的准确乘算。\end{multicols}\end{snippet}
% p015_s04_local1
\begin{snippet}{15.0}{10.5}{14}\SourceSans\setlength{\mathindent}{2bp}\(8\)\end{snippet}
% p015_s05
\begin{snippet}{505}{11.5}{14}\begin{multicols}{3}\raggedcolumns \textbf{三、2（1）：}\ 小括号中的\(24\)\allowbreak 与\(195\)\allowbreak 先相加，得到的\(219\)\allowbreak 是要乘\(15\)\allowbreak 的整个数。接着求这\(15\)\allowbreak 份的总量，最后才从\(4000\)\allowbreak 中减去这个积。不能先用\(4000\)\allowbreak 减\(15\)\allowbreak ，也不能把\(15\)\allowbreak 只乘括号中的\(24\)\allowbreak 。\par 简便提示（学完运算律后）：先按原顺序算得\(24+195=219\)\allowbreak 。把\(219\)\allowbreak 看成\(220-1\)\allowbreak ，先减\(15\times220=3300\)\allowbreak 是多减了\(15\)\allowbreak ，应再加回\(15\)\allowbreak 。\[\begin{aligned}&4000-15\times219\\={}&4000-15\times(220-1)\\={}&4000-(3300-15)\\={}&4000-3300+15\\={}&700+15\\={}&715\end{aligned}\]\end{multicols}\end{snippet}
% p015_s06
\begin{snippet}{505}{11.5}{14}\begin{multicols}{3}\raggedcolumns \textbf{三、2（2）：}\ 先求最内层的差\(84\)\allowbreak 减\(12\)\allowbreak ，再把所得的\(72\)\allowbreak 除以\(4\)\allowbreak 。中括号把这一步的商作为一个整体，它才是\(180\)\allowbreak 要除以的数。因而外层除法必须最后做，不能按从左往右的顺序先动\(180\)\allowbreak 。\end{multicols}\end{snippet}
% p015_s06_local1
\begin{snippet}{82.09676513671877}{10}{12.8}\SourceSans\setlength{\mathindent}{2bp}\(\displaystyle=180\div[72\div4]\)\end{snippet}
% p015_s07
\begin{snippet}{505}{11.5}{14}\begin{multicols}{3}\raggedcolumns \textbf{三、2（3）：}\ 先求小括号中\(60\)\allowbreak 与\(15\)\allowbreak 的差，再用\(160\)\allowbreak 减去整个差，得到中括号表示的数。中括号全部算好后，才用\(34\)\allowbreak 乘这个数。尤其不能把“减去\(60-15\)\allowbreak 的差”误算成连续减\(60\)\allowbreak 再减\(15\)\allowbreak ，那样会改变减去的数量。\par 简便提示（学完运算律后）：减去一个差，等于先减前数、再加后数。先用\(160-60\)\allowbreak 凑\(100\)\allowbreak ，再按乘法分配律计算；\(15\)\allowbreak 前应是加号。\[\begin{aligned}&34\times[160-(60-15)]\\={}&34\times[160-60+15]\\={}&34\times[100+15]\\={}&34\times100+34\times15\\={}&3400+510\\={}&3910\end{aligned}\]\end{multicols}\end{snippet}
% p015_s07_local1
\begin{snippet}{87.07808837890626}{10}{12.8}\SourceSans\setlength{\mathindent}{2bp}\(\displaystyle=34\times[160-45]\)\end{snippet}
% p015_s08
\begin{snippet}{505}{11.5}{14}\begin{multicols}{3}\raggedcolumns \textbf{三、2（4）：}\ 小括号中的\(295\)\allowbreak 减\(289\)\allowbreak 先得到除数\(6\)\allowbreak ，然后才能求\(288\)\allowbreak 除以\(6\)\allowbreak 的商。中括号算完以后，把得到的\(48\)\allowbreak 再乘\(17\)\allowbreak ，不能先乘右边的\(17\)\allowbreak 。最后的乘法可分别求\(48\)\allowbreak 的\(10\)\allowbreak 倍和\(7\)\allowbreak 倍后合并，帮助检查计算。\par\textbf{简便提示（学完运算律后）}\par 括号与除法照常求得48，最后可把48看成50少2，按乘法分配律同时算50乘17与2乘17。\[\begin{aligned}&48\times17\\={}&(50-2)\times17\\={}&50\times17-2\times17\\={}&850-34\\={}&816\end{aligned}\]\end{multicols}\end{snippet}
% p015_s08_local1
\begin{snippet}{82.1055847167969}{10}{12.8}\SourceSans\setlength{\mathindent}{2bp}\(\displaystyle=[288\div6]\times17\)\end{snippet}
% p015_s09
\begin{snippet}{505}{11.5}{14}\begin{multicols}{3}\raggedcolumns 先用三角板的两条直角边确定过\(A\)\allowbreak 的垂直方向，再把直尺的零刻度对准\(A\)\allowbreak ，沿这条垂线量到\(3\)\allowbreak 厘米处确定\(B\)\allowbreak 。过\(B\)\allowbreak 画平行线时，用直尺固定方向，平移三角板到作图边经过\(B\)\allowbreak 。完成后检查\(AB\)\allowbreak 确为\(3\)\allowbreak 厘米、与原线成直角，新线经过\(B\)\allowbreak 且保持与原线相同方向。\end{multicols}\end{snippet}
% p015_s10
\begin{snippet}{505}{11.5}{14}\begin{multicols}{3}\raggedcolumns 陆地上走垂直于河岸的最短线段，在河里游泳。\par “省力气”的关键条件是小鸭走陆路费劲，因此先让下水前和上岸后的两段陆地路程分别最短。小鸭家到近岸、小鸡家到另一岸，都应选对应的垂线段，并把两个垂足分别作为下水点、上岸点。再用线段连接两个垂足作为水中路线；不能只把两家直接相连，就认为费力的陆地部分也一定最短。\end{multicols}\end{snippet}
% p016_s01
\begin{snippet}{505}{11.5}{14}\begin{multicols}{3}\raggedcolumns 起点固定在\(A\)\allowbreak ，终点可以选在对岸，正好是点到直线的最短路线问题。过\(A\)\allowbreak 向对岸作垂线段，沿这条线段穿过人行横道，就走过了两条平行路边的垂直距离。检查路线确实经过\(A\)\allowbreak ，与对岸的交点处标直角；斜穿马路会多走一段路。\end{multicols}\end{snippet}
% p016_s02
\begin{snippet}{505}{11.5}{14}\begin{multicols}{3}\raggedcolumns 沿最短路线走过的路程，正好等于马路的宽。原方法先求厘米总数，再单独换成米；换单位不能只改单位名称。另一种分组思路是把每\(10\)\allowbreak 秒看成一段，先同时求每段的厘米数和段数：\par\noindent\(\displaystyle 60\times10=600\text{(厘米)},\allowbreak \qquad40\div10=4\text{(段)}\)\allowbreak \par \noindent\(\displaystyle 600\text{厘米}=(600\div100)\text{米}=6\text{米}\)\allowbreak \par 再用每段长度乘段数：\par\noindent\[\displaystyle \begin{aligned}&6\times4\\={}&24\text{(米)}\end{aligned}\]\par\end{multicols}\end{snippet}
% p016_s02_local1
\begin{snippet}{505}{10.0}{14}\SourceSans\setlength{\mathindent}{2bp}\(60\times40\)\end{snippet}
% p016_s02_local2
\begin{snippet}{505}{10.0}{14}\SourceSans\setlength{\mathindent}{2bp}\(=2400\text{(厘米)}\)\end{snippet}
% p016_s02_local3
\begin{snippet}{505}{10.0}{14}\SourceSans\setlength{\mathindent}{2bp}\(2400\text{厘米}=(2400\div100)\text{米}=24\text{米}\)\end{snippet}
% p016_s03
\begin{snippet}{505}{11.5}{14}\begin{multicols}{3}\raggedcolumns 先量出两条平行线间的垂直距离，把它作为正方形的一条边长，再按同样长度确定其余各边。这样上下两边落在给定平行线上，左右两边与它们垂直，正方形恰好撑满两线之间。边长若再增加就放不下，若减小虽能放下却不是最大的；还应检查四边等长、四角都是直角。\end{multicols}\end{snippet}
% p016_s04
\begin{snippet}{505}{11.5}{14}\begin{multicols}{3}\raggedcolumns 先连接\(A\)\allowbreak 、\(B\)\allowbreak 并延长，确定公路的位置；再从\(C\)\allowbreak 向公路作垂线段，量得题中所说的“\(C\)\allowbreak 点到公路的距离”。过\(C\)\allowbreak 作公路的平行线，只在\(C\)\allowbreak 的左下方向截取一段与刚才垂线段等长的线段，所得端点就是藏宝地。平行方向确定了位置所在的直线，等长确定了离\(C\)\allowbreak 多远，左下方再排除另一侧的端点，三个条件缺一不可。\end{multicols}\end{snippet}
% p017_s01
\begin{snippet}{505}{11.5}{14}\begin{multicols}{3}\raggedcolumns \(45\)\allowbreak 元是\(3\)\allowbreak 千克草莓合起来的钱，每千克的价钱则是把这笔总价平均分成\(3\)\allowbreak 份后的一份。因为题目求的是每\(1\)\allowbreak 千克的价钱，所以用总价除以千克数，结果用“元/千克”表示。单价、数量、总价要按各自含义区分，不能因为都带有数值就混淆。\end{multicols}\end{snippet}
% p017_s01_local1
\begin{snippet}{80}{9.5}{14}\SourceSans\setlength{\mathindent}{2bp}\(\text{(元/千克)}\)\end{snippet}
% p017_s02
\begin{snippet}{505}{11.5}{14}\begin{multicols}{3}\raggedcolumns \textbf{一、2（牛奶）：}\ 牛奶每箱\(65\)\allowbreak 元，\(195\)\allowbreak 元里含有几个\(65\)\allowbreak 元，就可以买几箱。这里知道每份的钱和总钱数，求份数要用除法；不是把\(195\)\allowbreak 平均分成\(65\)\allowbreak 箱。得到的数量单位用箱，还可用每箱价格乘箱数核对是否回到\(195\)\allowbreak 元。\end{multicols}\end{snippet}
% p017_s03
\begin{snippet}{505}{11.5}{14}\begin{multicols}{3}\raggedcolumns \textbf{一、2（鸡蛋）：}\ \(108\)\allowbreak 元是\(6\)\allowbreak 盒鸡蛋的总价，要求每一盒多少钱，就把\(108\)\allowbreak 元平均分成\(6\)\allowbreak 份。每份的钱就是单价，所以用总价除以盒数，单位写元/盒。这里的\(6\)\allowbreak 是数量，不能当成每盒的价钱。\end{multicols}\end{snippet}
% p017_s04
\begin{snippet}{505}{11.5}{14}\begin{multicols}{3}\raggedcolumns \textbf{一、2（奶茶）：}\ 每杯奶茶都要\(12\)\allowbreak 元，买\(5\)\allowbreak 杯就有\(5\)\allowbreak 份这样的价钱。求\(5\)\allowbreak 个\(12\)\allowbreak 元合起来多少，用每杯价格乘杯数，结果是总价，单位只用元。表中的元/杯、杯与元分别对应单价、数量、总价，要随所求的量选择单位。\end{multicols}\end{snippet}
% p017_s05
\begin{snippet}{505}{11.5}{14}\begin{multicols}{3}\raggedcolumns “每小时约行驶\(305\)\allowbreak 千米”说的是每\(1\)\allowbreak 小时行多少路，因此表示速度；写速度单位时，千米和时都要保留。按这个速度行\(8\)\allowbreak 小时，就是求\(8\)\allowbreak 份\(305\)\allowbreak 千米的总量，要用速度乘时间。题中速度带有“约”，由此得到的总路程也按约数理解。\par 用速度乘时间求全程：\[\begin{aligned}&305\times8\\={}&2440\text{(千米)}\end{aligned}\]\end{multicols}\end{snippet}
% p017_s06
\begin{snippet}{505}{11.5}{14}\begin{multicols}{3}\raggedcolumns 先统一长度单位：\par\noindent\(\displaystyle 1\text{千米}=1000\text{米}\)\allowbreak \par 这样才能与每分钟\(100\)\allowbreak 米的速度使用同一种长度单位。再看\(1000\)\allowbreak 米里面包含多少份\(100\)\allowbreak 米，每一份对应\(1\)\allowbreak 分钟，因此用路程除以速度求时间。不能直接用数值\(1\)\allowbreak 除以\(100\)\allowbreak ，那样把千米与米混在了一起。\end{multicols}\end{snippet}
% p017_s07
\begin{snippet}{505}{11.5}{14}\begin{multicols}{3}\raggedcolumns 每买一根要付\(8\)\allowbreak 元，先看\(245\)\allowbreak 元能组成多少个完整的\(8\)\allowbreak 元。买\(30\)\allowbreak 根后余下\(5\)\allowbreak 元，小于一根的价格，因此这笔余款不能再算成一根。若增加到\(31\)\allowbreak 根，需要\(248\)\allowbreak 元，已经超过\(245\)\allowbreak 元，所以求“最多”应舍去不足一根的部分。\end{multicols}\end{snippet}
% p017_s08
\begin{snippet}{505}{11.5}{14}\begin{multicols}{3}\raggedcolumns 两辆车行驶的时间不相同，不能只看\(360\)\allowbreak 千米比\(320\)\allowbreak 千米多就判断快慢。把各自路程平均分到每小时，大货车每小时行\(72\)\allowbreak 千米，冷藏车每小时行\(80\)\allowbreak 千米。比较相同时间内的路程，后者更多，才能确定谁的速度更快。\end{multicols}\end{snippet}
% p017_s09
\begin{snippet}{505}{11.5}{14}\begin{multicols}{3}\raggedcolumns 全部坐满表示有\(102\)\allowbreak 人买票，也就是售出\(102\)\allowbreak 张同价票。每张\(38\)\allowbreak 元，总收入就是\(102\)\allowbreak 份\(38\)\allowbreak 元，按单价乘数量计算。计算时把\(102\)\allowbreak 分成\(100\)\allowbreak 和\(2\)\allowbreak ，先算\(100\)\allowbreak 张与\(2\)\allowbreak 张的钱再合起来，不能漏掉多出的两张票款。\par 将整百张和另外两张的票款合并：\[\begin{aligned}&38\times102\\={}&38\times100+38\times2\\={}&3800+76\\={}&3876\text{(元)}\end{aligned}\]\end{multicols}\end{snippet}
% p018_s01
\begin{snippet}{505}{11.5}{14}\begin{multicols}{3}\raggedcolumns 第一空按盒计价，\(288\)\allowbreak 元平均分给\(8\)\allowbreak 盒；第二空按条计价，每盒的钱还要再平均分给\(3\)\allowbreak 条。先求每盒，再求每条：\[\begin{aligned}&288\div8\\={}&36\text{(元/盒)}\end{aligned}\]\[\begin{aligned}&36\div3\\={}&12\text{(元/条)}\end{aligned}\]也可先用每盒\(3\)\allowbreak 条乘\(8\)\allowbreak 盒，求出总条数，再把总价平均分到每一条：\[\begin{aligned}&288\div(3\times8)\\={}&288\div24\\={}&12\text{(元/条)}\end{aligned}\]两种单价的单位分别是元/盒、元/条，不能混用；每条价应小于每盒价，且按每条价、每盒条数、盒数依次乘回去，应还原总价\(288\)\allowbreak 元。\end{multicols}\end{snippet}
% p018_s02
\begin{snippet}{505}{11.5}{14}\begin{multicols}{3}\raggedcolumns \textbf{一、9（1）：}\ 知道同一段路的平均速度和行驶时间，求这段路程，要用速度乘时间，因为每小时行驶的千米数就是一份，行驶几小时就有几份。题面却只说高速路段的平均速度是\(106\)\allowbreak 千米/时，给出的\(22\)\allowbreak 小时是全程时间。只有全程这\(22\)\allowbreak 小时都适用该平均速度，才可据此求全程路程；若还含其他速度未知的路段，条件不足，不能把高速路段的速度直接套到全程。\end{multicols}\end{snippet}
% p018_s03
\begin{snippet}{505}{11.5}{14}\begin{multicols}{3}\raggedcolumns \textbf{一、9（2）（两个空）：}\ 乘数\(22\)\allowbreak 由\(2\)\allowbreak 个一和\(2\)\allowbreak 个十组成，竖式两行部分积分别计算这两部分。第②行用的是十位上的\(2\)\allowbreak ，实际表示\(20\)\allowbreak ，因此该行的\(212\)\allowbreak 向左错一位，表示\(2120\)\allowbreak 千米。在题设速度适用的前提下，它对应\(20\)\allowbreak 小时的路程；不能只看写出的数字\(2\)\allowbreak ，就把它理解为\(2\)\allowbreak 小时。\end{multicols}\end{snippet}
% p018_s04
\begin{snippet}{505}{11.5}{14}\begin{multicols}{3}\raggedcolumns 先把每个数与它所描述的事情对应起来：\(10\)\allowbreak 分钟说的是走路用了多久，\(15\)\allowbreak 元说的是三本练习本合起来花了多少钱，\(3\)\allowbreak 本说的是买了多少本。总价要与“每一本多少钱”的单价分开，题目并没有把\(15\)\allowbreak 元说成一本的价钱。用单位配合语句含义判断，比只记名词顺序更可靠。\end{multicols}\end{snippet}
% p018_s05
\begin{snippet}{505}{11.5}{14}\begin{multicols}{3}\raggedcolumns 第一问每瓶都要\(3\)\allowbreak 元，\(12\)\allowbreak 瓶就是\(12\)\allowbreak 份相同的价钱，求总价用每瓶价格乘瓶数。第二问要求骑手每分钟行多少米，要把\(5\)\allowbreak 分钟行的\(1400\)\allowbreak 米平均分成\(5\)\allowbreak 份，因此用路程除以时间。两个问题一个求总量，一个求每份量，应分别选择对应的运算。\end{multicols}\end{snippet}
% p018_s06
\begin{snippet}{505}{11.5}{14}\begin{multicols}{3}\raggedcolumns 图中甲、乙两地相距\(280\)\allowbreak 千米，货车用\(4\)\allowbreak 小时行完，把全程平均分成\(4\)\allowbreak 份就得到每小时的路程。速度的单位要同时说明长度和时间，所以要用千米/时。\(280\)\allowbreak 千米是总路程，不能直接当成每小时的速度；单独的时或千米也都不能完整表示速度。\end{multicols}\end{snippet}
% p018_s07
\begin{snippet}{505}{11.5}{14}\begin{multicols}{3}\raggedcolumns 速度说明每\(1\)\allowbreak 小时行多少千米，乘行驶的小时数，求出的就是这几小时合起来的路程。汽车\(5\)\allowbreak 小时行驶的千米数正好属于这种总量问题。求每支钢笔的钱要用总价除以支数，求每分钟走的米数要用路程除以时间，长方形周长则要依据边长计算。\end{multicols}\end{snippet}
% p018_s08
\begin{snippet}{505}{11.5}{14}\begin{multicols}{3}\raggedcolumns \textbf{二、5：}\ 指示牌中的\(203\)\allowbreak 千米描述当前位置与青岛之间还有多远，所以它表示待行驶的路程。这里没有“每小时”或“每分钟”，不能把它看成速度；时间也不能用千米来表示。还要注意这是“剩下的路程”，并非已经走过的路程。\end{multicols}\end{snippet}
% p018_s09
\begin{snippet}{505}{11.5}{14}\begin{multicols}{3}\raggedcolumns \textbf{二、6（A、C都正确）：}\ C先把\(6\)\allowbreak 本的总价\(18\)\allowbreak 元平均分开，求出一本的价钱，再求\(4\)\allowbreak 本的钱。A先求\(4\)\allowbreak 组、每组\(6\)\allowbreak 本的钱，这时共有\(24\)\allowbreak 本，再平均分成\(6\)\allowbreak 份，每份恰好是\(4\)\allowbreak 本的钱。两种列式都保持了同一单价，因此都成立；不能因为A先乘后除就判错，也不能把\(4\)\allowbreak 本误当成\(4\)\allowbreak 元与单价相加。\par C按“先求单价、再求总价”计算：\[\begin{aligned}={}&3\times4\\={}&12\text{(元)}\end{aligned}\]A按“先求四组的钱、再平均分成六份”计算：\[\begin{aligned}={}&72\div6\\={}&12\text{(元)}\end{aligned}\]\end{multicols}\end{snippet}
% p019_s01
\begin{snippet}{505}{11.5}{14}\begin{multicols}{3}\raggedcolumns \textbf{三、1（1）：}\ \(125\)\allowbreak 与\(8\)\allowbreak 是常用的凑整组合，合起来正好是一千。熟记这一组，遇到含\(125\)\allowbreak 的连乘时，也可以寻找因数\(8\)\allowbreak 先凑整。\end{multicols}\end{snippet}
% p019_s02
\begin{snippet}{505}{11.5}{14}\begin{multicols}{3}\raggedcolumns \textbf{三、1（2）：}\ 把\(40\)\allowbreak 看成\(4\)\allowbreak 个十，先让\(4\)\allowbreak 与\(25\)\allowbreak 相乘凑成一百，再乘剩下的\(10\)\allowbreak 。拆分后原来的\(40\)\allowbreak 没有改变，不能漏乘这个\(10\)\allowbreak 。\end{multicols}\end{snippet}
% p019_s03
\begin{snippet}{505}{11.5}{14}\begin{multicols}{3}\raggedcolumns \textbf{三、1（3）：}\ 可把\(74\)\allowbreak 拆成\(70\)\allowbreak 和\(4\)\allowbreak ，从\(103\)\allowbreak 先减整十数，再减个位数。也可以先减\(100\)\allowbreak 再补回\(26\)\allowbreak ，因为减\(100\)\allowbreak 比减\(74\)\allowbreak 多减了\(26\)\allowbreak 。\end{multicols}\end{snippet}
% p019_s04
\begin{snippet}{505}{11.5}{14}\begin{multicols}{3}\raggedcolumns \textbf{三、1（4）：}\ 先算乘法，\(0\)\allowbreak 与\(75\)\allowbreak 相乘的积仍是\(0\)\allowbreak ，再把它与\(9\)\allowbreak 相加。不能先把\(75\)\allowbreak 与\(9\)\allowbreak 相加，否则改变了原来的运算顺序。\end{multicols}\end{snippet}
% p019_s05
\begin{snippet}{505}{11.5}{14}\begin{multicols}{3}\raggedcolumns \textbf{三、1（5）：}\ 乘法要先于加法，所以先求\(57\)\allowbreak 的\(2\)\allowbreak 倍，再加上\(13\)\allowbreak 。这里的\(2\)\allowbreak 只与\(57\)\allowbreak 相乘，不能把\(13\)\allowbreak 也算进被乘的数里。\end{multicols}\end{snippet}
% p019_s06
\begin{snippet}{505}{11.5}{14}\begin{multicols}{3}\raggedcolumns \textbf{三、1（6）：}\ 可先把\(188\)\allowbreak 凑成\(200\)\allowbreak ：从\(423\)\allowbreak 中拿出\(12\)\allowbreak ，还剩\(411\)\allowbreak 。这样只是把加数中的一部分移给另一个加数，总和不变，再把\(200\)\allowbreak 与\(411\)\allowbreak 相加。\end{multicols}\end{snippet}
% p019_s07
\begin{snippet}{505}{11.5}{14}\begin{multicols}{3}\raggedcolumns \textbf{三、1（7）：}\ 把\(50\)\allowbreak 看成\(5\times10\)\allowbreak ，先用\(5\)\allowbreak 乘\(18\)\allowbreak ，再把所得的\(90\)\allowbreak 扩大到原来的\(10\)\allowbreak 倍。最后补上的一个\(0\)\allowbreak 来自乘\(10\)\allowbreak ，不能漏掉。\end{multicols}\end{snippet}
% p019_s08
\begin{snippet}{505}{11.5}{14}\begin{multicols}{3}\raggedcolumns \textbf{三、1（8）：}\ 乘\(5\)\allowbreak 可以想成乘\(10\)\allowbreak 后取一半，先得到\(670\)\allowbreak ，再把它平均分成两份。也可以把\(67\)\allowbreak 拆成\(60\)\allowbreak 与\(7\)\allowbreak ，分别乘\(5\)\allowbreak 后把两部分合起来。\end{multicols}\end{snippet}
% p019_s09
\begin{snippet}{505}{11.5}{14}\begin{multicols}{3}\raggedcolumns \textbf{三、1（9）：}\ 两次除法同级，应从左往右依次计算。也可以把\(2\)\allowbreak 和\(4\)\allowbreak 相乘作为总除数，表示把\(32\)\allowbreak 先分成\(2\)\allowbreak 份后，每份再分成\(4\)\allowbreak 份；不能改成除以\(2\div4\)\allowbreak 。\end{multicols}\end{snippet}
% p019_s10
\begin{snippet}{505}{11.5}{14}\begin{multicols}{3}\raggedcolumns \textbf{三、1（10）：}\ 先把\(34\)\allowbreak 平均分成\(2\)\allowbreak 份，所得的商才是要从\(56\)\allowbreak 中减去的数。减法必须等除法算完后再算，不能先求\(56\)\allowbreak 与\(34\)\allowbreak 的差。\end{multicols}\end{snippet}
% p019_s11
\begin{snippet}{505}{11.5}{14}\begin{multicols}{3}\raggedcolumns \textbf{三、2（1）：}\ 先处理括号，但括号里也有运算顺序：先算\(360\div15\)\allowbreak ，再与\(12\)\allowbreak 相加。整个括号的和是外层除法的除数，不能把\(360\)\allowbreak 分别除以括号中的两个数。\end{multicols}\end{snippet}
% p019_s12
\begin{snippet}{505}{11.5}{14}\begin{multicols}{3}\raggedcolumns \textbf{三、2（2）：}\ 最里面的小括号先算，\(103\)\allowbreak 与\(97\)\allowbreak 相加正好凑成\(200\)\allowbreak 。再用这个和减去\(131\)\allowbreak ，得到中括号的值，最后才用\(80\)\allowbreak 去乘它。\par 简便提示（学完运算律后）：括号按原顺序算出\(69\)\allowbreak 后，把它看成\(70-1\)\allowbreak ，先算\(70\)\allowbreak 份\(80\)\allowbreak ，再减去多算的\(1\)\allowbreak 份\(80\)\allowbreak 。\[\begin{aligned}&80\times69\\={}&80\times(70-1)\\={}&5600-80\\={}&5520\end{aligned}\]\end{multicols}\end{snippet}
% p019_s12_local1
\begin{snippet}{92.05938110351565}{10}{12.8}\SourceSans\setlength{\mathindent}{2bp}\(\displaystyle=80\times[200-131]\)\end{snippet}
% p019_s13
\begin{snippet}{505}{11.5}{14}\begin{multicols}{3}\raggedcolumns \textbf{三、2（3）：}\ 先算乘除部分，\(85\)\allowbreak 先乘\(4\)\allowbreak ，所得的积再除以\(20\)\allowbreak ；乘法、除法同级，顺序从左往右。最后从\(680\)\allowbreak 中减去整个乘除部分的值，不能先算\(680-85\)\allowbreak 。\par 简便提示（学完运算律后）：把\(20\)\allowbreak 拆成\(4\times5\)\allowbreak ，乘\(4\)\allowbreak 与除\(4\)\allowbreak 抵消，可直接算\(85\div5\)\allowbreak 。这只是乘除部分的等值简算，不能先算\(680-85\)\allowbreak 。\[\begin{aligned}&680-85\times4\div20\\={}&680-85\times4\div(4\times5)\\={}&680-85\div5\\={}&680-17\\={}&663\end{aligned}\]\end{multicols}\end{snippet}
% p019_s14
\begin{snippet}{505}{11.5}{14}\begin{multicols}{3}\raggedcolumns \textbf{三、2（4）：}\ 先算\(25\)\allowbreak 与\(11\)\allowbreak 的和，把括号看成一个完整除数。括号去掉后，除法和乘法同级，要先用\(396\)\allowbreak 除以\(36\)\allowbreak ，再乘\(12\)\allowbreak ；不能把\(36\times12\)\allowbreak 一起作除数。\end{multicols}\end{snippet}
% p019_s15
\begin{snippet}{505}{11.5}{14}\begin{multicols}{3}\raggedcolumns 这里按每架无人机都装满\(120\)\allowbreak 千克计算，\(35\)\allowbreak 架各运一次，就有\(35\)\allowbreak 个\(120\)\allowbreak 千克，用每架载质量乘架数。若把\(35\)\allowbreak 架分成\(30\)\allowbreak 架和\(5\)\allowbreak 架，两部分运量相加仍是同一总量，只是使乘法更方便。\end{multicols}\end{snippet}
% p019_s15_local1
\begin{snippet}{80}{10.0}{14}\SourceSans\setlength{\mathindent}{2bp}\(\text{(千克)}\)\end{snippet}
% p019_s16
\begin{snippet}{505}{11.5}{14}\begin{multicols}{3}\raggedcolumns \(600\)\allowbreak 元是\(25\)\allowbreak 本读本的总价，价格相同，平均分到每一本的钱就是单价，所以用总价除以本数。若把\(25\)\allowbreak 看成\(5\times5\)\allowbreak ，连续两次除以\(5\)\allowbreak 只是这道除法的简便算法，单位仍是每本的元数。\end{multicols}\end{snippet}
% p019_s16_local1
\begin{snippet}{80}{10.0}{14}\SourceSans\setlength{\mathindent}{2bp}\(\text{(元)}\)\end{snippet}
% p020_s01
\begin{snippet}{505}{11.5}{14}\begin{multicols}{3}\raggedcolumns 先按“每套价钱\(\times\)\allowbreak 套数”求实际货款，再从付出的\(900\)\allowbreak 元中扣除，剩下的才是找回的钱。另一种思路把\(900\)\allowbreak 元看成能买整\(20\)\allowbreak 套的钱，实际少买\(2\)\allowbreak 套，因此找回的正是这\(2\)\allowbreak 套的价钱。\par \[\begin{aligned}&45\times(900\div45-18)\\={}&45\times(20-18)\\={}&45\times2\\={}&90\text{(元)}\end{aligned}\]\end{multicols}\end{snippet}
% p020_s01_local1
\begin{snippet}{80}{10.0}{14}\SourceSans\setlength{\mathindent}{2bp}\(\text{(元)}\)\end{snippet}
% p020_s02
\begin{snippet}{505}{11.5}{14}\begin{multicols}{3}\raggedcolumns 答：返回时平均每小时行80千米。\par 原路返回说明单程路程没有改变，先用去时每小时的\(60\)\allowbreak 千米乘\(4\)\allowbreak 小时，确定这段路长。再把相同路程平均分到返回的\(3\)\allowbreak 小时，求返回速度；同一路程用时减少，速度应比去时快。\par 简便提示（学完运算律后）：数量关系仍是“去时速度\(\times\)\allowbreak 去时时间\(\div\)\allowbreak 返回时间”。在计算中先用\(60\div3\)\allowbreak 得整十数\(20\)\allowbreak ，再乘\(4\)\allowbreak 。\[\begin{aligned}&60\times4\div3\\={}&60\div3\times4\\={}&20\times4\\={}&80\text{(千米/时)}\end{aligned}\]\end{multicols}\end{snippet}
% p020_s02_local1
\begin{snippet}{80}{10.0}{14}\SourceSans\setlength{\mathindent}{2bp}\(\text{(千米/时)}\)\end{snippet}
% p020_s03
\begin{snippet}{505}{11.5}{14}\begin{multicols}{3}\raggedcolumns 答：两市相距540千米。\par 两车同时出发，都是行驶\(3\)\allowbreak 小时，相遇时两车走过的两段路正好接成两市之间的全程。可先把\(80\)\allowbreak 与\(100\)\allowbreak 相加，得到每小时共同接近的距离，再乘相遇时间；也可各自求出\(3\)\allowbreak 小时的路程后相加。\par \[\begin{aligned}&80\times3+100\times3\\={}&240+300\\={}&540\text{(千米)}\end{aligned}\]\end{multicols}\end{snippet}
% p020_s03_local1
\begin{snippet}{80}{9.4}{14}\SourceSans\setlength{\mathindent}{2bp}\(\text{(千米)}\)\end{snippet}
% p020_s04
\begin{snippet}{505}{11.5}{14}\begin{multicols}{3}\raggedcolumns 按\(2\)\allowbreak 名成人和\(1\)\allowbreak 名\(18\)\allowbreak 岁以下儿童计算，组团游把两种票价分别按人数合计。自由游的\(900\)\allowbreak 元只是单程票价，每人须算两程；住宿和吃饭、交通的两项\(120\)\allowbreak 元都要按\(5\)\allowbreak 天、\(3\)\allowbreak 人累计。也可先合一人的往返机票与五天食宿费，再乘人数，两种次序都要把全部费用算齐后与组团费用比较。\par \end{multicols}\[\begin{aligned}&900\times2\times3+(120+120)\times5\times3\\={}&1800\times3+240\times5\times3\\={}&5400+1200\times3\\={}&5400+3600\\={}&9000\text{(元)}\end{aligned}\]\begin{multicols}{3}\raggedcolumns 也可先求一人的全部费用：\[\begin{aligned}&900\times2+(120+120)\times5\\={}&1800+240\times5\\={}&1800+1200\\={}&3000\text{(元)}\end{aligned}\]\[\begin{aligned}&3000\times3\\={}&9000\text{(元)}\end{aligned}\]\(6200<9000\)\allowbreak ，相差\(9000-6200=2800\)\allowbreak 元。答：组团游更划算，少花2800元。\end{multicols}\end{snippet}
% p020_s04_local1
\begin{snippet}{80}{10.0}{14}\SourceSans\setlength{\mathindent}{2bp}\(\text{(元)}\)\end{snippet}
% p021_s01
\begin{snippet}{505}{11.5}{14}\begin{multicols}{3}\raggedcolumns \(96\)\allowbreak 元是\(3\)\allowbreak 支相同钢笔的总价，先平均分成\(3\)\allowbreak 份，求出一支的钱。每支价格不变，再用这个单价乘\(10\)\allowbreak ，才能求\(10\)\allowbreak 支的总价，不能把\(96\)\allowbreak 直接乘\(10\)\allowbreak 。\par \[\begin{aligned}&96\div3\times10\\={}&32\times10\\={}&320\text{(元)}\end{aligned}\]\end{multicols}\end{snippet}
% p021_s02
\begin{snippet}{505}{11.5}{14}\begin{multicols}{3}\raggedcolumns \textbf{一、2（已知条件）：}\ \(40\)\allowbreak 元是买完这批苹果一共付的钱，表示总价；\(8\)\allowbreak 千克表示买了多少，属于数量。要求每\(1\)\allowbreak 千克的钱，就是把总价平均分给这\(8\)\allowbreak 千克，因此用除法求单价。\end{multicols}\end{snippet}
% p021_s03
\begin{snippet}{505}{11.5}{14}\begin{multicols}{3}\raggedcolumns \textbf{一、2（单价和读法）：}\ 把\(40\)\allowbreak 元平均分给\(8\)\allowbreak 千克，得到每\(1\)\allowbreak 千克应付的钱。单价的单位要同时交代钱数和对应的质量，所以写“元/千克”，斜线读作“每”，不能只写“元”。\par \[\begin{aligned}&40\div8\\={}&5\text{(元/千克)}\end{aligned}\]\end{multicols}\end{snippet}
% p021_s04
\begin{snippet}{505}{11.5}{14}\begin{multicols}{3}\raggedcolumns \textbf{一、3（汽车）：}\ \(72\)\allowbreak 千米表示汽车\(1\)\allowbreak 小时行驶的路程，是单位时间内的路程，也就是速度。写速度单位时用“千米/时”，斜线表示“每”；读的时候先读路程单位，再读“每时”。\end{multicols}\end{snippet}
% p021_s05
\begin{snippet}{505}{11.5}{14}\begin{multicols}{3}\raggedcolumns \textbf{一、3（步行）：}\ \(65\)\allowbreak 米对应的是\(1\)\allowbreak 分钟，表示每分钟的路程，所以速度单位用“米/分”。读作“米每分”时，仍要保留“分”，不能误读成每小时走\(65\)\allowbreak 米。\end{multicols}\end{snippet}
% p021_s06
\begin{snippet}{505}{11.5}{14}\begin{multicols}{3}\raggedcolumns \textbf{一、4（步行速度）：}\ \(840\)\allowbreak 米是\(12\)\allowbreak 分钟的总路程，把它平均分给\(12\)\allowbreak 分钟，就得到每分钟走的米数。路程用米、时间用分，相除所得速度的单位应为“米/分”。\par \[\begin{aligned}&840\div12\\={}&70\text{(米/分)}\end{aligned}\]\end{multicols}\end{snippet}
% p021_s07
\begin{snippet}{505}{11.5}{14}\begin{multicols}{3}\raggedcolumns \textbf{一、4（骑车路程）：}\ \(280\)\allowbreak 米/分表示每\(1\)\allowbreak 分钟骑行\(280\)\allowbreak 米，连续骑\(5\)\allowbreak 分钟就是\(5\)\allowbreak 个这样的路程。用每分钟的路程乘分钟数求总路程，结果单位是米。\par \[\begin{aligned}&280\times5\\={}&1400\text{(米)}\end{aligned}\]\end{multicols}\end{snippet}
% p021_s08
\begin{snippet}{505}{11.5}{14}\begin{multicols}{3}\raggedcolumns \textbf{一、4（船行时间）：}\ 船每\(1\)\allowbreak 小时走\(35\)\allowbreak 千米，要求走完\(105\)\allowbreak 千米需要几个小时，就求其中有几个\(35\)\allowbreak 千米。用总路程除以每小时路程，所得时间单位是小时。\par \[\begin{aligned}&105\div35\\={}&3\text{(小时)}\end{aligned}\]\end{multicols}\end{snippet}
% p021_s09
\begin{snippet}{505}{11.5}{14}\begin{multicols}{3}\raggedcolumns \textbf{一、4（汽车速度）：}\ \(428\)\allowbreak 千米是汽车\(4\)\allowbreak 小时行驶的总路程，求平均每小时行驶多少，就把总路程平均分成\(4\)\allowbreak 份。因时间单位是小时，速度的单位应写“千米/时”，不能沿用步行一栏的“米/分”。\par \[\begin{aligned}&428\div4\\={}&107\text{(千米/时)}\end{aligned}\]\end{multicols}\end{snippet}
% p021_s10
\begin{snippet}{505}{11.5}{14}\begin{multicols}{3}\raggedcolumns “照这样的速度”表示每小时都按\(435\)\allowbreak 千米计算，\(12\)\allowbreak 小时便有\(12\)\allowbreak 个\(435\)\allowbreak 千米。用速度乘时间求路程，也可把\(12\)\allowbreak 小时分成\(10\)\allowbreak 小时和\(2\)\allowbreak 小时，分别求路程后相加。\par \[\begin{aligned}&435\times12\\={}&435\times(10+2)\\={}&4350+870\\={}&5220\text{(千米)}\end{aligned}\]\end{multicols}\end{snippet}
% p021_s11
\begin{snippet}{505}{11.5}{14}\begin{multicols}{3}\raggedcolumns 收到回声前，声音先传播到对面，再从对面返回，所以\(6\)\allowbreak 秒对应来回两段相等的路程。先用\(340\)\allowbreak 乘\(6\)\allowbreak 求往返总路程，再取一半才是山谷宽；也可先把\(6\)\allowbreak 秒平分成两个\(3\)\allowbreak 秒，用速度乘单程时间。不能把往返总路程直接当成宽度。\par 先求往返路程再取一半：\[\begin{aligned}&340\times6\div2\\={}&2040\div2\\={}&1020\text{(米)}\end{aligned}\]先求单程时间再求宽度：\[\begin{aligned}&340\times(6\div2)\\={}&340\times3\\={}&1020\text{(米)}\end{aligned}\]\end{multicols}\end{snippet}
% p021_s12
\begin{snippet}{505}{11.5}{14}\begin{multicols}{3}\raggedcolumns \textbf{一、7（每千克价格）：}\ \(16\)\allowbreak 元对应\(2\)\allowbreak 千克，单价就是把总价平均分给这\(2\)\allowbreak 千克。先求出每\(1\)\allowbreak 千克的钱，后面无论买多少千克，都使用同一个单价。\par \[\begin{aligned}&16\div2\\={}&8\text{(元/千克)}\end{aligned}\]\end{multicols}\end{snippet}
% p021_s13
\begin{snippet}{505}{11.5}{14}\begin{multicols}{3}\raggedcolumns \textbf{一、7（24元购买质量）：}\ 先由前一空确定每千克黄瓜的钱，\(24\)\allowbreak 元里包含几个这样的单价，就能买几千克。这里已知新总价和不变的单价，应求数量，不能把\(24\)\allowbreak 元乘单价。\par \[\begin{aligned}&24\div8\\={}&3\text{(千克)}\end{aligned}\]\end{multicols}\end{snippet}
% p022_s01
\begin{snippet}{505}{11.5}{14}\begin{multicols}{3}\raggedcolumns 按通常题意把\(2\)\allowbreak 分钟看作首次相遇时间：两人从同一点向相反方向绕跑道跑，首次相遇时两段路程正好合成一整圈。先合并每分钟\(65\)\allowbreak 米和\(85\)\allowbreak 米的路程，再乘共同用时，就是跑道周长。若不限定首次，相遇时可能已合跑若干圈，相遇次数未知，便不能唯一确定周长。\par 按首次相遇计算：\[\begin{aligned}&(65+85)\times2\\={}&150\times2\\={}&300\text{(米)}\end{aligned}\]\end{multicols}\end{snippet}
% p022_s02
\begin{snippet}{505}{11.5}{14}\begin{multicols}{3}\raggedcolumns 先统一时间单位：\par\noindent\(\displaystyle \text{半小时}=30\text{分钟}\)\allowbreak \par 把\(432\)\allowbreak 个字平均分到\(6\)\allowbreak 分钟，求每分钟打字数，再乘\(30\)\allowbreak ；也可看出\(30\)\allowbreak 分钟是\(6\)\allowbreak 分钟的\(5\)\allowbreak 倍，打字数也应变为\(5\)\allowbreak 倍。\par 先求每分钟打字数：\par\noindent\[\displaystyle \begin{aligned}&432\div6\times30\\={}&72\times30\\={}&2160\text{(个)}\end{aligned}\]\par 按时间倍数求：\par\noindent\[\displaystyle \begin{aligned}&432\times(30\div6)\\={}&432\times5\\={}&2160\text{(个)}\end{aligned}\]\par\end{multicols}\end{snippet}
% p022_s03
\begin{snippet}{505}{11.5}{14}\begin{multicols}{3}\raggedcolumns \(470\)\allowbreak 千米包括已行和未行两部分，剩下的\(30\)\allowbreak 千米还没有在这\(4\)\allowbreak 小时内行驶。先从全程中扣去未行路程，再把已行路程平均分到\(4\)\allowbreak 小时，才能求出每小时行驶多少。\par \[\begin{aligned}&(470-30)\div4\\={}&440\div4\\={}&110\text{(千米/时)}\end{aligned}\]\end{multicols}\end{snippet}
% p022_s04
\begin{snippet}{505}{11.5}{14}\begin{multicols}{3}\raggedcolumns \(70\)\allowbreak 千米/时表示每小时行\(70\)\allowbreak 千米，\(2\)\allowbreak 小时就是\(2\)\allowbreak 个\(70\)\allowbreak 千米，求总路程用乘法。价格类关系中的“总价、单价、数量”与这道行程题不对应，判断时要看清已知和所求的数量名称。\end{multicols}\end{snippet}
% p022_s05
\begin{snippet}{505}{11.5}{14}\begin{multicols}{3}\raggedcolumns 两种文具的单价和购买数量都不同，应先用每支\(5\)\allowbreak 元乘\(2\)\allowbreak 支，求钢笔款，再用每个\(8\)\allowbreak 元乘\(3\)\allowbreak 个，求文具盒款。两笔钱都属于总花费的组成部分，最后相加，不能把两种单价相加后随意乘一个数量。\par \[\begin{aligned}&5\times2+8\times3\\={}&10+24\\={}&34\text{(元)}\end{aligned}\]\end{multicols}\end{snippet}
% p022_s06
\begin{snippet}{505}{11.5}{14}\begin{multicols}{3}\raggedcolumns 出发和到达时刻的分钟数都是\(43\)\allowbreak ，从\(12\)\allowbreak 时\(43\)\allowbreak 分到\(14\)\allowbreak 时\(43\)\allowbreak 分正好经过\(2\)\allowbreak 小时。把\(406\)\allowbreak 千米平均分到这\(2\)\allowbreak 小时，就是平均速度；这里要求每小时的路程，应使用除法。\par \[\begin{aligned}&406\div2\\={}&203\text{(千米/时)}\end{aligned}\]\end{multicols}\end{snippet}
% p022_s07
\begin{snippet}{505}{11.5}{14}\begin{multicols}{3}\raggedcolumns 比较快慢应统一成相同的速度单位，题中多数已用“米/分”。先换长颈鹿的路程：\par\noindent\(\displaystyle 3\text{千米}=3000\text{米}\)\allowbreak \par 再除以\(3\)\allowbreak 分钟；得到的每分钟路程与另外三种动物的数值比较，数值最大者最快。\end{multicols}\end{snippet}
% p022_s08
\begin{snippet}{505}{11.5}{14}\begin{multicols}{3}\raggedcolumns 两人从桥的两端相向而行，同时走了\(6\)\allowbreak 分钟，相遇时各自路程之和正好等于桥长。每分钟两人共同走过\(70+60\)\allowbreak 米，所以先求速度和，再乘共同时间；不能用速度差求桥长。\par \[\begin{aligned}&(70+60)\times6\\={}&130\times6\\={}&780\text{(米)}\end{aligned}\]\end{multicols}\end{snippet}
% p022_s09
\begin{snippet}{505}{11.5}{14}\begin{multicols}{3}\raggedcolumns 先用每小时\(78\)\allowbreak 千米乘\(3\)\allowbreak 小时，确定服务区离甲地\(234\)\allowbreak 千米。全程\(700\)\allowbreak 千米分成三段，每段大约\(233\)\allowbreak 千米，因此服务区约在从甲到乙的三分之一处；按图估计位置时要从甲端量起。\end{multicols}\end{snippet}
% p023_s01
\begin{snippet}{505}{11.5}{14}\begin{multicols}{3}\raggedcolumns \textbf{三、1（1）：}\ 可以把\(34\)\allowbreak 拆成\(30\)\allowbreak 和\(4\)\allowbreak ，分别求\(5\)\allowbreak 份后相加。也可想成\(34\)\allowbreak 乘\(10\)\allowbreak 后取一半，这两种想法都保持了\(5\)\allowbreak 个\(34\)\allowbreak 的含义。\end{multicols}\end{snippet}
% p023_s02
\begin{snippet}{505}{11.5}{14}\begin{multicols}{3}\raggedcolumns \textbf{三、1（2）：}\ \(39\)\allowbreak 比\(40\)\allowbreak 少\(1\)\allowbreak ，先给\(564\)\allowbreak 加上\(40\)\allowbreak 比较方便，但这样多加了\(1\)\allowbreak 。再把多加的\(1\)\allowbreak 减掉，注意补偿的方向应是减。\end{multicols}\end{snippet}
% p023_s03
\begin{snippet}{505}{11.5}{14}\begin{multicols}{3}\raggedcolumns \textbf{三、1（3）：}\ 把\(21\)\allowbreak 拆成\(20\)\allowbreak 和\(1\)\allowbreak ，先求\(20\)\allowbreak 个\(105\)\allowbreak ，再加上\(1\)\allowbreak 个\(105\)\allowbreak 。第二部分是加\(105\)\allowbreak ，不是只加\(1\)\allowbreak ，这样才没有漏算一份。\end{multicols}\end{snippet}
% p023_s04
\begin{snippet}{505}{11.5}{14}\begin{multicols}{3}\raggedcolumns \textbf{三、1（4）：}\ 先求\(8\)\allowbreak 个\(75\)\allowbreak 是多少，再与前面的\(2\)\allowbreak 相加。\(2\)\allowbreak 不是乘法的因数，不能先把\(2\)\allowbreak 与\(8\)\allowbreak 凑成\(10\)\allowbreak 再乘。\par 简便提示（学完运算律后）：把乘数\(8\)\allowbreak 拆成\(2\times4\)\allowbreak ，用\(4\times75=300\)\allowbreak 凑整。最前面的\(2\)\allowbreak 仍是在乘积之外相加。\[\begin{aligned}&2+8\times75\\={}&2+2\times(4\times75)\\={}&2+2\times300\\={}&2+600\\={}&602\end{aligned}\]\end{multicols}\end{snippet}
% p023_s05
\begin{snippet}{505}{11.5}{14}\begin{multicols}{3}\raggedcolumns \textbf{三、1（5）：}\ 先算\(0\)\allowbreak 与\(2\)\allowbreak 的乘积，积为\(0\)\allowbreak 。给\(76\)\allowbreak 加\(0\)\allowbreak ，原数不变，不能先算\(76+0\)\allowbreak 后再乘\(2\)\allowbreak 。\end{multicols}\end{snippet}
% p023_s06
\begin{snippet}{505}{11.5}{14}\begin{multicols}{3}\raggedcolumns \textbf{三、1（6）：}\ 求商时可从乘法想起：\(19\)\allowbreak 乘\(30\)\allowbreak 得\(570\)\allowbreak ，\(703\)\allowbreak 减去它后还剩\(133\)\allowbreak ，正好是\(7\)\allowbreak 个\(19\)\allowbreak 。两部分合起来有\(30+7\)\allowbreak 个\(19\)\allowbreak ，可用乘法检验是否恰好回到\(703\)\allowbreak 。\end{multicols}\end{snippet}
% p023_s07
\begin{snippet}{505}{11.5}{14}\begin{multicols}{3}\raggedcolumns \textbf{三、1（7）：}\ 把\(18\)\allowbreak 看成\(20\)\allowbreak 少\(2\)\allowbreak ，先求\(20\)\allowbreak 个\(24\)\allowbreak ，再扣去多算的\(2\)\allowbreak 个\(24\)\allowbreak 。这里少的是两份\(24\)\allowbreak ，不能只从大积中减去\(2\)\allowbreak 。\end{multicols}\end{snippet}
% p023_s08
\begin{snippet}{505}{11.5}{14}\begin{multicols}{3}\raggedcolumns \textbf{三、1（8）：}\ 可把\(25\)\allowbreak 拆成\(20\)\allowbreak 和\(5\)\allowbreak ，先给\(327\)\allowbreak 加整十数，再加个位数。也可先拿出\(3\)\allowbreak 与\(327\)\allowbreak 凑成\(330\)\allowbreak ，剩下的\(22\)\allowbreak 仍要继续相加。\end{multicols}\end{snippet}
% p023_s09
\begin{snippet}{505}{11.5}{14}\begin{multicols}{3}\raggedcolumns \textbf{三、1（9）：}\ 先求\(8\div4\)\allowbreak 的商，整个商才是要从\(32\)\allowbreak 中减去的数。减法在除法之后，不能先算\(32-8\)\allowbreak 。\end{multicols}\end{snippet}
% p023_s10
\begin{snippet}{505}{11.5}{14}\begin{multicols}{3}\raggedcolumns \textbf{三、1（10）：}\ 把\(37\)\allowbreak 看成\(40\)\allowbreak 少\(3\)\allowbreak ，先从\(605\)\allowbreak 减\(40\)\allowbreak ，再加回多减的\(3\)\allowbreak 。也可以先减\(30\)\allowbreak 再减\(7\)\allowbreak ，两种口算都要留意跨十位的减法。\end{multicols}\end{snippet}
% p023_s11
\begin{snippet}{505}{11.5}{14}\begin{multicols}{3}\raggedcolumns \textbf{三、2（1）：}\ 括号内先做除法，用\(540\)\allowbreak 除以\(27\)\allowbreak 得到的商再与\(39\)\allowbreak 相加。整个括号的和最后乘\(17\)\allowbreak ，不能让\(17\)\allowbreak 只乘括号中的\(39\)\allowbreak 。\par 简便提示（学完运算律后）：括号照常先算得\(59\)\allowbreak 。最后把\(59\)\allowbreak 看成\(60-1\)\allowbreak ，先算\(60\)\allowbreak 份\(17\)\allowbreak ，再减去多算的一份\(17\)\allowbreak 。\[\begin{aligned}&59\times17\\={}&(60-1)\times17\\={}&1020-17\\={}&1003\end{aligned}\]\end{multicols}\end{snippet}
% p023_s12
\begin{snippet}{505}{11.5}{14}\begin{multicols}{3}\raggedcolumns \textbf{三、2（2）：}\ 先求括号内\(416\)\allowbreak 与\(536\)\allowbreak 的和，再将这个和除以\(34\)\allowbreak 。最后从\(75\)\allowbreak 中减去所得的商，不能把\(75\)\allowbreak 也提前放进除法计算。\end{multicols}\end{snippet}
% p023_s13
\begin{snippet}{505}{11.5}{14}\begin{multicols}{3}\raggedcolumns \textbf{三、2（3）：}\ 加号左边的\(180\)\allowbreak 乘\(3\)\allowbreak 与右边的\(728\)\allowbreak 除以\(8\)\allowbreak 互不依赖，可以在同一步分别完成。等两个结果都求出后再相加，不能按从左到右的顺序先把积与\(728\)\allowbreak 相加。\end{multicols}\end{snippet}
% p023_s14
\begin{snippet}{505}{11.5}{14}\begin{multicols}{3}\raggedcolumns \textbf{三、2（4）：}\ 先算最里面的差\(276-196\)\allowbreak ，把这个差作为\(480\)\allowbreak 的除数。再用\(156\)\allowbreak 除以中括号内的商；两层除法的除数不同，不能把整个过程误看成从左到右连续除。\end{multicols}\end{snippet}
% p023_s15
\begin{snippet}{505}{11.5}{14}\begin{multicols}{3}\raggedcolumns 两人同时出发到相遇，各走了\(3\)\allowbreak 分钟，两家之间的路程由乐乐和欢欢走的两段组成。先把每分钟\(220\)\allowbreak 米与\(260\)\allowbreak 米相加，求两人每分钟合走的路程，再乘\(3\)\allowbreak ；另一法分别求两人三分钟的路程后相加。\par \[\begin{aligned}&220\times3+260\times3\\={}&660+780\\={}&1440\text{(米)}\end{aligned}\]\end{multicols}\end{snippet}
% p023_s15_local1
\begin{snippet}{80}{9.4}{14}\SourceSans\setlength{\mathindent}{2bp}\(\text{(米)}\)\end{snippet}
% p023_s16
\begin{snippet}{505}{11.5}{14}\begin{multicols}{3}\raggedcolumns 两人这次从同一点同时向学校跑，走的是同一方向，欢欢每分钟比乐乐多跑\(40\)\allowbreak 米。经过\(5\)\allowbreak 分钟，累计多跑的路程就是两人间距，因此用速度差乘时间。也可各自求出\(5\)\allowbreak 分钟路程，再用较长的减去较短的；题目说明两人都未到校，才能按这\(5\)\allowbreak 分钟持续行进计算。\par \[\begin{aligned}&(260-220)\times5\\={}&40\times5\\={}&200\text{(米)}\end{aligned}\]也可分别求路程再相减：\[\begin{aligned}&260\times5-220\times5\\={}&1300-1100\\={}&200\text{(米)}\end{aligned}\]答：5分钟后两人相距200米。\end{multicols}\end{snippet}
% p024_s01
\begin{snippet}{505}{11.5}{14}\begin{multicols}{3}\raggedcolumns 答：两地相距368千米。\par 先把到达时刻改用二十四时计时法：\par\noindent\(\displaystyle \text{下午}1\text{时}=13\text{时}\)\allowbreak \par 从\(9\)\allowbreak 时到\(13\)\allowbreak 时共行驶\(4\)\allowbreak 小时。每小时行驶\(92\)\allowbreak 千米，求全程就求\(4\)\allowbreak 个\(92\)\allowbreak 千米，用速度乘时间；按中午分成\(3\)\allowbreak 小时和\(1\)\allowbreak 小时，只是核对同一段总用时。\end{multicols}\end{snippet}
% p024_s01_local1
\begin{snippet}{80}{10.0}{14}\SourceSans\setlength{\mathindent}{2bp}\(\text{(小时)。}\)\end{snippet}
% p024_s01_local2
\begin{snippet}{80}{10.0}{14}\SourceSans\setlength{\mathindent}{2bp}\(\text{(千米)}\)\end{snippet}
% p024_s02
\begin{snippet}{505}{11.5}{14}\begin{multicols}{3}\raggedcolumns 全路包含甲队已修、乙队已修和还没修的\(820\)\allowbreak 米三部分，未修部分应加回去。两队都修了\(130\)\allowbreak 天，可先把每天\(8\)\allowbreak 米和\(12\)\allowbreak 米合成共同进度，再求\(130\)\allowbreak 天合修长度；也可分别累计两队的长度后，再加\(820\)\allowbreak 米。\par \[\begin{aligned}&8\times130+12\times130+820\\={}&1040+1560+820\\={}&2600+820\\={}&3420\text{(米)}\end{aligned}\]\end{multicols}\end{snippet}
% p024_s02_local1
\begin{snippet}{80}{10.0}{14}\SourceSans\setlength{\mathindent}{2bp}\(\text{(米)}\)\end{snippet}
% p024_s03
\begin{snippet}{505}{11.5}{14}\begin{multicols}{3}\raggedcolumns 按照\(3\)\allowbreak 月的日销量计算，应先将\(248\)\allowbreak 台平均分到\(31\)\allowbreak 天，求每天销售的台数。第\(2\)\allowbreak 季度是\(4\)\allowbreak 、\(5\)\allowbreak 、\(6\)\allowbreak 月，共\(91\)\allowbreak 天，再用日销量乘天数；不能直接把\(3\)\allowbreak 月销量乘\(3\)\allowbreak ，因为这三个月比三个\(31\)\allowbreak 天少\(2\)\allowbreak 天。另一法先这样多算，再扣去\(2\)\allowbreak 天的销量。\par \[\begin{aligned}&248\times3-(248\div31)\times2\\={}&744-8\times2\\={}&744-16\\={}&728\text{(台)}\end{aligned}\]\par\textbf{简便提示（学完运算律后）}\par 计算季度天数时，两个30天可先合起来凑60，再加31天；这只是合计同样三个月的天数。\[\begin{aligned}&30+31+30\\={}&(30+30)+31\\={}&60+31\\={}&91\text{(天)}\end{aligned}\]\end{multicols}\end{snippet}
% p024_s03_local1
\begin{snippet}{80}{10.0}{14}\SourceSans\setlength{\mathindent}{2bp}\(\text{(台)}\)\end{snippet}
% p024_s04
\begin{snippet}{505}{11.5}{14}\begin{multicols}{3}\raggedcolumns 按题设的速度上限，\(8\)\allowbreak 分钟最多行驶\(8\)\allowbreak 个\(250\)\allowbreak 米，仍比目的地距离\(2400\)\allowbreak 米少，因而不能准时送到。也可以把\(2400\)\allowbreak 米平均分到\(8\)\allowbreak 分钟，求准时送达所需的平均速度，再和每分钟\(250\)\allowbreak 米的上限比较；比较的双方要用相同的速度单位。\par \[\begin{aligned}&2400\div8\\={}&300\text{(米/分)}\end{aligned}\]所需速度大于规定上限，结论与原位按最多路程比较相同。\end{multicols}\end{snippet}
% p024_s04_local1
\begin{snippet}{80}{10.0}{14}\SourceSans\setlength{\mathindent}{2bp}\(\text{(米)}\)\end{snippet}
% p025_s01
\begin{snippet}{505}{11.5}{14}\begin{multicols}{3}\raggedcolumns 把\(b+c\)\allowbreak 看成两部分，每一部分都要取\(a\)\allowbreak 份，因此\(b\)\allowbreak 和\(c\)\allowbreak 都必须分别与\(a\)\allowbreak 相乘。最后把两个积相加，表示的总量与先把两部分合起来再取\(a\)\allowbreak 份相同；不能漏乘任何一项。\end{multicols}\end{snippet}
% p025_s02
\begin{snippet}{505}{11.5}{14}\begin{multicols}{3}\raggedcolumns \textbf{一、2：}\ 竖式中\(42\)\allowbreak 的个位\(2\)\allowbreak 表示\(2\)\allowbreak ，十位\(4\)\allowbreak 表示\(40\)\allowbreak ，所以两行部分积分别对应\(45\)\allowbreak 乘\(2\)\allowbreak 和乘\(40\)\allowbreak 。第二行虽然写着\(180\)\allowbreak ，因向左错一位，实际表示\(1800\)\allowbreak ；两部分相加体现乘法分配律。\end{multicols}\end{snippet}
% p025_s03
\begin{snippet}{505}{11.5}{14}\begin{multicols}{3}\raggedcolumns \textbf{一、3：}\ 故障限制的是输入时不能按\(9\)\allowbreak 键，并不禁止屏幕计算出含\(9\)\allowbreak 的数。原位方法把\(99\)\allowbreak 看成\(100\)\allowbreak 少\(1\)\allowbreak ，先多算一份\(87\)\allowbreak 再减回。另可先输入\(88+11\)\allowbreak 并按等号键，得到\(99\)\allowbreak 后再乘\(87\)\allowbreak ；也可输入\(87\times33\times3\)\allowbreak ，用\(33\times3\)\allowbreak 代替\(99\)\allowbreak ，两种输入全程都不用\(9\)\allowbreak 键。\par 原位“先乘\(100\)\allowbreak 再减一份”的算式继续计算：\[\begin{aligned}={}&8700-87\\={}&8613\end{aligned}\]先加后乘的方法：\[\begin{aligned}&(88+11)\times87\\={}&99\times87\\={}&8613\end{aligned}\]先乘\(33\)\allowbreak 再乘\(3\)\allowbreak 的方法：\[\begin{aligned}&87\times33\times3\\={}&2871\times3\\={}&8613\end{aligned}\]\end{multicols}\end{snippet}
% p025_s04
\begin{snippet}{505}{11.5}{14}\begin{multicols}{3}\raggedcolumns \textbf{一、4（左上）：}\ 三个加数的位置仍是\(98\)\allowbreak 、\(74\)\allowbreak 、\(26\)\allowbreak ，只把后两个加数先结合，依据加法结合律。这样让\(74\)\allowbreak 与\(26\)\allowbreak 先凑成整百，再加\(98\)\allowbreak ，并没有增减任何一个加数。\end{multicols}\end{snippet}
% p025_s05
\begin{snippet}{505}{11.5}{14}\begin{multicols}{3}\raggedcolumns \textbf{一、4（左中）：}\ \(m\)\allowbreak 、\(21\)\allowbreak 、\(9\)\allowbreak 三个因数的位置没有交换，只让\(21\)\allowbreak 与\(9\)\allowbreak 先相乘，依据乘法结合律。括号内仍应是乘法，改成加法就改变了三个数连乘的意义。\end{multicols}\end{snippet}
% p025_s06
\begin{snippet}{505}{11.5}{14}\begin{multicols}{3}\raggedcolumns \textbf{一、4（左下）：}\ 先把\(380\)\allowbreak 平均分成\(5\)\allowbreak 份，再把每份平均分成\(2\)\allowbreak 份，相当于总共分成\(5\times2\)\allowbreak 份。因而括号里合并除数时应相乘，外面仍用除法，这也能先把除数凑成\(10\)\allowbreak 。\end{multicols}\end{snippet}
% p025_s07
\begin{snippet}{505}{11.5}{14}\begin{multicols}{3}\raggedcolumns \textbf{一、4（右上）：}\ 从\(612\)\allowbreak 中先减\(124\)\allowbreak ，又减\(76\)\allowbreak ，一共减掉的是这两个数的和。先让\(124\)\allowbreak 和\(76\)\allowbreak 凑成整百，括号内要用加法，括号外仍用减法。\end{multicols}\end{snippet}
% p025_s08
\begin{snippet}{505}{11.5}{14}\begin{multicols}{3}\raggedcolumns \textbf{一、4（右中）：}\ \(102\)\allowbreak 可以分成\(100\)\allowbreak 和\(2\)\allowbreak ，所以先给\(389\)\allowbreak 加整百数，再把剩下的\(2\)\allowbreak 加上。两部分都属于原来的加数，第二步仍是加，不能因为先加了\(100\)\allowbreak 而把\(2\)\allowbreak 减掉。\end{multicols}\end{snippet}
% p025_s09
\begin{snippet}{505}{11.5}{14}\begin{multicols}{3}\raggedcolumns \textbf{一、4（右下）：}\ \(a+b\)\allowbreak 中的两个加数都要各取\(9\)\allowbreak 份，因此分开后\(a\)\allowbreak 要乘\(9\)\allowbreak ，\(b\)\allowbreak 也要乘\(9\)\allowbreak 。这就是乘法分配律，两个积之间仍用加号连接。\end{multicols}\end{snippet}
% p025_s10
\begin{snippet}{505}{11.5}{14}\begin{multicols}{3}\raggedcolumns \textbf{一、5（左上）：}\ 两边都有\(50\)\allowbreak 个\(12\)\allowbreak ，只需比较另外一部分。左边的\(33\)\allowbreak 也要乘\(50\)\allowbreak ，表示\(50\)\allowbreak 个\(33\)\allowbreak ；右边只加\(1\)\allowbreak 个\(33\)\allowbreak ，所以左边较大，不能漏看括号外因数对\(33\)\allowbreak 的作用。\end{multicols}\end{snippet}
% p025_s11
\begin{snippet}{505}{11.5}{14}\begin{multicols}{3}\raggedcolumns \textbf{一、5（右上）：}\ 两边都先算\(132\)\allowbreak 减\(68\)\allowbreak ，得到相同的中间数。左边再减去\(12\)\allowbreak ，右边却加上\(12\)\allowbreak ，所以左边比右边少了两份\(12\)\allowbreak ，不必把两个长式分别算完再比较。\end{multicols}\end{snippet}
% p025_s12
\begin{snippet}{505}{11.5}{14}\begin{multicols}{3}\raggedcolumns \textbf{一、5（左中）：}\ 两边先用\(720\)\allowbreak 除以\(8\)\allowbreak ，所得的商相同且大于\(0\)\allowbreak 。接着左边乘\(6\)\allowbreak 会把这个商扩大，右边除以\(6\)\allowbreak 会把这个商缩小，因此左边较大；乘除同级，也都要先从左边算起。\end{multicols}\end{snippet}
% p025_s13
\begin{snippet}{505}{11.5}{14}\begin{multicols}{3}\raggedcolumns \textbf{一、5（右中）：}\ 左边三个因数连乘，右边只是把\(20\)\allowbreak 与\(8\)\allowbreak 先结合。因数和它们的位置都没有变化，依据乘法结合律，改变计算的结合顺序不会改变积。\end{multicols}\end{snippet}
% p025_s14
\begin{snippet}{505}{11.5}{14}\begin{multicols}{3}\raggedcolumns \textbf{一、5（左下）：}\ 两式的因数都是\(12\)\allowbreak 与\(56\)\allowbreak ，只是前后位置交换了。依据乘法交换律，交换两个因数的位置，积不变，无须把两个乘积都算出再比较。\end{multicols}\end{snippet}
% p025_s15
\begin{snippet}{505}{11.5}{14}\begin{multicols}{3}\raggedcolumns \textbf{一、5（右下）：}\ 两边都有\(85\)\allowbreak 个\(100\)\allowbreak ，只需比较剩余部分。左边还包含\(3\)\allowbreak 个\(85\)\allowbreak ，右边只加\(3\)\allowbreak ，因此左边较大；把\(103\)\allowbreak 拆成\(100+3\)\allowbreak 时，两部分都要乘\(85\)\allowbreak 。\end{multicols}\end{snippet}
% p025_s16
\begin{snippet}{505}{11.5}{14}\begin{multicols}{3}\raggedcolumns 两个乘积都有共同因数\(25\)\allowbreak ，可反用乘法分配律，先把圆形与三角形代表的数相加，再乘\(25\)\allowbreak 。已知两者的和是\(4\)\allowbreak ，直接用这个整体代入即可，不需要分别求出两个图形表示的数。\par 反用分配律后继续计算：\[\begin{aligned}={}&25\times(\bigcirc+\triangle)\\={}&25\times4\\={}&100\end{aligned}\]\end{multicols}\end{snippet}
% p025_s17
\begin{snippet}{505}{11.5}{14}\begin{multicols}{3}\raggedcolumns 正确计算中括号内的6应取8份，错式只取1份；含方框的部分相同。因此差就是少算的7份6，反用乘法分配律，无须求方框的值。\[\begin{aligned}&8\times6-6\\={}&6\times(8-1)\\={}&6\times7\\={}&42\end{aligned}\]\end{multicols}\end{snippet}
% p026_s01
\begin{snippet}{505}{11.5}{14}\begin{multicols}{3}\raggedcolumns 无论把\(32\)\allowbreak 拆成和、差还是积，拆分后的数值都必须仍是\(32\)\allowbreak 。A中的\(8+4\)\allowbreak 只有\(12\)\allowbreak ，少算了若干份\(125\)\allowbreak ；B的\(8\times4\)\allowbreak 、C的\(30+2\)\allowbreak 、D的\(40-8\)\allowbreak 都保留了\(32\)\allowbreak ，因此能表示原积。\end{multicols}\end{snippet}
% p026_s02
\begin{snippet}{505}{11.5}{14}\begin{multicols}{3}\raggedcolumns C把连续除以\(2\)\allowbreak 、\(5\)\allowbreak 合并为除以它们的积，先把除数凑成\(10\)\allowbreak ，符合除法的性质。A右边让\(7\)\allowbreak 也乘了\(54\)\allowbreak ，B从\(100\)\allowbreak 个\(16\)\allowbreak 减去的应是一整个\(16\)\allowbreak 而非\(1\)\allowbreak ，D把数字倒置不属于交换因数位置，均不能保证两边相等。\end{multicols}\end{snippet}
% p026_s03
\begin{snippet}{505}{11.5}{14}\begin{multicols}{3}\raggedcolumns \textbf{二、3：}\ 本题非唯一：C、D均可简算，见解析。\par 选D时，\(256\)\allowbreak 与\(244\)\allowbreak 先合成\(500\)\allowbreak ，利用连续减去两个数等于减去它们的和来简算。选C时，可先减\(283\)\allowbreak ，让\(783\)\allowbreak 先减到\(500\)\allowbreak ，再减\(256\)\allowbreak ，利用的是交换两次减去的先后而总共减去的数不变。两种都能先凑出\(500\)\allowbreak ，题干没有限定必须先合并减数，因此答案不唯一。\par 选D，先合并两个减数：\[\begin{aligned}&783-(256+244)\\={}&783-500\\={}&283\end{aligned}\]选C，先减\(283\)\allowbreak 凑整：\[\begin{aligned}&783-283-256\\={}&500-256\\={}&244\end{aligned}\]\end{multicols}\end{snippet}
% p026_s04
\begin{snippet}{505}{11.5}{14}\begin{multicols}{3}\raggedcolumns \textbf{二、4：}\ B图把长为\(2+3\)\allowbreak 、宽为\(m\)\allowbreak 的同一长方形分成两块，两块宽都仍是\(m\)\allowbreak 。整体求面积是总长乘宽，分块求面积则是长\(2\)\allowbreak 与长\(3\)\allowbreak 的两块面积相加；两种算法针对同一块面积，便能说明分配律等式。\end{multicols}\end{snippet}
% p026_s05
\begin{snippet}{505}{11.5}{14}\begin{multicols}{3}\raggedcolumns 双倍所得的\(2\)\allowbreak 份、\(4\)\allowbreak 份、\(8\)\allowbreak 份合起来正好是\(14\)\allowbreak 份，所以把这三部分的积相加就是\(23\)\allowbreak 乘\(14\)\allowbreak 。这里把\(14\)\allowbreak 拆成三个加数，再分别乘\(23\)\allowbreak 后相加，体现乘法分配律。\end{multicols}\end{snippet}
% p026_s06
\begin{snippet}{505}{11.5}{14}\begin{multicols}{3}\raggedcolumns ①把图形分成\(15\)\allowbreak 米乘\(32\)\allowbreak 米和\(15\)\allowbreak 米乘\(18\)\allowbreak 米两块；③把这两个宽同为\(15\)\allowbreak 米的长方形合并计算。④先补成边长为\(15+18\)\allowbreak 米和\(32\)\allowbreak 米的大长方形，再减去长\(18\)\allowbreak 米、宽\(32-15\)\allowbreak 米的缺角，也能得到原面积。②的前一项是面积、后一项\(18\)\allowbreak 只是长度，不能相加表示面积，因此只有另外三个算式正确。\end{multicols}\end{snippet}
% p026_s07
\begin{snippet}{505}{11.5}{14}\begin{multicols}{3}\raggedcolumns \textbf{三、1（1）：}\ \(45\)\allowbreak 的个位\(5\)\allowbreak 与\(55\)\allowbreak 的个位\(5\)\allowbreak 先凑成一个十，再与两个整十部分合起来正好成整百。熟悉这类互补的加数，可以直接凑整口算。\end{multicols}\end{snippet}
% p026_s08
\begin{snippet}{505}{11.5}{14}\begin{multicols}{3}\raggedcolumns \textbf{三、1（2）：}\ \(125\)\allowbreak 与\(8\)\allowbreak 是一组常见的凑整因数，乘积正好是一千。以后在连乘中遇到\(125\)\allowbreak ，可以先找是否能从其他因数中拆出\(8\)\allowbreak 。\end{multicols}\end{snippet}
% p026_s09
\begin{snippet}{505}{11.5}{14}\begin{multicols}{3}\raggedcolumns \textbf{三、1（3）：}\ 先交换\(7\)\allowbreak 与\(4\)\allowbreak 的位置，让\(5\)\allowbreak 和\(4\)\allowbreak 先相乘凑成\(20\)\allowbreak 。剩下的因数\(7\)\allowbreak 再与\(20\)\allowbreak 相乘，这里同时用了乘法交换律和结合律，不能漏掉\(7\)\allowbreak 。\end{multicols}\end{snippet}
% p026_s10
\begin{snippet}{505}{11.5}{14}\begin{multicols}{3}\raggedcolumns \textbf{三、1（4）：}\ \(25\)\allowbreak 与\(4\)\allowbreak 能凑成\(100\)\allowbreak ，可先把这两个因数结合。因数\(6\)\allowbreak 保留到下一步，再与整百数相乘；调整乘的顺序不会改变原积。\end{multicols}\end{snippet}
% p026_s11
\begin{snippet}{505}{11.5}{14}\begin{multicols}{3}\raggedcolumns \textbf{三、1（5）：}\ 可以想成\(4\)\allowbreak 个\(25\)\allowbreak ，先把每两个\(25\)\allowbreak 合成\(50\)\allowbreak ，再把两组\(50\)\allowbreak 凑成整百。记住\(4\)\allowbreak 和\(25\)\allowbreak 这组因数，有助于后面的乘除简算。\end{multicols}\end{snippet}
% p026_s12
\begin{snippet}{505}{11.5}{14}\begin{multicols}{3}\raggedcolumns \textbf{三、1（6）：}\ \(97\)\allowbreak 比\(100\)\allowbreak 少\(3\)\allowbreak ，先从\(123\)\allowbreak 减\(100\)\allowbreak 比较容易。这样多减了\(3\)\allowbreak ，必须再加回\(3\)\allowbreak ，不能再减\(3\)\allowbreak 。\end{multicols}\end{snippet}
% p026_s13
\begin{snippet}{505}{11.5}{14}\begin{multicols}{3}\raggedcolumns \textbf{三、1（7）：}\ 乘除同级，可以从左到右依次计算；其中先乘\(4\)\allowbreak 再除以\(4\)\allowbreak ，使数值回到原来的\(6\)\allowbreak 。最后仍要再乘末尾的\(6\)\allowbreak ，不能把两次出现的\(6\)\allowbreak 也消去。\end{multicols}\end{snippet}
% p026_s14
\begin{snippet}{505}{11.5}{14}\begin{multicols}{3}\raggedcolumns \textbf{三、1（8）：}\ 先求\(30\div5\)\allowbreak 的商，这个商才是从\(130\)\allowbreak 中减去的数。这里先除后减，不能按看到数字的先后先算\(130-30\)\allowbreak 。\end{multicols}\end{snippet}
% p027_s01
\begin{snippet}{505}{11.5}{14}\begin{multicols}{3}\raggedcolumns \textbf{三、2（1）：}\ 连续两次减去\(142\)\allowbreak 和\(58\)\allowbreak ，等于一次减去它们的和。两减数正好凑成\(200\)\allowbreak ，先合并后再从\(435\)\allowbreak 中减，比逐次退位减更方便；括号内应是加号。\end{multicols}\end{snippet}
% p027_s02
\begin{snippet}{505}{11.5}{14}\begin{multicols}{3}\raggedcolumns \textbf{三、2（2）：}\ 两个乘积都含有因数\(67\)\allowbreak ，分别表示\(68\)\allowbreak 份和\(32\)\allowbreak 份\(67\)\allowbreak 。把份数相加正好是\(100\)\allowbreak 份，反用乘法分配律即可，合并时共同因数\(67\)\allowbreak 不能漏掉。\end{multicols}\end{snippet}
% p027_s03
\begin{snippet}{505}{11.5}{14}\begin{multicols}{3}\raggedcolumns \textbf{三、2（3）：}\ \(98\)\allowbreak 比\(100\)\allowbreak 少\(2\)\allowbreak ，先把\(215\)\allowbreak 减去\(100\)\allowbreak 容易计算。先减\(100\)\allowbreak 相当于多减\(2\)\allowbreak ，所以要加回\(2\)\allowbreak ，这就是补偿的方向。\end{multicols}\end{snippet}
% p027_s04
\begin{snippet}{505}{11.5}{14}\begin{multicols}{3}\raggedcolumns \textbf{三、2（4）：}\ 前面的乘积表示\(99\)\allowbreak 份\(99\)\allowbreak ，末尾单独的\(99\)\allowbreak 还表示\(1\)\allowbreak 份\(99\)\allowbreak 。两部分合起来恰好是\(100\)\allowbreak 份\(99\)\allowbreak ，可反用分配律凑整百，不能把末尾的一项漏掉。\end{multicols}\end{snippet}
% p027_s05
\begin{snippet}{505}{11.5}{14}\begin{multicols}{3}\raggedcolumns \textbf{三、2（5）：}\ 把\(3600\)\allowbreak 先平均分成\(25\)\allowbreak 份，再将每份分成\(4\)\allowbreak 份，相当于一共平均分成\(25\times4\)\allowbreak 份。先把两个除数合成\(100\)\allowbreak ，再除整百数比较简便；括号内不能用除号。\end{multicols}\end{snippet}
% p027_s06
\begin{snippet}{505}{11.5}{14}\begin{multicols}{3}\raggedcolumns \textbf{三、2（6）：}\ 两旁的\(25\)\allowbreak 适合与\(4\)\allowbreak 凑\(100\)\allowbreak ，\(125\)\allowbreak 适合与\(8\)\allowbreak 凑\(1000\)\allowbreak ，所以把中间的\(32\)\allowbreak 拆成\(4\times8\)\allowbreak 。再用乘法结合律让它们分别配对，既保留了全部因数，又把计算化成整百数乘整千数。\end{multicols}\end{snippet}
% p027_s07
\begin{snippet}{505}{11.5}{14}\begin{multicols}{3}\raggedcolumns \textbf{三、2（7）：}\ 让括号里的\(25\)\allowbreak 和\(11\)\allowbreak 分别乘\(40\)\allowbreak ，得到的两个积相加，依据乘法分配律。\(25\)\allowbreak 与\(40\)\allowbreak 容易凑整，两个乘法互不依赖，可在同一步都算完后再相加；不能只让\(25\)\allowbreak 乘\(40\)\allowbreak 。\end{multicols}\end{snippet}
% p027_s08
\begin{snippet}{505}{11.5}{14}\begin{multicols}{3}\raggedcolumns \textbf{三、2（8）：}\ 两个要加的数\(357\)\allowbreak 和\(43\)\allowbreak 能凑成\(400\)\allowbreak ，两个要减的数\(156\)\allowbreak 和\(144\)\allowbreak 能凑成\(300\)\allowbreak 。先分别合并这两类数，再用总加数减去总减数；重组时每个数原来的加、减含义必须保留，不能把减数变成加数。\end{multicols}\end{snippet}
% p027_s09
\begin{snippet}{505}{11.5}{14}\begin{multicols}{3}\raggedcolumns \textbf{三、2（9）：}\ 原位把\(35\)\allowbreak 拆成\(7\times5\)\allowbreak ，连续除以\(7\)\allowbreak 和\(5\)\allowbreak ，便于先得到整十的商。另一法利用商不变性质，让被除数和除数同时乘\(2\)\allowbreak ，把除数\(35\)\allowbreak 化成\(70\)\allowbreak 再算；必须两者一起乘，不能只改除数。\par \[\begin{aligned}&(280\times2)\div(35\times2)\\={}&560\div70\\={}&8\end{aligned}\]\end{multicols}\end{snippet}
% p027_s10
\begin{snippet}{505}{11.5}{14}\begin{multicols}{3}\raggedcolumns 答：一共跳700下。\par 要求四天一共跳多少下，要把四天的次数全部相加。\(156\)\allowbreak 与\(144\)\allowbreak 能凑\(300\)\allowbreak ，\(195\)\allowbreak 与\(205\)\allowbreak 能凑\(400\)\allowbreak ，交换并结合这些加数不会改变总次数，计算时也不会遗漏任何一天。\end{multicols}\end{snippet}
% p027_s10_local1
\begin{snippet}{80}{10.0}{14}\SourceSans\setlength{\mathindent}{2bp}\(\text{(下)}\)\end{snippet}
% p027_s11
\begin{snippet}{505}{11.5}{14}\begin{multicols}{3}\raggedcolumns 答：一共有1200台。\par 每组25台，每区4组，共12区。原位主法先求每区台数，再乘区数，已先凑100。另一数量思路先用每区4组乘12区，表示全仓库的小组数；列式后仍用乘法结合律先算25与4，避免直接算25乘48。两条分析路径的层次不同，简算后的计算顺序可以相同。\par \textbf{另一数量思路：按总小组数列式，计算先凑100}\par \[\begin{aligned}&25\times(4\times12)\\={}&(25\times4)\times12\\={}&100\times12\\={}&1200\text{(台)}\end{aligned}\]答：这个仓库一共有1200台分拣机器人。\end{multicols}\end{snippet}
% p027_s11_local1
\begin{snippet}{80}{10.0}{14}\SourceSans\setlength{\mathindent}{2bp}\(\text{(台)}\)\end{snippet}
% p028_s01
\begin{snippet}{505}{11.5}{14}\begin{multicols}{3}\raggedcolumns 从总善款中扣去两笔费用，665与135的和正好是800，原位主法已简便。按购买先后列成连减式时，也应用减法的性质把两个减数合起来，避免先算1500减665。这与“先求总支出”是同一数量关系，不另算一种独立解法。\par \textbf{按购买顺序列式，仍先合并支出凑800}\par \[\begin{aligned}&1500-665-135\\={}&1500-(665+135)\\={}&1500-800\\={}&700\text{(元)}\end{aligned}\]答：还剩700元善款。\end{multicols}\end{snippet}
% p028_s01_local1
\begin{snippet}{80}{10.0}{14}\SourceSans\setlength{\mathindent}{2bp}\(\text{(元)}\)\end{snippet}
% p028_s02
\begin{snippet}{505}{11.5}{14}\begin{multicols}{3}\raggedcolumns 答：每个小队12名。\par 先分清25个中队与每中队4个小队。原位主法先按每中队4个小队乘25个中队，得到100个小队，再平均分配。另一路数量思路是先分给25个中队，再分给每中队4个小队；列成连除式后，按连除性质合并两个除数凑100，计算仍然简便。\par \textbf{另一数量思路：先分中队，再分小队}\par \[\begin{aligned}&1200\div25\div4\\={}&1200\div(25\times4)\\={}&1200\div100\\={}&12\text{(名)}\end{aligned}\]答：平均每个小队有12名少先队员。\end{multicols}\end{snippet}
% p028_s02_local1
\begin{snippet}{80}{10.0}{14}\SourceSans\setlength{\mathindent}{2bp}\(\text{(名)}\)\end{snippet}
% p028_s03
\begin{snippet}{505}{11.5}{14}\begin{multicols}{3}\raggedcolumns 两队合作16天，全长等于两队各自完成的路程之和。原位主法先合并每天的米数，53与47正好凑100。另一路从两队路程相加列式，每个乘法都按“每天米数乘天数”；再反用乘法分配律，避免分别计算两个较大的乘积。\par \textbf{另一数量思路：两队路程相加，计算先凑100}\par \[\begin{aligned}&53\times16+47\times16\\={}&(53+47)\times16\\={}&100\times16\\={}&1600\text{(米)}\end{aligned}\]答：配送通道全长1600米。\end{multicols}\end{snippet}
% p028_s03_local1
\begin{snippet}{80}{10.0}{14}\SourceSans\setlength{\mathindent}{2bp}\(\text{(米)}\)\end{snippet}
% p028_s04
\begin{snippet}{505}{11.5}{14}\begin{multicols}{3}\raggedcolumns 答：甲队多修96米。\par 两队施工天数同为16天，比较的是各自总路程之差。原位主法先求每天多修的6米，再乘16天，已很简便。另一路从两队总路程相减列式，仍可反用乘法分配律，先比较每天米数，避免分别求848米、752米。\par \textbf{另一数量思路：两队路程相减，先合并共同天数}\par \[\begin{aligned}&53\times16-47\times16\\={}&(53-47)\times16\\={}&6\times16\\={}&96\text{(米)}\end{aligned}\]答：甲队比乙队多修96米。\end{multicols}\end{snippet}
% p028_s04_local1
\begin{snippet}{80}{10.0}{14}\SourceSans\setlength{\mathindent}{2bp}\(\text{(米)}\)\end{snippet}
% p029_s01
\begin{snippet}{505}{11.5}{14}\begin{multicols}{3}\raggedcolumns 先观察哪些加数能凑成整百：164与536的和是700，因此先合并第二、三箱。其余第一箱的277件再加进去，三箱文具一箱也没有漏算。只改变加法的结合顺序，用的是加法结合律。\par \[\begin{aligned}&164+536\\={}&700\\&277+700\\={}&977\text{(件)}\end{aligned}\]\end{multicols}\end{snippet}
% p029_s02
\begin{snippet}{505}{11.5}{14}\begin{multicols}{3}\raggedcolumns 125与8相乘能凑1000，5与20相乘能凑100，先配成这两组，后面的乘法就只需处理整百、整千。把原来不相邻的因数换到一起用乘法交换律，先算每组再相乘用乘法结合律。两种律各有作用，不能只写其中一种。\par 两组分别为\(125\times8=1000\)\allowbreak 、\(5\times20=100\)\allowbreak ，再把两个整数组的积相乘。\end{multicols}\end{snippet}
% p029_s03
\begin{snippet}{505}{11.5}{14}\begin{multicols}{3}\raggedcolumns \textbf{一、3（左上）：}\ 从左边变到右边，260变为26是缩小到原来的十分之一，48变为480是扩大到原来的10倍。一个因数的缩小正好被另一个因数的扩大抵消，所以积不变。比较时无需分别做两次大乘法。\par 变化可写成\(260\div10=26\)\allowbreak 、\(48\times10=480\)\allowbreak ，扩大与缩小的倍数相同。\end{multicols}\end{snippet}
% p029_s04
\begin{snippet}{505}{11.5}{14}\begin{multicols}{3}\raggedcolumns \textbf{一、3（右上）：}\ 301比300多1，表示左边比300个24还多1个24。右边在300个24后面只加1，并没有再加1个24。相同部分抵消后比较24与1，就能判断左边较大。\par 分配后多出的两部分相差\(24-1=23\)\allowbreak ，所以左边较大。\end{multicols}\end{snippet}
% p029_s05
\begin{snippet}{505}{11.5}{14}\begin{multicols}{3}\raggedcolumns \textbf{一、3（左中）：}\ 两边都是从145里一共拿走18和82，差别只在先后顺序。连续减两次，可以把两个减数先合并成100，再一次减去。因此左右两边表示相同的剩余量。\par 先合并减数：\(18+82=100\)\allowbreak ，再求\(145-100=45\)\allowbreak 。\end{multicols}\end{snippet}
% p029_s06
\begin{snippet}{505}{11.5}{14}\begin{multicols}{3}\raggedcolumns \textbf{一、3（右中）：}\ 括号里的差乘50时，33和13都要乘50，表示从33个50中减去13个50。右边却只从33个50中减去13，减去的量小得多。起始量相同，左边减得多，剩下的数就较小。\par 左边实际减去\(13\times50=650\)\allowbreak ，而\(650>13\)\allowbreak ，所以左边剩下较少。\end{multicols}\end{snippet}
% p029_s07
\begin{snippet}{505}{11.5}{14}\begin{multicols}{3}\raggedcolumns \textbf{一、3（左下）：}\ 左边把17看成13与4两部分，每一部分都重复6次。右边分别求13的6倍和4的6倍再相加，计算的是同一总量。这正是乘法分配律，所以两边相等。\end{multicols}\end{snippet}
% p029_s08
\begin{snippet}{505}{11.5}{14}\begin{multicols}{3}\raggedcolumns \textbf{一、3（右下）：}\ 左边是把50除以2与5的积，也就是先除以2再除以5。右边没有括号，乘除同级要从左向右，除以2后又乘5。最后一步一个缩小、一个扩大，所以不能认为两式相等。\par 左边先得除数\(2\times5=10\)\allowbreak ，再算\(50\div10=5\)\allowbreak ；右边先得\(50\div2=25\)\allowbreak ，再算\(25\times5=125\)\allowbreak 。\end{multicols}\end{snippet}
% p029_s09
\begin{snippet}{505}{11.5}{14}\begin{multicols}{3}\raggedcolumns 先求一方队需要多少气球：每人2个，25人就是25份，再把一方队的用量乘8个方队。也可先求8个方队一共多少人，再按每人2个求总数。两条路径分别先求“每方队气球数”和“总人数”，都依据每份数乘份数。\par \[\begin{aligned}&2\times25\times8\\={}&50\times8\\={}&400\text{(个)}\end{aligned}\]另一路径先求\(25\times8=200\text{(人)}\)\allowbreak ，再求\(2\times200=400\text{(个)}\)\allowbreak 。\end{multicols}\end{snippet}
% p029_s10
\begin{snippet}{505}{11.5}{14}\begin{multicols}{3}\raggedcolumns 图中把98行补到100行，多补的是2行。先用每行18棵求100行的棵数，再减去同样每行18棵的2行，不能只减2棵。由整体减去补上的部分，体现乘法分配律。\par 把题中已列的括号式展开，继续算：\[\begin{aligned}={}&18\times100-18\times2\\={}&1800-36\\={}&1764\text{(棵)}\end{aligned}\]\end{multicols}\end{snippet}
% p030_s02
\begin{snippet}{505}{11.5}{14}\begin{multicols}{3}\raggedcolumns 先把67拆成66加1，把182拆成181加1，这样两次分配后都能出现181与66的乘积。第一式还多181个1，第二式还多66个1，因此只比较181与66即可。两组因数的和同为248，181与67的差较小，这个较大的积也说明：正数的和相同时，两数的差越小，积越大。\par 两组因数的和都为\(181+67=182+66=248\)\allowbreak ；共同部分之外，左边比右边多\(181-66=115\)\allowbreak 。\end{multicols}\end{snippet}
% p030_s03
\begin{snippet}{505}{11.5}{14}\begin{multicols}{3}\raggedcolumns 先确认两组因数的和相同，才能使用上一题发现的规律。右边把9999减少1，同时把2026增加1，使两因数更接近，因此右边的积较大。也可把9998与2026的乘积作为共同部分，只比较另加的2026与9998来验证。\par 两组因数的和都为12025，差分别为\(9999-2026=7973\)\allowbreak 、\(9998-2027=7971\)\allowbreak ，右边差较小。共同部分之外，右边又比左边多\(9998-2026=7972\)\allowbreak ，可相互核验。\end{multicols}\end{snippet}
% p030_s04
\begin{snippet}{505}{11.5}{14}\begin{multicols}{3}\raggedcolumns \textbf{一、7（例式）：}\ 先选一个非零除数，再选两个都便于除以它的加数。左边把两部分先合起来平均分，右边把两部分分别平均分后再合起来，每份得到的总量相同。仿写时两部分必须除以同一个数，除数不能为0。\end{multicols}\end{snippet}
% p030_s05
\begin{snippet}{505}{11.5}{14}\begin{multicols}{3}\raggedcolumns \textbf{一、7（字母规律）：}\ 把两部分分别记作\(a,\allowbreak b\)\allowbreak ，把共同的除数记作\(c\)\allowbreak ，就能用字母概括前面的数量关系。括号内两个加数都必须分别除以\(c\)\allowbreak ，然后把两个商相加。还要写明\(c\ne0\)\allowbreak ，因为0不能作除数。\end{multicols}\end{snippet}
% p030_s06
\begin{snippet}{505}{11.5}{14}\begin{multicols}{3}\raggedcolumns 99可以看作100少1，把53分别乘100和1，再求两个积的差。这里改变的是“一个数乘差”的计算形式，两个部分都与53相乘，所以依据是乘法分配律。因数并没有单纯调换位置，也不是只改变连乘的分组。\end{multicols}\end{snippet}
% p030_s07
\begin{snippet}{505}{11.5}{14}\begin{multicols}{3}\raggedcolumns 先按100个27计算，比原来的99个27多算了1个27。要还原原积，应减掉27，而不是只减1或减100。检查选项时看它是否同时保留了整百积，并扣回这一份。\end{multicols}\end{snippet}
% p030_s08
\begin{snippet}{505}{11.5}{14}\begin{multicols}{3}\raggedcolumns 先用乘法结合律，把\(a\)\allowbreak 与\(b\)\allowbreak 的积看成一个整体。这个整体再乘18得到216，反求整体应做除法。已知总积及一个因数，另一个因数等于总积除以这个因数。\par \(216\div18=12\)\allowbreak ，所以前两个因数的积是12。\end{multicols}\end{snippet}
% p030_s09
\begin{snippet}{505}{11.5}{14}\begin{multicols}{3}\raggedcolumns 两式中的\(a\)\allowbreak 没有变化，括号里的另一加数由6变为9，多了3。整个括号还要乘5，所以增加的不是3本身，而是3的5倍。无论\(a\)\allowbreak 取什么数，这一增加量都不变。\par 先求多出的部分\(9-6=3\)\allowbreak ，再求\(5\times3=15\)\allowbreak 。\end{multicols}\end{snippet}
% p030_s10
\begin{snippet}{505}{11.5}{14}\begin{multicols}{3}\raggedcolumns 先按图认数量：15与24是相邻两块地的宽，32是两块地共同的长。两宽相加后乘共同的长，求合并后长方形面积。豆角区域宽20没有进入算式，所以不在所求总量中，选D。若继续求面积，把39看成40少1：\[\begin{aligned}&(15+24)\times32\\={}&39\times32\\={}&(40-1)\times32\\={}&40\times32-1\times32\\={}&1280-32\\={}&1248\text{(平方米)}\end{aligned}\]\end{multicols}\end{snippet}
% p030_s11
\begin{snippet}{505}{11.5}{14}\begin{multicols}{3}\raggedcolumns 图②把同宽的两个长方形合为一个整体，图③把每排的6个与4个合并，都是共同重复3份。①只是三段线段的长度相加，没有3份的含义；④书与笔的件数不同，不能共同乘3。判断时要让每个数与图中同一种数量一一对应。\end{multicols}\end{snippet}
% p031_s01
\begin{snippet}{505}{11.5}{14}\begin{multicols}{3}\raggedcolumns \textbf{三、1（1）：}\ 28再加172正好凑成两个整百。口算时可先看个位8与2凑成10，再把进位合入十位、百位，不需要另列竖式。\end{multicols}\end{snippet}
% p031_s02
\begin{snippet}{505}{11.5}{14}\begin{multicols}{3}\raggedcolumns \textbf{三、1（2）：}\ 连续减去45和55，等于一共减去它们的和。两减数刚好凑100，先合并后再从136中减，比逐次退位更方便。\end{multicols}\end{snippet}
% p031_s03
\begin{snippet}{505}{11.5}{14}\begin{multicols}{3}\raggedcolumns \textbf{三、1（3）：}\ 连续除以25再除以2，可以合并为除以二者的积。25与2先凑成50，200里有几个50就是所求的商。这里是连续两次除法，不能把后一除改成乘。\end{multicols}\end{snippet}
% p031_s04
\begin{snippet}{505}{11.5}{14}\begin{multicols}{3}\raggedcolumns \textbf{三、1（4）：}\ 式中乘法、除法同级，又没有括号，应按从左到右的顺序算。先把160平均分成2份，再把所得的一份数扩大到5倍，不能先把2和5合成除数。\end{multicols}\end{snippet}
% p031_s05
\begin{snippet}{505}{11.5}{14}\begin{multicols}{3}\raggedcolumns \textbf{三、1（5）：}\ 把101看成100加1，表示12的100倍再加1个12。用乘法分配律拆开后，先算整百倍，再把多出来的一份补上，口算更方便。\end{multicols}\end{snippet}
% p031_s06
\begin{snippet}{505}{11.5}{14}\begin{multicols}{3}\raggedcolumns \textbf{三、1（6）：}\ 139与要减去的39能先凑成100，再加65。调换加减的计算顺序时，39仍必须保留减号，65仍保留加号，才能使数值不变。\end{multicols}\end{snippet}
% p031_s07
\begin{snippet}{505}{11.5}{14}\begin{multicols}{3}\raggedcolumns \textbf{三、1（7）：}\ 后面的37就是1个37，前面是99个37。把这两部分合起来恰好是100个37，运用乘法分配律逆用即可口算。\end{multicols}\end{snippet}
% p031_s08
\begin{snippet}{505}{11.5}{14}\begin{multicols}{3}\raggedcolumns \textbf{三、1（8）：}\ 整数乘除可整理次序：先乘8、再除8会还原，余下两个5仍都是乘数。把8除以8与5乘5同一步算完；不要把末尾的乘5误改为除5。\[\begin{aligned}&8\times5\div8\times5\\={}&(8\div8)\times(5\times5)\\={}&1\times25\\={}&25\end{aligned}\]\end{multicols}\end{snippet}
% p031_s09
\begin{snippet}{505}{11.5}{14}\begin{multicols}{3}\raggedcolumns \textbf{三、2（1）：}\ 先观察能凑整的搭配：138与62凑200，293与107凑400。交换加数位置后分别结合，使用了加法交换律和结合律；所有加数仍各用一次。\end{multicols}\end{snippet}
% p031_s10
\begin{snippet}{505}{11.5}{14}\begin{multicols}{3}\raggedcolumns \textbf{三、2（2）：}\ 从574连续减去189、211，表示一共减去两数的和。两减数先相加正好400，之后只需从原数减去整百数，这就是简便处。\end{multicols}\end{snippet}
% p031_s11
\begin{snippet}{505}{11.5}{14}\begin{multicols}{3}\raggedcolumns \textbf{三、2（3）：}\ 括号中的25乘8能直接凑200，因而适合用乘法分配律。把25与7都分别乘8，再把两个积相加；不能只把25乘8而漏掉7的8倍。\end{multicols}\end{snippet}
% p031_s12
\begin{snippet}{505}{11.5}{14}\begin{multicols}{3}\raggedcolumns \textbf{三、2（4）：}\ 连乘时先找能凑整的因数：250与40相乘能凑10000。先交换57与40的位置，再把250与40结合，分别使用乘法交换律和结合律。最后乘57时只需处理整万数。\end{multicols}\end{snippet}
% p031_s13
\begin{snippet}{505}{11.5}{14}\begin{multicols}{3}\raggedcolumns \textbf{三、2（5）：}\ 连续两次平均分，可以把两个除数先相乘。16与5凑80，再求720里有几个80，避免先求45后再除5的中间计算。这是除法性质，除数都不能为0。\end{multicols}\end{snippet}
% p031_s14
\begin{snippet}{505}{11.5}{14}\begin{multicols}{3}\raggedcolumns \textbf{三、2（6）：}\ 102接近100，把它拆成100加2。分别求45的100倍和2倍后相加，整百倍的部分便于口算。拆成和后两部分都要乘45，所用的是乘法分配律。\end{multicols}\end{snippet}
% p031_s15
\begin{snippet}{505}{11.5}{14}\begin{multicols}{3}\raggedcolumns 两种算法求的是同一个长方形的面积。\par 两种表示要说明数的是同一批物品：9行里，每行的12棵与8棵可以分开计数，也可以先合成每行20棵再计数。面积图中两块共用高9，分别求两块面积再相加，与先合并成整体求面积相同。关键是两个部分都重复9次，才能合并出共同因数9。\end{multicols}\end{snippet}
% p031_s16
\begin{snippet}{505}{11.5}{14}\begin{multicols}{3}\raggedcolumns 要求全书页数，必须把两天已看的页数和还没看的页数全部合起来。为了简便，先把178与122凑成300，再加第二天看的63页，依据是加法交换律、结合律。按阅读先后依次相加也能核验总数，只是凑整计算更省步骤。\end{multicols}\end{snippet}
% p031_s16_local1
\begin{snippet}{80}{10.0}{14}\SourceSans\setlength{\mathindent}{2bp}\(\text{(页)}\)\end{snippet}
% p032_s01
\begin{snippet}{505}{11.5}{14}\begin{multicols}{3}\raggedcolumns 两类物品各有120个，可以每种取1个配成一组。原位主法先求一组价格，58与42刚好凑100。另一路按“单价乘数量”分别列出两类收入再相加；反用乘法分配律后也先凑100，避免先算6960与5040。这是两种数量分析路径归到同一简算结构。\par \textbf{另一数量思路：两类收入相加，计算先凑100}\par \[\begin{aligned}&58\times120+42\times120\\={}&(58+42)\times120\\={}&100\times120\\={}&12000\text{(元)}\end{aligned}\]答：这些摆件和挂件一共卖了12000元。\end{multicols}\end{snippet}
% p032_s01_local1
\begin{snippet}{80}{10.0}{14}\SourceSans\setlength{\mathindent}{2bp}\(\text{(元)}\)\end{snippet}
% p032_s02
\begin{snippet}{505}{11.5}{14}\begin{multicols}{3}\raggedcolumns 每套4本，每箱25套，先按每份本数乘套数，得到一箱图书的本数，再乘12箱。另一种方法先求12箱共有多少套，再用每套4本乘总套数。两条路径改变的是先求“每箱本数”还是“总套数”，图书总量不变。第二种列式仍可按乘法结合律先凑4乘25等于100；只是计算次序优化，不增列新的解法。\par \textbf{另法：先求总套数}\par \[\begin{aligned}&4\times(25\times12)\\={}&(4\times25)\times12\\={}&100\times12\\={}&1200\text{(本)}\end{aligned}\]\end{multicols}\end{snippet}
% p032_s02_local1
\begin{snippet}{80}{10.0}{14}\SourceSans\setlength{\mathindent}{2bp}\(\text{(本)}\)\end{snippet}
% p032_s03
\begin{snippet}{505}{11.5}{14}\begin{multicols}{3}\raggedcolumns 答：每组分10本。\par 先使用上一问求出的图书总数，再确定平均分给多少个读书小组。方法一先求15个班共有多少组，用总本数除以总组数；方法二先平均分到15个班，再把每班所得平均分给8个组。平均分的层次不同，最终每组所得相同。\par \textbf{另法：先分到班，再分到组}\par \[\begin{aligned} &1200\div15\div8\\={}&80\div8\\={}&10\text{(本)}\end{aligned}\]\end{multicols}\end{snippet}
% p032_s03_local1
\begin{snippet}{80}{10.0}{14}\SourceSans\setlength{\mathindent}{2bp}\(\text{(本)}\)\end{snippet}
% p032_s04
\begin{snippet}{505}{11.5}{14}\begin{multicols}{3}\raggedcolumns 答：劳动试验田的面积是1080平方米。\par 横向分割时，下方长方形是41乘18，上方是18乘19；交换第二项的因数，反用乘法分配律，把41与19凑成60，原位主法已展示。纵分时先同时求出左块总高37与右块宽23，再用共同的18合并成60乘18。补形时，大长方形中的41乘19包含缺角；扣去缺角后，这一层只留下18乘19，因此也可化为横分式。三种是不同的几何思路，简算变形不另算第四种方法。\par \textbf{另法一：纵向分割，37与23凑60}\par \end{multicols}\[\begin{aligned}&18\times(19+18)+(41-18)\times18\\={}&18\times37+23\times18\\={}&(37+23)\times18\\={}&60\times18\\={}&1080\text{(平方米)}\end{aligned}\]\begin{multicols}{3}\raggedcolumns 答：劳动试验田的面积是1080平方米。\par \textbf{另法二：补形后减缺角，先消去相同面积}\par \end{multicols}\[\begin{aligned}&41\times(19+18)-(41-18)\times19\\={}&41\times19+41\times18-(41\times19-18\times19)\\={}&41\times18+18\times19\\={}&(41+19)\times18\\={}&60\times18\\={}&1080\text{(平方米)}\end{aligned}\]\begin{multicols}{3}\raggedcolumns 答：劳动试验田的面积是1080平方米。\end{multicols}\end{snippet}
% p032_s04_local1
\begin{snippet}{289}{10.5}{12.8}\SourceSans\setlength{\mathindent}{2bp}\[\begin{aligned}&41\times18+18\times19\\={}&(41+19)\times18\\={}&60\times18\\={}&1080\text{(平方米)}\end{aligned}\]答：劳动试验田的面积是1080平方米。\end{snippet}
% p033_s01
\begin{snippet}{505}{11.5}{14}\begin{multicols}{3}\raggedcolumns \textbf{一、1（锐角）：}\ 三角形三个内角合起来是\(180^\circ\)\allowbreak 。如果有两个角都是直角或钝角，它们的和已经不少于\(180^\circ\)\allowbreak ，第三个角就没有度数可分了。因此至少两个角必须小于\(90^\circ\)\allowbreak ，也就是锐角。\end{multicols}\end{snippet}
% p033_s02
\begin{snippet}{505}{11.5}{14}\begin{multicols}{3}\raggedcolumns \textbf{一、1（钝角）：}\ 一个钝角已大于\(90^\circ\)\allowbreak ，若再有一个钝角，两角的和就超过三角形的内角和。由此最多只能有一个钝角。有钝角三角形这样的实例，说明“最多一个”可以达到。\end{multicols}\end{snippet}
% p033_s03
\begin{snippet}{505}{11.5}{14}\begin{multicols}{3}\raggedcolumns \textbf{一、2（内角和）：}\ 把三角形的三个内角拼在一起，可以拼成一个平角。平角是\(180^\circ\)\allowbreak ，所以填的是三个内角的总和，不是某一个内角的度数。这个总和与三角形大小、形状无关。\end{multicols}\end{snippet}
% p033_s04
\begin{snippet}{505}{11.5}{14}\begin{multicols}{3}\raggedcolumns \textbf{一、2（三边关系）：}\ 两边相接后必须能够弯折，与第三边围成封闭图形。如果两边的和等于第三边，三条边只能排成一条直线；若还小于第三边，就连不到一起。因此这里必须用严格的“大于”。\end{multicols}\end{snippet}
% p033_s05
\begin{snippet}{505}{11.5}{14}\begin{multicols}{3}\raggedcolumns 先看它只有一组对边平行，仍然属于梯形。题中一条腰与两底垂直，腰的两端各形成一个直角，因此按这一特征称为直角梯形。\end{multicols}\end{snippet}
% p033_s06
\begin{snippet}{505}{11.5}{14}\begin{multicols}{3}\raggedcolumns 要使切开的两部分都是梯形，分割线连接长方形的一组对边，且不能平行于另一组边。每部分都保留一条原长方形的侧边，这条边同时垂直于两底。因此两个小梯形都是直角梯形。\end{multicols}\end{snippet}
% p033_s07
\begin{snippet}{505}{11.5}{14}\begin{multicols}{3}\raggedcolumns \textbf{一、5（三角形）：}\ 找高对应的底，先看虚线末端的直角记号。高必须与底所在直线垂直，图中被垂直的那条边标为6厘米。水平放置的9厘米边并不与这条高垂直，不能因为画在下方就把它当作对应底。\end{multicols}\end{snippet}
% p033_s08
\begin{snippet}{505}{11.5}{14}\begin{multicols}{3}\raggedcolumns \textbf{一、5（平行四边形）：}\ 底已指定为\(AB\)\allowbreak ，要找从它的对边到\(AB\)\allowbreak 所在直线的垂线段；按直角标记，对应的是标4厘米的\(BD\)\allowbreak 。标3厘米的垂线对应另一组底，不能混用。原图5、4、3厘米不能同时满足所示垂直关系，这里只依标记识别底高，不据冲突数值推算其他长度。\end{multicols}\end{snippet}
% p033_s09
\begin{snippet}{505}{11.5}{14}\begin{multicols}{3}\raggedcolumns \textbf{一、6（等腰三角形）：}\ 等腰三角形的两个底角相等，所以先从内角和中减去顶角，再把剩余度数平均分给两个底角。求出每个底角也为\(60^\circ\)\allowbreak 后，可见三个角全是锐角，按角分为锐角三角形。若按边分，它还属于等边三角形，两种分类标准要区分。\par \[\begin{aligned}&(180^\circ-60^\circ)\div2\\={}&120^\circ\div2\\={}&60^\circ\end{aligned}\]\end{multicols}\end{snippet}
% p033_s10
\begin{snippet}{505}{11.5}{14}\begin{multicols}{3}\raggedcolumns \textbf{一、6（直角三角形）：}\ 直角已经占去\(90^\circ\)\allowbreak ，两个锐角合起来只剩\(90^\circ\)\allowbreak 。其中一个是\(45^\circ\)\allowbreak ，另一个也只能是\(45^\circ\)\allowbreak 。两个角相等的三角形有两条相等的边，所以按边分还是等腰三角形。\par 两个锐角共\(90^\circ\)\allowbreak ，因此\(90^\circ-45^\circ=45^\circ\)\allowbreak 。\end{multicols}\end{snippet}
% p033_s11
\begin{snippet}{505}{11.5}{14}\begin{multicols}{3}\raggedcolumns 周长是沿图形四周走一圈的总长度。平行四边形两组对边分别相等，5厘米与3厘米各要算两次，可以先求相邻两边的和再乘2。题中给的是边长，不是底和高。\par \[\begin{aligned}&(5+3)\times2\\={}&8\times2\\={}&16\text{(厘米)}\end{aligned}\]\end{multicols}\end{snippet}
% p033_s12
\begin{snippet}{505}{11.5}{14}\begin{multicols}{3}\raggedcolumns 等腰三角形有两条相等的腰，但题目没有指明4厘米是腰还是底，要分情况检查。若4厘米作腰，两腰的和小于9厘米，不能围成；因此只能把9厘米作为腰。保留的4、9、9又满足两短边和大于最长边，再求三边总长才有依据。\par 两腰取4厘米时，\(4+4=8<9\)\allowbreak ，应排除。两腰取9厘米时，\(4+9>9\)\allowbreak ，可围成。\[\begin{aligned}&9\times2+4\\={}&18+4\\={}&22\text{(厘米)}\end{aligned}\]\end{multicols}\end{snippet}
% p034_s01
\begin{snippet}{505}{11.5}{14}\begin{multicols}{3}\raggedcolumns 两个已知锐角的和超过\(90^\circ\)\allowbreak ，从\(180^\circ\)\allowbreak 中扣去这部分，余下的第三个角就小于\(90^\circ\)\allowbreak 。原有两角也都小于\(90^\circ\)\allowbreak ，所以三个角全部是锐角。判断三角形类型必须同时考虑第三角，不能只看已有两个锐角。\end{multicols}\end{snippet}
% p034_s02
\begin{snippet}{505}{11.5}{14}\begin{multicols}{3}\raggedcolumns \textbf{一、10（80度是顶角）：}\ 先把\(80^\circ\)\allowbreak 看作顶角，两个未知角就是相等的底角。用三角形内角和扣去顶角后，剩下的度数由两底角平均分。下一种情况还须把\(80^\circ\)\allowbreak 改看作底角，不能漏掉。\par \[\begin{aligned}&(180^\circ-80^\circ)\div2\\={}&100^\circ\div2\\={}&50^\circ\end{aligned}\]\end{multicols}\end{snippet}
% p034_s03
\begin{snippet}{505}{11.5}{14}\begin{multicols}{3}\raggedcolumns \textbf{一、10（80度是底角）：}\ 若\(80^\circ\)\allowbreak 是底角，另一底角也必须是\(80^\circ\)\allowbreak ，先把两个底角合起来。余下的顶角用内角和减去两个底角的和求出。和上一种情形相比，分清已知角的位置，另外两角才能列全。\par \[\begin{aligned}&180^\circ-80^\circ\times2\\={}&180^\circ-160^\circ\\={}&20^\circ\end{aligned}\]\end{multicols}\end{snippet}
% p034_s04
\begin{snippet}{505}{11.5}{14}\begin{multicols}{3}\raggedcolumns 升降台抬高或降低时，连接杆的长度保持不变，改变的是杆件之间的角度。平行四边形在边长不变时仍可改变形状，因此能实现升降。三角形三边确定后形状稳定，不适合这种伸缩变化。\end{multicols}\end{snippet}
% p034_s05
\begin{snippet}{505}{11.5}{14}\begin{multicols}{3}\raggedcolumns 等腰三角形的周长包括一条底和两条相等的腰。先扣去已经知道的12厘米底边，得到的是两腰合计的长度，再平均分成2份，才是一条腰的长度。不能把扣底后的总长直接当成腰长。\par \[\begin{aligned}&(48-12)\div2\\={}&36\div2\\={}&18\text{(厘米)}\end{aligned}\]\end{multicols}\end{snippet}
% p034_s06
\begin{snippet}{505}{11.5}{14}\begin{multicols}{3}\raggedcolumns 先用三角形内角和扣去两个\(70^\circ\)\allowbreak ，确认第三角的度数。两个已知角相等，说明两条对应边相等，所以按边分为等腰；求出的第三角与另外两角都小于\(90^\circ\)\allowbreak ，所以按角分为锐角。两个条件要同时满足选项。\par \[\begin{aligned}&180^\circ-70^\circ-70^\circ\\={}&110^\circ-70^\circ\\={}&40^\circ\end{aligned}\]\end{multicols}\end{snippet}
% p034_s07
\begin{snippet}{505}{11.5}{14}\begin{multicols}{3}\raggedcolumns 从四边形一个顶点连一条对角线，可把它分成两个三角形，因而内角和是两个\(180^\circ\)\allowbreak 。把三个已知内角先合计，再从\(360^\circ\)\allowbreak 中扣除，就得到最后一个角。不要套用一个三角形的内角和。\par \[\begin{aligned}&360^\circ-110^\circ-110^\circ-50^\circ\\={}&360^\circ-270^\circ\\={}&90^\circ\end{aligned}\]\end{multicols}\end{snippet}
% p034_s08
\begin{snippet}{505}{11.5}{14}\begin{multicols}{3}\raggedcolumns 从这个固定顶点出发，它的两条不经过该顶点的对边都可以分别作底。对每一条对边所在直线，都能从该点作出一条垂线段，因此共有两条。垂足落在边的延长线上时，同样是相应底边上的高。\end{multicols}\end{snippet}
% p034_s09
\begin{snippet}{505}{11.5}{14}\begin{multicols}{3}\raggedcolumns 框架拉伸时，各根杆的长度没有改变，只是杆件夹角改变。周长计算四条边的总长度，所以变成长方形后仍与原来相等。形状变化不能直接推断周长变化。\end{multicols}\end{snippet}
% p034_s10
\begin{snippet}{505}{11.5}{14}\begin{multicols}{3}\raggedcolumns 题目中的“任意”要求无论取哪两个三角形都能拼成，但大小或形状不同的三角形不能保证这一点。取两个完全一样的三角形，沿对应边反向拼合，才可以保证拼出平行四边形。给出一组不同的三角形作反例，就足以否定原说法。\end{multicols}\end{snippet}
% p034_s11
\begin{snippet}{505}{11.5}{14}\begin{multicols}{3}\raggedcolumns 三角形至少有两个锐角，所以仅有这个条件还不能确定第三角。可取两个\(30^\circ\)\allowbreak 锐角，剩余角为\(120^\circ\)\allowbreak ，就得到钝角三角形。这个反例满足题设却不满足结论，因而判断错误。\par 反例三个角满足\(30^\circ+30^\circ+120^\circ=180^\circ\)\allowbreak ，其中第三角是钝角。\end{multicols}\end{snippet}
% p034_s12
\begin{snippet}{505}{11.5}{14}\begin{multicols}{3}\raggedcolumns 取两个完全一样的梯形，把其中一个转半周，使对应的一条腰重合拼接。两梯形的上、下底分别接成新图形的一组平行对边，另外两条腰也成为一组平行对边。拼成的四边形符合平行四边形的特征，所以原判断成立。\end{multicols}\end{snippet}
% p035_s01
\begin{snippet}{505}{11.5}{14}\begin{multicols}{3}\raggedcolumns 高要与所选的底一一对应，三角形有三条边，就有三条高。直角三角形的两条直角边可以互作底和高；钝角三角形有两条高需要画到底边的延长线上。高不全在内部，不表示这些高不存在。\end{multicols}\end{snippet}
% p035_s02
\begin{snippet}{505}{11.5}{14}\begin{multicols}{3}\raggedcolumns \textbf{四、1（平行四边形）：}\ 先认准标“底”的边，把三角尺的一条直角边靠住它，另一条直角边对准对边上的一点。沿尺画到垂足，得到连接两条平行线的垂线段，画虚线并标直角。不能把斜边当高，同一底上可作出的各条高长度相等。\end{multicols}\end{snippet}
% p035_s03
\begin{snippet}{505}{11.5}{14}\begin{multicols}{3}\raggedcolumns \textbf{四、1（三角形）：}\ 本图指定的底是上方斜边，先找这条边对面的顶点。用三角尺检验垂直方向，从该顶点向斜底画到垂足，并在交点标出直角。高的方向跟底有关，不一定沿纸面的竖直方向。\end{multicols}\end{snippet}
% p035_s04
\begin{snippet}{505}{11.5}{14}\begin{multicols}{3}\raggedcolumns \textbf{四、2（1）：}\ 先画横跨3个间隔的底边，再在与底相隔2个间隔的平行格线上画同样长的对边。把对应端点连接，使两组对边分别平行。高要按两条平行线之间的垂直距离数2格，斜边的长度不能代替高。\end{multicols}\end{snippet}
% p035_s05
\begin{snippet}{505}{11.5}{14}\begin{multicols}{3}\raggedcolumns \textbf{四、2（2）：}\ 先选相隔4格的两条平行横线安放上下底，高就确定为4厘米。上底居中，下底左右伸出的格数相同，可保证两腰相等。画好后同时检查“两底平行”和“两腰相等”，才是符合要求的等腰梯形。\end{multicols}\end{snippet}
% p035_s06
\begin{snippet}{505}{11.5}{14}\begin{multicols}{3}\raggedcolumns \textbf{四、2（3）：}\ 底边先取3个间隔，第三个顶点放在与底相隔3行、又超出底端外侧的位置。这样对应的高落在底的延长线上，垂直数3格；靠近该顶点一侧的底角大于直角。要分别检查底长、高和钝角三个条件。\end{multicols}\end{snippet}
% p035_s07
\begin{snippet}{505}{11.5}{14}\begin{multicols}{3}\raggedcolumns 圆规开口先截取\(AB\)\allowbreak 的长度，此后画两段弧时保持不变。以两端点为圆心画弧，交点到两端的距离都等于\(AB\)\allowbreak ，连接后就有三条相等的边。保留交点附近的弧线，才能看出两条新边是按同一长度作出的。\end{multicols}\end{snippet}
% p035_s08
\begin{snippet}{505}{11.5}{14}\begin{multicols}{3}\raggedcolumns 答：破碎角是67度。\par 玻璃虽然破损，原来的三个内角总和仍是\(180^\circ\)\allowbreak 。两个保留下来的内角已经知道，用总和扣去它们即可求缺角；可连续减，也可先把两已知角合并再减。最后把所求角与两已知角相加，应回到\(180^\circ\)\allowbreak 。\end{multicols}\end{snippet}
% p036_s01
\begin{snippet}{505}{11.5}{14}\begin{multicols}{3}\raggedcolumns 答：较大锐角72度。\par 直角已占\(90^\circ\)\allowbreak ，两个锐角共有\(90^\circ\)\allowbreak 。把小角看成1份，大角就是4份，合起来5份，先平均分得一份的度数。方法一取4份求大角；方法二从两个锐角的总度数中扣去1份小角，留下的就是大角。\par 小角看作一份，大角四份，共五份。\par \textbf{另法：用直角减去小角}\par \[\begin{aligned}&90^\circ-90^\circ\div(1+4)\\={}&90^\circ-90^\circ\div5\\={}&90^\circ-18^\circ\\={}&72^\circ\end{aligned}\]\end{multicols}\end{snippet}
% p036_s02
\begin{snippet}{505}{11.5}{14}\begin{multicols}{3}\raggedcolumns 两条腰各由26厘米的木条加固，等腰梯形的两腰相等，所以先求两腰总长。周长中除两腰外还包括上底、下底，扣去两腰与已知上底后，剩下的才是下底。可连续扣除，也可先合计三条已知边再扣除，依据都是四边长度之和等于周长。\end{multicols}\end{snippet}
% p036_s02_local1
\begin{snippet}{80}{10.0}{14}\SourceSans\setlength{\mathindent}{2bp}\(\text{(厘米)}\)\end{snippet}
% p036_s03
\begin{snippet}{505}{11.5}{14}\begin{multicols}{3}\raggedcolumns 两侧三角形中，相等的小角均为36度。\par 从同一顶点连对角线把五边形分成3个三角形，先求全部内角的总度数，再平均分给5个相等内角。两侧三角形各有两个相等的小角，扣去\(108^\circ\)\allowbreak 后要再平均分成2份。方法一从A处整个内角中减去左右两个小角；方法二看中间的三角形，先用C、D各自的内角扣去旁边小角，再用三角形内角和求中间角。\par 五边形从一个顶点连对角线，可分成三个三角形。两侧小角的求法补为：\[\begin{aligned}&(180^\circ-108^\circ)\div2\\={}&72^\circ\div2\\={}&36^\circ\end{aligned}\]\textbf{另法：看中间的三角形ACD}\par 在C、D处的角各为\(108^\circ-36^\circ=72^\circ\)\allowbreak 。\[\begin{aligned}&180^\circ-72^\circ-72^\circ\\={}&108^\circ-72^\circ\\={}&36^\circ\end{aligned}\]\end{multicols}\end{snippet}
% p037_s01
\begin{snippet}{505}{11.5}{14}\begin{multicols}{3}\raggedcolumns 第一问先用内角和扣去两个已知角，所得第三角超过\(90^\circ\)\allowbreak ，因此按角分为钝角三角形。第二问三个内角相等，把\(180^\circ\)\allowbreak 平均分成3份，每角都是\(60^\circ\)\allowbreak 。三个等角对应三条等边，按边分为等边三角形；三角都是锐角，按角分又是锐角三角形。\par \[\begin{aligned}&180^\circ-37^\circ-49^\circ\\={}&143^\circ-49^\circ\\={}&94^\circ\end{aligned}\]三个角相等时，\(180^\circ\div3=60^\circ\)\allowbreak 。\end{multicols}\end{snippet}
% p037_s02
\begin{snippet}{505}{11.5}{14}\begin{multicols}{3}\raggedcolumns \textbf{一、2：}\ 撑杆与窗框原有的两段边围成一个三角形。三条边的长度一旦固定，三角形的形状也就固定，不能任意晃动。因此这种装置利用的是三角形的稳定性。\end{multicols}\end{snippet}
% p037_s03
\begin{snippet}{505}{11.5}{14}\begin{multicols}{3}\raggedcolumns \textbf{一、3：}\ 篱笆总长就是梯形的周长，先扣去上底4米和下底8米。剩下的是两条腰合起来的长度，等腰梯形的两腰相等，还要平均分成2份。只减两底而不除以2，求到的是两腰总长。\par \[\begin{aligned}&(24-4-8)\div2\\={}&12\div2\\={}&6\text{(米)}\end{aligned}\]\end{multicols}\end{snippet}
% p037_s04
\begin{snippet}{505}{11.5}{14}\begin{multicols}{3}\raggedcolumns \textbf{一、4：}\ 一个平行四边形由两条横向平行线和两条斜向平行线围成。从三条横线中选上下两条有3种，从三条斜线中选左右两条也有3种，逐一配对即可不重不漏。还可按大小复查：4个小的、4个由两个小图组成的、1个整体。\par 按选边计数是\(3\times3=9\)\allowbreak 个；按大小复查是\(4+4+1=9\)\allowbreak 个。\end{multicols}\end{snippet}
% p037_s05
\begin{snippet}{505}{11.5}{14}\begin{multicols}{3}\raggedcolumns \textbf{一、5：}\ 分割改变了图形的大小，却没有改变每块都是三角形这一事实。任意一个小三角形的三个内角仍能拼成一个平角，内角和都是\(180^\circ\)\allowbreak 。原三角形内角和不能平均分给两个小三角形。\end{multicols}\end{snippet}
% p037_s06
\begin{snippet}{505}{11.5}{14}\begin{multicols}{3}\raggedcolumns 延长上底后，它与下底同为6厘米，而且两边仍然平行，连接对应端点形成平行四边形。缩短3厘米后，上底长度变成0，原来两个上端点合成一个顶点，图形只剩三条边。先跟踪上底长度和端点怎样变化，再判断图形名称。\par 延长后上底为\(3+3=6\)\allowbreak 厘米；缩短后为\(3-3=0\)\allowbreak 厘米。\end{multicols}\end{snippet}
% p037_s07
\begin{snippet}{505}{11.5}{14}\begin{multicols}{3}\raggedcolumns 第三边必须小于已知两边的和，否则无法首尾相接围成。它还必须大于两边的差，保证较短边与第三边的和超过最长边。因此范围是大于3厘米而小于19厘米，取整厘米数时两端都不能取等号。\par 两边和为\(8+11=19\)\allowbreak 厘米，两边差为\(11-8=3\)\allowbreak 厘米，因此\(3<\text{第三边}<19\)\allowbreak 。取整厘米数，最短为4厘米，最长为18厘米。\end{multicols}\end{snippet}
% p037_s08
\begin{snippet}{505}{11.5}{14}\begin{multicols}{3}\raggedcolumns 先分别把三个已知点当作两邻边的公共顶点，按对边平行且相等补出第四点；检查格图范围后，只保留图内的一点。数梯形时先找与横底平行的另一底，排除原点及组成平行四边形的点，再检查斜向两底的可能。格点列自左、排自上数：平行四边形补第2列第2排；梯形补第2排第1、3、4、6、7列，以及第7列第4排。斜放的最后一种也只有一组对边平行，不能漏数。\par 梯形补点的两类合计为\(5+1=6\)\allowbreak 种。\end{multicols}\end{snippet}
% p038_s01
\begin{snippet}{505}{11.5}{14}\begin{multicols}{3}\raggedcolumns 把一个三角形绕一条边的中点转半周，与另一个完全相同的三角形沿这条对应边拼合。外边界的两组对边分别平行，必定构成平行四边形。要拼成长方形或正方形，还需原三角形具备特定角度、边长条件，题目没有给出。\end{multicols}\end{snippet}
% p038_s02
\begin{snippet}{505}{11.5}{14}\begin{multicols}{3}\raggedcolumns 平行四边形的对边分别相等，因此7米的边有两条，6米的边也有两条。周长要把四边全部合计，可以先求7米与6米的和再乘2。只算两条邻边会漏掉另外一半。\par \[\begin{aligned}&(7+6)\times2\\={}&13\times2\\={}&26\text{(米)}\end{aligned}\]\end{multicols}\end{snippet}
% p038_s03
\begin{snippet}{505}{11.5}{14}\begin{multicols}{3}\raggedcolumns 等腰直角三角形有一个直角，余下两个角的和是\(90^\circ\)\allowbreak 。这两个角相等，必须把\(90^\circ\)\allowbreak 平均分成2份，得到一个底角。直角是顶角，不能把它填作底角。\par \[\begin{aligned}&(180^\circ-90^\circ)\div2\\={}&90^\circ\div2\\={}&45^\circ\end{aligned}\]\end{multicols}\end{snippet}
% p038_s04
\begin{snippet}{505}{11.5}{14}\begin{multicols}{3}\raggedcolumns 伸缩门拉开与合上时，杆件长度基本不变，而相邻杆件的夹角改变。平行四边形在四边长度固定时仍能改变形状，正适合这种拉伸过程。因此对应的是“易变形”的特性。\end{multicols}\end{snippet}
% p038_s05
\begin{snippet}{505}{11.5}{14}\begin{multicols}{3}\raggedcolumns 每组先找最长边，只要两条较短边的和大于它，另外两次检验也能满足。B组两短边之和恰好等于最长边，只能排成直线，不能围成三角形。其余三组两短边之和都超过最长边，所以只有B不符合。\par 逐组检查：A有\(6+6>6\)\allowbreak ，B有\(3+3=6\)\allowbreak ，C有\(7+8>9\)\allowbreak ，D有\(6+10>12\)\allowbreak 。只有B没有满足严格大于。\end{multicols}\end{snippet}
% p038_s06
\begin{snippet}{505}{11.5}{14}\begin{multicols}{3}\raggedcolumns 梯形的高是两底所在平行线之间的垂直距离。分成的两个梯形仍把原平行四边形的同一组对边作为底所在直线，因此它们的高必然相等。分割位置可以变化，面积、周长和两底之和却随之变化，不能保证相等。\end{multicols}\end{snippet}
% p038_s07
\begin{snippet}{505}{11.5}{14}\begin{multicols}{3}\raggedcolumns 最小角已有\(46^\circ\)\allowbreak ，另外两个角都不会比它小。即使取另外两角中的一个也为\(46^\circ\)\allowbreak ，最大角也只剩\(88^\circ\)\allowbreak ，仍小于直角。因此三个角都小于\(90^\circ\)\allowbreak ，可确定为锐角三角形。\par \(180^\circ-46^\circ-46^\circ=88^\circ<90^\circ\)\allowbreak ，这是最大角能达到的上限。\end{multicols}\end{snippet}
% p038_s08
\begin{snippet}{505}{11.5}{14}\begin{multicols}{3}\raggedcolumns 判断特殊关系时，看后一类图形是否满足前一类的全部条件。长方形两组对边分别平行，符合平行四边形的定义；正方形四个角都是直角，也符合长方形的定义。正方形多了四边相等的条件，不妨碍它同时属于长方形。\end{multicols}\end{snippet}
% p038_s09
\begin{snippet}{505}{11.5}{14}\begin{multicols}{3}\raggedcolumns 要否定“一定比腰短”，只需找到高与腰相等的例子。直角梯形中，垂直于两底的那条腰本身就是高，两者长度相等。因此原说法把所有梯形都包括进去，结论不成立。\end{multicols}\end{snippet}
% p038_s10
\begin{snippet}{505}{11.5}{14}\begin{multicols}{3}\raggedcolumns 高不只可以从顶点画，还可以从一条边上的其他点向对边画垂线段。这样的点有无数个，所以平行四边形的高也有无数条。对应同一底的高虽然长度相等，但位置不同，不能只数四个顶点处的线段。\end{multicols}\end{snippet}
% p038_s11
\begin{snippet}{505}{11.5}{14}\begin{multicols}{3}\raggedcolumns 两个梯形完全一样，才能把对应的一条腰严密拼合。把其中一个转半周后拼在另一个旁边，原来的上、下底接成一组平行边，另外两腰构成另一组平行边。所以可得到平行四边形，原判断正确。\end{multicols}\end{snippet}
% p038_s12
\begin{snippet}{505}{11.5}{14}\begin{multicols}{3}\raggedcolumns 剪掉一个角后，要重新看剩下图形有几条边，不能从原内角和直接减去这个角。沿连接两边的线截去顶角，剩下会成为四边形，内角和是\(360^\circ\)\allowbreak ；剪线若经过另一顶点，剩下也可能仍是三角形。图形的内角总和由图形类别决定，不是把剪掉的角度直接扣除。\end{multicols}\end{snippet}
% p038_s13
\begin{snippet}{505}{11.5}{14}\begin{multicols}{3}\raggedcolumns 把一个三角形放大，只会改变边长，三个角的大小不会因此增加。任意三角形的内角和都固定为\(180^\circ\)\allowbreak ，所以大三角形与小三角形的内角和相同。图形面积大不代表角度和大。\end{multicols}\end{snippet}
% p038_s14
\begin{snippet}{505}{11.5}{14}\begin{multicols}{3}\raggedcolumns 梯形必须是“只有一组对边平行”的四边形。平行四边形也具有一组对边平行，但它的另一组对边同样平行，便不是这里所说的梯形。原条件没有限制另一组边，因此不足以保证结论。\end{multicols}\end{snippet}
% p039_s01
\begin{snippet}{505}{11.5}{14}\begin{multicols}{3}\raggedcolumns \textbf{四、1（1）：}\ 等腰三角形的顶点应与底边两端距离相等，所以先找到给定底边的中点。经过中点沿垂直于底的方向选一点，再与底的两端连接，两腰便对应相同的格数。最后从顶点画到底边中点的垂线段，标出直角；本图底竖放，高就沿横向画。\end{multicols}\end{snippet}
% p039_s02
\begin{snippet}{505}{11.5}{14}\begin{multicols}{3}\raggedcolumns \textbf{四、1（2）：}\ 先画只有一组对边平行的梯形，再从上底的一端向下底作线段，使它平行于另一侧的腰。这样一块同时具有两组平行对边，是平行四边形；余下三条线段围成三角形。完成后检查分割线确实接到下底，且与所选腰平行。\end{multicols}\end{snippet}
% p039_s03
\begin{snippet}{505}{11.5}{14}\begin{multicols}{3}\raggedcolumns 先按线段\(a\)\allowbreak 的长度画底边，两端点分别作圆心。以\(b,\allowbreak c\)\allowbreak 为半径画弧，弧的交点同时满足到两底端距离分别为\(b,\allowbreak c\)\allowbreak ，连接后就满足三边条件。每画一段弧时圆规开口要固定，交点附近的弧线应保留，不能凭目测选第三顶点。\end{multicols}\end{snippet}
% p039_s04
\begin{snippet}{505}{11.5}{14}\begin{multicols}{3}\raggedcolumns 问题示例：这个三角形按角分，是什么三角形？\par 问题要围绕上一问实际画出的三角形，并且有可以观察、测量或推理的数学内容。可以从角的分类、边的关系或高的位置提出问题。题目要求“不用解答”，因此写出一个清楚、合理的问题即可，不必再计算结果。\end{multicols}\end{snippet}
% p039_s05
\begin{snippet}{505}{11.5}{14}\begin{multicols}{3}\raggedcolumns 答：等边三角形的边长是12厘米。\par 同一根铁丝改围图形，总长度不变。按正方形每边9厘米乘4条边列出总长，再平均分成3条边；计算时利用乘除运算的性质，先算9除以3得3，再乘4。另一种理解是把每条原边都平均分成3小份，四条原边共12小份，每条新边取4小份。两种理解得到同一个简算式，原位已经写出，不在旁边重抄。\end{multicols}\end{snippet}
% p039_s05_local1
\begin{snippet}{359}{10.5}{12.8}\SourceSans\setlength{\mathindent}{2bp}\[\begin{aligned}&9\times4\div3\\={}&9\div3\times4\\={}&3\times4\\={}&12\text{(厘米)}\end{aligned}\]答：等边三角形的边长是12厘米。\end{snippet}
% p040_s01
\begin{snippet}{505}{11.5}{14}\begin{multicols}{3}\raggedcolumns 关系：\ensuremath{\angle}4 = \ensuremath{\angle}1 + \ensuremath{\angle}2，三角形的一个外角等于不相邻的两个内角之和。\par 先分清\(\angle3\)\allowbreak 是三角形内角，\(\angle4\)\allowbreak 与它拼成一条直线上的平角。方法一用内角和扣去两个已知角求\(\angle3\)\allowbreak ，再用平角扣去\(\angle3\)\allowbreak 求\(\angle4\)\allowbreak 。方法二比较同为\(180^\circ\)\allowbreak 的三角形内角和与平角，扣去共同的\(\angle3\)\allowbreak ，就能说明\(\angle4=\angle1+\angle2\)\allowbreak 。这个关系由已学的两个角度和推出。\par 比较内角和与底边平角，减去共同的\(\angle3\)\allowbreak ，可先求\(\angle4\)\allowbreak ：\[\begin{aligned}&\angle4\\={}&75^\circ+60^\circ\\={}&135^\circ\\&\angle3\\={}&180^\circ-135^\circ\\={}&45^\circ\end{aligned}\]\end{multicols}\end{snippet}
% p040_s02
\begin{snippet}{505}{11.5}{14}\begin{multicols}{3}\raggedcolumns 11 \ensuremath{-}2 = 9，其余两段只能长4厘米和5厘米。\par 第一刀剪下2厘米后，剩下两段总长9厘米。要与2厘米的边围成三角形，余下两段的差必须小于2；整厘米数且和为9，只能配成4与5。把4厘米或5厘米从2厘米刻度起量，便得到两个备选剪口；也可逐个列出3到10厘米的剪口，检查两短边之和是否大于最长边。图上两处是备选，实际只选一处再剪。\end{multicols}\end{snippet}
% p040_s03
\begin{snippet}{505}{11.5}{14}\begin{multicols}{3}\raggedcolumns 先利用平行四边形的对边相等，把切线长度对应到梯形的斜腰。切出的直角三角形另一条边就是原梯形的直角腰，所以这两边之和已有10厘米。剩下的水平边等于下底比上底多出的3厘米，把这三边合计即可，不必分别求两条腰的长度。\end{multicols}\end{snippet}
% p040_s03_local1
\begin{snippet}{80}{9.5}{14}\SourceSans\setlength{\mathindent}{2bp}\(\text{(厘米)}\)\end{snippet}
% p041_s01
\begin{snippet}{505}{11.5}{14}\begin{multicols}{3}\raggedcolumns 从最后得到72向前倒推：9乘括号内的商等于72，先用72除以9求这个商。再按“被除数等于商乘除数”，把所求商乘4，得到方框中的数。代回时仍要先算括号，再乘9检查。\par 倒推两步为\(72\div9=8\)\allowbreak 、\(8\times4=32\)\allowbreak ；代回检验，括号中的商先得8，再乘9正好72。\end{multicols}\end{snippet}
% p041_s02
\begin{snippet}{505}{11.5}{14}\begin{multicols}{3}\raggedcolumns 在同一平面内，两条不同直线如果延长后相交，位置关系就是相交；始终不相交，才是互相平行。垂直是相交的一种特殊情况，还必须形成直角。不能把所有相交的直线都叫垂直。\end{multicols}\end{snippet}
% p041_s03
\begin{snippet}{505}{11.5}{14}\begin{multicols}{3}\raggedcolumns \textbf{一、3（1）：}\ 先找能凑整的加数：357需要43凑400，168需要32凑200。把这两组分别放入括号，所有加数各用一次，和就不变。调换位置用加法交换律，改变先算哪组用加法结合律。\par 两组分别凑成\(357+43=400\)\allowbreak 、\(168+32=200\)\allowbreak 。\end{multicols}\end{snippet}
% p041_s04
\begin{snippet}{505}{11.5}{14}\begin{multicols}{3}\raggedcolumns \textbf{一、3（2）：}\ 括号外的125要分别乘括号里的80和8，不能只乘其中一个加数。展开后仍是这两部分的积相加，对应乘法分配律。80与8的次序保持与题中式子对应即可。\end{multicols}\end{snippet}
% p041_s05
\begin{snippet}{505}{11.5}{14}\begin{multicols}{3}\raggedcolumns \textbf{一、3（3）：}\ 单独的\(a\)\allowbreak 表示1个\(a\)\allowbreak ，可写成\(a\times1\)\allowbreak 。这样前后两部分分别是99个\(a\)\allowbreak 与1个\(a\)\allowbreak ，可以合并为共同因数\(a\)\allowbreak 乘份数的和。括号内填的是99与1，而不是把\(a\)\allowbreak 再填进去。\end{multicols}\end{snippet}
% p041_s06
\begin{snippet}{505}{11.5}{14}\begin{multicols}{3}\raggedcolumns 平行四边形的另外两条边分别与6厘米、4厘米的邻边相等。把相邻两边先相加，再乘2，四边就都计入周长。这里求周长，只用边长，不需要高。\par \[\begin{aligned}&(6+4)\times2\\={}&10\times2\\={}&20\text{(厘米)}\end{aligned}\]\end{multicols}\end{snippet}
% p041_s07
\begin{snippet}{505}{11.5}{14}\begin{multicols}{3}\raggedcolumns 等边三角形的三条边长度完全相同。把周长27厘米平均分给这三条边，就得到一条边长。题目要求的是边长，所以用除以3来平均分。\par \(27\div3=9\text{(厘米)}\)\allowbreak 。\end{multicols}\end{snippet}
% p041_s08
\begin{snippet}{505}{11.5}{14}\begin{multicols}{3}\raggedcolumns ①右边分配后比左边多出3个125，所以左边较小。②两式都有9个36，另加的2个36比2大，所以左边较大；③两个加数都乘45，符合分配律，两边相等。④左边除以积，等于先除2再除5，右边却是除2后乘5，因而左边较小。逐式认清相同部分和变化部分，可少做完整计算。\par ①相同部分之外，右边多\(125\times3=375\)\allowbreak ；②比较\(36\times2=72>2\)\allowbreak 。④左边除数先得\(2\times5=10\)\allowbreak ，商为\(600\div10=60\)\allowbreak ；右边先得\(600\div2=300\)\allowbreak ，再乘5得1500。\end{multicols}\end{snippet}
% p041_s09
\begin{snippet}{505}{11.5}{14}\begin{multicols}{3}\raggedcolumns 直角三角形两锐角的和是\(90^\circ\)\allowbreak ，一个是\(45^\circ\)\allowbreak ，另一个也为\(45^\circ\)\allowbreak ，故有两条等边。拼成大三角形后，图形仍是三角形，内角和仍为\(180^\circ\)\allowbreak ，不能把两个小三角形的内角和直接相加。拼成正方形则看四个直角，总和为\(360^\circ\)\allowbreak 。\par 另一个锐角为\(90^\circ-45^\circ=45^\circ\)\allowbreak ；正方形内角和为\(90^\circ\times4=360^\circ\)\allowbreak 。\end{multicols}\end{snippet}
% p042_s01
\begin{snippet}{505}{11.5}{14}\begin{multicols}{3}\raggedcolumns 每箱24千克，按每袋8千克，先求每箱可分3袋，再按每箱袋数乘6箱。原式中的先乘6再除8，可利用乘除运算的性质改为先除8再乘6，避免先算总质量144千克。收入按每袋12元乘18袋，箱数、袋数与质量不能混用。\par \[\begin{aligned}&24\times6\div8\\={}&24\div8\times6\\={}&3\times6\\={}&18\text{(袋)}\\&12\times18\\={}&216\text{(元)}\end{aligned}\]\end{multicols}\end{snippet}
% p042_s02
\begin{snippet}{505}{11.5}{14}\begin{multicols}{3}\raggedcolumns 为有序计数，可把下方从左到右4个交点记作\(A,\allowbreak B,\allowbreak C,\allowbreak D\)\allowbreak ，上方3个交点记作\(E,\allowbreak F,\allowbreak G\)\allowbreak 。三角形是\(BFC,\allowbreak FCG,\allowbreak CGD\)\allowbreak ；平行四边形是\(ABFE,\allowbreak BCGF,\allowbreak ACGE\)\allowbreak ，它们各有两组平行边。梯形是\(ACFE,\allowbreak ADGE,\allowbreak BDGF,\allowbreak CDGF\)\allowbreak ，逐个检查都只有一组对边平行，较大的组合也不能漏掉。按类别列出后再数个数，可避免把平行四边形重复计入梯形。\end{multicols}\end{snippet}
% p042_s03
\begin{snippet}{505}{11.5}{14}\begin{multicols}{3}\raggedcolumns 括号里的\(8\times9\)\allowbreak 是一个整体除数，表示把720连续按8份、再按9份平均分。去括号后应保持两次都是除法，即依次除以8、9。若把后一项改为乘9，就改变了平均分的含义。\end{multicols}\end{snippet}
% p042_s04
\begin{snippet}{505}{11.5}{14}\begin{multicols}{3}\raggedcolumns 先用分配律把括号里的\(a\)\allowbreak 和4都分别乘25。两式都有\(25\times a\)\allowbreak ，这一部分相同，只需比较4个25与1个25。它们相差3个25，所以差与\(a\)\allowbreak 取值无关。\par \(25\times4=100\)\allowbreak ，两式的差就是\(100-25=75\)\allowbreak 。\end{multicols}\end{snippet}
% p042_s05
\begin{snippet}{505}{11.5}{14}\begin{multicols}{3}\raggedcolumns 先画出\(b\)\allowbreak ，再把\(a,\allowbreak c\)\allowbreak 都画成与它成直角的不同直线，可以看到二者方向相同。同一平面内，垂直于同一条直线的两条不同直线互相平行。这里按题意将\(a,\allowbreak c\)\allowbreak 看作不同直线来判断。\end{multicols}\end{snippet}
% p042_s06
\begin{snippet}{505}{11.5}{14}\begin{multicols}{3}\raggedcolumns 各组只需先检查两条最短边的和能否超过最长边。A的和等于最长边，B、D的和小于最长边，均不能围成。C的5与6合起来超过10，三边才可以闭合为三角形。\par A：\(4+5=9\)\allowbreak ；B：\(3+4=7<8\)\allowbreak ；C：\(5+6=11>10\)\allowbreak ；D：\(2+2=4<5\)\allowbreak 。\end{multicols}\end{snippet}
% p042_s07
\begin{snippet}{505}{11.5}{14}\begin{multicols}{3}\raggedcolumns 最小角是\(50^\circ\)\allowbreak ，说明另外两个角也至少有\(50^\circ\)\allowbreak 。为给最大角留下尽可能多的度数，先从内角和扣去两个\(50^\circ\)\allowbreak ，最多只剩\(80^\circ\)\allowbreak 。最大角都小于直角，三个角便全是锐角。\par 最大角至多\(180^\circ-50^\circ-50^\circ=80^\circ<90^\circ\)\allowbreak 。\end{multicols}\end{snippet}
% p042_s08
\begin{snippet}{505}{11.5}{14}\begin{multicols}{3}\raggedcolumns ①把25与4先结合凑100，再乘3；②125与8直接凑1000；③乘10时各数位向左进一位；④先算100除以10，再从200中减。⑤360里有6个60；⑥11乘20可先乘2再乘10；⑦先把括号内的和凑20；⑧先求括号内的积，或连续除以8、2。每小题先认清运算顺序，再选凑整搭配。\end{multicols}\end{snippet}
% p042_s09
\begin{snippet}{505}{11.5}{14}\begin{multicols}{3}\raggedcolumns \textbf{三、2（1）：}\ 先算最内层小括号，19减14得5；再把16与5相乘凑成整十除数80，直接用640除以80。这里先求易算的积，比拆成先除16再除5更简便。除数是16与5的积，不是它们的和。\end{multicols}\end{snippet}
% p042_s09_local1
\begin{snippet}{176.445529683431}{10.5}{12.8}\SourceSans\setlength{\mathindent}{2bp}\[\begin{aligned}={}&640\div[16\times5]\\={}&640\div80\\={}&8\end{aligned}\]\end{snippet}
% p042_s10
\begin{snippet}{505}{11.5}{14}\begin{multicols}{3}\raggedcolumns \textbf{三、2（2）：}\ 两个25都能与4凑100，把32拆成\(4\times8\)\allowbreak 后，可使前一个25配4、后一个25配8，分别凑100和200。运用乘法结合律调整分组，最后算两个整百数的积。这里拆成的是积，应保留连乘，不能按加法分配。\end{multicols}\end{snippet}
% p042_s11
\begin{snippet}{505}{11.5}{14}\begin{multicols}{3}\raggedcolumns \textbf{三、2（3）：}\ 最后的47是一份47，把它看成\(1\times47\)\allowbreak 。99份与1份合并成100份，再乘每份47，依据是乘法分配律的逆用。共同因数47保持不变，只把份数相加。\end{multicols}\end{snippet}
% p042_s12
\begin{snippet}{505}{11.5}{14}\begin{multicols}{3}\raggedcolumns \textbf{三、2（4）：}\ 98比100少2，先按100个35求积，会多算2个35。再减去这2份，便还原原来的98份。把乘数拆成差后，两部分都乘35，使用乘法分配律来凑整。\end{multicols}\end{snippet}
% p043_s01
\begin{snippet}{505}{11.5}{14}\begin{multicols}{3}\raggedcolumns 画垂线时，一条直角边先靠住已知直线，平移到另一条直角边经过\(A\)\allowbreak ，沿它画线并标直角。画平行线时，用直尺固定三角尺的移动方向，平推到对应边经过\(A\)\allowbreak 再画。最后检查两条所画直线都经过\(A\)\allowbreak ，分别满足直角与不相交的关系。\end{multicols}\end{snippet}
% p043_s02
\begin{snippet}{505}{11.5}{14}\begin{multicols}{3}\raggedcolumns 先画5个间隔长的底，再在垂直相隔3格的另一条平行格线上画同样长的对边。连接对应端点后，两组对边都要分别平行。高从上边一点垂直画到底边，长度按垂直方向数3格并标直角，不能用斜边代替。\end{multicols}\end{snippet}
% p043_s03
\begin{snippet}{505}{11.5}{14}\begin{multicols}{3}\raggedcolumns 实际天数是8减2，得6天；同一条路的总长不变，原位列式按原计划每天米数乘8天，再除以6。利用乘除运算的性质先算120除以6得20，再乘8，避免先求960。另一思路保留实际6天原有的每天120米，把原来2天的任务平均加进来；计算增加部分时也先把120除以6得20。提前完成，每天米数应比120增加，可检查结果方向。\par \textbf{另一数量思路：把原两天任务平均加到实际六天}\par \[\begin{aligned}&120+120\times2\div(8-2)\\={}&120+120\times2\div6\\={}&120+120\div6\times2\\={}&120+20\times2\\={}&120+40\\={}&160\text{(米)}\end{aligned}\]答：实际平均每天修160米。\end{multicols}\end{snippet}
% p043_s03_local1
\begin{snippet}{268}{10.5}{12.8}\SourceSans\setlength{\mathindent}{2bp}\[\begin{aligned}&120\times8\div(8-2)\\={}&120\times8\div6\\={}&120\div6\times8\\={}&20\times8\\={}&160\text{(米)}\end{aligned}\]答：实际平均每天修160米。\end{snippet}
% p043_s04
\begin{snippet}{505}{11.5}{14}\begin{multicols}{3}\raggedcolumns 题中未说明6厘米还是8厘米为腰，等腰三角形有两条相等边，因此必须分两种情况。分别把6或8重复一次，先用两短边和大于最长边检验，两个分支都能成立。最后各求一次三边总长，就得到两种可能周长。\end{multicols}\end{snippet}
% p043_s04_local1
\begin{snippet}{80}{9.5}{14}\SourceSans\setlength{\mathindent}{2bp}\(\text{(厘米)}\)\end{snippet}
% p043_s04_local2
\begin{snippet}{80}{9.5}{14}\SourceSans\setlength{\mathindent}{2bp}\(\text{(厘米)}\)\end{snippet}
% p044_s01
\begin{snippet}{505}{11.5}{14}\begin{multicols}{3}\raggedcolumns 答：等边三角形的边长是12厘米。\par 同一根铁丝改围图形，总长度不变。按正方形每边9厘米乘4条边列出总长，再平均分成3条边；计算时利用乘除运算的性质，先算9除以3得3，再乘4。另一种理解是把每条原边都平均分成3小份，四条原边共12小份，每条新边取4小份。两种理解得到同一个简算式，原位已经写出，不在旁边重抄。\end{multicols}\end{snippet}
% p044_s01_local1
\begin{snippet}{359}{10.5}{12.8}\SourceSans\setlength{\mathindent}{2bp}\[\begin{aligned}&9\times4\div3\\={}&9\div3\times4\\={}&3\times4\\={}&12\text{(厘米)}\end{aligned}\]答：等边三角形的边长是12厘米。\end{snippet}
% p044_s02
\begin{snippet}{505}{11.5}{14}\begin{multicols}{3}\raggedcolumns “买5送1”是付5瓶的钱实际得到6瓶，因此先按6瓶分组，不能按5瓶分组。方法一每组付费5瓶，用每组付费瓶数乘组数。方法二每组赠1瓶，先由组数求赠送瓶数，再从实际得到的24瓶中扣除。\par 每满六瓶赠一瓶。\par \textbf{另法：先求赠品瓶数，再扣去}\par \[\begin{aligned}&24-24\div6\\={}&24-4\\={}&20\text{(瓶)}\end{aligned}\]\end{multicols}\end{snippet}
% p044_s02_local1
\begin{snippet}{67.43834838867188}{9}{12.8}\SourceSans\setlength{\mathindent}{2bp}\(\displaystyle=5\times[24\div6]\)\end{snippet}
% p044_s02_local2
\begin{snippet}{80}{9.5}{14}\SourceSans\setlength{\mathindent}{2bp}\(\text{(瓶)}\)\end{snippet}
% p044_s03
\begin{snippet}{505}{11.5}{14}\begin{multicols}{3}\raggedcolumns 上一问求得付费20瓶，原位主法按每瓶6元乘20瓶，已经简便。另一路先按24瓶全价列式，再扣赠品价；每6瓶赠1瓶，赠品为4瓶。两项价格都有每瓶6元这个共同因数，反用乘法分配律，先把24减4凑20，不必分别求144与24。\par \textbf{另一数量思路：全价减赠品价，先凑20瓶}\par \[\begin{aligned}&6\times24-6\times(24\div6)\\={}&6\times24-6\times4\\={}&6\times(24-4)\\={}&6\times20\\={}&120\text{(元)}\end{aligned}\]答：一共要付120元。\end{multicols}\end{snippet}
% p044_s03_local1
\begin{snippet}{80}{9.5}{14}\SourceSans\setlength{\mathindent}{2bp}\(\text{(元)}\)\end{snippet}
% p044_s04
\begin{snippet}{505}{11.5}{14}\begin{multicols}{3}\raggedcolumns 答：两个年级一共可以挽救31棵大树。\par 五年级比四年级的2倍少50千克，应先乘2再减50，得到500千克。主法先按总回收量列式，再把275与500分别按每25千克换成树数，同一步得到11与20，不必合成775再除25。另一思路先把两年级看作四年级的3份，再扣去少回收50千克对应的棵数；也可分别计算两个年级的棵数再相加。三条是数量分析路径，简算后的相同计算结构不另算新方法。\par \textbf{另法一：按三份推算，再扣少回收部分}\par \[\begin{aligned}&275\div25\times3-50\div25\\={}&11\times3-2\\={}&33-2\\={}&31\text{(棵)}\end{aligned}\]答：两个年级一共可以挽救31棵大树。\par \textbf{另法二：两个年级分别换算棵数}\par \end{multicols}\[\begin{aligned}&275\div25+(275\times2-50)\div25\\={}&11+(550-50)\div25\\={}&11+500\div25\\={}&11+20\\={}&31\text{(棵)}\end{aligned}\]\begin{multicols}{3}\raggedcolumns 答：两个年级一共可以挽救31棵大树。\end{multicols}\end{snippet}
% p044_s04_local1
\begin{snippet}{356}{10.5}{12.8}\SourceSans\setlength{\mathindent}{2bp}\[\begin{aligned}&[275+(275\times2-50)]\div25\\={}&[275+(550-50)]\div25\\={}&[275+500]\div25\\={}&275\div25+500\div25\\={}&11+20\\={}&31\text{(棵)}\end{aligned}\]答：两个年级一共可以挽救31棵大树。\end{snippet}
% p045_s01
\begin{snippet}{505}{11.5}{14}\begin{multicols}{3}\raggedcolumns 要求先减法，就先用小括号把370与295括成一个差。要求再加法，就把200与这个差的和再作为整体括起来。最后的除以5要留在整个括号外，才能保证最后一步是除法。\end{multicols}\end{snippet}
% p045_s02
\begin{snippet}{505}{11.5}{14}\begin{multicols}{3}\raggedcolumns 48元是每1平方米的费用，总面积有96份1平方米，按每份费用乘份数求总费用。为了简算，可先按100平方米计算，再扣掉多算4平方米的费用。扣的是4份48元，不能只减4元。\par \[\begin{aligned}&48\times(100-4)\\={}&4800-192\\={}&4608\text{(元)}\end{aligned}\]\end{multicols}\end{snippet}
% p045_s03
\begin{snippet}{505}{11.5}{14}\begin{multicols}{3}\raggedcolumns 整时时，分针总指向12，时针只会指向整点数字。垂直表示两针形成直角，相当于相隔3个大格。与12分别相隔3个大格的整点是3和9。\par 相隔3个大格时，夹角为\(30^\circ\times3=90^\circ\)\allowbreak 。\end{multicols}\end{snippet}
% p045_s04
\begin{snippet}{505}{11.5}{14}\begin{multicols}{3}\raggedcolumns 题目使用加法交换律，表示两个加数只交换位置，数值不变。左边末位2的数应与右边十位5的数对应，另一个十位4的数应与末位8的数对应。先拼出两个相同加数，再用和为100核验，就能确定全部方框。\par 两个加数对应为52与48，再核验\(52+48=100\)\allowbreak 。\end{multicols}\end{snippet}
% p045_s05
\begin{snippet}{505}{11.5}{14}\begin{multicols}{3}\raggedcolumns \textbf{一、5（1）：}\ 175与225能凑400，264与136也能凑400，先把这两对分别结合。为让不相邻的加数配对，使用加法交换律；先求每对的和，再求总和，使用结合律。填入运算符时始终都是加法，各数各用一次。\par 两组分别为\(175+225=400\)\allowbreak 、\(264+136=400\)\allowbreak 。\end{multicols}\end{snippet}
% p045_s06
\begin{snippet}{505}{11.5}{14}\begin{multicols}{3}\raggedcolumns \textbf{一、5（2）：}\ 右边共同因数是\(a\)\allowbreak ，所以左边两个积都应含有\(a\)\allowbreak 。单独的一份写成\(a\times1\)\allowbreak ，与99份合起来是99加1份，因此括号内填加号。按乘法分配律逆用检查，括号展开后必须与左边的两项完全对应。\end{multicols}\end{snippet}
% p045_s07
\begin{snippet}{505}{11.5}{14}\begin{multicols}{3}\raggedcolumns 每小时印300份，一天工作8小时，一周只算5个工作日。先按“每小时份数\(\times\)\allowbreak 每天小时数\(\times\)\allowbreak 工作天数”列式，再用乘法结合律，先把8与5相乘凑成整十数40，求出一周总工作时间：\[\begin{aligned}&300\times8\times5\\={}&300\times(8\times5)\\={}&300\times40\\={}&12000\text{(份)}\end{aligned}\]原空格后已有“份”，只填12000。\end{multicols}\end{snippet}
% p045_s08
\begin{snippet}{505}{11.5}{14}\begin{multicols}{3}\raggedcolumns ①两式都有5个38，另加的6个38比6大；②把21拆成7与3的积，除以积等于连续除，所以相等。③按200个78求积后，左边需扣78，右边只扣1，左边较小；④共同因数74不变，只比较20与12。⑤共同因数65不变，比较24与80；⑥分别算出535与735，再按大小判断。每式都先认清加、乘、连续除的不同结构。\par 相应的简算依据为：①\(38\times6=228>6\)\allowbreak ；②\(21=7\times3\)\allowbreak ；④\(10\times2=20>12\)\allowbreak ；⑤\(4\times20=80>24\)\allowbreak 。⑥分别按\((100+7)\times5=500+35=535\)\allowbreak 、\((100+5)\times7=700+35=735\)\allowbreak 比较。\end{multicols}\end{snippet}
% p045_s09
\begin{snippet}{505}{11.5}{14}\begin{multicols}{3}\raggedcolumns 同一组白线沿着相同方向排列，把它们所在直线向两端延长也不会相交。这样的线段可以看作互相平行。观察的是白线之间的位置关系，与白线长短无关。\end{multicols}\end{snippet}
% p045_s10
\begin{snippet}{505}{11.5}{14}\begin{multicols}{3}\raggedcolumns 等腰三角形两个底角相等，把每个底角看成1份，顶角就是4份。三个角合起来一共6份，必须用\(180^\circ\)\allowbreak 平均分成6份才能求一个底角。这里的倍数关系只针对一个底角，不能漏算另一个底角的1份。\par \[\begin{aligned}&180^\circ\div(1+1+4)\\={}&180^\circ\div6\\={}&30^\circ\end{aligned}\]\end{multicols}\end{snippet}
% p046_s01
\begin{snippet}{505}{11.5}{14}\begin{multicols}{3}\raggedcolumns 按图示标注填写。\par 指定底后，要从对面顶点向底所在直线找垂线段，不能按线段看起来长短来选。底\(AC\)\allowbreak 对应\(BD\)\allowbreak ，底\(AB\)\allowbreak 对应与它成直角的\(BC\)\allowbreak ；D到\(BC\)\allowbreak 的垂线则是三角形\(BCD\)\allowbreak 以\(BC\)\allowbreak 为底的高。原标\(DB=1.8\)\allowbreak 厘米，却小于D到\(BC\)\allowbreak 的垂线段1.9厘米，违反垂线段最短。因此只按原直角标记识别所填底高，冲突尺寸应核实，不再据此推算。\end{multicols}\end{snippet}
% p046_s02
\begin{snippet}{505}{11.5}{14}\begin{multicols}{3}\raggedcolumns 连续除以25再除以4，可以合成除以二者的积，所以D符合除法性质。A、B把乘和除混淆了，去括号后不能把本来应继续除的量改成乘。C分配时漏掉了8与25相乘，两个加数都应受到同样的乘法作用。\end{multicols}\end{snippet}
% p046_s03
\begin{snippet}{505}{11.5}{14}\begin{multicols}{3}\raggedcolumns 每扇窗4块玻璃，每层18扇，共25层，按“每份数\(\times\)\allowbreak 份数”先列\(4\times18\times25\)\allowbreak 。选项B的\(25\times18\times4\)\allowbreak 只是交换因数位置，积不变，所以选B。计算时用乘法交换律和结合律，把4与25凑成100：\[\begin{aligned}&4\times18\times25\\={}&(4\times25)\times18\\={}&100\times18\\={}&1800\text{(块)}\end{aligned}\]\end{multicols}\end{snippet}
% p046_s04
\begin{snippet}{505}{11.5}{14}\begin{multicols}{3}\raggedcolumns 先挑出每组的两条短边，与最长边比较。A的3与4之和恰好7，只能成直线；B、D的两短边之和还不够最长边，都不能围成。C的两条6厘米边合起来超过10厘米，满足围成三角形的条件。\par 具体比较为\(3+4=7\)\allowbreak 、\(2+3=5<6\)\allowbreak 、\(6+6=12>10\)\allowbreak 、\(2+2=4<5\)\allowbreak 。\end{multicols}\end{snippet}
% p046_s05
\begin{snippet}{505}{11.5}{14}\begin{multicols}{3}\raggedcolumns 左边先把两个图形代表的数相加，再乘34。右边把每个加数都分别乘34，然后将两个积相加，算的是相同总量。由“和乘一个数”变成“分别相乘再相加”，正是乘法分配律。\end{multicols}\end{snippet}
% p046_s06
\begin{snippet}{505}{11.5}{14}\begin{multicols}{3}\raggedcolumns 原梯形只有一组对边平行。若沿同一条切线分成两个平行四边形，两块各自另一组对边也要平行，拼合后的外边界就会具有两组平行对边，与梯形条件不符。其他选项可分别通过斜切、连对角线或平行一腰切割得到，因此只排除这一类。\end{multicols}\end{snippet}
% p046_s07
\begin{snippet}{505}{11.5}{14}\begin{multicols}{3}\raggedcolumns ①12乘9可按12乘10再减一份12；②152乘4可分成150与2分别乘4；③600除以25可把被除数、除数同时乘4，化成除以100。④先算括号；⑤先求23与6的积再从397中减；⑥125与8凑1000；⑦24乘5可先乘10再减半。⑧连续除法从左到右，先除以80再除以6；⑨把49、31分别看成50、30；⑩把61、95分别看成60、100。后两题求的是估计值，要写约等号，不能与精确计算混同。\end{multicols}\end{snippet}
% p046_s08
\begin{snippet}{505}{11.5}{14}\begin{multicols}{3}\raggedcolumns \textbf{三、2（1）：}\ 先算括号中的除法，确定360平均分成18份后每份多少。再用12乘所得商，最后才从480中扣除这个积。式子按既定的括号与运算顺序计算即可，不因结果是整数就称为简算。\end{multicols}\end{snippet}
% p046_s09
\begin{snippet}{505}{11.5}{14}\begin{multicols}{3}\raggedcolumns \textbf{三、2（2）：}\ 先处理最内层小括号，把63与177相加凑成240。再计算中括号内的除法，得到一个完整因数，最后才与45相乘。每去掉一层括号，都应先完成这一层指定的运算。\end{multicols}\end{snippet}
% p046_s09_local1
\begin{snippet}{87.07808837890627}{10}{12.8}\SourceSans\setlength{\mathindent}{2bp}\(\displaystyle=45\times[240\div24]\)\end{snippet}
% p046_s10
\begin{snippet}{505}{11.5}{14}\begin{multicols}{3}\raggedcolumns \textbf{三、2（3）：}\ 末尾减去的136就是减去1份136，前面共有101份。提出共同的136，把份数101减1，正好凑成100份。用的是乘法分配律逆用，能避免先算较大的积再作减法。\end{multicols}\end{snippet}
% p046_s10_local1
\begin{snippet}{157.0}{10.5}{12.8}\SourceSans\setlength{\mathindent}{2bp}\(\begin{aligned}={}&136\times(101-1)\\={}&136\times100\\={}&13600\end{aligned}\)\end{snippet}
% p046_s11
\begin{snippet}{505}{11.5}{14}\begin{multicols}{3}\raggedcolumns \textbf{三、2（4）：}\ 主解把48拆成8与6的积，让125先与8凑成1000，再乘6，依据是乘法结合律。另一拆法把48看作50少2，先算50个125，再扣掉多算的2个125，用的是乘法分配律。两条路径都利用具体的凑整结构，而不只是按原顺序求积。\par \textbf{另一拆法：用乘法分配律}\par \[\begin{aligned}={}&125\times(50-2)\\={}&125\times50-125\times2\\={}&6250-250\\={}&6000\end{aligned}\]\end{multicols}\end{snippet}
% p047_s01
\begin{snippet}{505}{11.5}{14}\begin{multicols}{3}\raggedcolumns 从点\(A\)\allowbreak 到公路有许多条连线，其中与公路边缘垂直的线段最短。用三角尺一条直角边靠住公路近侧边缘，平移使另一条边经过\(A\)\allowbreak ，再画到垂足并标直角。终点应落在近侧公路边缘，不能继续穿到中心线，否则多出一段路程。\end{multicols}\end{snippet}
% p047_s02
\begin{snippet}{505}{11.5}{14}\begin{multicols}{3}\raggedcolumns 圆规开口取给定\(AB\)\allowbreak 的长度，以A、B分别为圆心画弧，过程中开口保持不变。两弧交点到A、B的距离都等于\(AB\)\allowbreak ，连接后便有三条等边。保留作图弧，若两弧尚未相交应延长弧线，不能随意改变半径。\end{multicols}\end{snippet}
% p047_s03
\begin{snippet}{505}{11.5}{14}\begin{multicols}{3}\raggedcolumns 方法一先看\(138^\circ\)\allowbreak 与\(\angle2\)\allowbreak 合成平角，求出\(\angle2\)\allowbreak 后，再用三角形内角和扣去两个已知角。方法二把同为\(180^\circ\)\allowbreak 的平角与三角形内角和比较，扣去共同的\(\angle2\)\allowbreak ，便知\(138^\circ\)\allowbreak 是\(\angle1\)\allowbreak 与\(61^\circ\)\allowbreak 的和。因此可直接从这个和中减去\(61^\circ\)\allowbreak ，求出未知角。\par 比较底边平角与三角形内角和，减去共同的\(\angle2\)\allowbreak ，所求角与\(61^\circ\)\allowbreak 的和就是\(138^\circ\)\allowbreak 。\[\begin{aligned}&138^\circ-61^\circ\\={}&77^\circ\end{aligned}\]\end{multicols}\end{snippet}
% p047_s04
\begin{snippet}{505}{11.5}{14}\begin{multicols}{3}\raggedcolumns 答：相距3300千米。\par 图中以机场为共同起点，两架飞机向相反方向飞，相距路程等于各自飞行路程的和。方法一先合计一小时两机共飞的路程，再乘共同飞行的2小时；方法二分别按“速度乘时间”求路程后相加。不能求路程差，因为两机不在机场同一侧。\par \textbf{另法：分别求路程，再相加}\par \[\begin{aligned}&800\times2+850\times2\\={}&1600+1700\\={}&3300\text{(千米)}\end{aligned}\]\end{multicols}\end{snippet}
% p047_s04_local1
\begin{snippet}{80}{9.5}{14}\SourceSans\setlength{\mathindent}{2bp}\(\text{(千米)}\)\end{snippet}
% p048_s01
\begin{snippet}{505}{11.5}{14}\begin{multicols}{3}\raggedcolumns 题目把100克看成一份，先求200克中有多少份100克。每份约含62毫克，按每份含量乘份数求出摄入量，再与“超过100毫克”比较。这里质量单位克与含量单位毫克各有含义，不需要把200克直接化为毫克来相乘。\end{multicols}\end{snippet}
% p048_s01_local1
\begin{snippet}{80}{9.5}{14}\SourceSans\setlength{\mathindent}{2bp}\(\text{(毫克)}\)\end{snippet}
% p048_s02
\begin{snippet}{505}{11.5}{14}\begin{multicols}{3}\raggedcolumns 三个社团的报名人数全部相加。用加法交换律和结合律，先把128与172凑成300，再加275。交换、结合加数都不改变总人数。\end{multicols}\end{snippet}
% p048_s02_local1
\begin{snippet}{80}{9.5}{14}\SourceSans\setlength{\mathindent}{2bp}\(\text{(人)}\)\end{snippet}
% p048_s03
\begin{snippet}{505}{11.5}{14}\begin{multicols}{3}\raggedcolumns “编程比另外两个社团的总和少”，要以航模和书法的合计为较大数。方法一先把这两社团人数相加，再减编程人数；另一方法先从编程人数中扣去书法人数，剩下部分再与航模人数比较。后一方法相当于先给比较双方都扣去同一个书法人数，差保持不变。\par \textbf{另法：先求编程比书法多多少}\par \[\begin{aligned}&128-(275-172)\\={}&128-103\\={}&25\text{(人)}\end{aligned}\]\end{multicols}\end{snippet}
% p048_s03_local1
\begin{snippet}{80}{9.5}{14}\SourceSans\setlength{\mathindent}{2bp}\(\text{(人)}\)\end{snippet}
% p048_s04
\begin{snippet}{505}{11.5}{14}\begin{multicols}{3}\raggedcolumns 国庆节在10月，先按表选择旺季票价，再区分90名学生与20位老师。主解按两类各自的“每人票价乘人数”求费用，合计后从3000元中扣除；也可按两次购票顺序连续扣款。另一方法先让110人都按儿童价计费，再给20位老师补上每人10元的差价，补的只属于老师。总费用求准后，剩余金额仍用原有钱数减实际支出。\par \textbf{另法：全按儿童票，再补老师差价}\par \end{multicols}\[\begin{aligned}&3000-20\times(90+20)-(30-20)\times20\\={}&3000-20\times110-10\times20\\={}&3000-2200-200\\={}&800-200\\={}&600\text{(元)}\end{aligned}\]\end{snippet}
% p048_s04_local1
\begin{snippet}{264}{10.5}{12.8}\SourceSans\setlength{\mathindent}{2bp}国庆节按旺季票价。\par \[\begin{aligned}&3000-(20\times90+30\times20)\\={}&3000-(1800+600)\\={}&3000-2400\\={}&600\text{(元)}\end{aligned}\]答：还剩600元。\end{snippet}
% p049_s01
\begin{snippet}{505}{11.5}{14}\begin{multicols}{3}\raggedcolumns 读数时把整数部分与小数部分分开：30按整数读法读，小数点后6、4逐个读。计数单位看最后写到的数位，本数末位是百分位，所以每个单位是\(0.01\)\allowbreak 。再把十、十分位上的数量也换成百分之一，合计就是共有多少个这样的单位。\par 3个十相当于3000个百分之一，6个十分之一相当于60个百分之一，合计\(3000+60+4=3064\)\allowbreak 个。\end{multicols}\end{snippet}
% p049_s02
\begin{snippet}{505}{11.5}{14}\begin{multicols}{3}\raggedcolumns 8个十分之一先换成80个百分之一，与另7个百分之一合起来就是87个百分之一，组成0.87。1可以看成100个百分之一，用总单位个数减去已有单位个数，便知道还差多少个。第二空要求的是百分之一的个数，不是直接填差的小数。\par 统一按百分之一计数，8个十分之一就是80个百分之一，先算\(80+7=87\)\allowbreak 个。1有100个百分之一，还差\(100-87=13\)\allowbreak 个。\end{multicols}\end{snippet}
% p049_s03
\begin{snippet}{505}{11.5}{14}\begin{multicols}{3}\raggedcolumns 改成“万”按10000合成1份，改成“亿”按100000000合成1份，准确改写不改变大小：\par\noindent\end{multicols}\[\displaystyle \begin{aligned}168000&=(168000\div10000)\text{万}=16.8\text{万}\\995000000&=(995000000\div100000000)\text{亿}=9.95\text{亿}\end{aligned}\]\begin{multicols}{3}\raggedcolumns \par 再按指定数位取近似数：\par\noindent\[\displaystyle \begin{aligned}16.8\text{万}&\approx17\text{万}\\9.95\text{亿}&\approx10.0\text{亿}\end{aligned}\]\par 改写用等号，取近似数才用约等号；10.0的末尾0不可省。\end{multicols}\end{snippet}
% p049_s04
\begin{snippet}{505}{11.5}{14}\begin{multicols}{3}\raggedcolumns 先看所换的量及单位。长度按10，质量按1000，面积按100的相应进率换算；复名数保留原整数部分，只换小单位部分：\par\noindent\end{multicols}\[\displaystyle \begin{aligned}3\text{分米}&=(3\div10)\text{米}=0.3\text{米}\\0.06\text{吨}&=(0.06\times1000)\text{千克}=60\text{千克}\\5.06\text{吨}&=5\text{吨}60\text{千克}\\32\text{克}&=(32\div1000)\text{千克}=0.032\text{千克}\\3\text{平方分米}&=(3\div100)\text{平方米}=0.03\text{平方米}\\5\text{平方米}3\text{平方分米}&=(5+0.03)\text{平方米}=5.03\text{平方米}\end{aligned}\]\begin{multicols}{3}\raggedcolumns \par 原空后已有单位，不重复填写单位。\end{multicols}\end{snippet}
% p049_s05
\begin{snippet}{505}{11.5}{14}\begin{multicols}{3}\raggedcolumns \textbf{一、5：}\ 这些数的整数部分都是0，先比十分位；十分位相同再比百分位，仍相同才比千分位。可在小数末尾补0统一写成三位小数，这不会改变数值，却方便对齐比较。先排十分位为0的两数，再排十分位为4、5的两组，组内依次逐位比较。\end{multicols}\end{snippet}
% p049_s06
\begin{snippet}{505}{11.5}{14}\begin{multicols}{3}\raggedcolumns \textbf{一、6：}\ 先保留84002的数字顺序，再把小数点放在不同位置试读。整数部分按整数规则读零，小数部分却要逐个读出每一个0；由此排除读两个零或一个零也不读的写法。保留的写法应读作“八百四十点零二”，恰好读一个零。\end{multicols}\end{snippet}
% p049_s07
\begin{snippet}{505}{11.5}{14}\begin{multicols}{3}\raggedcolumns ①把6.02末尾补0后比百分位；②小数末尾添0不改变大小。③与⑤先统一单位：\par\noindent\end{multicols}\[\displaystyle \begin{aligned}50\text{克}&=(50\div1000)\text{千克}=0.05\text{千克}\\9\text{千克}50\text{克}&=(9+0.05)\text{千克}=9.05\text{千克}\\0.57\text{米}&=(0.57\times100)\text{厘米}=57\text{厘米}\end{aligned}\]\begin{multicols}{3}\raggedcolumns \par 因此9.05千克小于9.5千克，而0.57米与57厘米相等。④比较百分位1与7；⑥比较整数部分8与4。\end{multicols}\end{snippet}
% p050_s01
\begin{snippet}{505}{11.5}{14}\begin{multicols}{3}\raggedcolumns 近似数写到百分位，所以原三位小数要看千分位来取舍。求最大值，从7.20后面添最大的可舍去数字4；再添5就会进到7.21。求最小值，从7.19后面添最小的能进位数字5；若添4便只能约成7.19。两个端点分别从“舍去”和“进位”检查，才能把范围找全。\par 两侧再多取一位检查：\(7.194\approx7.19\)\allowbreak 、\(7.205\approx7.21\)\allowbreak ，都不能约成7.20，说明所选两端已到边界。\end{multicols}\end{snippet}
% p050_s02
\begin{snippet}{505}{11.5}{14}\begin{multicols}{3}\raggedcolumns 先辨认每幅图的一个单位有多大，再数有多少份。A每颗珠是百分之一，B每小格是百分之一，C把0.1分成5份，每间隔是0.02，因此三者都能表示六个百分之一。D却把整个1只分10份，涂6份表示六个十分之一，比六个百分之一大10倍。\par C把10个百分之一平均分5份，每格有\(10\div5=2\)\allowbreak 个百分之一；点在第3格处，表示\(2\times3=6\)\allowbreak 个百分之一。\end{multicols}\end{snippet}
% p050_s03
\begin{snippet}{505}{11.5}{14}\begin{multicols}{3}\raggedcolumns 小数末尾添0或去0，只改变写法，不改变大小。B的0都接在小数末尾，可直接去掉；其他选项的0位于整数部分或小数中间。把这些0删掉会使后面的数字占到不同数位，因此数值改变。\end{multicols}\end{snippet}
% p050_s04
\begin{snippet}{505}{11.5}{14}\begin{multicols}{3}\raggedcolumns “最接近”要比较各数到9的距离，在9两边的数都要检查。可以先把四个数与9统一按千分之一计数，再比较相差的单位个数。差的个数越少，距离就越短，不能只看原数比9大还是小。\par 统一按千分之一计数，9有9000个，四个数依次有8900、9010、9001、8890个。距离分别用\(9000-8900=100\)\allowbreak 、\(9010-9000=10\)\allowbreak 、\(9001-9000=1\)\allowbreak 、\(9000-8890=110\)\allowbreak 个千分之一表示，最小的只有1个。\end{multicols}\end{snippet}
% p050_s05
\begin{snippet}{505}{11.5}{14}\begin{multicols}{3}\raggedcolumns 先确定中间落点在1.24米与1.53米之间，再看它更靠近右端1.53米。若方框填3，1.33到左端更近，与图不符；填4得到的数则离右端更近。只判断在两数之间还不能唯一选出数字，必须结合图示的相对距离。\par 按百分之一米计距离：若填3，两端距离为\(133-124=9\)\allowbreak 、\(153-133=20\)\allowbreak ；若填4，则为\(143-124=19\)\allowbreak 、\(153-143=10\)\allowbreak 。后一种更靠近右端，与图相符。\end{multicols}\end{snippet}
% p050_s06
\begin{snippet}{505}{11.5}{14}\begin{multicols}{3}\raggedcolumns 先把0.6写成0.60，与两位小数按位比较。若十分位填0到5，先在十分位就比6小；若填6，十分位相同，但百分位8比0大，已不符合。更大的7、8、9也都不能取，因此可取数字要从0开始数。\par 可填\(0,\allowbreak 1,\allowbreak 2,\allowbreak 3,\allowbreak 4,\allowbreak 5\)\allowbreak ，从0数到5共有\(5-0+1=6\)\allowbreak 个。\end{multicols}\end{snippet}
% p050_s07
\begin{snippet}{505}{11.5}{14}\begin{multicols}{3}\raggedcolumns 四张数字卡各用一次，整数部分只能选2，其他数字开头都不可能四舍五入成3。要由2进到3，十分位至少为5，所以只能选6或9。确定这两位后，余下两个数字各有两种交换次序，按两类各两种列全；写在末尾的0也要保留，才能符合三位小数的要求。\par 两类各有两种，\(2\times2=4\)\allowbreak 个；列全为\(2.609,\allowbreak 2.690,\allowbreak 2.906,\allowbreak 2.960\)\allowbreak 。\end{multicols}\end{snippet}
% p050_s08
\begin{snippet}{505}{11.5}{14}\begin{multicols}{3}\raggedcolumns 先扩大100倍再把小数点左移一位，后一变化相当于缩小到原来的十分之一。两次合起来相当于扩大10倍，因此反求原数要把39.25除以10。也可以倒着逐步还原：先右移一位撤销最后一步，再除以100撤销最初扩大。\par 综合变化倍数为\(100\div10=10\)\allowbreak ，所以原数为\(39.25\div10=3.925\)\allowbreak 。\end{multicols}\end{snippet}
% p051_s01
\begin{snippet}{505}{11.5}{14}\begin{multicols}{3}\raggedcolumns ①乘10右移小数点一位；②乘1000右移三位，原小数位不足时在末尾补0；③除以10左移一位。④除以100左移两位；⑤先右移两位再左移一位，合起来右移一位；⑥先左移两位再右移三位，也合起来右移一位。小数点移动时数字的先后次序保持不变，整数部分不足一位要补0占位。\end{multicols}\end{snippet}
% p051_s02
\begin{snippet}{505}{11.5}{14}\begin{multicols}{3}\raggedcolumns 按题面顺序，前两式分别缩小到十分之一、扩大100倍，确定除以10与乘100；第三式已给除以10，要反向扩大10倍求被除数。第四式扩大10倍；第五式可把未知数除以10求6.7，也可把未知数乘10求6.7，两种都符合题意。第六式已给除以1000，反向把0.735扩大1000倍即可。题面未限定方框必须为整数，所以第5式还可填\(0.67\times10=6.7\)\allowbreak 。\end{multicols}\end{snippet}
% p051_s03
\begin{snippet}{505}{11.5}{14}\begin{multicols}{3}\raggedcolumns 2.5与10.08的下界分别取整数部分，上界各多1。6.94与35.618同样先找夹住它的两个相邻整数，但题面用大于号，应把较大整数写左边、较小整数写右边。99.2与275.2用小于号，按从小到大把两个相邻整数放在两端。检查两端必须相差1，且中间的小数确实落在它们之间。\end{multicols}\end{snippet}
% p051_s04
\begin{snippet}{505}{11.5}{14}\begin{multicols}{3}\raggedcolumns 0.7按十分之一计数，在十等份圆里涂7份；0.39按百分之一计数，在百格里涂39格。1.43先涂一个完整的1，余下0.43再涂百格中的43格。两个用来表示1的整体大小必须一致，否则整块与小格的份量就不对应。\end{multicols}\end{snippet}
% p051_s05
\begin{snippet}{505}{11.5}{14}\begin{multicols}{3}\raggedcolumns 先认出相邻两个整数之间分成10格，每格0.1，4.4、4.7分别从4起数4、7格，5.0就在5，5.4从5起数4格。5.65需要继续细分：在5.6与5.7之间再看百分位，位于这一格正中。点要落实在数轴上，再在旁边标数，不能把5.65移到5.6或5.7的整刻度上。\end{multicols}\end{snippet}
% p051_s06
\begin{snippet}{505}{11.5}{14}\begin{multicols}{3}\raggedcolumns 48元是10千克西瓜的总价，要求每1千克多少钱，就要把总价平均分成10份。除以10时，小数点向左移动一位，所得数的单位仍为元。再把每千克价格乘10千克，应回到原来的48元，可用来核验。\end{multicols}\end{snippet}
% p051_s06_local1
\begin{snippet}{80}{9.5}{14}\SourceSans\setlength{\mathindent}{2bp}\(\text{(元)}\)\end{snippet}
% p052_s01
\begin{snippet}{505}{11.5}{14}\begin{multicols}{3}\raggedcolumns 先把总质量化成克，与“每100克”的单位统一：\par\noindent\(\displaystyle 10\text{千克}=10000\text{克}\)\allowbreak \par 方法一求总西瓜质量包含多少份100克，再用每份出汁81.5克乘份数；方法二先求1千克含有几份100克，算出1千克的出汁量，再乘10份。第二法在计算时可先把两个10结合成100，一次移动小数点；两种方法的中间结果都是克，最后还要除以1000化成千克；计算只需乘10、100与单位换算。\par \textbf{另法：先求每千克出汁量，再乘总质量}\par \noindent\(\displaystyle 1\text{千克}=1000\text{克}\)\allowbreak \par \noindent\[\displaystyle \begin{aligned}&81.5\times(1000\div100)\times10\\={}&81.5\times10\times10\\={}&81.5\times(10\times10)\\={}&81.5\times100\\={}&8150\text{(克)}\end{aligned}\]\par \noindent\(\displaystyle 8150\text{克}=(8150\div1000)\text{千克}=8.15\text{千克}\)\allowbreak \par\end{multicols}\end{snippet}
% p052_s01_local1
\begin{snippet}{80}{9.0}{14}\SourceSans\setlength{\mathindent}{2bp}\(\text{(克)}\)\end{snippet}
% p052_s01_local2
\begin{snippet}{505}{9.0}{14}\SourceSans\setlength{\mathindent}{2bp}\(8150\text{克}=(8150\div1000)\text{千克}=8.15\text{千克}\)\end{snippet}
% p052_s02
\begin{snippet}{505}{11.5}{14}\begin{multicols}{3}\raggedcolumns 老师、亮亮买成人票，东东半票，明明免票。\par 先把厘米统一化成米，各自除以100：\par\noindent\[\displaystyle \begin{aligned}135\text{厘米}&=1.35\text{米}\\119\text{厘米}&=1.19\text{米}\\158\text{厘米}&=1.58\text{米}\end{aligned}\]\par 再逐人对照票价区间：东东半票、明明免票，老师与亮亮成人票。方法一分别计算成人票和半票后相加，免票者不计钱。方法二把一张半票的钱看成1份，每张成人票就是2份，两张成人票与一张半票共5份，再按每份钱乘份数。\par \textbf{另法：把半票钱看成一份}\par \noindent\[\displaystyle \begin{aligned}&(132\div2)\times(2\times2+1)\\={}&66\times(4+1)\\={}&66\times5\\={}&330\text{(元)}\end{aligned}\]\par\end{multicols}\end{snippet}
% p052_s02_local1
\begin{snippet}{80}{9.0}{14}\SourceSans\setlength{\mathindent}{2bp}\(\text{(元)}\)\end{snippet}
% p052_s03
\begin{snippet}{505}{11.5}{14}\begin{multicols}{3}\raggedcolumns 方法一先求一块瓷砖的面积，再按每行18块、共5行求总块数，用每块面积乘块数。方法二先用一块的长乘每行块数、一块的宽乘行数，分别求整面墙的长和宽，再求长方形面积。两条路径都先使用厘米作整数计算，所得面积是平方厘米；换成平方米必须按面积进率10000去除。\par \textbf{另法：先求墙的长、宽}\par\noindent\[\displaystyle \begin{aligned}&(40\times18)\times(30\times5)\\={}&720\times150\\={}&108000\text{(平方厘米)}\end{aligned}\]\par \noindent\end{multicols}\[\displaystyle \begin{aligned}1\text{平方米}&=10000\text{平方厘米}\\108000\text{平方厘米}&=(108000\div10000)\text{平方米}\\&=10.8\text{平方米}\end{aligned}\]\begin{multicols}{3}\raggedcolumns \par\end{multicols}\end{snippet}
% p052_s03_local1
\begin{snippet}{80}{9.0}{14}\SourceSans\setlength{\mathindent}{2bp}\(\text{(平方厘米)}\)\end{snippet}
% p052_s03_local2
\begin{snippet}{505}{9.0}{14}\SourceSans\setlength{\mathindent}{2bp}\(108000\text{平方厘米}=(108000\div10000)\text{平方米}=10.8\text{平方米}\)\end{snippet}
% p052_s04
\begin{snippet}{505}{11.5}{14}\begin{multicols}{3}\raggedcolumns 三张数字卡各用一次，一位小数表示小数点后只能放一张卡，另外两张组成十位和个位。按十位分别为1、3、6分类，每类再交换剩余两张卡，就能不重不漏地列全。比较时先排十位，同十位再排个位，只有前面都相同才比较十分位。\end{multicols}\end{snippet}
% p052_s04_local1
\begin{snippet}{80}{9.0}{14}\SourceSans\setlength{\mathindent}{2bp}\(\text{(个)}\)\end{snippet}
% p053_s01
\begin{snippet}{505}{11.5}{14}\begin{multicols}{3}\raggedcolumns 先把1写成一位小数1.0，这样与要求的小数位数一致。相邻一位小数之间都相差0.1，向前少一个0.1，向后多一个0.1。不能把“相邻”理解成相邻整数。\par 按十分之一计数，1就是10个这样的单位，前后各取一个单位：\(10-1=9\)\allowbreak 、\(10+1=11\)\allowbreak ，分别组成0.9与1.1。\end{multicols}\end{snippet}
% p053_s02
\begin{snippet}{505}{11.5}{14}\begin{multicols}{3}\raggedcolumns 0.618最后写到千分位，所以以千分之一为一个计数单位，三个小数位合起来可看成618个这样的单位。改成以万分之一计数时，要在小数末尾添0，写到第四位小数。添0不改变大小，却把每个计数单位细分成原来的十分之一，所以单位个数也相应变为原来的10倍。\par 千分之一再细分成10个万分之一，因此单位个数变为\(618\times10=6180\)\allowbreak 。\end{multicols}\end{snippet}
% p053_s03
\begin{snippet}{505}{11.5}{14}\begin{multicols}{3}\raggedcolumns 计数器先看珠子在哪一数位，百分位4颗珠表示4个百分之一。圆图先看整体平均分成10份，涂4份表示4个十分之一。两者要统一成同一种计数单位，再比较单位个数，不能只看数字都是4就认为相等。\par 都按百分之一计数，圆的4个十分之一就是40个百分之一，计数器只有4个百分之一。相差\(40-4=36\)\allowbreak 个百分之一，即0.36。\end{multicols}\end{snippet}
% p053_s04
\begin{snippet}{505}{11.5}{14}\begin{multicols}{3}\raggedcolumns 计数单位取决于当前写法的最后一个数位，三位小数写到千分位，两位小数写到百分位。比较大小则看各数位代表的实际量，小数末尾添0并没有增加数值。因此“计数单位相同”和“数值相同”要按不同依据分别配对。\end{multicols}\end{snippet}
% p053_s05
\begin{snippet}{505}{11.5}{14}\begin{multicols}{3}\raggedcolumns 先按数位安放数字：十位放3，十分位放8，千分位放5，没有给出的个位与百分位用0占位。读整数部分时，末尾0不读；小数部分则逐个数字读，百分位的0要读出。数一数哪些0按这两种规则省读，就能求第二空。\end{multicols}\end{snippet}
% p053_s06
\begin{snippet}{505}{11.5}{14}\begin{multicols}{3}\raggedcolumns 甲原来就是小数，末尾添两个0只改变写法；乙原来是整数，末尾添两个0使它扩大100倍。先比较原来38与3.8相差的倍数，再把乙新增的100倍计入。两个数都添0，并不表示它们同时按相同比例变化。\par 原有倍数可由\(3.8\times10=38\)\allowbreak 看出；乙再扩大100倍后，相比甲的倍数变为\(10\times100=1000\)\allowbreak 。\end{multicols}\end{snippet}
% p053_s07
\begin{snippet}{505}{11.5}{14}\begin{multicols}{3}\raggedcolumns 求最小数要从高位往低位逐位选最小的可行数字。个位先取0，十分位就不能再取0，最小取1；百分位再取与前面不同的最小数字2。最后由数字和9补出千分位，并检查它与前三位均不同，所以没有漏掉更小且合条件的选择。\par 千分位由\(9-0-1-2=6\)\allowbreak 确定；数字0、1、2、6互不相同，数位条件与数字和条件都满足。\end{multicols}\end{snippet}
% p053_s08
\begin{snippet}{505}{11.5}{14}\begin{multicols}{3}\raggedcolumns 一组答案为“厘米、米；克、千克”。前式进率100，后式进率1000：\par\noindent\end{multicols}\[\displaystyle \begin{aligned}45\text{厘米}&=(45\div100)\text{米}=0.45\text{米}\\3540\text{克}&=(3540\div1000)\text{千克}=3.540\text{千克}\end{aligned}\]\begin{multicols}{3}\raggedcolumns \par 题目未限定单位种类，原有其他合理组合也分别列出：\par\noindent\[\displaystyle \begin{aligned}45\text{分}&=0.45\text{元}\\45\text{平方厘米}&=0.45\text{平方分米}\\45\text{平方分米}&=0.45\text{平方米}\\45\text{公顷}&=0.45\text{平方千米}\\3540\text{毫米}&=3.540\text{米}\\3540\text{米}&=3.540\text{千米}\\3540\text{千克}&=3.540\text{吨}\end{aligned}\]\par 保留题面给定的3.540写法。两边必须是同一种量，不能把长度与质量混配。这里只在原单位空填写单位，不重复数值。\end{multicols}\end{snippet}
% p053_s09
\begin{snippet}{505}{11.5}{14}\begin{multicols}{3}\raggedcolumns 先把1吨化成1000千克，与已知的10千克采用同一单位。求出其中有多少份10千克，再用每份出油4.5千克乘份数，得到总出油千克数。最后按1000千克为1吨再换算，第二空不能直接抄第一空的数。\par \noindent\(\displaystyle 1\text{吨}=1000\text{千克}\)\allowbreak \par \noindent\[\displaystyle \begin{aligned}&4.5\times(1000\div10)\\={}&4.5\times100\\={}&450\text{(千克)}\end{aligned}\]\par \noindent\(\displaystyle 450\text{千克}=(450\div1000)\text{吨}=0.45\text{吨}\)\allowbreak \par\end{multicols}\end{snippet}
% p054_s01
\begin{snippet}{505}{11.5}{14}\begin{multicols}{3}\raggedcolumns 每次都先看\(A\)\allowbreak 与\(B\)\allowbreak 之间10个间隔表示多少量，再除以10求一小格。第一种\(A,\allowbreak B\)\allowbreak 相差1，每格0.1，C从A向右数7格；第二种相差0.1，每格0.01，D从B向右数2格。第三种仍相差0.1，C从A向右数7个0.01。字母点位置没有变，但已知数值改变后，每格代表的大小也会改变。\par 第一种每格\(1\div10=0.1\)\allowbreak ，从1向右数7格组成1.7。后两种每格\(0.1\div10=0.01\)\allowbreak ，从2.4再数2个百分之一组成2.42，从48.6再数7个百分之一组成48.67。\end{multicols}\end{snippet}
% p054_s02
\begin{snippet}{505}{11.5}{14}\begin{multicols}{3}\raggedcolumns 先判断数字9在十分位，表示9个十分之一米。按米化分米的进率10换算：\par\noindent\(\displaystyle 0.9\text{米}=(0.9\times10)\text{分米}=9\text{分米}\)\allowbreak \par 所以选9分米，不能只读数字而不看数位与单位。\end{multicols}\end{snippet}
% p054_s03
\begin{snippet}{505}{11.5}{14}\begin{multicols}{3}\raggedcolumns A保留两位小数时要看再后一位，即千分位；B从3.45起就会四舍五入到3.5，并非最小3.50。C可比较0.99与1.0，位数更多的小数反而较小，因此位数不能直接决定大小。D中的整数部分可以是0，也仍然存在，故只有这一项成立。\par B的反例可写为\(3.45\approx3.5\)\allowbreak （保留一位小数），C的反例为\(0.99<1.0\)\allowbreak 。\end{multicols}\end{snippet}
% p054_s04
\begin{snippet}{505}{11.5}{14}\begin{multicols}{3}\raggedcolumns 只解读题中已经给出的\(6\times0.001\)\allowbreak ：它表示6个千分之一。6回答“有几个”，是新的计数单位的个数；0.001回答“每一个是多少”，是新的计数单位。不要把个数与单位互换，也无需另用尚未学习的一般小数乘法计算方法。\end{multicols}\end{snippet}
% p054_s05
\begin{snippet}{505}{11.5}{14}\begin{multicols}{3}\raggedcolumns 短跑用时越少名次越靠前，所以浩浩必须比另外两位用时都少。明明比亮亮慢，只需先保证浩浩小于15.47秒；十分位填4时，再比较百分位6与7，仍能胜出。若十分位再增到5，就比15.47大，因此4已经是可填的最大值。\par 试最大的可行数字：\(15.46<15.47\)\allowbreak ；再增加一位到5，\(15.56>15.47\)\allowbreak ，便不能得第一。\end{multicols}\end{snippet}
% p054_s06
\begin{snippet}{505}{11.5}{14}\begin{multicols}{3}\raggedcolumns 厘米化米除以100，原有8米不变，只把6厘米换算后合并：\par\noindent\end{multicols}\[\displaystyle \begin{aligned}6\text{厘米}&=(6\div100)\text{米}=0.06\text{米}\\8\text{米}6\text{厘米}&=(8+0.06)\text{米}=8.06\text{米}\\560\text{厘米}&=(560\div100)\text{米}=5.6\text{米}\end{aligned}\]\begin{multicols}{3}\raggedcolumns \par 统一成米后相加，正确算式为\(8.06+5.6\)\allowbreak 。\end{multicols}\end{snippet}
% p054_s07
\begin{snippet}{505}{11.5}{14}\begin{multicols}{3}\raggedcolumns ①除以100左移小数点两位；②乘1000右移三位；③除以10左移一位；④先合并变化倍数，\(100\div10=10\)\allowbreak ，所以直接算\(45\times10=450\)\allowbreak ，小数点只右移一位。⑤除以10左移一位，十分位不足时补0；⑥除以1000左移三位；⑦除以100左移两位；⑧先合并变化倍数，\(1000\div10=100\)\allowbreak ，直接算\(3.2\times100=320\)\allowbreak ，小数点只右移两位。小数点移动只改变各位数字表示的大小，数字顺序不能改变。\end{multicols}\end{snippet}
% p055_s01
\begin{snippet}{505}{11.5}{14}\begin{multicols}{3}\raggedcolumns 精确到百分位，保留末位是百分位，应直接看原数的千分位。前一数千分位3小于5，舍去后面各位；后一数千分位7大于5，百分位进1。不能先取三位小数，再连续四舍五入，否则可能改变结果。\end{multicols}\end{snippet}
% p055_s02
\begin{snippet}{505}{11.5}{14}\begin{multicols}{3}\raggedcolumns “万”与“亿”分别是10000和100000000，改写时将小数点相应左移四位、八位，数的大小不变，所以用等号。保留一位小数时再看百分位，第二个数要连续进位。结果中的\(255.0\)\allowbreak 仍须保留末尾0，表示保留到十分位，不能省成255。\par\noindent\[\displaystyle \begin{aligned}327900&=32.79\text{万}\\25499270000&=254.9927\text{亿}\end{aligned}\]\par\end{multicols}\end{snippet}
% p055_s03
\begin{snippet}{505}{11.5}{14}\begin{multicols}{3}\raggedcolumns \textbf{三、3（1）：}\ 整数部分相同，先比较十分位；右边十分位是6，左边填7、8、9都能使它更大。题目要求“最大”，就在符合条件的数字中取9。小数位数多并不表示数更大。\end{multicols}\end{snippet}
% p055_s04
\begin{snippet}{505}{11.5}{14}\begin{multicols}{3}\raggedcolumns \textbf{三、3（2）：}\ 两数整数部分相同，左边十分位是0，因此右边十分位不能填1或更大的数字。填0后，再比较百分位，右边6小于左边7，正好符合严格小于的要求。这里0既是允许填写的数字，也是唯一符合条件的数字。\end{multicols}\end{snippet}
% p055_s05
\begin{snippet}{505}{11.5}{14}\begin{multicols}{3}\raggedcolumns \textbf{三、3（3）：}\ 先按1000的进率把克化成千克：\par\noindent\(\displaystyle 3240\text{克}=(3240\div1000)\text{千克}=3.24\text{千克}\)\allowbreak \par 右边填3时，\(3.15<3.24\)\allowbreak ；若填4或更大，整数部分已超过左边，所以最大填3。\end{multicols}\end{snippet}
% p055_s06
\begin{snippet}{505}{11.5}{14}\begin{multicols}{3}\raggedcolumns \textbf{三、3（4）：}\ 先把左边化成千克：\par\noindent\(\displaystyle 6.25\text{吨}=(6.25\times1000)\text{千克}=6250\text{千克}\)\allowbreak \par 千位相同，比较百位。填2时6225小于6250；填3时6325已经超过6250，所以最大填2。\end{multicols}\end{snippet}
% p055_s07
\begin{snippet}{505}{11.5}{14}\begin{multicols}{3}\raggedcolumns \textbf{三、3（5）：}\ 先统一成米：\par\noindent\(\displaystyle 15\text{厘米}=(15\div100)\text{米}=0.15\text{米}\)\allowbreak \par 两边个位和百分位相同，十分位填大于1的数字都符合条件，再取最大的9。\end{multicols}\end{snippet}
% p055_s08
\begin{snippet}{505}{11.5}{14}\begin{multicols}{3}\raggedcolumns \textbf{三、3（6）：}\ 小数末尾添0大小不变，可以把左边看作\(21.50\)\allowbreak 。右边十分位若也填5，百分位9比0大，反而超过左边；因此十分位必须小于5。符合条件的数字中，4最大。\end{multicols}\end{snippet}
% p055_s09
\begin{snippet}{505}{11.5}{14}\begin{multicols}{3}\raggedcolumns 先看每个完整正方形都表示1，两数的整数部分都是1，所以都先涂满一个整方形。\(1.4\)\allowbreak 剩下4个十分之一，按十分格涂4份；\(1.07\)\allowbreak 剩下7个百分之一，应在百格图中涂7小格。不要把\(0.07\)\allowbreak 误涂成7条十分格，涂色边界要落在格线上。\end{multicols}\end{snippet}
% p055_s10
\begin{snippet}{505}{11.5}{14}\begin{multicols}{3}\raggedcolumns \textbf{四、2（1）：}\ 先确定刻度单位：相邻两大刻度相差\(0.1\)\allowbreak ，中间分成10小格，每小格就是\(0.01\)\allowbreak 。箭头从\(1.4\)\allowbreak 再向右数5小格，读出\(1.45\)\allowbreak 米。最后把这个数与表中各人的身高对应，不能只凭箭头位置猜姓名。\end{multicols}\end{snippet}
% p055_s11
\begin{snippet}{505}{11.5}{14}\begin{multicols}{3}\raggedcolumns \textbf{四、2（2）：}\ 沿用每小格表示\(0.01\)\allowbreak 米的刻度。乐乐的\(1.32\)\allowbreak 应从\(1.3\)\allowbreak 向右数2小格；王华的\(1.19\)\allowbreak 应在\(1.2\)\allowbreak 左边1小格。箭头或点须准确对准刻度，再在相应位置标姓名。\end{multicols}\end{snippet}
% p055_s12
\begin{snippet}{505}{11.5}{14}\begin{multicols}{3}\raggedcolumns \textbf{四、2（3）：}\ 第三名必须低于已知较高的两人，又高于第四名。可可和乐乐分别为\(1.45\)\allowbreak 、\(1.32\)\allowbreak 米，王华为\(1.19\)\allowbreak 米，因此冬冬须夹在\(1.32\)\allowbreak 与\(1.19\)\allowbreak 之间。结合百分位固定为4，十分位只能填2；填1会低于王华，填3会高于乐乐。\end{multicols}\end{snippet}
% p055_s13
\begin{snippet}{505}{11.5}{14}\begin{multicols}{3}\raggedcolumns 现在的边长只相当于原边长的一百分之一，还原时应扩大100倍，不能再缩小。先求原来的边长，再按“边长×边长”求正方形面积；边长扩大100倍并不表示面积只扩大100倍。\[\begin{aligned}&0.03\times100\\={}&3\text{(米)}\end{aligned}\]\[\begin{aligned}&3\times3\\={}&9\text{(平方米)}\end{aligned}\]答：原菜地的面积是9平方米。\end{multicols}\end{snippet}
% p056_s01
\begin{snippet}{505}{11.5}{14}\begin{multicols}{3}\raggedcolumns 题目要以千克作答，先按1000的进率把总质量换算：\par\noindent\(\displaystyle 3.75\text{吨}=(3.75\times1000)\text{千克}=3750\text{千克}\)\allowbreak \par 方法一，先用“每车件数×车辆数”求总件数，再以总质量除以总件数，得到每件质量：\par\noindent\[\displaystyle \begin{aligned}&3750\div(20\times5)\\={}&3750\div100\\={}&37.5\text{(千克)}\end{aligned}\]\par 方法二，仍按先分给5辆车、再按每车20件平均分来分析。计算时连续除以5和20，等于一次除以它们的积；5与20先凑成100，可直接用小数点移动规律计算：\par\noindent\[\displaystyle \begin{aligned}&3750\div5\div20\\={}&3750\div(5\times20)\\={}&3750\div100\\={}&37.5\text{(千克)}\end{aligned}\]\par 答：平均每件重37.5千克。\par 先把件数写清：\par\noindent\(\displaystyle 1\text{万件}=10000\text{件}\)\allowbreak \par 同样重的快递有10000件，按“每件质量×件数”求总质量，即把每件质量扩大10000倍：\par\noindent\[\displaystyle \begin{aligned}&37.5\times10000\\={}&375000\text{(千克)}\end{aligned}\]\par 答：1万件重375000千克。\end{multicols}\end{snippet}
% p056_s02
\begin{snippet}{505}{11.5}{14}\begin{multicols}{3}\raggedcolumns \textbf{五、3（1）：}\ 本问只采用每个人三轮中的最高分，不把三轮相加。先分别找出94、\(93.5\)\allowbreak 、\(89.25\)\allowbreak ，再在这三个数中比较大小。取分范围要对三人一致，不能给某人取总分、给另一人取单轮分。\end{multicols}\end{snippet}
% p056_s03
\begin{snippet}{505}{11.5}{14}\begin{multicols}{3}\raggedcolumns \textbf{五、3（2）：}\ 本问的排名依据变为三轮总分，必须把每人的三轮都计入。可以先把能凑成整数的两轮分数结合，再加第三轮，减少小数相加的负担。1号的第一、三轮可先合成105分；下面分别求三人的总分：1号\[\begin{aligned}&28.75+94+76.25\\={}&(28.75+76.25)+94\\={}&105+94\\={}&199\text{(分)}\end{aligned}\]2号\[\begin{aligned}&89+33+93.5\\={}&122+93.5\\={}&215.5\text{(分)}\end{aligned}\]3号\[\begin{aligned}&89.25+75.75+82\\={}&165+82\\={}&247\text{(分)}\end{aligned}\]三人的计分范围相同后，再比较\(247>215.5>199\)\allowbreak 。\end{multicols}\end{snippet}
% p056_s04
\begin{snippet}{505}{11.5}{14}\begin{multicols}{3}\raggedcolumns \textbf{五、3（3）：}\ 要让2号夺冠，须提出一条对三人都使用的计分规则，并用数据检查。此题可以有不同规则，例如可规定：各去掉本人的最低分，再比较余下两轮的总分。1号\[\begin{aligned}&94+76.25\\={}&170.25\text{(分)}\end{aligned}\]2号\[\begin{aligned}&89+93.5\\={}&182.5\text{(分)}\end{aligned}\]3号\[\begin{aligned}&89.25+82\\={}&171.25\text{(分)}\end{aligned}\]\(182.5>171.25>170.25\)\allowbreak ，这条规则能使2号夺冠。去掉最低分后，比较的是各自保留两轮的总分，不是各自挑不同轮数。这个例子满足2号第一，但不表示这是唯一可能的规则。\end{multicols}\end{snippet}
% p056_s05
\begin{snippet}{505}{11.5}{14}\begin{multicols}{3}\raggedcolumns A比D高且比B矮，可先排出\(B>A>D\)\allowbreak ；C又比D矮，所以完整顺序为\(B>A>D>C\)\allowbreak 。把四个身高也按从大到小排列，再按相同名次逐一配对。\(138.8\)\allowbreak 与\(138.5\)\allowbreak 整数部分相同，要看十分位，不能颠倒D与C。答：A为141厘米，B为142.6厘米，C为138.5厘米，D为138.8厘米。\end{multicols}\end{snippet}
% p057_s01
\begin{snippet}{505}{11.5}{14}\begin{multicols}{3}\raggedcolumns 只有相同计数单位的个数才能直接相加减，所以应先对齐相同数位。小数点对齐后，个位、十分位、百分位也就分别对齐了。不能按各数末尾对齐，因为一位小数的末位与两位小数的末位不在同一数位。\end{multicols}\end{snippet}
% p057_s02
\begin{snippet}{505}{11.5}{14}\begin{multicols}{3}\raggedcolumns 第一式把\(4.9\)\allowbreak 移到前面，再将\(6.3\)\allowbreak 与\(13.7\)\allowbreak 结合，依据是加法交换律和结合律，两数还能先凑成20。第二式连续减去两数，相当于一次减去它们的和，括号内应使用加号。不能把两个减号直接改成括号内的减号。\end{multicols}\end{snippet}
% p057_s03
\begin{snippet}{505}{11.5}{14}\begin{multicols}{3}\raggedcolumns 一个加数增加5.9，另一个不变，原和应增加一次5.9，不用分别求两个加数。把5.9凑成6，先加整数，再减去多加的0.1：\[\begin{aligned}&36.8+5.9\\={}&36.8+6-0.1\\={}&42.8-0.1\\={}&42.7\end{aligned}\]\end{multicols}\end{snippet}
% p057_s04
\begin{snippet}{505}{11.5}{14}\begin{multicols}{3}\raggedcolumns 同样是20个，计数单位不同，表示的数也不同：20个十分之一合2，20个百分之一合\(0.2\)\allowbreak 。先把两组分别写成小数，再相加求和、以较大数减较小数求差。不能把20个\(0.01\)\allowbreak 也看成2。\end{multicols}\end{snippet}
% p057_s05
\begin{snippet}{505}{11.5}{14}\begin{multicols}{3}\raggedcolumns \textbf{一、5（1）：}\ 先比较相邻两项，\(4.0\)\allowbreak 比\(3.2\)\allowbreak 多\(0.8\)\allowbreak ，\(4.8\)\allowbreak 又比\(4.0\)\allowbreak 多\(0.8\)\allowbreak ，说明每次增加\(0.8\)\allowbreak 。从\(4.8\)\allowbreak 继续依次加\(0.8\)\allowbreak 填写两个空。填好后再检查相邻各项的差是否一致。\end{multicols}\end{snippet}
% p057_s06
\begin{snippet}{505}{11.5}{14}\begin{multicols}{3}\raggedcolumns \textbf{一、5（2）：}\ 前两次都是后一项比前一项少\(1.2\)\allowbreak ，因此沿用每次减\(1.2\)\allowbreak 的规律。第一空从\(7.7\)\allowbreak 往后减，第二空再从第一空往后减。不要把“递减”误用成加法，也不要两空都从\(7.7\)\allowbreak 减一次。\end{multicols}\end{snippet}
% p057_s07
\begin{snippet}{505}{11.5}{14}\begin{multicols}{3}\raggedcolumns \textbf{一、5（3）：}\ \(9.5\)\allowbreak 与\(3.5\)\allowbreak 、\(15.5\)\allowbreak 与\(9.5\)\allowbreak 都相差6，因此相邻两项依次加6。空格在\(15.5\)\allowbreak 与\(27.5\)\allowbreak 之间，应比前者多6，也比后者少6。用左右两个方向检查，可确认没有漏掉一项。\end{multicols}\end{snippet}
% p057_s08
\begin{snippet}{505}{11.5}{14}\begin{multicols}{3}\raggedcolumns 扩大到原来的10倍，先把\(0.46\)\allowbreak 的小数点向右移动一位，得到\(4.6\)\allowbreak 。题目问“比原数大多少”，还须用扩大后的数减去原数。\(4.6\)\allowbreak 是扩大后的总量，不是增加量。\end{multicols}\end{snippet}
% p057_s09
\begin{snippet}{505}{11.5}{14}\begin{multicols}{3}\raggedcolumns 甲、乙组成总和，已知总和与甲数后，余下部分就是乙数。因此用\(26.3\)\allowbreak 减\(8.6\)\allowbreak 来求，不能把两数再相加。求得乙数后，可将它与甲数相加，检查是否回到题给的和。\end{multicols}\end{snippet}
% p057_s10
\begin{snippet}{505}{11.5}{14}\begin{multicols}{3}\raggedcolumns “甲比乙多\(3.7\)\allowbreak ”说明甲是较大数，乙是较小数。由较大数去掉多出的部分，才能还原乙，所以用\(9.2\)\allowbreak 减\(3.7\)\allowbreak 。不要看到“多”字就直接用加法，要先判断所求的是谁。\end{multicols}\end{snippet}
% p057_s11
\begin{snippet}{505}{11.5}{14}\begin{multicols}{3}\raggedcolumns \(4.8\)\allowbreak 是一位小数，计数单位是\(0.1\)\allowbreak 。整数4里有40个\(0.1\)\allowbreak ，再加小数部分的8个，共有48个；5里有50个同样的单位。比较48个与50个，还须添2个，不能把\(0.2\)\allowbreak 误填成添的个数。\end{multicols}\end{snippet}
% p057_s12
\begin{snippet}{505}{11.5}{14}\begin{multicols}{3}\raggedcolumns \textbf{一、10（1）：}\ 先分别求左右两个算式的值，左边两个加数可先凑成10，右边用12减\(1.7\)\allowbreak 。再比较10与\(10.3\)\allowbreak ，看整数部分相同后的十分位。圆圈中比较的是整个算式的值，不能只比较某一个加数。\end{multicols}\end{snippet}
% p057_s13
\begin{snippet}{505}{11.5}{14}\begin{multicols}{3}\raggedcolumns \textbf{一、10（2）：}\ 把\(5.99\)\allowbreak 看成比6少\(0.01\)\allowbreak ，与\(1.08\)\allowbreak 相加就比\(7.08\)\allowbreak 少\(0.01\)\allowbreak 。因此左边略小，不必把小数位数的多少当作大小依据。也可直接算出左边\(7.07\)\allowbreak ，再按数位比较。\end{multicols}\end{snippet}
% p057_s14
\begin{snippet}{505}{11.5}{14}\begin{multicols}{3}\raggedcolumns \textbf{一、10（3）：}\ 两边的运算不同，应先分别求值：左边减去6，右边将\(4.8\)\allowbreak 与\(1.9\)\allowbreak 相加。得到\(7.5\)\allowbreak 与\(6.7\)\allowbreak 后，整数部分已能确定大小。不要只比较题面最前面的\(13.5\)\allowbreak 与\(4.8\)\allowbreak 。\end{multicols}\end{snippet}
% p057_s15
\begin{snippet}{505}{11.5}{14}\begin{multicols}{3}\raggedcolumns \textbf{一、10（4）：}\ 先把两组加数分别求和，再比较\(13.2\)\allowbreak 与\(12.9\)\allowbreak 。虽然右边第一个加数较大，另一加数却更小，不能据一个加数决定总和大小。比较总和时先看整数部分。\end{multicols}\end{snippet}
% p057_s16
\begin{snippet}{505}{11.5}{14}\begin{multicols}{3}\raggedcolumns 两根木条原长相加时，重叠的\(3.6\)\allowbreak 厘米计了两次，拼好总长只计一次。因此先给29厘米补回这段重叠，得到两根原长之和，再减去已知的一根。不能从29直接减\(18.2\)\allowbreak ，否则会把另一根的重叠部分漏掉。\[\begin{aligned}&29+3.6-18.2\\={}&32.6-18.2\\={}&14.4\text{(厘米)}\end{aligned}\]\end{multicols}\end{snippet}
% p058_s01
\begin{snippet}{505}{11.5}{14}\begin{multicols}{3}\raggedcolumns 明日温度以今日\(19.5\)\allowbreak 摄氏度为起点，还要升高\(3.8\)\allowbreak 摄氏度。因此把原温度与升高量相加，求出的才是明日温度。列式时把两个小数的小数点对齐。\end{multicols}\end{snippet}
% p058_s02
\begin{snippet}{505}{11.5}{14}\begin{multicols}{3}\raggedcolumns 所求数与9.5相加组成30，应由和减去已知加数。把9.5凑成10，先减整数，再补回多减的0.5：\[\begin{aligned}&30-9.5\\={}&30-10+0.5\\={}&20+0.5\\={}&20.5\end{aligned}\]填回检验：\(20.5+9.5=30\)\allowbreak ，所以选C。\end{multicols}\end{snippet}
% p058_s03
\begin{snippet}{505}{11.5}{14}\begin{multicols}{3}\raggedcolumns 角对应元的十分位，分对应元的百分位，先分别换算：\par\noindent\[\displaystyle \begin{aligned}5\text{元}3\text{角}&=5.30\text{元}\\8\text{角}7\text{分}&=0.87\text{元}\end{aligned}\]\par 再按相同数位相加，\(5.30+0.87=6.17\)\allowbreak 元。不能把8角当8元，也不能把3角放到百分位。\end{multicols}\end{snippet}
% p058_s04
\begin{snippet}{505}{11.5}{14}\begin{multicols}{3}\raggedcolumns 一位小数末尾添0后，可以与两位小数按相同数位相加。若按竖式暂不化简、保留两位的写法，可选B；但去掉结果末尾的0后，可能是一位小数，甚至成为整数。比如\(1.2+0.30=1.50\)\allowbreak ，化简后为\(1.5\)\allowbreak ，所以B需要上述写法约定，不能作为无条件结论。\end{multicols}\end{snippet}
% p058_s05
\begin{snippet}{505}{11.5}{14}\begin{multicols}{3}\raggedcolumns 总共可减的数量没有变，减掉的部分却增加了\(3.6\)\allowbreak ，剩下的差就减少\(3.6\)\allowbreak 。变化方向与减数相反，因此选“减少”。不能把减数增加误认为差也增加。\end{multicols}\end{snippet}
% p058_s06
\begin{snippet}{505}{11.5}{14}\begin{multicols}{3}\raggedcolumns \(9.99\)\allowbreak 离10只差\(0.01\)\allowbreak ，可先把它凑成10，方便与原数相加。这一步比原来多加了\(0.01\)\allowbreak ，所以最后必须减回\(0.01\)\allowbreak 。补偿的符号要与“多加”相反。\end{multicols}\end{snippet}
% p058_s07
\begin{snippet}{505}{11.5}{14}\begin{multicols}{3}\raggedcolumns 路程相同，“小明比小雨慢0.9秒”表示小明多用0.9秒，小雨应用17.8秒减0.9秒。把0.9凑成1，多减的0.1再加回：\[\begin{aligned}&17.8-0.9\\={}&17.8-1+0.1\\={}&16.8+0.1\\={}&16.9\text{(秒)}\end{aligned}\]因此选B。\end{multicols}\end{snippet}
% p058_s08
\begin{snippet}{505}{11.5}{14}\begin{multicols}{3}\raggedcolumns 口算仍要按相同数位相加减：\(8.88-0.8\)\allowbreak 中的\(0.8\)\allowbreak 是8个十分之一，不能当成8个百分之一。连减\(0.36\)\allowbreak 、\(0.64\)\allowbreak 可合成减1，\(0.77\)\allowbreak 与\(0.23\)\allowbreak 、\(4.65\)\allowbreak 与\(5.35\)\allowbreak 也分别能凑整。整数减小数可先想成末尾添0的小数，防止退位时漏位。\end{multicols}\end{snippet}
% p058_s09
\begin{snippet}{505}{11.5}{14}\begin{multicols}{3}\raggedcolumns 列竖式时小数点对齐，\(4.6\)\allowbreak 、50可分别看作\(4.60\)\allowbreak 、\(50.00\)\allowbreak ，添0不改变大小。加法满十进1；\(20.05\)\allowbreak 和\(50.00\)\allowbreak 的减法遇到0不够减时，要逐级向高位退1，并记录借位后的数字。带星号题将差与减数相加，应还原50。\end{multicols}\end{snippet}
% p058_s09_local1
\begin{snippet}{37.0}{10.5}{14}\SourceSans\setlength{\mathindent}{2bp}\(17.55\)\end{snippet}
% p058_s09_local2
\begin{snippet}{37.0}{10.5}{14}\SourceSans\setlength{\mathindent}{2bp}\(6.25\)\end{snippet}
% p058_s09_local3
\begin{snippet}{37.0}{10.5}{14}\SourceSans\setlength{\mathindent}{2bp}\(12.16\)\end{snippet}
% p058_s09_local4
\begin{snippet}{37.0}{10.5}{14}\SourceSans\setlength{\mathindent}{2bp}\(33.66\)\end{snippet}
% p059_s01
\begin{snippet}{505}{11.5}{14}\begin{multicols}{3}\raggedcolumns \textbf{三、3（1）：}\ 减去括号中两个数的和，可以改为分别减去这两个数。先减\(12.46\)\allowbreak ，小数部分正好抵消而得到整数40，再减\(8.3\)\allowbreak 更方便。改变的是两个减数的先后次序，被减数始终不动。\end{multicols}\end{snippet}
% p059_s01_local1
\begin{snippet}{195.00000000000003}{10.5}{12.8}\SourceSans\setlength{\mathindent}{2bp}\[\begin{aligned}={}&52.46-12.46-8.3\\={}&40-8.3\\={}&31.7\end{aligned}\]\end{snippet}
% p059_s02
\begin{snippet}{505}{11.5}{14}\begin{multicols}{3}\raggedcolumns \textbf{三、3（2）：}\ \(10.1\)\allowbreak 由10与\(0.1\)\allowbreak 组成，减去它等于先减10，再减\(0.1\)\allowbreak 。这样每次只处理一个数位，更便于口算。拆开的两部分都是减数，后一项不能改成加\(0.1\)\allowbreak 。\end{multicols}\end{snippet}
% p059_s02_local1
\begin{snippet}{217.0}{10.5}{12.8}\SourceSans\setlength{\mathindent}{2bp}\[\begin{aligned}={}&86.72-10-0.1\\={}&76.72-0.1\\={}&76.62\end{aligned}\]\end{snippet}
% p059_s03
\begin{snippet}{505}{11.5}{14}\begin{multicols}{3}\raggedcolumns \textbf{三、3（3）：}\ 连续减去两个数，可以一次减去这两个数的和。观察到\(25.62\)\allowbreak 与\(74.38\)\allowbreak 的整数、小数部分合起来正好是100，先凑整再减。添括号后两个减数之间要改用加号，才能保持原来的差不变。\end{multicols}\end{snippet}
% p059_s03_local1
\begin{snippet}{196.00000000000003}{10.5}{12.8}\SourceSans\setlength{\mathindent}{2bp}\[\begin{aligned}={}&250-(25.62+74.38)\\={}&250-100\\={}&150\end{aligned}\]\end{snippet}
% p059_s04
\begin{snippet}{505}{11.5}{14}\begin{multicols}{3}\raggedcolumns \textbf{三、3（4）：}\ 移动一个数时，要连同它前面的加减符号一起移动。\(+6.9\)\allowbreak 与\(-6.9\)\allowbreak 先抵消为0，余下两个\(+1.3\)\allowbreak 再合并。不能把最后的\(+1.3\)\allowbreak 误看成跟在前一减号后也要减去。\end{multicols}\end{snippet}
% p059_s04_local1
\begin{snippet}{232.0}{10.5}{12.8}\SourceSans\setlength{\mathindent}{2bp}\[\begin{aligned}={}&(6.9-6.9)+(1.3+1.3)\\={}&0+2.6\\={}&2.6\end{aligned}\]\end{snippet}
% p059_s05
\begin{snippet}{505}{11.5}{14}\begin{multicols}{3}\raggedcolumns 方法一，“比另一种多多少”求的是两种续航的差，以较长的燃油车续航减去纯电动车续航：\[\begin{aligned}&596.5-425.8\\={}&170.7\text{(千米)}\end{aligned}\]方法二，把减数凑成整数426，先减整更方便；这样多减了0.2千米，最后须加回：\[\begin{aligned}&596.5-426+0.2\\={}&170.5+0.2\\={}&170.7\text{(千米)}\end{aligned}\]答：燃油汽车的单次续航多170.7千米。\end{multicols}\end{snippet}
% p059_s06
\begin{snippet}{505}{11.5}{14}\begin{multicols}{3}\raggedcolumns 方法一，本月比上月多，先在上月节水量上加1.15吨求本月，再把两个月相加：\[\begin{aligned}&3.65+(3.65+1.15)\\={}&3.65+4.80\\={}&8.45\text{(吨)}\end{aligned}\]方法二，把本月拆成“与上月同样多的一份”和“多出的1.15吨”。两个月就共有两份上月量，再加一次多出的量：\[\begin{aligned}&3.65+3.65+1.15\\={}&7.30+1.15\\={}&8.45\text{(吨)}\end{aligned}\]答：两个月共节约8.45吨。\end{multicols}\end{snippet}
% p059_s07
\begin{snippet}{505}{11.5}{14}\begin{multicols}{3}\raggedcolumns 方法一，两个驻留时长相比较，以较长时长减去较短时长，所得差就是多出的天数：\[\begin{aligned}&186.5-92.8\\={}&93.7\text{(天)}\end{aligned}\]方法二，将92.8凑成100后先减整百数；100比92.8多7.2，因此要把多减的7.2天补回：\[\begin{aligned}&186.5-100+7.2\\={}&86.5+7.2\\={}&93.7\text{(天)}\end{aligned}\]答：新一代空间站单次驻留多93.7天。\end{multicols}\end{snippet}
% p060_s01
\begin{snippet}{505}{11.5}{14}\begin{multicols}{3}\raggedcolumns 方法一，多找的9.3元不属于正确找零，应先从26.7元中扣除，得到应找的钱；再用付款减正确找零求售价：\[\begin{aligned}&100-(26.7-9.3)\\={}&100-17.4\\={}&82.6\text{(元)}\end{aligned}\]方法二，直接用100减错误找零，会使算出的售价少9.3元，因此再加9.3元补回。两种方法都是把多找的钱还原到售价里：\[\begin{aligned}&100-26.7+9.3\\={}&73.3+9.3\\={}&82.6\text{(元)}\end{aligned}\]答：围巾82.6元。\end{multicols}\end{snippet}
% p060_s02
\begin{snippet}{505}{11.5}{14}\begin{multicols}{3}\raggedcolumns 方法一，短画卷比长画卷少1.35米，先用长卷减相差量，求短卷；再与长卷相加：\[\begin{aligned}&3.6+(3.6-1.35)\\={}&3.6+2.25\\={}&5.85\text{(米)}\end{aligned}\]方法二，假想两幅都长3.6米，相加后只在短卷那一份多算了1.35米，所以减去一次即可：\[\begin{aligned}&3.6+3.6-1.35\\={}&7.2-1.35\\={}&5.85\text{(米)}\end{aligned}\]答：两幅画卷共长5.85米。\end{multicols}\end{snippet}
% p060_s03
\begin{snippet}{505}{11.5}{14}\begin{multicols}{3}\raggedcolumns 方法一，各买一件，两种商品的总价就是两个单价的和，按相同数位相加：\[\begin{aligned}&49.8+26.35\\={}&76.15\text{(元)}\end{aligned}\]方法二，49.8离50只差0.2，可先凑整相加，再减去多加的0.2元，保持总价不变：\[\begin{aligned}&50+26.35-0.2\\={}&76.35-0.2\\={}&76.15\text{(元)}\end{aligned}\]答：一共需要76.15元。\end{multicols}\end{snippet}
% p060_s04
\begin{snippet}{505}{11.5}{14}\begin{multicols}{3}\raggedcolumns 方法一，三种商品各买一件，先把三项单价相加。洗衣液与纸巾的价格正好凑成40元，先结合计算更方便：\[\begin{aligned}&49.8+26.35+13.65\\={}&49.8+(26.35+13.65)\\={}&49.8+40\\={}&89.8\text{(元)}\end{aligned}\]\(89.8<100\)\allowbreak ，钱够。\[\begin{aligned}&100-89.8\\={}&10.2\text{(元)}\end{aligned}\]方法二，从100元中先扣奶钱，再一次扣掉另外两种商品的合计价。扣完仍有余款，说明钱够，余款就是所求：\[\begin{aligned}&100-49.8-(26.35+13.65)\\={}&50.2-40\\={}&10.2\text{(元)}\end{aligned}\]答：钱够，还剩10.2元。\end{multicols}\end{snippet}
% p061_s01
\begin{snippet}{505}{11.5}{14}\begin{multicols}{3}\raggedcolumns 整数部分从右往左的首个计数单位是1，小数部分从左往右的首个计数单位是\(0.1\)\allowbreak 。整数计数单位越往左越大，小数计数单位越往右越小，所以这两个单位分别符合“最小”和“最大”。再对1与\(0.1\)\allowbreak 求和、求差。\end{multicols}\end{snippet}
% p061_s02
\begin{snippet}{505}{11.5}{14}\begin{multicols}{3}\raggedcolumns 先用原数减去剩下的\(0.05\)\allowbreak ，求出实际去掉的部分是\(0.40\)\allowbreak 。\(0.40\)\allowbreak 等于40个\(0.01\)\allowbreak ，所求是单位的个数，应填40。不要把去掉的数\(0.40\)\allowbreak 直接当作“几个”。\end{multicols}\end{snippet}
% p061_s03
\begin{snippet}{505}{11.5}{14}\begin{multicols}{3}\raggedcolumns 第一空以7.8为基准，“少2.81”用减法。把2.81拆成2.8与0.01，先减2.8可以凑成整数5：\[\begin{aligned}&7.8-2.81\\={}&7.8-2.8-0.01\\={}&5-0.01\\={}&4.99\end{aligned}\]第二空以2.36为基准，“多1.4”用加法，\(2.36+1.4=3.76\)\allowbreak 。两空的基准不同，不能共用一个原数。\end{multicols}\end{snippet}
% p061_s04
\begin{snippet}{505}{11.5}{14}\begin{multicols}{3}\raggedcolumns 小数点向右移动一位，相当于把原数扩大到10倍，先得到75。题目问增加了多少，应再从75中减去原来的\(7.5\)\allowbreak 。移动后的数与增加的部分是两个不同的量。\end{multicols}\end{snippet}
% p061_s05
\begin{snippet}{505}{11.5}{14}\begin{multicols}{3}\raggedcolumns 第一式缺被减数，用“差＋减数”还原；第二式缺一个加数，用“和－另一个加数”求。第三式缺的是减数，应以“被减数－差”求，不能也把差加上。填好后逐空代回原式，检查等号两边是否相等。\end{multicols}\end{snippet}
% p061_s06
\begin{snippet}{505}{11.5}{14}\begin{multicols}{3}\raggedcolumns \textbf{一、6（1）：}\ 两边被减数都是\(4.3\)\allowbreak ，只需比较减数。\(2.5\)\allowbreak 大于\(0.25\)\allowbreak ，左边减掉更多，剩下的差反而更小。这里不能根据减数较大就判左边的算式较大。\end{multicols}\end{snippet}
% p061_s07
\begin{snippet}{505}{11.5}{14}\begin{multicols}{3}\raggedcolumns \textbf{一、6（2）：}\ 把两组加数分别拆成整数部分与小数部分，整数部分都是3与9，小数部分都是\(0.1\)\allowbreak 与\(0.84\)\allowbreak 。各部分相同，只是配到的加数不同，因此两个和相等。本小问是加法，不能套用“减数越大差越小”的理由。\end{multicols}\end{snippet}
% p061_s08
\begin{snippet}{505}{11.5}{14}\begin{multicols}{3}\raggedcolumns \textbf{一、6（3）：}\ 先看括号：\(3.8\)\allowbreak 加\(1.6\)\allowbreak 的结果大于\(3.8\)\allowbreak 减\(1.6\)\allowbreak 的结果。两边再同样用\(7.8\)\allowbreak 去减，左边因减数较大而得到较小的差。必须先把括号看作一个整体，不能只比较括号外的数。\end{multicols}\end{snippet}
% p061_s09
\begin{snippet}{505}{11.5}{14}\begin{multicols}{3}\raggedcolumns 先统一单位，不能把8厘米当0.8米：\par\noindent\(\displaystyle 8\text{厘米}=(8\div100)\text{米}=0.08\text{米}\)\allowbreak \par 求较小的小芳成绩，用\(1.25-0.08=1.17\)\allowbreak 米，结果应小于小明的1.25米。原空只填1.17。\end{multicols}\end{snippet}
% p061_s10
\begin{snippet}{505}{11.5}{14}\begin{multicols}{3}\raggedcolumns 两位小数的计数单位是\(0.01\)\allowbreak ，不是\(0.1\)\allowbreak 。已知原数加一个这样的单位才得到\(1.2\)\allowbreak ，求原数须反过来减\(0.01\)\allowbreak 。把\(1.2\)\allowbreak 写成\(1.20\)\allowbreak ，就能按百分位清楚地减去。\end{multicols}\end{snippet}
% p061_s11
\begin{snippet}{505}{11.5}{14}\begin{multicols}{3}\raggedcolumns 甲数比乙数多\(0.75\)\allowbreak ，因此乙数应先用甲数减\(0.75\)\allowbreak 求出。再把甲、乙两数相加，才得到题目要的和。若先按两个甲数相加，最后也只能减去一次\(0.75\)\allowbreak ，因为只有乙的那一份较少。\[\begin{aligned}&10.9+(10.9-0.75)\\={}&10.9+10.15\\={}&21.05\end{aligned}\]\end{multicols}\end{snippet}
% p061_s12
\begin{snippet}{505}{11.5}{14}\begin{multicols}{3}\raggedcolumns \textbf{一、10（1）：}\ 已知\(1.2\)\allowbreak 比\(0.5\)\allowbreak 多\(0.7\)\allowbreak ，按每次加\(0.7\)\allowbreak 的规律填写。第一空应介于\(1.2\)\allowbreak 与\(2.6\)\allowbreak 之间，再由\(2.6\)\allowbreak 继续加\(0.7\)\allowbreak 求第二空。检查各相邻项的差都应相同。\end{multicols}\end{snippet}
% p061_s13
\begin{snippet}{505}{11.5}{14}\begin{multicols}{3}\raggedcolumns \textbf{一、10（2）：}\ 前面的相邻两项每次减少\(0.03\)\allowbreak ，即少3个百分之一。因此从\(5.13\)\allowbreak 接着减，先到\(5.10\)\allowbreak ，再到\(5.07\)\allowbreak 。\(5.10\)\allowbreak 虽可化简为\(5.1\)\allowbreak ，但按两位小数写能更清楚地看出规律。\[\begin{aligned}&5.13-0.03\\={}&5.10\\&5.10-0.03\\={}&5.07\end{aligned}\]\end{multicols}\end{snippet}
% p062_s01
\begin{snippet}{505}{11.5}{14}\begin{multicols}{3}\raggedcolumns 先分清学校与两家的位置关系，不能只画一种情形。学校在两家之间时，两段到校距离首尾相接，两家最远相距，用\(1.36+0.78\)\allowbreak 求；两家在学校同侧时，两段路程有重合，最近相距，用\(1.36-0.78\)\allowbreak 求。前页的最远一空与本页的最近一空分别对应求和、求差，填好后最近距离应小于最远距离。\[\begin{aligned}&1.36+0.78\\={}&2.14\text{(千米，最远)}\\&1.36-0.78\\={}&0.58\text{(千米，最近)}\end{aligned}\]\end{multicols}\end{snippet}
% p062_s02
\begin{snippet}{505}{11.5}{14}\begin{multicols}{3}\raggedcolumns 方法一，倒油前后桶重不变，连桶质量减少的部分就是倒出的半桶油；把两份半桶油的质量相加得到整桶油，再用原连桶质量减去油重求桶重。\[\begin{aligned}&46.8-24.5\\={}&22.3\text{(千克)}\end{aligned}\]\[\begin{aligned}&22.3+22.3\\={}&44.6\text{(千克)}\end{aligned}\]\[\begin{aligned}&46.8-44.6\\={}&2.2\text{(千克)}\end{aligned}\]方法二，把剩下的“半桶油连桶”算两份，油合起来正好是一整桶，却算进了两个桶；再减去原来的“整桶油连桶”，就只剩一个空桶。\[\begin{aligned}&24.5+24.5-46.8\\={}&49-46.8\\={}&2.2\text{(千克)}\end{aligned}\]\[\begin{aligned}&46.8-2.2\\={}&44.6\text{(千克)}\end{aligned}\]两空按“桶重、原有油重”的顺序填，不能颠倒。\end{multicols}\end{snippet}
% p062_s03
\begin{snippet}{505}{11.5}{14}\begin{multicols}{3}\raggedcolumns 第二个加数的百分位是6，第一个加数可看成百分位0，所以和的百分位仍为6，先排除A。再用最小的两个加数\(20.4\)\allowbreak 与\(9.06\)\allowbreak 求得最小和\(29.46\)\allowbreak ，B比它还小，也不可能。C可由\(21.4\)\allowbreak 与\(9.86\)\allowbreak 相加得到，既要排除，也要确认留下的选项确实可实现。\end{multicols}\end{snippet}
% p062_s04
\begin{snippet}{505}{11.5}{14}\begin{multicols}{3}\raggedcolumns 先沿小数点找到数位，9与5都在十分位，分别表示9个、5个十分之一。相加的14是14个十分之一，其中10个十分之一可进成1个一，余4个十分之一。选项问的是相加时的计数单位，不能只看9加5等于14就选14个一。\end{multicols}\end{snippet}
% p062_s05
\begin{snippet}{505}{11.5}{14}\begin{multicols}{3}\raggedcolumns \(9.9\)\allowbreak 可以拆成\(9+0.9\)\allowbreak ，也可以凑成\(10-0.1\)\allowbreak ，所以A与C都保持原和不变。B把9.9变成了\(10+0.1\)\allowbreak ，实际多加了\(0.2\)\allowbreak 。本问要找“不同”的算式，不能选到正确的等值变形。\end{multicols}\end{snippet}
% p062_s06
\begin{snippet}{505}{11.5}{14}\begin{multicols}{3}\raggedcolumns 和先因一个加数增加而多\(3.57\)\allowbreak ，又因另一个加数减少而少\(4.5\)\allowbreak 。减少量更大，二者抵消后还少\(4.5-3.57=0.93\)\allowbreak 。因此既要确定变化量，也要确定变化方向。\end{multicols}\end{snippet}
% p062_s07
\begin{snippet}{505}{11.5}{14}\begin{multicols}{3}\raggedcolumns “接近9”要比较与9的距离，9两边的数都可能最接近。三数与9分别相差\(0.02\)\allowbreak 、\(0.01\)\allowbreak 、\(0.1\)\allowbreak ，相差最少的最接近。不能因为\(9.02\)\allowbreak 大于9就优先选它，也不能只比较三个数本身的大小。\end{multicols}\end{snippet}
% p062_s08
\begin{snippet}{505}{11.5}{14}\begin{multicols}{3}\raggedcolumns ①可用一个整数加一个小数得到小数作反例；②同样的\(6.4\)\allowbreak 抵消后，还剩两个\(3.6\)\allowbreak ，并不是0。③必须对齐小数点，末位不一定属于同一数位；④被减数和减数同时增加相同的量，多出的部分互相抵消，差不变。逐条判断后，再数正确说法的个数。\end{multicols}\end{snippet}
% p062_s09
\begin{snippet}{505}{11.5}{14}\begin{multicols}{3}\raggedcolumns 一位小数与两位小数相加减时，仍须按数位对齐，如\(6.76-0.6\)\allowbreak 只减6个十分之一。\(2.36\)\allowbreak 与\(1.64\)\allowbreak 的小数部分能凑成1，可先利用凑整；\(4.2-4\)\allowbreak 的整数部分减完后是0，个位0与小数点要写全。整数加小数不会改变小数部分原来的数位。\end{multicols}\end{snippet}
% p063_s01
\begin{snippet}{505}{11.5}{14}\begin{multicols}{3}\raggedcolumns 先利用小数末尾添0的性质，把\(8.7\)\allowbreak 、\(5.1\)\allowbreak 、12分别写成两位小数，再对齐小数点。\(5.10\)\allowbreak 、\(12.00\)\allowbreak 做减法遇到不够减时，要逐级退位，借来1个十分之一就是10个百分之一。两道带星号题分别用和减去一个加数、差加上减数来验算。\end{multicols}\end{snippet}
% p063_s01_local1
\begin{snippet}{37.0}{10.5}{14}\SourceSans\setlength{\mathindent}{2bp}\(21.06\)\end{snippet}
% p063_s01_local2
\begin{snippet}{37.0}{10.5}{14}\SourceSans\setlength{\mathindent}{2bp}\(1.51\)\end{snippet}
% p063_s01_local3
\begin{snippet}{37.0}{10.5}{14}\SourceSans\setlength{\mathindent}{2bp}\(9.75\)\end{snippet}
% p063_s02
\begin{snippet}{505}{11.5}{14}\begin{multicols}{3}\raggedcolumns \textbf{三、3（1）：}\ 本式没有可直接配成整数的两数，按从左到右计算也可以。原红式把\(+6.3\)\allowbreak 与\(-3.7\)\allowbreak 结合，是保留符号的等值重排，先求两者的差，再与35相加。不要把\(35-3.7+6.3\)\allowbreak 误变成35减去两数的和。\end{multicols}\end{snippet}
% p063_s02_local1
\begin{snippet}{182.0}{10.5}{12.8}\SourceSans\setlength{\mathindent}{2bp}\[\begin{aligned}={}&35+(6.3-3.7)\\={}&35+2.6\\={}&37.6\end{aligned}\]\end{snippet}
% p063_s03
\begin{snippet}{505}{11.5}{14}\begin{multicols}{3}\raggedcolumns \textbf{三、3（2）：}\ \(9.98\)\allowbreak 比10少\(0.02\)\allowbreak ，减10比原来多减了\(0.02\)\allowbreak 。因此先算27减10，再把多减的\(0.02\)\allowbreak 加回。补偿时应加回，不能再减一次。\end{multicols}\end{snippet}
% p063_s03_local1
\begin{snippet}{232.0}{10.5}{12.8}\SourceSans\setlength{\mathindent}{2bp}\[\begin{aligned}={}&27-10+0.02\\={}&17+0.02\\={}&17.02\end{aligned}\]\end{snippet}
% p063_s04
\begin{snippet}{505}{11.5}{14}\begin{multicols}{3}\raggedcolumns \textbf{三、3（3）：}\ 减去括号中的和，等于依次减去两个加数。交换两个减数的顺序，先让\(45.35\)\allowbreak 减\(5.35\)\allowbreak ，恰好得到40，再减\(7.2\)\allowbreak 。这样利用相同的小数部分凑整，且两个减数的符号都不能改变。\end{multicols}\end{snippet}
% p063_s04_local1
\begin{snippet}{208.00000000000003}{10.5}{12.8}\SourceSans\setlength{\mathindent}{2bp}\[\begin{aligned}={}&45.35-5.35-7.2\\={}&40-7.2\\={}&32.8\end{aligned}\]\end{snippet}
% p063_s05
\begin{snippet}{505}{11.5}{14}\begin{multicols}{3}\raggedcolumns \textbf{三、3（4）：}\ 加减混合时，移动各数必须带着原来的符号。先把两个加数\(23.8\)\allowbreak 、\(36.2\)\allowbreak 凑成60，再把两个减数\(3.9\)\allowbreak 、\(5.1\)\allowbreak 合成9。两个减数合并到括号里应用加号，括号外只保留一个减号。\end{multicols}\end{snippet}
% p063_s05_local1
\begin{snippet}{279.0}{10.5}{12.8}\SourceSans\setlength{\mathindent}{2bp}\[\begin{aligned}={}&(23.8+36.2)-(3.9+5.1)\\={}&60-9\\={}&51\end{aligned}\]\end{snippet}
% p063_s06
\begin{snippet}{505}{11.5}{14}\begin{multicols}{3}\raggedcolumns 方法一，“十七号比十二号缩短19.5小时”说明十二号用时更长。以已知的较短时间加缩短量，才能还原较长时间：\[\begin{aligned}&9+19.5\\={}&28.5\text{(小时)}\end{aligned}\]方法二，先把19.5凑成20与9相加，再减去多加的0.5小时；这是凑整后补偿的计算思路：\[\begin{aligned}&9+20-0.5\\={}&29-0.5\\={}&28.5\text{(小时)}\end{aligned}\]答：神舟十二号返航大约用了28.5小时。\end{multicols}\end{snippet}
% p063_s07
\begin{snippet}{505}{11.5}{14}\begin{multicols}{3}\raggedcolumns 竹竿总长由入泥、水中、露出三段组成，水深只对应水中的一段。先统一单位：\par\noindent\[\displaystyle \begin{aligned}16\text{厘米}&=0.16\text{米}\\54\text{厘米}&=0.54\text{米}\end{aligned}\]\par 用米计算时，按连减性质先合并水深以外两段，0.16与0.54正好凑成0.7：\par\noindent\[\displaystyle \begin{aligned}&2-0.16-0.54\\={}&2-(0.16+0.54)\\={}&2-0.7\\={}&1.3\text{(米)}\end{aligned}\]\par 也可统一成厘米：\par\noindent\(\displaystyle 2\text{米}=200\text{厘米}\)\allowbreak \par 同样先合并两段：\par\noindent\[\displaystyle \begin{aligned}&200-(16+54)\\={}&200-70\\={}&130\text{(厘米)}\end{aligned}\]\par \noindent\(\displaystyle 130\text{厘米}=1.3\text{米}\)\allowbreak \par 两种单位写法的数量关系相同。答：水池的水深是1.3米。\end{multicols}\end{snippet}
% p064_s01
\begin{snippet}{505}{11.5}{14}\begin{multicols}{3}\raggedcolumns 方法一，华山比衡山高，先用衡山海拔加0.85千米求华山；泰山又比华山低0.62千米，因此接着减：\[\begin{aligned}&1.3+0.85-0.62\\={}&2.15-0.62\\={}&1.53\text{(千米)}\end{aligned}\]方法二，0.85千米中去掉华山高出泰山的0.62千米，余下就是泰山高出衡山的部分，再加衡山海拔：\[\begin{aligned}&1.3+(0.85-0.62)\\={}&1.3+0.23\\={}&1.53\text{(千米)}\end{aligned}\]答：东岳泰山约高1.53千米。\end{multicols}\end{snippet}
% p064_s02
\begin{snippet}{505}{11.5}{14}\begin{multicols}{3}\raggedcolumns 方法一，两根原绳的长度相加时，重叠接头被计了两次；接好后的实际总长只计一次。两者相减，正好去掉多计的一次接头：\[\begin{aligned}&2.35+2.85-4.7\\={}&5.20-4.70\\={}&0.50\text{(米)}\end{aligned}\]方法二，先从接后总长扣去第二根整段，剩下的是第一根未重叠部分；再从第一根原长中扣去它，得到接头长：\[\begin{aligned}&2.35-(4.7-2.85)\\={}&2.35-1.85\\={}&0.50\text{(米)}\end{aligned}\]答：接头处长0.5米。\end{multicols}\end{snippet}
% p064_s03
\begin{snippet}{505}{11.5}{14}\begin{multicols}{3}\raggedcolumns 方法一，10月比9月少，先用9月销量减0.26万箱求10月销量，再加回9月销量求两月合计：\[\begin{aligned}&1.8+(1.8-0.26)\\={}&1.8+1.54\\={}&3.34\text{(万箱)}\end{aligned}\]方法二，先假设两月各售1.8万箱，两份合计后，只需扣掉10月少售的一次0.26万箱。全过程的单位都是万箱，不要中途漏掉“万”：\[\begin{aligned}&1.8+1.8-0.26\\={}&3.60-0.26\\={}&3.34\text{(万箱)}\end{aligned}\]答：两个月共售出3.34万箱。\end{multicols}\end{snippet}
% p064_s04
\begin{snippet}{505}{11.5}{14}\begin{multicols}{3}\raggedcolumns 方法一，38.5元里含有多找的10元，先扣除多找部分，才是正确找零；付款减去正确找零就是书价：\[\begin{aligned}&50-(38.5-10)\\={}&50-28.5\\={}&21.5\text{(元)}\end{aligned}\]方法二，付款减错误找零，会把书价少算10元，须再加回10元。多找使售价被低估，所以补偿应使用加法：\[\begin{aligned}&50-38.5+10\\={}&11.5+10\\={}&21.5\text{(元)}\end{aligned}\]答：这本书21.5元。\end{multicols}\end{snippet}
% p065_s01
\begin{snippet}{505}{11.5}{14}\begin{multicols}{3}\raggedcolumns 逐个寻找一条可能的对称轴，再设想沿它对折，核对两侧的轮廓与内部结构是否都重合。第1、3、5、6、7、10个符合条件；其余图有斜边、波浪或长短臂不能配对，不能仅凭“看起来平衡”就打勾。判断的是轴对称，图形转半圈能重合也不等于能对折重合。\end{multicols}\end{snippet}
% p065_s02
\begin{snippet}{505}{11.5}{14}\begin{multicols}{3}\raggedcolumns 按建筑结构填；含牌匾文字的整图不完全对称。\par 本空把门坊的建筑结构作为观察对象，左右柱、梁等构件互相对应，可沿中间竖线对折，所以通常填“轴对称”。观察范围须说清：牌匾上的四个汉字并不会因建筑左右对称而变成对称文字。若把文字也纳入整个平面图案，就不能断言整幅图能沿这条线完全重合。\end{multicols}\end{snippet}
% p065_s03
\begin{snippet}{505}{11.5}{14}\begin{multicols}{3}\raggedcolumns \textbf{一、3（1）：}\ 在三幅船上选同一个顶点，如船头尖，依照给定中间图比较它的前后位置。第一步同列向下跨5格，第二步同排向右跨4格，船身朝向始终不变。数格时数相邻格线之间的距离，不能数图形之间空着几格，也不能把起点格点算作第一格。\end{multicols}\end{snippet}
% p065_s04
\begin{snippet}{505}{11.5}{14}\begin{multicols}{3}\raggedcolumns \textbf{一、3（2）：}\ 固定用鱼嘴尖作对应点，先比较第一幅与中间幅，再比较中间幅与最后一幅。两段移动依次是向上5格、向左5格，其他顶点也应移动相同的方向和格数。不能把两个不同部位当作一对对应点来数格。\end{multicols}\end{snippet}
% p065_s05
\begin{snippet}{505}{11.5}{14}\begin{multicols}{3}\raggedcolumns 对称轴是直线。一般长方形有经过两组对边中点的2条对称轴；正方形再加上两条对角线所在的直线，共4条。等边三角形有分别经过顶点和对边中点的3条对称轴。\end{multicols}\end{snippet}
% p066_s01
\begin{snippet}{505}{11.5}{14}\begin{multicols}{3}\raggedcolumns 基本单元是圆，把同样大小的圆移到不同位置，就能组成五环的轮廓，因此选择平移。数对称轴时要看五个圆的整体排列，而不能按单个圆有无数条轴来填。按题图只看圆环轮廓，只有中间一条竖直轴能让上排、下排的圆分别配对；沿它画虚线。\end{multicols}\end{snippet}
% p066_s02
\begin{snippet}{505}{11.5}{14}\begin{multicols}{3}\raggedcolumns \textbf{一、6（1）：}\ 比较房顶的朝向，依次为上、右、空白、左，说明每次顺时针转\(90^\circ\)\allowbreak 。因此第三幅房顶朝下，其他各边也应随之转动一个直角。画时保持边长与顶点连接顺序，不能只改屋顶方向而把房身留在原位。\end{multicols}\end{snippet}
% p066_s03
\begin{snippet}{505}{11.5}{14}\begin{multicols}{3}\raggedcolumns \textbf{一、6（2）：}\ 每幅折箭头都比前一幅顺时针转了一个直角，第四格继续沿用相同方向与角度。应同时转动两条臂和箭头，得到竖直向上、水平向右的形状。长短臂的位置关系也随旋转改变，不能把图形翻面代替旋转。\end{multicols}\end{snippet}
% p066_s04
\begin{snippet}{505}{11.5}{14}\begin{multicols}{3}\raggedcolumns \textbf{一、7（①）：}\ 7. ①先向上平移1格，再向右平移1格。 ②绕A点逆时针旋转90\ensuremath{^\circ}，再向上平移1格。\par 比较图二①与图一对应部分，用同一个格点作起点和终点。先向上移1格，使高低位置吻合，再向右移1格，使左右位置吻合。平移中各点都要作同样的移动，马身的形状与朝向不改变。\end{multicols}\end{snippet}
% p066_s05
\begin{snippet}{505}{11.5}{14}\begin{multicols}{3}\raggedcolumns \textbf{一、7（②）：}\ ②先要改正朝向，再调整位置：以标出的A为固定点，逆时针转\(90^\circ\)\allowbreak ，之后整体上移1格。旋转时A不动，其他点到A的距离保持不变；平移时所有点都上移1格。两种动作的依据不同，不能把旋转中心随意换到图形中间。\end{multicols}\end{snippet}
% p066_s06
\begin{snippet}{505}{11.5}{14}\begin{multicols}{3}\raggedcolumns 观察活动部件怎样运动：键盘托架、抽屉抽出时整体沿直线同向移动，方向不变，是平移。订书机活动部分、旋转闸、门环都绕固定处转动，是旋转。先逐个分类，再数平移的个数，不能把所有“位置变了”的运动都算成平移。\end{multicols}\end{snippet}
% p066_s07
\begin{snippet}{505}{11.5}{14}\begin{multicols}{3}\raggedcolumns 将甲图与乙图的对应顶点、边长逐一比较，再设想绕一个固定点转动。第一、二组能在保持形状、大小的同时改变朝向而重合，第三组表现为同朝向移位。选择要同时包含前两组，不能把位置改变本身当成旋转的充分依据。\end{multicols}\end{snippet}
% p066_s08
\begin{snippet}{505}{11.5}{14}\begin{multicols}{3}\raggedcolumns 风车的每一片不仅有外轮廓，还有黑白两部分，基本单元应选完整的黑白拼块。让它绕公共中心旋转，才能依次取得不同朝向并保持黑白位置对应。单靠平移不改变朝向，不能形成题图风叶的排列。\end{multicols}\end{snippet}
% p067_s01
\begin{snippet}{505}{11.5}{14}\begin{multicols}{3}\raggedcolumns 把折边当作对称轴，设想把每个选项展开成完整图形，再比较四个尖角的方向与长短。C展开后能使各尖角成对对应，轮廓与题图吻合。不能只看到一个尖角相似就选，还须核对其他各处的凹口和边缘。\end{multicols}\end{snippet}
% p067_s02
\begin{snippet}{505}{11.5}{14}\begin{multicols}{3}\raggedcolumns 分针正常走动是顺时针，从6依次经过7、8到9，共跨3个大格。钟面一周\(360^\circ\)\allowbreak 平均分12大格，每格\(30^\circ\)\allowbreak ，3格就是\(90^\circ\)\allowbreak 。要数经过的格数，不能把6、7、8、9四个数字当作4格。\end{multicols}\end{snippet}
% p067_s03
\begin{snippet}{505}{11.5}{14}\begin{multicols}{3}\raggedcolumns \textbf{三、1（第1图）：}\ 先找左半图的各个拐点，沿水平格线数它们到竖直轴的距离，在右侧找等距点。轴上的点保持不动，再按原图顺序把对应点连接。最后检查对应点连线是否垂直于轴，并被轴平分，不能只把轮廓画得大致相像。\end{multicols}\end{snippet}
% p067_s04
\begin{snippet}{505}{11.5}{14}\begin{multicols}{3}\raggedcolumns \textbf{三、1（第2图）：}\ 先补外轮廓，再对内部小三角形的三个顶点分别找对称点。竖直轴两侧的对应点必须在同一水平格线上，离轴的格数相等。内部三角形也属于图形，漏掉它就不能使两半完全重合。\end{multicols}\end{snippet}
% p067_s05
\begin{snippet}{505}{11.5}{14}\begin{multicols}{3}\raggedcolumns \textbf{三、1（第3图）：}\ 外侧斜边、底边以及内部横线都要补齐，可先确定每条线段两端的对称点。轴上的端点不动，其余点沿水平格线在另一侧取等距位置。连线后检查左右斜边是否成镜像、内部横线是否两侧等长。\end{multicols}\end{snippet}
% p067_s06
\begin{snippet}{505}{11.5}{14}\begin{multicols}{3}\raggedcolumns \textbf{三、2（第1图）：}\ 三片相同叶片围绕同一中心均匀排列，一条轴经过一片叶片的中线时，另两片应互相重合。对三片叶片分别这样检查，可得到3条不同的轴。画轴用虚线，并穿过公共中心，不能把相邻叶片之间的任意直线都算进去。\end{multicols}\end{snippet}
% p067_s07
\begin{snippet}{505}{11.5}{14}\begin{multicols}{3}\raggedcolumns \textbf{三、2（第2图）：}\ 既要让外圆环重合，也要让斜带重合，才是整个标志的对称轴。沿斜带中线的一条轴符合条件；过圆心且垂直于斜带的另一条轴也符合条件。任意直径虽是圆的对称轴，却不一定是这个组合图案的对称轴。\end{multicols}\end{snippet}
% p067_s08
\begin{snippet}{505}{11.5}{14}\begin{multicols}{3}\raggedcolumns \textbf{三、2（第3图）：}\ 把扇形的两条半径互相对折，折痕须经过圆心，并平分两半径夹角。沿这条直线对折时，两条半径和两段弧都能重合，所以只有1条轴。其他半径不能同时平分图形，不应多画。\end{multicols}\end{snippet}
% p067_s09
\begin{snippet}{505}{11.5}{14}\begin{multicols}{3}\raggedcolumns \textbf{三、2（第4图）：}\ 先找连接相对顶点的直线，三对相对顶点给出3条轴。再找连接相对边中点的直线，另外有3条，均经过中心。两类轴互不重复，共6条，不能只画过顶点的一类。\end{multicols}\end{snippet}
% p067_s10
\begin{snippet}{505}{11.5}{14}\begin{multicols}{3}\raggedcolumns 第一步固定O，把各顶点到O的线段顺时针转\(90^\circ\)\allowbreak ，长度保持不变，再按顺序连点。第二步的对象是刚旋转得到的图，将其所有顶点向右移3格，包括刚才与O重合的那个顶点。题目原标记O仍作参照，两步顺序不能随意颠倒。\end{multicols}\end{snippet}
% p068_s01
\begin{snippet}{505}{11.5}{14}\begin{multicols}{3}\raggedcolumns 三角形以B为中心逆时针转，长方形以D为中心顺时针转，两次旋转各有自己的中心。分别固定B、D，再按直角方向逐点定位，所有边长保持不变。拼合后才检查整个“屋顶与矩形”能否沿竖直中线对折，不能把原来的整个组合绕同一点一起转。\end{multicols}\end{snippet}
% p068_s02
\begin{snippet}{505}{11.5}{14}\begin{multicols}{3}\raggedcolumns \textbf{三、3（3）填空：}\ 把小三角形的一个顶点与凹口中的对应顶点配对，先数水平方向相差9格，再数竖直方向相差2格。需要先向右，再向上，其他顶点也应得到同样的位移。检验移动后是否恰好填满凹口，既不重叠也不留下缝隙。\end{multicols}\end{snippet}
% p068_s03
\begin{snippet}{505}{11.5}{14}\begin{multicols}{3}\raggedcolumns \textbf{三、3（3）作图：}\ 先按上一问把小三角形填入凹口，确认得到完整长方形。再固定O，将这个长方形整体逆时针转\(90^\circ\)\allowbreak ，原水平长边转成竖直长边，原竖直短边也随之转向。逐顶点保持到O的距离，所得图形应在O左上方，不能只旋转小三角形。\end{multicols}\end{snippet}
% p068_s04
\begin{snippet}{505}{11.5}{14}\begin{multicols}{3}\raggedcolumns \textbf{三、4（1）：}\ 共同顶点B始终不动，因而B是旋转中心。按题图弧箭头采用“虚线到实线”的读法，方向为逆时针；A的对应点须保持到B的距离，标在右侧括号中。若改以标有字母的实线为原图、虚线为结果，起止互换后应读成顺时针，所以必须交代采用的读图约定。\end{multicols}\end{snippet}
% p068_s05
\begin{snippet}{505}{11.5}{14}\begin{multicols}{3}\raggedcolumns \textbf{三、4（2）：}\ 长方形共同顶点D不动，中心应是D。沿弧箭头从虚线到实线看为顺时针，B的对应点在右下括号中，并与B到D等距。若以实线为原图、虚线为结果，方向会反为逆时针；本图也要沿用前问的箭头约定，不能混用两种读法。\end{multicols}\end{snippet}
% p068_s06
\begin{snippet}{505}{11.5}{14}\begin{multicols}{3}\raggedcolumns 先绕原来的旋转中心逆时针旋转90\ensuremath{^\circ}。 再向左平移6格，即可回到原来的位置。\par 还原要从最后一步开始撤销，动作的顺序也须倒过来。先绕原来同一个中心逆时针转\(90^\circ\)\allowbreak ，抵消最后的顺时针旋转；再向左移6格，抵消最初的右移。若先左移再反转，图形相对旋转中心的位置已经改变，一般不能回到原处。\end{multicols}\end{snippet}
% p069_s01
\begin{snippet}{505}{11.5}{14}\begin{multicols}{3}\raggedcolumns 轴对称的关键是沿一条直线对折后两部分完全重合，这条折痕所在的直线才叫对称轴。对折时一对对应点会合到一起，因此它们到轴的距离相等，连线与轴垂直。不能只满足离轴一样远却错开高低位置，两个条件需要同时符合。\end{multicols}\end{snippet}
% p069_s02
\begin{snippet}{505}{11.5}{14}\begin{multicols}{3}\raggedcolumns 平移时所有点向同一方向移动相同的距离，所以对应边长与连接方式不变。图形的形状、大小保持原样，只是位置改变，朝向也不改变。可找两组对应点检查它们的移动是否同向、等距，避免把转动误认为平移。\end{multicols}\end{snippet}
% p069_s03
\begin{snippet}{505}{11.5}{14}\begin{multicols}{3}\raggedcolumns 滑动解锁和缆车沿轨道整体移动，题图中朝向保持不变，按平移标△。摩天轮整体绕中心转，冰箱门绕门轴转，按旋转标○。本题判断的是指定物体的运动方式，不能转而判断图片是不是轴对称。\end{multicols}\end{snippet}
% p069_s04
\begin{snippet}{505}{11.5}{14}\begin{multicols}{3}\raggedcolumns 均匀刻度下，同样增加2千克，指针应再顺时针转同样的\(90^\circ\)\allowbreak 。此时总共放了4千克，指针已从0转过两个直角；全部拿走要反向回到0，所以逆时针转\(180^\circ\)\allowbreak 。最后一空是从4千克位置回零的总角度，不能仍填一个直角。\end{multicols}\end{snippet}
% p069_s05
\begin{snippet}{505}{11.5}{14}\begin{multicols}{3}\raggedcolumns A、B是一对对称点，对称轴把AB平分，所以A到轴的距离是整条AB的一半。AB横跨4格且长4厘米，可知每格1厘米，再数出C、F各离轴3格。D的对应点应在轴另一侧同一条垂线上的等距位置，沿图核对到E，不能只按字母顺序配对。\end{multicols}\end{snippet}
% p070_s01
\begin{snippet}{505}{11.5}{14}\begin{multicols}{3}\raggedcolumns \textbf{一、6：}\ 先把一周\(360^\circ\)\allowbreak 平均分成12大格，得每大格\(30^\circ\)\allowbreak 。从2到5依次跨3大格，正常转动方向是顺时针，因此总角度为3个\(30^\circ\)\allowbreak 。应数2到3、3到4、4到5这三段，不能数成4格。\end{multicols}\end{snippet}
% p070_s02
\begin{snippet}{505}{11.5}{14}\begin{multicols}{3}\raggedcolumns \textbf{一、7：}\ 比较同一顶点B与B′，相差2厘米。厘米化毫米乘10：\par\noindent\(\displaystyle 2\text{厘米}=(2\times10)\text{毫米}=20\text{毫米}\)\allowbreak \par 平移距离是对应点之间的距离，不是两幅图之间的空隙；原空只填20。\end{multicols}\end{snippet}
% p070_s03
\begin{snippet}{505}{11.5}{14}\begin{multicols}{3}\raggedcolumns \textbf{二、1：}\ “都有”要求每一种三角形都成立，举出一个反例就能否定。不等边三角形通常不能沿任何直线对折重合，而等边三角形才有3条对称轴。因此不能把等边三角形的性质推广到所有三角形。\end{multicols}\end{snippet}
% p070_s04
\begin{snippet}{505}{11.5}{14}\begin{multicols}{3}\raggedcolumns 平移的要求就是每个点同向、等距地移动，图形各部分不能各走各的距离。因此任意两点走过的路径长度一样，判断为正确。与绕中心旋转相比，这里不存在离中心远的点走得更长的问题。\end{multicols}\end{snippet}
% p070_s05
\begin{snippet}{505}{11.5}{14}\begin{multicols}{3}\raggedcolumns 旋转时可把图形看成一块完整的硬纸片，绕中心转动不会拉长边，也不会改变角的大小。改变的是位置或朝向，原来的形状与大小仍然保持。因此说形状会改变是错误的。\end{multicols}\end{snippet}
% p070_s06
\begin{snippet}{505}{11.5}{14}\begin{multicols}{3}\raggedcolumns 对角线是线段，对称轴是直线，原句应判错。正确说法：正方形每条对角线所在的直线都是它的对称轴。\end{multicols}\end{snippet}
% p070_s06_local1
\begin{snippet}{9.800000000000011}{10.5}{12.8}\SourceSans\setlength{\mathindent}{2bp}\(\times\)\end{snippet}
% p070_s07
\begin{snippet}{505}{11.5}{14}\begin{multicols}{3}\raggedcolumns 先把给出的一半关于虚线反射，轴两侧对应端点取等距位置。补成后，上下两底仍互相平行，左右两腰互为镜像而等长，符合等腰梯形的条件。不能只凭左右对称判断，还须确认上下两底平行。\end{multicols}\end{snippet}
% p070_s08
\begin{snippet}{505}{11.5}{14}\begin{multicols}{3}\raggedcolumns 看的是每个完整的组合图，轴必须让所有圆同时重合。同心两圆的任何过圆心直线都满足条件，轴最多；不同心且内切的一大一小两圆，只有连两圆心的那条线能使整体重合，轴最少。不能只看其中一个圆而给所有选项都填无数条。\end{multicols}\end{snippet}
% p070_s09
\begin{snippet}{505}{11.5}{14}\begin{multicols}{3}\raggedcolumns 长方形、圆和半圆都能找出使两侧对折后完全重合的直线；半圆的轴经过圆心且垂直于直径。题图的一般平行四边形虽然两组对边分别平行，但沿直线对折通常不能重合。因此“有平行边”与“轴对称”不是同一个条件。\end{multicols}\end{snippet}
% p070_s10
\begin{snippet}{505}{11.5}{14}\begin{multicols}{3}\raggedcolumns 把折边当成对称轴，将每个选项反射展开，再对照题图的完整外轮廓。A能使两侧的凸起、凹口与题图逐处对应。剪纸会把边缘一起复制，不能只核对中间某一段相似就决定选项。\end{multicols}\end{snippet}
% p070_s11
\begin{snippet}{505}{11.5}{14}\begin{multicols}{3}\raggedcolumns “任意一条边上的高所在的直线”都要成为对称轴，说明三组对应边分别相等，因此是等边三角形。一般等腰三角形只有1条对称轴，不满足“任意一条边”的要求。\end{multicols}\end{snippet}
% p070_s12
\begin{snippet}{505}{11.5}{14}\begin{multicols}{3}\raggedcolumns 三项条件必须同时符合：四条边、没有平行边、恰有一条对称轴。A的两侧边沿中间竖线可以配对，而且没有一组平行边；其他选项分别有平行边或不能对折重合。不能只看它有一条对称轴就忽略另外两个条件。\end{multicols}\end{snippet}
% p070_s13
\begin{snippet}{505}{11.5}{14}\begin{multicols}{3}\raggedcolumns 先把N关于M、N之间的竖直界线作镜像，确定M处应有的轮廓和朝向。再寻找能够只靠平移放入M的选项，A的朝向与所需镜像一致。平移只改变位置，不能把其他朝向的图形翻面或旋转后再说它符合要求。\end{multicols}\end{snippet}
% p071_s01
\begin{snippet}{505}{11.5}{14}\begin{multicols}{3}\raggedcolumns C中三个弧形箭头形状、大小相同，均围绕一个共同中心排列。把一个箭头绕中心转\(120^\circ\)\allowbreak ，能到下一个位置，继续同向旋转就得到全图。检验时既要看外轮廓，也要看箭头指向是否随转动一起对应。\end{multicols}\end{snippet}
% p071_s02
\begin{snippet}{505}{11.5}{14}\begin{multicols}{3}\raggedcolumns \textbf{四、1：}\ O在两次旋转中都固定，先分别保留OA、OB的长度定位两端。顺时针转直角时，左端A到O上方、右端B到O下方；逆时针转时上下方向相反。两种结果都与原线段垂直，但若OA与OB不等，不能把两个端点画成离O一样远。\end{multicols}\end{snippet}
% p071_s03
\begin{snippet}{505}{11.5}{14}\begin{multicols}{3}\raggedcolumns \textbf{四、2（1）：}\ 此图的轴是水平虚线，因此对应点应沿竖直格线寻找，在轴下方取与原顶点同样远的位置。轴上的点保持不动，其余各拐点定位后按原顺序连线。两侧斜边的两个端点都要分别反射，不能只画一条看起来接近的斜线。\end{multicols}\end{snippet}
% p071_s04
\begin{snippet}{505}{11.5}{14}\begin{multicols}{3}\raggedcolumns \textbf{四、2（2）：}\ 先选船尖等一个顶点向下数7格，再将船体、桅杆、旗帜的其余顶点都同样下移7格。按原来的连接顺序重连后，图形形状、大小和朝向都不变。不能只移动船体而把上方桅杆或旗帜留在旧位置。\end{multicols}\end{snippet}
% p071_s05
\begin{snippet}{505}{11.5}{14}\begin{multicols}{3}\raggedcolumns \textbf{四、2（3）：}\ 固定A不动，先把从A指向各顶点的方向逆时针转一个直角。原在A左方的点应到A下方，原在左上方的点应到左下方，各点到A的距离保持不变。定位完三个顶点再连边，不能绕三角形自身中心转。\end{multicols}\end{snippet}
% p071_s06
\begin{snippet}{505}{11.5}{14}\begin{multicols}{3}\raggedcolumns \textbf{四、3：}\ 以O为共同中心，将同一个三角形依次转过\(90^\circ\)\allowbreak 、\(180^\circ\)\allowbreak 、\(270^\circ\)\allowbreak ，补成其余三片风叶。也可连续三次向同一方向转\(90^\circ\)\allowbreak ，所得四片的位置相同。每片的边长、形状都与原图一致，且相邻朝向相差一个直角，不能把其中一片改成镜像。\end{multicols}\end{snippet}
% p071_s07
\begin{snippet}{505}{11.5}{14}\begin{multicols}{3}\raggedcolumns \textbf{四、4：}\ 镜面使钟面左右对应，正上方12不变，右下方4对应到左下方8。两根指针须分别作镜像：长分针仍指12，短时针转到8，共同端点仍在钟心。不能交换长短针，也不能把整个钟面简单转半圈。\end{multicols}\end{snippet}
% p072_s01
\begin{snippet}{505}{11.5}{14}\begin{multicols}{3}\raggedcolumns \textbf{五、1（图②）：}\ 图①显示每走到下一格，就把同一图形顺时针转\(90^\circ\)\allowbreak ，不是沿直线翻折。按左上、右上、右下、左下的顺序，Y形的三条臂整体转动，右上开口朝右，左下开口朝左。各臂长度与交点位置都要保持对应，不能只转开口方向。\end{multicols}\end{snippet}
% p072_s02
\begin{snippet}{505}{11.5}{14}\begin{multicols}{3}\raggedcolumns \textbf{五、1（图③）：}\ 沿用每次顺时针转一个直角的规律，F形从左上到右下共转\(180^\circ\)\allowbreak ，到左下共转\(270^\circ\)\allowbreak 。长竖笔、顶横笔、中横笔各自随图一起转动，长短关系保持。F不左右对称，若只翻面会得到不同图形，要特别核对中横笔的位置。\end{multicols}\end{snippet}
% p072_s03
\begin{snippet}{505}{11.5}{14}\begin{multicols}{3}\raggedcolumns 第一幅两鱼沿中线对折能重合，属于对称；第二幅同样的动物绕共同中心改变朝向，属于旋转。第三幅动物保持同样朝向，只在不同位置重复出现，属于平移。区分时同时看位置和朝向，不能把所有重复图案都归为同一种现象。\end{multicols}\end{snippet}
% p072_s04
\begin{snippet}{505}{11.5}{14}\begin{multicols}{3}\raggedcolumns \textbf{五、3（1）：}\ 先确定通过花圃中间横边的水平格线为轴，再把现有下半部各拐点沿竖直方向反射到上方。每对点距轴的格数相等，连线顺序与原轮廓相同。完成后检查上下能否对折重合，房屋图中的每格2米比例不改变。\end{multicols}\end{snippet}
% p072_s05
\begin{snippet}{505}{11.5}{14}\begin{multicols}{3}\raggedcolumns \textbf{五、3（2）：}\ 图中每格代表2米，实际距离要先除以2，得到下移6格、右移2格。再按题目顺序让厨房的每个顶点作相同位移，最后连接并标名称。不能直接下移12格、右移4格，否则实际移动距离会加倍。\end{multicols}\end{snippet}
% p072_s06
\begin{snippet}{505}{11.5}{14}\begin{multicols}{3}\raggedcolumns \textbf{五、3（3）：}\ C是原车库左下角，旋转时该点不动。逆时针转直角后，原来向上的长边转为向左，原来向右的短边转为向上，所以新车库位于C左上方。边长不变，只是长短边朝向互换，画成横3格、竖2格后再标名称。\end{multicols}\end{snippet}
% p072_s07
\begin{snippet}{505}{11.5}{14}\begin{multicols}{3}\raggedcolumns 先尝试竖直、水平和斜向的轴，查看原有涂色格哪些位置缺少对应格；每种方法都只能补一格。按网格左上角起数，三种补位分别是第4行第3列、第2行第5列、第4行第6列，对应竖轴、横轴、斜轴。补完逐格检查，反射后每个涂色格都要有伙伴，不能只让外形看起来相似；三种位置互不重复。\end{multicols}\end{snippet}
% p073_s01
\begin{snippet}{505}{11.5}{14}\begin{multicols}{3}\raggedcolumns 以112个为基准，105个少7个，119个多7个，多、少正好抵消。把119个中的7个移给105个，三人就同为112个，总数不变，所以平均数是112个。\[\begin{aligned}&112+(7-7)\div3\\={}&112+0\div3\\={}&112+0\\={}&112\text{(个)}\end{aligned}\]\end{multicols}\end{snippet}
% p073_s02
\begin{snippet}{505}{11.5}{14}\begin{multicols}{3}\raggedcolumns 甲瓶比乙瓶多100克，从甲倒给乙时，甲减少多少，乙就同时增加多少。因此只需移差的一半50克，两瓶便相等，不是把全部差100克都倒过去。再把总质量平均分成两份，可核对每瓶都是450克。\[\begin{aligned}&(500-400)\div2\\={}&100\div2\\={}&50\text{(克，倒出量)}\end{aligned}\]\[\begin{aligned}&(500+400)\div2\\={}&900\div2\\={}&450\text{(克，每瓶)}\end{aligned}\]\end{multicols}\end{snippet}
% p073_s03
\begin{snippet}{505}{11.5}{14}\begin{multicols}{3}\raggedcolumns 求四次观测温度的平均数，除数是4。用加法交换律、结合律，把8与24、15与17分别配成32；两个相同的和更易算。这里估计的是题目规定的四时刻平均温度。\[\begin{aligned}&(8+15+24+17)\div4\\={}&[(8+24)+(15+17)]\div4\\={}&[32+32]\div4\\={}&64\div4\\={}&16\text{(摄氏度)}\end{aligned}\]\end{multicols}\end{snippet}
% p073_s04
\begin{snippet}{505}{11.5}{14}\begin{multicols}{3}\raggedcolumns 平均数乘个数还原三个数的总和，再减去已知两数。连续减两个数，等于减去它们的和；先把24与16凑成40。乘法与括号内加法互相独立，在同一步算出。\[\begin{aligned}&30\times3-24-16\\={}&30\times3-(24+16)\\={}&90-40\\={}&50\end{aligned}\]\end{multicols}\end{snippet}
% p073_s05
\begin{snippet}{505}{11.5}{14}\begin{multicols}{3}\raggedcolumns 两队参加人数不同，比较总投中数会受到人数影响，应比较平均每人投中数。分别把102个分给6人、112个分给7人，再比较17个与16个。乙队总数虽多，平均每人的成绩却较少，因此要按人均水平判断。\[\begin{aligned}&102\div6\\={}&17\text{(个，甲队)}\\&112\div7\\={}&16\text{(个，乙队)}\end{aligned}\]\end{multicols}\end{snippet}
% p073_s06
\begin{snippet}{505}{11.5}{14}\begin{multicols}{3}\raggedcolumns 题目问“平均每天”，应先把上、中、下旬销量合并成全月总销量，再除以五月的31天。三个旬是三个记录时段，并不是3天，且下旬比前两旬多1天。不能将三旬销量相加后直接除以3当作每天销量。\[\begin{aligned}&(98+125+149)\div31\\={}&(223+149)\div31\\={}&372\div31\\={}&12\text{(部)}\end{aligned}\]\end{multicols}\end{snippet}
% p073_s07
\begin{snippet}{505}{11.5}{14}\begin{multicols}{3}\raggedcolumns 以80为基准，甲、乙的总量正好是2份80；若丙也为80，三数平均仍为80。丙实际少30，应把这30平均从3份中扣去，所以先用30除以3得10，再从80减10，避免先求210。\[\begin{aligned}&80-(80-50)\div3\\={}&80-30\div3\\={}&80-10\\={}&70\end{aligned}\]原“先求总量再平均”的数量思路也成立，但不能把代表两份的80和代表一份的50直接平均。\end{multicols}\end{snippet}
% p073_s08
\begin{snippet}{505}{11.5}{14}\begin{multicols}{3}\raggedcolumns \textbf{一、8：}\ 两根原长相加时，重叠4厘米被算了两遍，实际总长中只能算一遍。因此先求两根原长之和，再减去一次重叠长。也可保留第一根20厘米，只加第二根未重叠的16厘米，得到相同结果。\[\begin{aligned}&20\times2-4\\={}&40-4\\={}&36\text{(厘米)}\end{aligned}\]\end{multicols}\end{snippet}
% p073_s09
\begin{snippet}{505}{11.5}{14}\begin{multicols}{3}\raggedcolumns \textbf{一、9：}\ 设想把高条多出的部分移给矮条，直到三条高度一样，总高度保持不变。所选平均线应使线上多出的总长度与线下不足的总长度大致相等，图中③最接近这个平衡。不能只选中间那个人的身高，也不能只看某一条与虚线的距离。\end{multicols}\end{snippet}
% p073_s10
\begin{snippet}{505}{11.5}{14}\begin{multicols}{3}\raggedcolumns \textbf{一、10：}\ 填\(26\)\allowbreak 。简便思路：留下3个数平均各比20少2，去掉的数应比20多3份2。\[\begin{aligned}&20+(20-18)\times3\\={}&20+2\times3\\={}&20+6\\={}&26\end{aligned}\]另可用总和之差，同时计算两个乘法：\[\begin{aligned}&20\times4-18\times3\\={}&80-54\\={}&26\end{aligned}\]\end{multicols}\end{snippet}
% p074_s01
\begin{snippet}{505}{11.5}{14}\begin{multicols}{3}\raggedcolumns 平均数把总量平均分给全班，不能小于人人至少读的28本，也不能大于任何人都不超过的42本。先用这个范围排除25、45、20，只留下35。本题问“可能”，只需判断候选值是否落在合理范围内。\end{multicols}\end{snippet}
% p074_s02
\begin{snippet}{505}{11.5}{14}\begin{multicols}{3}\raggedcolumns 女生平均55下高于男生平均45下，能说明本次女生的整体水平较好。平均数不表示每个人都达到同一成绩，因此不能推出每个男生都低于女生，也不能排除某男生高于55下。全班平均还要知道两类人数，不能在人数未知时直接平均45与55。\end{multicols}\end{snippet}
% p074_s03
\begin{snippet}{505}{11.5}{14}\begin{multicols}{3}\raggedcolumns 书法与机器人两项人数相加时，同时参加的16人被计了两次，须减掉一次，得到至少参加一项的45人。只有题目另约定全班人人至少参加一项时，这45人才等于全班人数，可选C。若有人两项都不参加，还应另加这部分人数，题面现有条件不能无条件确定全班总人数。\[\begin{aligned}&29+32-16\\={}&61-16\\={}&45\text{(人，至少参加一项)}\end{aligned}\]\end{multicols}\end{snippet}
% p074_s04
\begin{snippet}{505}{11.5}{14}\begin{multicols}{3}\raggedcolumns 按当前题面15千克，\(39>15\)\allowbreak ，加入后平均体重提高，应选C。书后A与当前数据不符；若原平均为45千克，\(39<45\)\allowbreak ，加入后就会变轻，选A。判断增减只需比较新队员体重与原平均，不必先算总量。\end{multicols}\end{snippet}
% p074_s05
\begin{snippet}{505}{11.5}{14}\begin{multicols}{3}\raggedcolumns “平均每台500多元”应理解为大于500元、小于600元。6台总价也应大于3000元、小于3600元，将四个选项放进这个范围检查，只有3300元符合。反过来除以6，平均550元也符合“500多元”。\end{multicols}\end{snippet}
% p074_s06
\begin{snippet}{505}{11.5}{14}\begin{multicols}{3}\raggedcolumns 第三条终点正落在平均线上，所以平均长度就是14厘米。第二条比平均多6厘米，移多补少时这6厘米应补给第一条，第一条就比14厘米少6厘米。不能把20与14直接平均，因为实际共有三条丝带。\[\begin{aligned}&14-(20-14)\\={}&14-6\\={}&8\text{(厘米)}\end{aligned}\]\end{multicols}\end{snippet}
% p075_s01
\begin{snippet}{505}{11.5}{14}\begin{multicols}{3}\raggedcolumns \textbf{方法一：分组凑整求总和}\par 用加法交换律、结合律，把70与130、240与60配对。互相独立的分组在同一步计算。\end{multicols}\[\begin{aligned}&(70+120+170+240+190+130+60)\div7\\={}&[(70+130)+(240+60)+(120+170+190)]\div7\\={}&[200+300+(290+190)]\div7\\={}&[500+480]\div7\\={}&980\div7\\={}&140\text{(厘米)}\end{aligned}\]\begin{multicols}{3}\raggedcolumns 答：这条河的平均水深是140厘米。\par\textbf{方法二：以100厘米为基准移多补少}\par 多出20、70、140、90、30，少30、40；利用加减运算的性质，把100与160、90减70分别括成一组。\end{multicols}\[\begin{aligned}&100+[(70+30)+(20+140)+90-(30+40)]\div7\\={}&100+[100+160+90-70]\div7\\={}&100+[(100+160)+(90-70)]\div7\\={}&100+[260+20]\div7\\={}&100+280\div7\\={}&100+40\\={}&140\text{(厘米)}\end{aligned}\]\begin{multicols}{3}\raggedcolumns 答：这条河的平均水深是140厘米。\end{multicols}\end{snippet}
% p075_s02
\begin{snippet}{505}{11.5}{14}\begin{multicols}{3}\raggedcolumns 有危险。最大水深240厘米> 200厘米，深处超过它的身高。\par 平均数只反映七处数据的整体水平，不表示每一处都同样深。判断能否过河，应检查各处水深是否超过身高，尤其不能忽略最深处，不能因为平均水深较浅就断定处处安全。\end{multicols}\end{snippet}
% p075_s02_local1
\begin{snippet}{505}{11}{14}\SourceSans\setlength{\mathindent}{2bp}\(2\text{米}=200\text{厘米}\)\par 有危险。最大水深\(240\text{厘米}>200\text{厘米}\)，深处超过它的身高。\end{snippet}
% p075_s03
\begin{snippet}{505}{11.5}{14}\begin{multicols}{3}\raggedcolumns \textbf{总路程括成整体，再除以总时间}\par 来回总路程按每程720米乘2程列式，并用括号表示整体。简算时利用\(20=2\times10\)\allowbreak 和乘除运算的性质，先算\(2\div2\)\allowbreak ，再除以10。\[\begin{aligned}&(720\times2)\div(12+8)\\={}&(720\times2)\div20\\={}&720\times(2\div2)\div10\\={}&720\times1\div10\\={}&720\div10\\={}&72\text{(米/分)}\end{aligned}\]答：来回平均每分钟走72米。\end{multicols}\end{snippet}
% p075_s04
\begin{snippet}{505}{11.5}{14}\begin{multicols}{3}\raggedcolumns 原位主式以85分为基准：87多2、83少2，88多3、82少3，两对偏差正好抵消；只剩97多12与80少5，净多7分平均补给7份。\par\textbf{另法：分组求总分再平均}\par 用加法交换律、结合律分组，再用总分除以7；各组互不依赖的计算同一步完成。\end{multicols}\[\begin{aligned}&(80+87+83+85+88+97+82)\div7\\={}&[(87+83)+(88+82)+(80+85+97)]\div7\\={}&[170+170+(165+97)]\div7\\={}&[340+262]\div7\\={}&602\div7\\={}&86\text{(分)}\end{aligned}\]\begin{multicols}{3}\raggedcolumns 答：七位裁判评分的平均分是86分。\end{multicols}\end{snippet}
% p075_s04_local1
\begin{snippet}{189}{9.0}{12.8}\SourceSans\setlength{\mathindent}{2bp}\[\begin{aligned}&85+(12-5)\div7\\={}&85+7\div7\\={}&85+1\\={}&86\text{(分)}\end{aligned}\]\end{snippet}
% p075_s05
\begin{snippet}{505}{11.5}{14}\begin{multicols}{3}\raggedcolumns \textbf{方法一：从总分扣除最高、最低分}\par 去掉2个评分，总分与份数都要相应减少；分子先算602减97，分母同时算7减2。\[\begin{aligned}&(602-97-80)\div(7-2)\\={}&(505-80)\div5\\={}&425\div5\\={}&85\text{(分)}\end{aligned}\]答：去掉一个最高分和一个最低分后，平均分是85分。\par\textbf{方法二：移多补少}\par 从87分移2分给83分，从88分移3分给82分，原来的85分不动，五个评分都成85分。按每份85分、5份求平均，利用乘除运算性质先算5除以5。\[\begin{aligned}&85\times5\div5\\={}&85\times(5\div5)\\={}&85\times1\\={}&85\text{(分)}\end{aligned}\]答：平均分是85分。\end{multicols}\end{snippet}
% p075_s06
\begin{snippet}{505}{11.5}{14}\begin{multicols}{3}\raggedcolumns 第二个结果较合理：去掉一个最高分和一个最低分，可减小极端评分对结果的影响。\par 第一个平均数采用7个评分，第二个采用去掉两端评分后剩下的5个，数据范围与份数必须相对应。两者都是各自数据的正确平均数，选择哪一个作代表值取决于采用的评分规则。这里比较的是结果的代表性，不能把未选作代表值的那个平均数说成计算错误。\end{multicols}\end{snippet}
% p076_s01
\begin{snippet}{505}{11.5}{14}\begin{multicols}{3}\raggedcolumns 每名学生有3位老师评分，必须沿同一学生所在列取数。补未知评分可用移多补少，求平均可先配对凑整或取基准数；以下按表中五个空依次说明。\par 四（1）班4号：另外两次评分比平均67分别少2、1，李老师应比67多3；利用加法结合律先合并2与1。\[\begin{aligned}&67+(67-65)+(67-66)\\={}&67+2+1\\={}&67+(2+1)\\={}&67+3\\={}&70\text{(分)}\end{aligned}\]四（1）班3号：以95分为基准，96多1分，94少1分，移多补少恰好抵消。\end{multicols}\[\begin{aligned}&95+[(96-95)-(95-94)]\div3\\={}&95+[1-1]\div3\\={}&95+0\div3\\={}&95+0\\={}&95\text{(分)}\end{aligned}\]\begin{multicols}{3}\raggedcolumns 四（2）班5号：88与92先凑180，原式已易算，保持。\[\begin{aligned}&90\times3-(88+92)\\={}&270-180\\={}&90\text{(分)}\end{aligned}\]四（2）班3号：一次95恰等于平均，另一次93少2，刘老师应比95多2。\[\begin{aligned}&95+(95-93)\\={}&95+2\\={}&97\text{(分)}\end{aligned}\]四（2）班4号：以83为基准，86多出的3分平均分给3份。\[\begin{aligned}&83+(86-83)\div3\\={}&83+3\div3\\={}&83+1\\={}&84\text{(分)}\end{aligned}\]填完后分别求该列三次评分的平均，检查与表中平均分相符。\end{multicols}\end{snippet}
% p076_s02
\begin{snippet}{505}{11.5}{14}\begin{multicols}{3}\raggedcolumns 统计对象是每名学生最后的平均得分，而不是把三位老师的评分各算一人。按等级界限依次归类，每人只记一次：92分进“优”，68分进“达标”。为使分段互不重叠，末组“68分以下”在本表按不足68分处理，填后每班四类人数之和都应为6。\end{multicols}\end{snippet}
% p076_s03
\begin{snippet}{505}{11.5}{14}\begin{multicols}{3}\raggedcolumns 两班优秀人数相同，都是2人；四（2）班良好人数更多，且没有待达标学生，整体情况较好。\par 两班各抽查6人，可直接比较同等级人数。判断整体情况应综合各等级分布，不能只凭优秀人数相同就认定两班整体相同。结论限于本次抽查，不能据此断定每名学生的情况，也不能直接推及未被抽查的学生。\end{multicols}\end{snippet}
% p077_s01
\begin{snippet}{505}{11.5}{14}\begin{multicols}{3}\raggedcolumns 第7天准备量取前6天的平均，除数是6。先配成三对34块，再把6拆成3乘2，利用乘除运算的性质先算3除以3，避免先求总量102。\end{multicols}\[\begin{aligned}&(16+15+20+18+14+19)\div6\\={}&[(16+18)+(15+19)+(20+14)]\div6\\={}&[34+34+34]\div6\\={}&[34\times3]\div(3\times2)\\={}&34\times(3\div3)\div2\\={}&34\times1\div2\\={}&34\div2\\={}&17\text{(块)}\end{aligned}\]\end{snippet}
% p077_s02
\begin{snippet}{505}{11.5}{14}\begin{multicols}{3}\raggedcolumns 第一季度给总量，第二季度给月平均量，先按每月145台乘3个月求第二季度总销量，再求半年总量。\[\begin{aligned}&237+145\times3\\={}&237+435\\={}&672\text{(台)}\end{aligned}\]两个季度各3个月，半年总量除以总月数，不能把237与145直接平均。\[\begin{aligned}&672\div(3+3)\\={}&672\div6\\={}&112\text{(台)}\end{aligned}\]\end{multicols}\end{snippet}
% p077_s03
\begin{snippet}{505}{11.5}{14}\begin{multicols}{3}\raggedcolumns 连续五个自然数，中间数两边各有两个数，左右对称的两数相加都等于中间数的2倍。移多补少后五份都等于中间数，因此平均数26就是中间的那个数。最小与最大分别离中间数2，不是离中间数5。\end{multicols}\end{snippet}
% p077_s04
\begin{snippet}{505}{11.5}{14}\begin{multicols}{3}\raggedcolumns 填\(14\)\allowbreak 。已知两个数比平均12合计少2，第三个数应比12多2。\[\begin{aligned}&12+(12-11)\times2\\={}&12+1\times2\\={}&12+2\\={}&14\end{aligned}\]另可用总和之差，同时计算两个乘法：\[\begin{aligned}&12\times3-11\times2\\={}&36-22\\={}&14\end{aligned}\]\end{multicols}\end{snippet}
% p077_s05
\begin{snippet}{505}{11.5}{14}\begin{multicols}{3}\raggedcolumns 填\(52\)\allowbreak 。把两人相差的8次分成两半，较多的一方减4就成为平均数。\[\begin{aligned}&56-8\div2\\={}&56-4\\={}&52\text{(次)}\end{aligned}\]另一方法先求总次数再平均：\[\begin{aligned}&[56+(56-8)]\div2\\={}&[56+48]\div2\\={}&104\div2\\={}&52\text{(次)}\end{aligned}\]\end{multicols}\end{snippet}
% p077_s06
\begin{snippet}{505}{11.5}{14}\begin{multicols}{3}\raggedcolumns 以155厘米为基准：156多1、154少1，151少4、159多4，两对多、少部分分别抵消。只剩160比155多5厘米，平均补到5人，每人多1厘米。\[\begin{aligned}&155+(160-155)\div5\\={}&155+5\div5\\={}&155+1\\={}&156\text{(厘米)}\end{aligned}\]\par\textbf{另法：分组求总身高再平均}\par 两对身高分别合为310厘米，与160厘米合计后除以5。\end{multicols}\[\begin{aligned}&(156+154+151+159+160)\div5\\={}&[(156+154)+(151+159)+160]\div5\\={}&[310+310+160]\div5\\={}&[620+160]\div5\\={}&780\div5\\={}&156\text{(厘米)}\end{aligned}\]\begin{multicols}{3}\raggedcolumns 平均数反映整体水平，不表示每位队员都高156厘米。\end{multicols}\end{snippet}
% p077_s07
\begin{snippet}{505}{11.5}{14}\begin{multicols}{3}\raggedcolumns B队平均157厘米、共5人，总身高是5份157厘米。已知四人身高先用150与160凑310，159与157合计316；总身高的乘法与两个分组相互独立，同一步算出。剩下的才是5号身高；填回可验平均仍为157厘米。\end{multicols}\[\begin{aligned}&157\times5-(150+159+157+160)\\={}&157\times5-[(150+160)+(159+157)]\\={}&785-[310+316]\\={}&785-626\\={}&159\text{(厘米)}\end{aligned}\]\end{snippet}
% p077_s08
\begin{snippet}{505}{11.5}{14}\begin{multicols}{3}\raggedcolumns 交换的是A队154厘米与B队159厘米的队员，人数仍各5人。总身高的增减先平均分给5人，再分别加到或从原平均身高中减去。\par A队：\[\begin{aligned}&156+(159-154)\div5\\={}&156+5\div5\\={}&156+1\\={}&157\text{(厘米)}\end{aligned}\]B队：\[\begin{aligned}&157-(159-154)\div5\\={}&157-5\div5\\={}&157-1\\={}&156\text{(厘米)}\end{aligned}\]交换后157厘米大于156厘米，填A。\end{multicols}\end{snippet}
% p078_s01
\begin{snippet}{505}{11.5}{14}\begin{multicols}{3}\raggedcolumns 三人总身高为\(136\times3=408\)\allowbreak 厘米；若3人全能乘坐，即使每人只按140厘米计，也至少需\(140\times3=420\)\allowbreak 厘米。\(408<420\)\allowbreak ，故3人全能坐不可能。\((136,\allowbreak 136,\allowbreak 136)\)\allowbreak 、\((146,\allowbreak 131,\allowbreak 131)\)\allowbreak 、\((146,\allowbreak 146,\allowbreak 116)\)\allowbreak 三组身高总和均为408厘米，分别实现0、1、2人能坐。\end{multicols}\end{snippet}
% p078_s01_local1
\begin{snippet}{80}{12}{12.8}\SourceSans\setlength{\mathindent}{2bp}{\color{black}\(\mathrm{136}\)}\end{snippet}
% p078_s01_local2
\begin{snippet}{80}{12}{12.8}\SourceSans\setlength{\mathindent}{2bp}{\color{black}\(\mathrm{140}\)}\end{snippet}
% p078_s01_local3
\begin{snippet}{200}{12}{12.8}\SourceSans\setlength{\mathindent}{2bp}{\color{black}（每题2分，共12分）}\end{snippet}
% p078_s02
\begin{snippet}{505}{11.5}{14}\begin{multicols}{3}\raggedcolumns 45名报名者中有10人两项都没参加，所以真正至少参加一项的只有35人。两项名单人数相加为38，比分开去重后的35多出的3人，就是被重复计一次的两项都参加者。不能直接从38减45，否则遗漏了未参加的10人。\[\begin{aligned}&21+17-(45-10)\\={}&38-35\\={}&3\text{(人)}\end{aligned}\]\end{multicols}\end{snippet}
% p078_s03
\begin{snippet}{505}{11.5}{14}\begin{multicols}{3}\raggedcolumns 最大数不能小于平均数45，其余两数均严格大于40，所以最大数要严格小于下面所得上界：\[\begin{aligned}&45\times3-(40+40)\\={}&135-80\\={}&55\end{aligned}\]只有49落在45至55之间，且可由49、43、43实现。40不包括在“大于40”内。\end{multicols}\end{snippet}
% p078_s04
\begin{snippet}{505}{11.5}{14}\begin{multicols}{3}\raggedcolumns 甲比平均年龄多6岁，乙比平均年龄少6岁。把甲多出的6岁补给乙，两人都等于平均年龄，丙也必须等于平均年龄，因此选B。\end{multicols}\end{snippet}
% p078_s04_local1
\begin{snippet}{22}{11.5}{12.8}\SourceSans\setlength{\mathindent}{2bp}\(\mathrm{B}\)\end{snippet}
% p078_s05
\begin{snippet}{505}{11.5}{14}\begin{multicols}{3}\raggedcolumns 以22摄氏度为基准比较四点的左右偏差。第二天高出的部分约能补足低下的部分，移多补少后平均约为22摄氏度，选B。第一天偏高，第三、第四天偏低；不能只凭某一个点接近22摄氏度来判断。\end{multicols}\end{snippet}
% p078_s05_local1
\begin{snippet}{200}{10.5}{12.8}\SourceSans\setlength{\mathindent}{2bp}{\color{black}\input{p078_update/assets/p078_temperatures_photo.tikz}}\end{snippet}
% p078_s06
\begin{snippet}{505}{11.5}{14}\begin{multicols}{3}\raggedcolumns 94分是两项的平均分，要先乘2还原秧歌、剪纸两项总分。三项总分281分减去这188分，余下才是草编分数。不能从281直接减94，因为94只是一份平均水平，并非原两项总分。\[\begin{aligned}&281-94\times2\\={}&281-188\\={}&93\text{(分)}\end{aligned}\]\end{multicols}\end{snippet}
% p078_s06_local1
\begin{snippet}{80}{12}{12.8}\SourceSans\setlength{\mathindent}{2bp}{\color{black}秧歌}\end{snippet}
% p079_s01
\begin{snippet}{505}{11.5}{14}\begin{multicols}{3}\raggedcolumns \textbf{方法一：总额除以组数}\par 问的是“平均每组”，把总捐款平均分成3份。\[\begin{aligned}&(178+224+198)\div3\\={}&(402+198)\div3\\={}&600\div3\\={}&200\text{(元)}\end{aligned}\]答：平均每组捐款200元。\par \textbf{方法二：移多补少}\par 从第二组的224元中移22元给第一组，再移2元给第三组；总钱数不变。三组移补后的金额分别为：\[178+22=200\text{(元)}\]\[\begin{aligned}&224-22-2\\={}&202-2\\={}&200\text{(元)}\end{aligned}\]\[198+2=200\text{(元)}\]三组恰好一样多。答：平均每组捐款200元。\end{multicols}\end{snippet}
% p079_s02
\begin{snippet}{505}{11.5}{14}\begin{multicols}{3}\raggedcolumns \textbf{求总金额和总人数}\par 问的是“平均每人”，用总捐款除以总人数。三组人数不同，不能把各组人均捐款直接平均。总金额与总人数的两条计算支路在每一步同时计算。\end{multicols}\[\begin{aligned}&(178+224+198)\div(9+11+10)\\={}&(402+198)\div(20+10)\\={}&600\div30\\={}&20\text{(元)}\end{aligned}\]\begin{multicols}{3}\raggedcolumns 答：平均每人捐款20元。\end{multicols}\end{snippet}
% p079_s03
\begin{snippet}{505}{11.5}{14}\begin{multicols}{3}\raggedcolumns \textbf{方法一：分组凑整后求平均}\end{multicols}\[\begin{aligned}&(123+145+156+172+184)\div5\\={}&[(156+184)+(123+145+172)]\div5\\={}&[340+(268+172)]\div5\\={}&[340+440]\div5\\={}&780\div5\\={}&156\text{(克)}\end{aligned}\]\begin{multicols}{3}\raggedcolumns 答：抽出的土豆平均每个重156克。\par \textbf{方法二：以150克为基准移多补少}\par 多出6、22、34克，不足27、5克；先把6与34凑40；不足的27与5合为32。\end{multicols}\[\begin{aligned}&150+[(6+34)+22-(27+5)]\div5\\={}&150+[40+22-32]\div5\\={}&150+[62-32]\div5\\={}&150+30\div5\\={}&150+6\\={}&156\text{(克)}\end{aligned}\]\begin{multicols}{3}\raggedcolumns 答：平均每个重156克。\end{multicols}\end{snippet}
% p079_s04
\begin{snippet}{505}{11.5}{14}\begin{multicols}{3}\raggedcolumns 上一问只称了5个土豆，用这5个的平均质量156克代表整袋的大致水平。再按“每个平均质量×40个”估计总质量，最后把克按1000的进率化成千克。因为是抽样估计，不能说每个土豆都重156克，也不能把所得估计数当作整袋实称的精确值。\end{multicols}\end{snippet}
% p079_s04_local1
\begin{snippet}{80}{9.0}{14}\SourceSans\setlength{\mathindent}{2bp}\(\text{(克)}\)\end{snippet}
% p079_s04_local2
\begin{snippet}{505}{9}{14}\SourceSans\setlength{\mathindent}{2bp}\(6240\text{克}=(6240\div1000)\text{千克}=6.24\text{千克}\)\end{snippet}
% p079_s05
\begin{snippet}{505}{11.5}{14}\begin{multicols}{3}\raggedcolumns 原方法先用89乘4还原前四次总分，加入第五次94分后，再除以5。另法把第五次也先看成89分，此时五份都按89分平均；实际第五次多出的5分，应补给全部5次，而不是只分给原四次：\[\begin{aligned}&89+(94-89)\div5\\={}&89+5\div5\\={}&89+1\\={}&90\text{(分)}\end{aligned}\]\end{multicols}\end{snippet}
% p079_s05_local1
\begin{snippet}{80}{9.0}{14}\SourceSans\setlength{\mathindent}{2bp}\(\text{(分)}\)\end{snippet}
% p080_s01
\begin{snippet}{505}{11.5}{14}\begin{multicols}{3}\raggedcolumns 先将共同出现的王晴、赵磊、周敏放在重叠区，另外两边分别放各自只参加一项的人，每人只写一次。方法一，两张名单相加时共同的3人被计了两次，要减掉一次重复：\[\begin{aligned}&6+5-3\\={}&11-3\\={}&8\text{(人)}\end{aligned}\]方法二，分别数“只参加环保、两项都参加、只参加编程”三块，人数互不重复，可以直接相加。两侧单独区域须先从各自名单人数里减掉共同的3人：\[\begin{aligned}&(6-3)+3+(5-3)\\={}&3+3+2\\={}&6+2\\={}&8\text{(人)}\end{aligned}\]答：一共有8人参加活动。\end{multicols}\end{snippet}
% p080_s02
\begin{snippet}{505}{11.5}{14}\begin{multicols}{3}\raggedcolumns 先按性别分别划记，再按表头四区间归类，每人只记一次。边界4.2归第一栏，4.3和4.6归第二栏，4.7和4.9归第三栏，5.0归最后一栏，避免漏记或重记。男生先将2、3、5凑成10，再加11得21人；女生将8与12先凑20：\[\begin{aligned}&1+2+8+12\\={}&(1+2)+(8+12)\\={}&3+20\\={}&23\text{(人)}\end{aligned}\]纵向男女相加得到各栏总人数，横向两行合计为44人，两种方向应相符。\end{multicols}\end{snippet}
% p080_s03
\begin{snippet}{505}{11.5}{14}\begin{multicols}{3}\raggedcolumns 求全班人数，可把男生21人与女生23人相加，也可把“总计”行四个视力段人数相加，两种结果都应为44。判断哪个范围人数最多，要比较全班总计的四个数3、5、13、23。不能只看男生或女生单独一行，也不能把最大的视力数值误当成人数最多的范围。\end{multicols}\end{snippet}
% p080_s04
\begin{snippet}{505}{11.5}{14}\begin{multicols}{3}\raggedcolumns 要注意用眼卫生，保持正确读写姿势，保证户外活动时间，并定期检查视力。\par 题目把5.0及以上规定为正常，因此低于5.0只包括前三栏，边界5.0不能算进去。可用全班44人减正常23人，也可将前三栏3、5、13相加，两种方法都得到21人。两种计数方法应得到相同人数，可互相检查是否漏记边界。\end{multicols}\end{snippet}
% p081_s01
\begin{snippet}{505}{11.5}{14}\begin{multicols}{3}\raggedcolumns 第一空用12根的总价除以12，求每根跳绳的单价。买90根的总价再用“单价×根数”求，最后以带的钱减应付总价，得到找零。2600元不是90根的总价，不能直接除以90求单价。\[\begin{aligned}&300\div12\\={}&25\text{(元，每根)}\end{aligned}\]\[\begin{aligned}&2600-25\times90\\={}&2600-2250\\={}&350\text{(元)}\end{aligned}\]\end{multicols}\end{snippet}
% p081_s02
\begin{snippet}{505}{11.5}{14}\begin{multicols}{3}\raggedcolumns 把16看成\(4\times4\)\allowbreak ，其中一个4与括号内的“除以4”抵消，整个原式相当于所填数乘4。因此用512除以4，直接求出所填数。这是整数乘除关系的整理，不用小数乘法。\[\begin{aligned}&512\div(16\div4)\\={}&512\div4\\={}&128\end{aligned}\]代回检验：\(16\times(128\div4)=16\times32=512\)\allowbreak 。\end{multicols}\end{snippet}
% p081_s03
\begin{snippet}{505}{11.5}{14}\begin{multicols}{3}\raggedcolumns 3个十写在十位，6个十分之一写在十分位，9个千分之一写在千分位，个位与百分位都没有单位须用0占位。读小数部分时按顺序逐字读，百分位的0也要读出。保留两位小数应看千分位9并向百分位进1，不能直接把末尾9删去。\end{multicols}\end{snippet}
% p081_s04
\begin{snippet}{505}{11.5}{14}\begin{multicols}{3}\raggedcolumns 改写成亿作单位，只是换一种表示方式，应把原数除以100000000，大小不变，用等号。省略亿位后的尾数则是在取近似数，应看千万位3，小于5就舍去，用约等号。两个空的目的不同，不能都填整数或都用约等号。\[331000000=3.31\text{亿}\]\[331000000\approx3\text{亿}\]\end{multicols}\end{snippet}
% p081_s05
\begin{snippet}{505}{11.5}{14}\begin{multicols}{3}\raggedcolumns 正确算法先减25再乘12，相当于从12份原数中减去12份25；错算只减去了一份25。共同的“原数×12”部分相同，不用知道三角形代表多少。比较两式实际减去的300与25，错式少减275，所以得数比正确结果多275。\[\begin{aligned}&25\times12-25\\={}&300-25\\={}&275\end{aligned}\]\end{multicols}\end{snippet}
% p081_s06
\begin{snippet}{505}{11.5}{14}\begin{multicols}{3}\raggedcolumns 每分钟跑240米，6分钟就是6份240米，先用“速度×时间”求已跑路程。再用1500米全程减去这段路程，余下才是离终点的距离。求完应为正数且比全程短，不能把1440米直接填入剩余距离。\[\begin{aligned}&1500-240\times6\\={}&1500-1440\\={}&60\text{(米)}\end{aligned}\]\end{multicols}\end{snippet}
% p081_s07
\begin{snippet}{505}{11.5}{14}\begin{multicols}{3}\raggedcolumns \textbf{一、7（面积）：}\ 相邻面积单位的进率是100，1平方米包含100个平方分米。大单位化小单位应乘100，即把520扩大100倍。不能套用米与分米之间的进率10，面积换算要考虑两个方向。\par\noindent\(\displaystyle 520\text{平方米}=52000\text{平方分米}\)\allowbreak \par\end{multicols}\end{snippet}
% p081_s08
\begin{snippet}{505}{11.5}{14}\begin{multicols}{3}\raggedcolumns \textbf{一、7（钱数）：}\ 将角和分分别化成元：\par\noindent\(\displaystyle 1\text{角}=0.1\text{元}\)\allowbreak \par \noindent\(\displaystyle 1\text{分}=0.01\text{元}\)\allowbreak \par 把587元、\(0.3\)\allowbreak 元、\(0.05\)\allowbreak 元合并，角对应十分位、分对应百分位。\par\noindent\(\displaystyle 587\text{元}3\text{角}5\text{分}=587.35\text{元}\)\allowbreak \par 不能把3角5分写成\(0.035\)\allowbreak 元，或把元数也跟着缩小。\end{multicols}\end{snippet}
% p081_s09
\begin{snippet}{505}{11.5}{14}\begin{multicols}{3}\raggedcolumns \textbf{一、7（质量一）：}\ \par\noindent\(\displaystyle 1000\text{克}=1\text{千克}\)\allowbreak \par 先从6260克中取出6个1000克。剩下260克不足1千克，作为第二个单位的零头保留。\par\noindent\(\displaystyle 6260\text{克}=6\text{千克}260\text{克}\)\allowbreak \par 因此两空分别写整千克数与余下克数，不能把余数再当成小数千克填在克的位置。\end{multicols}\end{snippet}
% p081_s10
\begin{snippet}{505}{11.5}{14}\begin{multicols}{3}\raggedcolumns \textbf{一、7（长度）：}\ \par\noindent\(\displaystyle 1\text{米}=100\text{厘米}\)\allowbreak \par 由厘米化米是小单位化大单位，应除以100。把\(329.4\)\allowbreak 的小数点向左移动两位，注意原小数位仍保留。\par\noindent\(\displaystyle 329.4\text{厘米}=3.294\text{米}\)\allowbreak \par 不能只移动一位，也不能向右移使数反而更大。\end{multicols}\end{snippet}
% p081_s11
\begin{snippet}{505}{11.5}{14}\begin{multicols}{3}\raggedcolumns \textbf{一、7（质量二）：}\ 先将\(275.09\)\allowbreak 吨分成整数部分275吨与小数部分\(0.09\)\allowbreak 吨。整数吨数直接保留，小数吨按下列进率换算：\par\noindent\(\displaystyle 1\text{吨}=1000\text{千克}\)\allowbreak \par \noindent\(\displaystyle 0.09\text{吨}=90\text{千克}\)\allowbreak \par 所以\par\noindent\(\displaystyle 275.09\text{吨}=275\text{吨}90\text{千克}\)\allowbreak \par\(0.09\)\allowbreak 吨不是9千克，换算时不能漏掉进率中的一个10。\end{multicols}\end{snippet}
% p081_s12
\begin{snippet}{505}{11.5}{14}\begin{multicols}{3}\raggedcolumns 第一空是左移两位后的原数，要倒过来右移两位，等于把\(86.27\)\allowbreak 扩大100倍。第二空是扩大1000倍前的数，须把7460缩小到原来的一千分之一。两空都是还原原数，运算方向应与题中已经发生的变化相反。\[\begin{aligned}&86.27\times100\\={}&8627\\&7460\div1000\\={}&7.46\end{aligned}\]\end{multicols}\end{snippet}
% p081_s13
\begin{snippet}{505}{11.5}{14}\begin{multicols}{3}\raggedcolumns 原数只有两位小数，保留一位时由百分位决定舍去还是进位。最小时，十分位原为4，百分位至少5才能进到\(236.5\)\allowbreak ，故下界为\(236.45\)\allowbreak ；最大时十分位为5，百分位至多4才不进位，故上界为\(236.54\)\allowbreak 。再往外的\(236.44\)\allowbreak 、\(236.55\)\allowbreak 分别会得到\(236.4\)\allowbreak 、\(236.6\)\allowbreak ，都不符合。\end{multicols}\end{snippet}
% p081_s14
\begin{snippet}{505}{11.5}{14}\begin{multicols}{3}\raggedcolumns 题目按当前实际在馆的图书本数计算，借出且尚未归还的部分应先从原有量中扣去。新进图书再增加到剩余量上，所以按“原有－借出＋新进”计算。三项都以万册为单位，可直接运算，结果也应保留万册。\[\begin{aligned}&6.26-1.3+2.25\\={}&4.96+2.25\\={}&7.21\text{(万册)}\end{aligned}\]\end{multicols}\end{snippet}
% p082_s01
\begin{snippet}{505}{11.5}{14}\begin{multicols}{3}\raggedcolumns 先把两个6放到指定数位上：百位上的6表示6个百，百分位上的6表示6个百分之一。十位、个位和十分位都没有相应的数，要写0占位；尤其不能漏掉小数点后的十分位0。\end{multicols}\end{snippet}
% p082_s02
\begin{snippet}{505}{11.5}{14}\begin{multicols}{3}\raggedcolumns 602比600多2，把它拆成\(600+2\)\allowbreak ，就是把602个32分成600个32与2个32。两部分都要与32相乘，最后再把积相加；A少算一个32，C把乘积错成了和。\end{multicols}\end{snippet}
% p082_s03
\begin{snippet}{505}{11.5}{14}\begin{multicols}{3}\raggedcolumns 判断“最接近”要比较各数到287的差，不能只看整数部分。三个差依次为\(0.001,\allowbreak 0.01,\allowbreak 0.1\)\allowbreak ，最小的差表示两数之间距离最小；小数位数多并不表示差更大。\end{multicols}\end{snippet}
% p082_s04
\begin{snippet}{505}{11.5}{14}\begin{multicols}{3}\raggedcolumns 先右移两位、再左移三位，抵消后相当于向左移一位，得到的数是原数的十分之一。现在要从结果倒推原数，就应把327.06的小数点向右移一位；倒推时移动方向与合并后的方向相反。\end{multicols}\end{snippet}
% p082_s05
\begin{snippet}{505}{11.5}{14}\begin{multicols}{3}\raggedcolumns 小括号在中括号里面，所以先算\(12+34\)\allowbreak ，得到除数。中括号内还须先除后减，算完中括号才与外面的24相乘；不能因为乘号写在前面就先算乘法。\end{multicols}\end{snippet}
% p082_s06
\begin{snippet}{505}{11.5}{14}\begin{multicols}{3}\raggedcolumns “24个32相加”表示每份32、共24份，应列\(32\times24\)\allowbreak ，不能列成\(24+32\)\allowbreak 。除数是“12的22倍”这整个数量，要用括号把\(12\times22\)\allowbreak 括起来；不加括号就会变成先除12再乘22，改变题意。\end{multicols}\end{snippet}
% p082_s07
\begin{snippet}{505}{11.5}{14}\begin{multicols}{3}\raggedcolumns 两个成绩的单位不同，先统一单位：\par\noindent\(\displaystyle 4\text{厘米}=0.04\text{米}\)\allowbreak \par 再与1.23米比较。王兵的成绩比张宇高，因此求张宇时要从王兵的成绩中减去差，不能反而加上4厘米。\end{multicols}\end{snippet}
% p082_s08
\begin{snippet}{505}{11.5}{14}\begin{multicols}{3}\raggedcolumns \textbf{三、1（1）—（4）：}\ （1）先把\(2\times5\)\allowbreak 结合成10，再与37相乘，能直接得到整十数。（2）101是100多1，先求25的100倍，再加1个25。（3）（4）把6.7、8.6分别看成6.70、8.60，与另一个加数的小数点对齐，十分位与十分位相加。\end{multicols}\end{snippet}
% p082_s09
\begin{snippet}{505}{11.5}{14}\begin{multicols}{3}\raggedcolumns \textbf{三、1（5）—（8）：}\ （5）32与18相加正好是50，连续减去这两数可改成一次减去它们的和。（6）十分位相加满十时向个位进1；（7）182可写成182.0，十分位不够减须退位。（8）把55.4看成55.40，与9.46按同数位相加，百分位的6不能加到十分位上。\end{multicols}\end{snippet}
% p082_s10
\begin{snippet}{505}{11.5}{14}\begin{multicols}{3}\raggedcolumns \textbf{三、2（1）：}\ 观察小数部分，233.2与26.8能凑成260，56.4与83.6能凑成140。利用加法交换律、结合律把这两组数配在一起，先各自凑整，再相加更简便。\end{multicols}\end{snippet}
% p082_s10_local1
\begin{snippet}{198.233}{10.5}{12.8}\SourceSans\setlength{\mathindent}{2bp}\[\begin{aligned}={}&(233.2+26.8)+(56.4+83.6)\\={}&260+140\\={}&400\end{aligned}\]\end{snippet}
% p082_s11
\begin{snippet}{505}{11.5}{14}\begin{multicols}{3}\raggedcolumns \textbf{三、2（2）：}\ 63.46与24.54的小数部分相加正好凑成1，两个减数的和是整数88。用连续减去两个数等于减去它们的和，把减法合并后再算；合并减数时括号内应写加号。\end{multicols}\end{snippet}
% p082_s11_local1
\begin{snippet}{164.567}{10.5}{12.8}\SourceSans\setlength{\mathindent}{2bp}\[\begin{aligned}={}&200-(63.46+24.54)\\={}&200-88\\={}&112\end{aligned}\]\end{snippet}
% p082_s12
\begin{snippet}{505}{11.5}{14}\begin{multicols}{3}\raggedcolumns \textbf{三、2（3）：}\ 102接近整百数，把它拆成100与2，分别求87的100倍和2倍。用乘法分配律把两部分的积相加；不能把多出来的2只加一次，也不能只乘100。\end{multicols}\end{snippet}
% p082_s12_local1
\begin{snippet}{143.233}{10.5}{12.8}\SourceSans\setlength{\mathindent}{2bp}\[\begin{aligned}={}&87\times(100+2)\\={}&87\times100+87\times2\\={}&8700+174\\={}&8874\end{aligned}\]\end{snippet}
% p083_s01
\begin{snippet}{505}{11.5}{14}\begin{multicols}{3}\raggedcolumns \textbf{三、2（4）：}\ 乘除部分中的36能被9整除，可先用36除以9，再与43相乘，计算更方便。这里调整的是乘积除以9的算法，外面的540仍减去整个乘除结果；不能先算\(540-36\)\allowbreak 。\end{multicols}\end{snippet}
% p083_s01_local1
\begin{snippet}{152.733}{10.5}{12.8}\SourceSans\setlength{\mathindent}{2bp}\[\begin{aligned}={}&540-(36\div9)\times43\\={}&540-4\times43\\={}&540-172\\={}&368\end{aligned}\]\end{snippet}
% p083_s02
\begin{snippet}{505}{11.5}{14}\begin{multicols}{3}\raggedcolumns \textbf{三、2（5）：}\ 32可以拆成\(4\times8\)\allowbreak ，使25与4配对成100，250与8配对成2000。利用乘法交换律、结合律重新分组，所有因数都恰好用一次；分组后的两积仍要相乘，不能改成相加。\end{multicols}\end{snippet}
% p083_s02_local1
\begin{snippet}{141.4}{10.5}{12.8}\SourceSans\setlength{\mathindent}{2bp}\[\begin{aligned}={}&(25\times4)\times(8\times250)\\={}&100\times2000\\={}&200000\end{aligned}\]\end{snippet}
% p083_s03
\begin{snippet}{505}{11.5}{14}\begin{multicols}{3}\raggedcolumns \textbf{三、2（6）：}\ 先计算最里面的减法，得到16，再计算中括号中的\(96\div16\)\allowbreak 。最后才用246除以中括号的结果；不能把原式改成连续除以96和16，因为除数是一个完整的商。\end{multicols}\end{snippet}
% p083_s03_local1
\begin{snippet}{216.067}{10.5}{12.8}\SourceSans\setlength{\mathindent}{2bp}\[\begin{aligned}={}&246\div[96\div16]\\={}&246\div6\\={}&41\end{aligned}\]\end{snippet}
% p083_s04
\begin{snippet}{505}{11.5}{14}\begin{multicols}{3}\raggedcolumns 两车都行驶3小时，尚未相遇。可先把每小时两车接近的路程合并再乘时间，也可分别列出两车路程，再用乘法分配律提取共同的3；最后都要加上尚余193千米。\par 另一列式可这样化为原位主式：\[\begin{aligned}&95\times3+75\times3+193\\={}&(95+75)\times3+193\end{aligned}\]后续共用原位的完整简便脱式。\end{multicols}\end{snippet}
% p083_s05
\begin{snippet}{505}{11.5}{14}\begin{multicols}{3}\raggedcolumns 大米比食用油少花23.9元，食用油就是79.5元再多23.9元，两种费用合计相当于两个79.5元再加23.9元。用加法结合律先把两个79.5凑成整数159，再加23.9；这是同一数量关系的简便整理，不另算一种解法。5千克和5升是包装信息，本题无需计算单价。\end{multicols}\end{snippet}
% p083_s05_local1
\begin{snippet}{189}{10.5}{12.8}\SourceSans\setlength{\mathindent}{2bp}\[\begin{aligned}&(79.5+79.5)+23.9\\={}&159+23.9\\={}&182.9\text{(元)}\end{aligned}\]\end{snippet}
% p083_s06
\begin{snippet}{505}{11.5}{14}\begin{multicols}{3}\raggedcolumns 四年级人数是158人的2倍再少68，比较时要从362中减去这一整体。减去“316少68”，等于先从362减316再补68。按减去一个差的性质去括号，再把362与68结合，先凑成430；430与316两项可同一步算出。不要把这一整理再算作另一种独立解法。\end{multicols}\end{snippet}
% p083_s06_local1
\begin{snippet}{179}{10.5}{12.8}\SourceSans\setlength{\mathindent}{2bp}\[\begin{aligned}&362-(158\times2-68)\\={}&362+68-158\times2\\={}&430-316\\={}&114\text{(人)}\end{aligned}\]\end{snippet}
% p084_s01
\begin{snippet}{505}{11.5}{14}\begin{multicols}{3}\raggedcolumns 两种水果各35箱，可把一箱葡萄与一箱橙子配成一组，先求每组价钱。也可分别按单价乘箱数列出总价，再用乘法分配律把55与45凑成100。\par 另一列式可这样化为原位主式：\[\begin{aligned}&55\times35+45\times35\\={}&(55+45)\times35\end{aligned}\]后续共用原位的完整简便脱式。\end{multicols}\end{snippet}
% p084_s01_local1
\begin{snippet}{80}{9.0}{14}\SourceSans\setlength{\mathindent}{2bp}\(\text{(元)}\)\end{snippet}
% p084_s02
\begin{snippet}{505}{11.5}{14}\begin{multicols}{3}\raggedcolumns 两种水果箱数相同，每箱价差乘35箱就是总价差。也可先列两种总价之差，再用乘法分配律把55减45凑成10。\par 另一列式可这样化为原位主式：\[\begin{aligned}&55\times35-45\times35\\={}&(55-45)\times35\end{aligned}\]后续共用原位的完整简便脱式。\end{multicols}\end{snippet}
% p084_s02_local1
\begin{snippet}{80}{9.0}{14}\SourceSans\setlength{\mathindent}{2bp}\(\text{(元)}\)\end{snippet}
% p084_s03
\begin{snippet}{505}{11.5}{14}\begin{multicols}{3}\raggedcolumns 表中1014元是桌椅合计，先扣除6张课桌的费用，剩下的才是12把椅子的费用，再平均到每把椅子。另法把12把椅子和6张桌子分成6组，每组恰好一桌两椅；一组总价减125元后，还要除以2才是每把椅子的价钱。\par 另法：每组一桌配两椅，先求一组椅子的价钱。\[\begin{aligned}&(1014\div6-125)\div2\\={}&(169-125)\div2\\={}&44\div2\\={}&22\text{(元)}\end{aligned}\]\end{multicols}\end{snippet}
% p084_s03_local1
\begin{snippet}{80}{9.0}{14}\SourceSans\setlength{\mathindent}{2bp}\(\text{(元)}\)\end{snippet}
% p084_s04
\begin{snippet}{505}{11.5}{14}\begin{multicols}{3}\raggedcolumns 两人的钱合起来正好够买一套，说明王军所缺的9.3元恰好是张兵已有的钱，张兵所缺的14.7元也是王军已有的钱。因此两笔缺额相加就是一套书的价钱，再平均分给价钱相同的两本书；不能把缺额之差当成书价。\end{multicols}\end{snippet}
% p084_s04_local1
\begin{snippet}{80}{9.0}{14}\SourceSans\setlength{\mathindent}{2bp}\(\text{(元)}\)\end{snippet}
% p085_s01
\begin{snippet}{505}{11.5}{14}\begin{multicols}{3}\raggedcolumns 给已知直线画平行线时，可以在它两侧选取不同距离，因而画法没有固定的条数限制。第二空固定了直线外的同一点，过这个点与已知直线成直角的方向只有一个；不要把“任意位置”与“过指定点”混淆。\end{multicols}\end{snippet}
% p085_s02
\begin{snippet}{505}{11.5}{14}\begin{multicols}{3}\raggedcolumns 邻边是共享一个顶点的两条边，长方形每个顶点处都是直角，所以要看成直角的关系。对边是互不相邻的两条边，在同一平面内延长也不相交；一组边竖着画、另一组横着画，都不影响这些关系。\end{multicols}\end{snippet}
% p085_s03
\begin{snippet}{505}{11.5}{14}\begin{multicols}{3}\raggedcolumns 先固定题目选定的这一条底，从对边上不同的点都能向这条底作垂线段，因此可作无数条高。这些高都夹在同一组平行线之间，长度相等；不能换另一条边作底后，把两种底对应的高混在一起比较。\end{multicols}\end{snippet}
% p085_s04
\begin{snippet}{505}{11.5}{14}\begin{multicols}{3}\raggedcolumns \textbf{一、4：}\ 先给点命名以便有次序地找图形：上排从左到右为\(A,\allowbreak B,\allowbreak C,\allowbreak D\)\allowbreak ，下排为\(E,\allowbreak F,\allowbreak G\)\allowbreak 。先数三条已画线段围成的三角形，再检查四边形的两组对边关系，不能只数最小图形。\par 三角形：\(ABE,\allowbreak BEF,\allowbreak BCF,\allowbreak CFG,\allowbreak CDG\)\allowbreak 。\par 平行四边形：\(ABFE,\allowbreak ACGE,\allowbreak BCFE,\allowbreak BCGF\)\allowbreak 。\par 梯形：\(ACFE,\allowbreak ADGE,\allowbreak BCGE,\allowbreak BDGE,\allowbreak BDGF,\allowbreak CDGF\)\allowbreak ；这里梯形按只有一组对边平行判断，不能把平行四边形再计入。\end{multicols}\end{snippet}
% p085_s05
\begin{snippet}{505}{11.5}{14}\begin{multicols}{3}\raggedcolumns \textbf{一、5：}\ 用任意两边之和大于第三边来限制第三边：它必须大于\(9-5=4\)\allowbreak 厘米，又小于\(9+5=14\)\allowbreak 厘米。题目要求整厘米数，就在4与14之间取整数，最靠近两端的可取5和13；等于4或14只能排成一条线，不能围成三角形。\end{multicols}\end{snippet}
% p085_s06
\begin{snippet}{505}{11.5}{14}\begin{multicols}{3}\raggedcolumns \textbf{一、6：}\ 先认准所问的底，再找对面顶点到这条底所在直线的垂线段：BC对应AD，AB对应AC。AC与AB在A处成直角，因此第二空可直接读AC的长度；AD原标末位尚未核清，2.4与2.6的疑读条件应保留，不能用未学的面积法反推原标。\end{multicols}\end{snippet}
% p085_s07
\begin{snippet}{505}{11.5}{14}\begin{multicols}{3}\raggedcolumns \textbf{一、7（1）：}\ 本小问研究的是方向盘本身，要看它是否围绕固定中心转动。驾驶时方向盘各部分绕中心改变方向，所以按转动的特征判断；汽车整体沿公路前进是另一个研究对象。\end{multicols}\end{snippet}
% p085_s08
\begin{snippet}{505}{11.5}{14}\begin{multicols}{3}\raggedcolumns \textbf{一、7（2）：}\ 把国旗看作一个整体，升旗时各点沿旗杆向上移动相同的距离，朝向保持不变。按题中理想的上升运动判断，不把风吹旗面摆动混入这次运动。\end{multicols}\end{snippet}
% p085_s09
\begin{snippet}{505}{11.5}{14}\begin{multicols}{3}\raggedcolumns \textbf{一、7（3）：}\ 按题意把拖地动作理解为拖布沿直线来回推动，拖布整体向同一方向移动，朝向不变。现实中转动拖布是另一种动作，判断时要依据题中所研究的直线推动过程。\end{multicols}\end{snippet}
% p085_s10
\begin{snippet}{505}{11.5}{14}\begin{multicols}{3}\raggedcolumns \textbf{一、7（4）：}\ 这里研究车轮绕车轴转动，各部分都围绕同一个中心运动，符合旋转的特征。自行车整体向前行驶与车轮本身的转动要分别观察，不能用整辆车的位置变化代替对车轮的判断。\end{multicols}\end{snippet}
% p085_s11
\begin{snippet}{505}{11.5}{14}\begin{multicols}{3}\raggedcolumns 对称轴是直线。正方形的4条对称轴是两条对角线所在的直线，以及两组对边中点连线所在的直线；等边三角形有3条，等腰梯形有经过两底中点的1条。\end{multicols}\end{snippet}
% p085_s12
\begin{snippet}{505}{11.5}{14}\begin{multicols}{3}\raggedcolumns 末题后两空按分针顺时针转90度理解；原角度数字不清。\par 原钟面分针指向12、时针指向3，可以先读出原来的整时。后续旋转的直角个数未核清，不能据此确定唯一时刻；仅按“一个直角”理解时，\(90\div30=3\)\allowbreak (大格)，\(5\times3=15\)\allowbreak (分钟)，这是有条件的计算。\end{multicols}\end{snippet}
% p086_s01
\begin{snippet}{505}{11.5}{14}\begin{multicols}{3}\raggedcolumns 先跟踪原图三格横排的位置，顺时针转直角后，这三格变成竖排。再看原来左下伸出的那格，它应转到竖排的上端左侧；核对这个拐角位置，能避免把旋转误作翻面。\end{multicols}\end{snippet}
% p086_s02
\begin{snippet}{505}{11.5}{14}\begin{multicols}{3}\raggedcolumns 直线a、b都与c成直角，因此a、b相对于c有相同的方向。题目已经限定三条直线在同一平面内，可据“同一平面内垂直于同一直线的两条直线互相平行”判断；不能因为都出现直角就说a、b也互相垂直。\end{multicols}\end{snippet}
% p086_s03
\begin{snippet}{505}{11.5}{14}\begin{multicols}{3}\raggedcolumns 40度没有说明是顶角还是底角，必须分两种情况，不能只用一种图形判断。若它是顶角，余下的140度平均分给两个相等的底角；若它是底角，先确定另一个底角也是40度，再用180度减去两个底角。\end{multicols}\end{snippet}
% p086_s04
\begin{snippet}{505}{11.5}{14}\begin{multicols}{3}\raggedcolumns 框架被拉动时，四条边仍是原来的四根边，长度没有改变，只是相邻边的夹角变了。周长是四边长度的总和，所以可以直接根据边长是否改变来比较；图形变斜或变矮并不改变这次比较的结果。\end{multicols}\end{snippet}
% p086_s05
\begin{snippet}{505}{11.5}{14}\begin{multicols}{3}\raggedcolumns 逐个按原图完整轮廓尝试对折，只有两部分完全重合才算轴对称。箭头、十字形和圆形可找到对称方向，而一般平行四边形沿对角线或中线对折都不能重合；不要把它当成长方形。\end{multicols}\end{snippet}
% p086_s06
\begin{snippet}{505}{11.5}{14}\begin{multicols}{3}\raggedcolumns 工具标识的外框和内部图案都是判断对象，应一起沿同一条直线对折检查。C中的闪电与轮廓不能同时重合，不能因为外框近似对称就把整个标识判成轴对称；判断依照保留的原字形，不另换简化图案。\end{multicols}\end{snippet}
% p086_s07
\begin{snippet}{505}{11.5}{14}\begin{multicols}{3}\raggedcolumns \textbf{三、1：}\ 题目同时固定了直线外的一个点和已知直线，在同一平面内，过该点保持与已知直线平行的方向只有一个。虽然能画许多条与已知直线平行的线，但其中只有一条经过这个指定点。\end{multicols}\end{snippet}
% p086_s08
\begin{snippet}{505}{11.5}{14}\begin{multicols}{3}\raggedcolumns \textbf{三、2：}\ 判断平行必须同时满足“在同一平面内”和“不相交”，本句只给出不相交，条件不足。比如盒子上不在同一平面内的两条棱所在直线也可能不相交，却不能称为平行；因此不能漏掉同一平面的限制。\end{multicols}\end{snippet}
% p086_s09
\begin{snippet}{505}{11.5}{14}\begin{multicols}{3}\raggedcolumns \textbf{三、3：}\ 原底角末位有\(26^\circ\)\allowbreak 或\(25^\circ\)\allowbreak 两种读法，两种读数都小于\(45^\circ\)\allowbreak 。等腰三角形的两个底角相等，所以底角和均小于90度，余下的顶角就大于90度；判断相同，但不能把某一种疑读计算出的顶角当作已核清数值。\end{multicols}\end{snippet}
% p086_s10
\begin{snippet}{505}{11.5}{14}\begin{multicols}{3}\raggedcolumns \textbf{三、4：}\ 旋转是图形围绕一个中心转动，对应边的长度和对应角的大小都保持不变。可以把原图想象为硬纸片转动：朝向和位置可以改变，形状、大小却仍与原图相同。\end{multicols}\end{snippet}
% p086_s11
\begin{snippet}{505}{11.5}{14}\begin{multicols}{3}\raggedcolumns \textbf{三、5：}\ 两条直线相交成直角，就是互相垂直，本题给出的四个角都是直角，已经满足这个条件。判断依据是相交处的角度，不是线条画成水平、竖直还是倾斜。\end{multicols}\end{snippet}
% p086_s12
\begin{snippet}{505}{11.5}{14}\begin{multicols}{3}\raggedcolumns \textbf{三、6：}\ “等底等高”只限制了底和高，不能保证两梯形的腰长、倾斜方向和形状完全相同，因此不能保证拼接后的两组对边都平行。教材能保证拼成平行四边形的条件是两个完全相同的梯形，并沿合适的腰拼接；不能把两种条件等同。\end{multicols}\end{snippet}
% p087_s01
\begin{snippet}{505}{11.5}{14}\begin{multicols}{3}\raggedcolumns \textbf{四、1（1）：}\ 沿虚线轴逐个数顶点到轴的水平格数，在轴另一侧的同一水平线上取等距点，再按原轮廓的顺序连接。距轴2格、3格的顶点，另一侧也分别距轴2格、3格；轴上两个端点不移动，不能把整幅半图直接平移过去。\end{multicols}\end{snippet}
% p087_s02
\begin{snippet}{505}{11.5}{14}\begin{multicols}{3}\raggedcolumns \textbf{四、1（2）：}\ 先以O为固定中心，把各顶点到O的连线顺时针转直角，顶点与O的距离保持不变。确定旋转后的三角形后，再把它的每个顶点向左数3格并顺次连接；第二步必须从旋转图起数，不能从原图起数。\end{multicols}\end{snippet}
% p087_s03
\begin{snippet}{505}{11.5}{14}\begin{multicols}{3}\raggedcolumns \textbf{四、2（1）：}\ 先把三角形的三个顶点各向上数3格，再从所得位置各向右数5格。按原来的顶点顺序连接，所画图形的边长和朝向都应与原图一致；数的是移动的格数，不是把起点也算一格。\end{multicols}\end{snippet}
% p087_s04
\begin{snippet}{505}{11.5}{14}\begin{multicols}{3}\raggedcolumns \textbf{四、2（2）：}\ O是本次旋转中心，保持原位；原来从O向左的3格线段，顺时针转直角后应朝上3格。另一个顶点也围绕O按同一方向转动，再连接三个对应点；这一小问针对原图，不接着上一问的平移图旋转。\end{multicols}\end{snippet}
% p087_s05
\begin{snippet}{505}{11.5}{14}\begin{multicols}{3}\raggedcolumns \textbf{四、2（3）：}\ 以竖直虚线为轴，原图左侧两个顶点距轴都是3格，应在轴右侧同一水平线上各取距轴3格的对应点。O在轴上保持不动，连接对应点补成三角形；本问从原图作对称，不接着前两次操作变换。\end{multicols}\end{snippet}
% p087_s06
\begin{snippet}{505}{11.5}{14}\begin{multicols}{3}\raggedcolumns 先认准每个图中标“底”的边：三角形和平行四边形的底在下方，梯形所标的底在上方。用三角尺一条直角边靠底所在直线，移动到另一条直角边通过所选顶点，再画虚线垂线段。三角形从对面顶点作高，平行四边形从上边顶点作高，梯形从下底端点向上底作高，垂足处均标直角；斜边、斜腰不能直接当高。\end{multicols}\end{snippet}
% p087_s07
\begin{snippet}{505}{11.5}{14}\begin{multicols}{3}\raggedcolumns \textbf{五、1（1）：}\ 王庄与加油站都用固定的位置点表示，所求路线的两端必须分别落在这两个点上。直接连接它们得到线段，依据两点之间线段最短；不需要为了与公路成直角而改变终点。\end{multicols}\end{snippet}
% p087_s08
\begin{snippet}{505}{11.5}{14}\begin{multicols}{3}\raggedcolumns \textbf{五、1（2）：}\ 到公路时，终点可以选在公路上，应从王庄向靠近它的公路边线作垂线段。依据点到直线的垂线段最短，垂足处要标直角；本题的终点不是加油站，不能照用上一问的连线。\end{multicols}\end{snippet}
% p088_s01
\begin{snippet}{505}{11.5}{14}\begin{multicols}{3}\raggedcolumns 等腰梯形两腰相等。\textbf{方法一：先扣除两底，再平均分给两腰}\[\begin{aligned}&(68-(8+16))\div2\\={}&(68-24)\div2\\={}&44\div2\\={}&22\text{(厘米)}\end{aligned}\]答：一条腰长22厘米。\par\textbf{方法二：从一半周长扣除两底和的一半}\[\begin{aligned}&68\div2-(8+16)\div2\\={}&34-24\div2\\={}&34-12\\={}&22\text{(厘米)}\end{aligned}\]答：一条腰长22厘米。\end{multicols}\end{snippet}
% p088_s02
\begin{snippet}{505}{11.5}{14}\begin{multicols}{3}\raggedcolumns 24厘米这条边比邻边短2厘米，所以邻边应在24厘米上加2厘米，再把一组邻边的和乘2。另法先把四边都按24厘米计，只有两条较长边各少算2厘米，合计补2份2厘米；不能给四条边都加2。\par 另法：四边先都按较短边算，再补两条邻边各多的长度。\[\begin{aligned}&24\times4+2\times2\\={}&96+4\\={}&100\text{(厘米)}\end{aligned}\]\end{multicols}\end{snippet}
% p088_s02_local1
\begin{snippet}{72.05178222656252}{9}{12.8}\SourceSans\setlength{\mathindent}{2bp}\(\displaystyle=[24+26]\times2\)\end{snippet}
% p088_s02_local2
\begin{snippet}{80}{9.0}{14}\SourceSans\setlength{\mathindent}{2bp}\(\text{(厘米)}\)\end{snippet}
% p088_s03
\begin{snippet}{505}{11.5}{14}\begin{multicols}{3}\raggedcolumns 同一根铁丝重新围图形，总长不变，先按正方形边长6厘米乘4条边列总长，再平均分给等边三角形的3条边。计算时6能被3整除，可先算6除以3得2，再乘4；这是乘除性质的简便整理，不能把原来的6厘米只除以3而漏乘4。\end{multicols}\end{snippet}
% p088_s03_local1
\begin{snippet}{319}{10.5}{12.8}\SourceSans\setlength{\mathindent}{2bp}\[\begin{aligned}&6\times4\div3\\={}&6\div3\times4\\={}&2\times4\\={}&8\text{(厘米)}\end{aligned}\]\end{snippet}
% p089_s01
\begin{snippet}{505}{11.5}{14}\begin{multicols}{3}\raggedcolumns 把一组数据的总量平均分成与数据个数同样多的份，每份的数就是平均数。它反映整组的整体水平，可能与任何一个原数据都不相等，不能把平均数理解为某个人实际的成绩。\end{multicols}\end{snippet}
% p089_s02
\begin{snippet}{505}{11.5}{14}\begin{multicols}{3}\raggedcolumns 简便法，以150为基准移多补少：\[\begin{aligned}&150+(160-150)\div2\\={}&150+10\div2\\={}&150+5\\={}&155\text{(个)}\end{aligned}\]另法，求总量再平均：\[\begin{aligned}&(160+150)\div2\\={}&310\div2\\={}&155\text{(个)}\end{aligned}\]第二空仅小涵增加，小玉保持160个，应补足全部差：\(160-150=10\)\allowbreak (个)。\end{multicols}\end{snippet}
% p089_s03
\begin{snippet}{505}{11.5}{14}\begin{multicols}{3}\raggedcolumns 用加法结合律把12与8、7与13分别配成20：\[\begin{aligned}&(12+8+7+13)\div4\\={}&[(12+8)+(7+13)]\div4\\={}&[20+20]\div4\\={}&40\div4\\={}&10\end{aligned}\]4个数据，除数仍为4。\end{multicols}\end{snippet}
% p089_s04
\begin{snippet}{505}{11.5}{14}\begin{multicols}{3}\raggedcolumns 移多补少：二班少5本，一班多3本，三班须多2本。利用加减运算的性质，先算5减3。\[\begin{aligned}&20+(20-15)-(23-20)\\={}&20+5-3\\={}&20+(5-3)\\={}&20+2\\={}&22\text{(本)}\end{aligned}\]另法，先求三个班总本数，再减两班之和：\[\begin{aligned}&20\times3-(23+15)\\={}&60-38\\={}&22\text{(本)}\end{aligned}\]\end{multicols}\end{snippet}
% p089_s05
\begin{snippet}{505}{11.5}{14}\begin{multicols}{3}\raggedcolumns 连续三个自然数相邻相差1，把最大数多出的1移给最小数，三数就都等于中间数，总和没有改变。因此平均数12就是中间数，最大数应比它多1；不要把12误认成最大数。\end{multicols}\end{snippet}
% p089_s06
\begin{snippet}{505}{11.5}{14}\begin{multicols}{3}\raggedcolumns 合并两组的总次数与总人数，两个括号可同一步算：\[\begin{aligned}&(96+156)\div(3+4)\\={}&252\div7\\={}&36\text{(个)}\end{aligned}\]两组人数不同，不能把32与39直接平均。\end{multicols}\end{snippet}
% p089_s07
\begin{snippet}{505}{11.5}{14}\begin{multicols}{3}\raggedcolumns 只去掉99和92，剩5项。\textbf{方法一：以96分为基准}98多2分，两个95各少1分，余两项正好96分。连续减去两个数，等于减去它们的和。\[\begin{aligned}&96+(2-1-1)\div5\\={}&96+[2-(1+1)]\div5\\={}&96+[2-2]\div5\\={}&96+0\div5\\={}&96+0\\={}&96\text{(分)}\end{aligned}\]\textbf{方法二：分组求总分}用加法交换律、结合律：\end{multicols}\[\begin{aligned}&(98+95+96+95+96)\div5\\={}&[98+(96+96)+(95+95)]\div5\\={}&[98+192+190]\div5\\={}&[290+190]\div5\\={}&480\div5\\={}&96\text{(分)}\end{aligned}\]\begin{multicols}{3}\raggedcolumns 答：最后得分是96分。\end{multicols}\end{snippet}
% p089_s08
\begin{snippet}{505}{11.5}{14}\begin{multicols}{3}\raggedcolumns 一组深度的平均数应在最浅1.1米和最深2.7米之间，先用这个范围逐项排除。由于这里已经有不同的最浅、最深值，整体平均不会等于任一端点，0.8米又比最浅处还小，因此只有中间的数符合“可能”。\end{multicols}\end{snippet}
% p089_s09
\begin{snippet}{505}{11.5}{14}\begin{multicols}{3}\raggedcolumns 全班平均体重35千克只表示整体水平，不能据此要求每个同学都恰好35千克。小萱比平均数轻3千克，可以由其他同学较重的部分补足，所以32千克是可能的，却不能断定她一定是这个体重。\end{multicols}\end{snippet}
% p090_s01
\begin{snippet}{505}{11.5}{14}\begin{multicols}{3}\raggedcolumns 要求平均的是四次测得的气温8、15、26、19，要先求这四个温度的和，再除以4次。3时、8时、15时、21时说明观测时间，不是气温值；一天有24小时也不表示这里有24个观测数据。\end{multicols}\end{snippet}
% p090_s02
\begin{snippet}{505}{11.5}{14}\begin{multicols}{3}\raggedcolumns 原选项均不成立。由从低到高排列得\(81\le a\le87\)\allowbreak 。先求其他四项和，按加法交换律、结合律分组：\[\begin{aligned}&73+81+87+95\\={}&(73+87)+(81+95)\\={}&160+176\\={}&336\end{aligned}\]最小平均数：\[\begin{aligned}&(336+81)\div5\\={}&417\div5\\={}&83.4\text{(下)}\end{aligned}\]最大平均数：\[\begin{aligned}&(336+87)\div5\\={}&423\div5\\={}&84.6\text{(下)}\end{aligned}\]因此平均数在83.4至84.6下之间；72、85、95、96均不符合。题目未要求取整，不能自行把84.6约为85。\end{multicols}\end{snippet}
% p090_s03
\begin{snippet}{505}{11.5}{14}\begin{multicols}{3}\raggedcolumns 原数为20，原选项均不成立。平均数增加10，6个数的总和就增加60；只有一个数改动，故从80扣去这60：\[\begin{aligned}&80-(60-50)\times6\\={}&80-10\times6\\={}&80-60\\={}&20\end{aligned}\]\end{multicols}\end{snippet}
% p090_s04
\begin{snippet}{505}{11.5}{14}\begin{multicols}{3}\raggedcolumns \textbf{三、1（1）—（4）：}\ 这些数都能用十分之一的个数来理解：4个十分之一加5个十分之一，或从10个十分之一中减7个十分之一。（3）十分位合成10个时要进成1；（4）2个十分之一不够减8个十分之一，要从个位退1再减。\end{multicols}\end{snippet}
% p090_s05
\begin{snippet}{505}{11.5}{14}\begin{multicols}{3}\raggedcolumns \textbf{三、1（5）—（8）：}\ （1）把0.4与1.2的十分位相加；（2）（4）两数的小数部分相同，减去后小数部分为0，可直接看整数部分的差。（3）3个十分之一加8个十分之一满十进1，进上去的1还要加到个位上。\end{multicols}\end{snippet}
% p090_s06
\begin{snippet}{505}{11.5}{14}\begin{multicols}{3}\raggedcolumns \textbf{三、2（1）：}\ 先把15.05与2.63的小数点对齐，确保百分位的5与3相加，十分位的0与6相加。按位相加时这一题不需要进位，十位上的1也要保留，不能把两数按末位对齐。\end{multicols}\end{snippet}
% p090_s06_local1
\begin{snippet}{24.25}{10}{14}\SourceSans\setlength{\mathindent}{2bp}\(17.68\)\end{snippet}
% p090_s07
\begin{snippet}{505}{11.5}{14}\begin{multicols}{3}\raggedcolumns \textbf{三、2（2）：}\ 把2.5写成等值的2.50，小数点对齐后才能逐位相减。百分位先算6减0，十分位0不够减5，要从个位9退1，把它换成10个十分之一；个位随后用剩下的8减2。\end{multicols}\end{snippet}
% p090_s07_local1
\begin{snippet}{24.25}{10}{14}\SourceSans\setlength{\mathindent}{2bp}\(6.56\)\end{snippet}
% p090_s08
\begin{snippet}{505}{11.5}{14}\begin{multicols}{3}\raggedcolumns \textbf{三、2（3）：}\ 5.7可在末尾补0写成5.70，便于与12.34对齐百分位。十分位7加3满十进1，十分位要写0，个位相加时再加上进来的1；结果中间的0不能漏写。\end{multicols}\end{snippet}
% p090_s08_local1
\begin{snippet}{24.75}{10}{14}\SourceSans\setlength{\mathindent}{2bp}\(18.04\)\end{snippet}
% p090_s09
\begin{snippet}{505}{11.5}{14}\begin{multicols}{3}\raggedcolumns \textbf{三、2（4）及验算：}\ 先把12.8补写成12.80，百分位0不够减4，要从十分位8退1后再算。带星号还须用差加减数验算，原竖式中的11.56与1.24相加应恢复12.80，这才与原被减数12.8相同。\end{multicols}\end{snippet}
% p090_s09_local1
\begin{snippet}{34.25}{10}{14}\SourceSans\setlength{\mathindent}{2bp}\(11.56\)\end{snippet}
% p091_s01
\begin{snippet}{505}{11.5}{14}\begin{multicols}{3}\raggedcolumns 第3、4项字形有缺损，不能认定原题已核清。以下只按96、88、92、88计算。\textbf{方法一：分组求总分再平均}先按加法结合律分成两组，才能同一步求两个和：\[\begin{aligned}&(96+88+92+88)\div4\\={}&[(96+88)+(92+88)]\div4\\={}&[184+180]\div4\\={}&364\div4\\={}&91\text{(分)}\end{aligned}\]答：按这组数据，平均成绩是91分。\par\textbf{方法二：以88分为基准移多补少}\end{multicols}\[\begin{aligned}&88+[(96-88)+(92-88)]\div4\\={}&88+[8+4]\div4\\={}&88+12\div4\\={}&88+3\\={}&91\text{(分)}\end{aligned}\]\begin{multicols}{3}\raggedcolumns 答：按这组数据，平均成绩是91分。\par 若缺损两项实际为82、18，则：\[\begin{aligned}&(96+88+82+18)\div4\\={}&[(96+88)+(82+18)]\div4\\={}&[184+100]\div4\\={}&284\div4\\={}&71\text{(分)}\end{aligned}\]答：仅按该备选读法，平均成绩是71分；仍须核清原题。\end{multicols}\end{snippet}
% p091_s02
\begin{snippet}{505}{11.5}{14}\begin{multicols}{3}\raggedcolumns 男生8人与女生12人的人数不同，应分别用平均身高乘人数，求出两组身高总和，再除以20人。另法先把20人都按132厘米计算，12名女生每人多出的10厘米合计120厘米，再平均补给20人；不能把132和142直接平均。\par 另法：以132厘米为基准，女生多出的总身高再平均给20人。\end{multicols}\[\begin{aligned} &132+(142-132)\times12\div(8+12)\\={}&132+10\times12\div20\\={}&132+120\div20\\={}&132+6\\={}&138\text{(厘米)}\end{aligned}\]\end{snippet}
% p091_s02_local1
\begin{snippet}{80}{9.0}{14}\SourceSans\setlength{\mathindent}{2bp}\(\text{(厘米)}\)\end{snippet}
% p091_s03
\begin{snippet}{505}{11.5}{14}\begin{multicols}{3}\raggedcolumns 前2天给的是每天平均本数，须乘2才能得到这两天总数；后6天464本已经是总数，不再乘6。合并8天的784本后再求平均；求出总量后，也可与每天100本比较作复核，8天总共少16本，把差额平均到8天再从100中扣除。\par 基准复核：沿用主法已求出的总数784本，与每天100本比较。\[\begin{aligned}&100-(100\times8-784)\div8\\={}&100-(800-784)\div8\\={}&100-16\div8\\={}&100-2\\={}&98\text{(本)}\end{aligned}\]\end{multicols}\end{snippet}
% p091_s03_local1
\begin{snippet}{80}{9.0}{14}\SourceSans\setlength{\mathindent}{2bp}\(\text{(本)}\)\end{snippet}
% p091_s04
\begin{snippet}{505}{11.5}{14}\begin{multicols}{3}\raggedcolumns 平均数有重影，以下仅按2500毫升作条件计算。\textbf{方法一：先合并甲、乙的量}用减法性质合并两个减数：\[\begin{aligned}&2500\times3-(2700+2200)\\={}&7500-4900\\={}&2600\text{(毫升)}\end{aligned}\]答：若平均数是2500毫升，丙的肺活量是2600毫升。\par\textbf{方法二：移多补少}乙少300毫升，甲多200毫升，丙应在平均数上补足100毫升：\end{multicols}\[\begin{aligned}&2500+[(2500-2200)-(2700-2500)]\\={}&2500+[300-200]\\={}&2500+100\\={}&2600\text{(毫升)}\end{aligned}\]\begin{multicols}{3}\raggedcolumns 答：若平均数是2500毫升，丙的肺活量是2600毫升。\par 若原平均数实际为3000毫升，则：\[\begin{aligned}&3000\times3-(2700+2200)\\={}&9000-4900\\={}&4100\text{(毫升)}\end{aligned}\]答：若平均数是3000毫升，丙的肺活量是4100毫升。须核清原题后确定唯一答案。\end{multicols}\end{snippet}
% p092_s01
\begin{snippet}{505}{11.5}{14}\begin{multicols}{3}\raggedcolumns 女生按原表7、12、1、1读取；末列淡字暂按1。\par 合并表时，先把女生、男生各型号人数填到相应的行和列，再逐列相加求该型号总人数。最后分别核对各行合计与全班合计是否相同，不能把“同一型号”与“同一性别”加混。女生特大号的淡字暂按1人，若读数改变，该列总计、女生合计、全班合计必须同步修改。\par 暂按女生特大号1人核对：女生\(7+12+1+1=21\)\allowbreak (人)，男生\(3+12+7+2=24\)\allowbreak (人)，全班\(21+24=45\)\allowbreak (人)；逐列合计再核对\(10+24+8+3=45\)\allowbreak (人)。\end{multicols}\end{snippet}
% p092_s02
\begin{snippet}{505}{11.5}{14}\begin{multicols}{3}\raggedcolumns 女生最少人数要把小号7、中号12、大号1与尚未核清的特大号人数逐项比较，不能只找一个格子就停止。特大号若0人则它最少，若1人则与大号并列最少，若2人或更多则大号最少。男生各格不受该淡字影响，可直接比较3、12、7、2找出最多型号。\end{multicols}\end{snippet}
% p093_s01
\begin{snippet}{505}{11.5}{14}\begin{multicols}{3}\raggedcolumns 第一空先把60米除以1000换成千米，再与原有7千米合并，整数部分7不再另行缩小：\par\noindent\(\displaystyle 60\text{米}=0.06\text{千米}\)\allowbreak \par \noindent\(\displaystyle 7\text{千米}60\text{米}=(7+0.06)\text{千米}=7.06\text{千米}\)\allowbreak \par 第二空是面积单位换算，不能沿用长度单位的进率：\par\noindent\(\displaystyle 1\text{平方米}=10000\text{平方厘米}\)\allowbreak \par \noindent\(\displaystyle 800\text{平方厘米}=(800\div10000)\text{平方米}=0.08\text{平方米}\)\allowbreak \par\end{multicols}\end{snippet}
% p093_s02
\begin{snippet}{505}{11.5}{14}\begin{multicols}{3}\raggedcolumns 把大数改用“亿”作单位，要除以一亿，也就是把小数点向左移8位：\par\noindent\(\displaystyle 240960000\text{公顷}=2.4096\text{亿公顷}\)\allowbreak \par 保留两位小数看千分位9，要向百分位进1；两空后原来只印“公顷”，填空中还须带上“亿”，才能保持原面积的大小。\end{multicols}\end{snippet}
% p093_s03
\begin{snippet}{505}{11.5}{14}\begin{multicols}{3}\raggedcolumns 先按个位、十分位、百分位、千分位的顺序放入2、7、0、7，百分位的0起占位作用，不能省掉。题目要求三位小数，最低位是千分位，所以它的计数单位是一个千分之一，而不是看某个非零数字大小来确定。\end{multicols}\end{snippet}
% p093_s04
\begin{snippet}{505}{11.5}{14}\begin{multicols}{3}\raggedcolumns 题目问“可能”，可以分别尝试两次都沿同一方向折，以及横、竖方向各折一次。前者的折痕方向相同，后者的两种折痕成直角；应把这两种位置关系都考虑进去，不能只依据自己的一种折法。\end{multicols}\end{snippet}
% p093_s05
\begin{snippet}{505}{11.5}{14}\begin{multicols}{3}\raggedcolumns 先按48厘米竹条的根数分类，避免把同一组边长的不同摆放重复计算；原题未要求两种长度都用，等边三角形也属于等腰三角形。根数足够时四种组合如下：\par \((24,\allowbreak 24,\allowbreak 24)\)\allowbreak ：\(24\times3=72\)\allowbreak (厘米)；\par \((24,\allowbreak 24,\allowbreak 48)\)\allowbreak ：\(24+24=48\)\allowbreak ，不能围成；\par \((24,\allowbreak 48,\allowbreak 48)\)\allowbreak ：\(48\times2+24=96+24=120\)\allowbreak (厘米)；\par \((48,\allowbreak 48,\allowbreak 48)\)\allowbreak ：\(48\times3=144\)\allowbreak (厘米)。\par 若另要求两种长度都用，才只有120厘米；原题未说明各长竹条的根数，不能把这项附加条件当成原有条件。\end{multicols}\end{snippet}
% p093_s06
\begin{snippet}{505}{11.5}{14}\begin{multicols}{3}\raggedcolumns 700克与1.8千克单位不同，先统一单位：\par\noindent\(\displaystyle 700\text{克}=0.7\text{千克}\)\allowbreak \par 白鸭比黑鸭重，先在黑鸭质量上加0.7求出白鸭，再把两只鸭的质量相加；不能把0.7直接当作白鸭的质量。\par \noindent\[\displaystyle \begin{aligned}&1.8+(1.8+0.7)\\={}&1.8+2.5\\={}&4.3\text{(千克)}\end{aligned}\]\par\end{multicols}\end{snippet}
% p093_s07
\begin{snippet}{505}{11.5}{14}\begin{multicols}{3}\raggedcolumns （1）左边的两个因数8与4合为32，右边对应24，可在同乘125的情况下比较。（2）连续除以3、2相当于除以它们的积6，不能当成除以它们的和5。（3）连减125、75可合并为减去它们的和。（4）减98改成减100时多减2，应再补2，题中右边却又减2，因此两式不同。\end{multicols}\end{snippet}
% p093_s08
\begin{snippet}{505}{11.5}{14}\begin{multicols}{3}\raggedcolumns 最多人数对应完全没有重复演员，此时把两项人数直接相加。最少人数对应36名舞蹈演员全部都在42名合唱演员中，此时总人数只看较大的一项；不能把两项都参加的人重复计算。\end{multicols}\end{snippet}
% p093_s09
\begin{snippet}{505}{11.5}{14}\begin{multicols}{3}\raggedcolumns 按角给三角形分类，先把最大的角与90度比较，150度已经是钝角。有一个钝角即可确定这一类别，不必再求另外两个角；也不能因其余两个角较小就改判为锐角三角形。\end{multicols}\end{snippet}
% p094_s01
\begin{snippet}{505}{11.5}{14}\begin{multicols}{3}\raggedcolumns 先按每个数的原写法读，整数部分按整数读法，小数点后各位却要逐个读出，包括末尾的0。A有“点零零”，B末位也要读零，C、D同样含需读出的零，所以原选项均不符合；把900.60先化简为900.6会改变本题的读数条件。\end{multicols}\end{snippet}
% p094_s02
\begin{snippet}{505}{11.5}{14}\begin{multicols}{3}\raggedcolumns 32可拆成\(8\times4\)\allowbreak ，125与8凑1000，25与4凑100，两组积再相乘。这样每个因数只用一次，运用了乘法交换律、结合律；A把32用了两遍，B、C则把应有的乘号写成加号。\par 两组各凑整后，\(1000\times100=100000\)\allowbreak ；凑整过程中只改变因数的分组，没有增加或减少因数。\end{multicols}\end{snippet}
% p094_s03
\begin{snippet}{505}{11.5}{14}\begin{multicols}{3}\raggedcolumns 连接后的70厘米只把重合段算一次，而两根各自长度相加会把这5厘米算两次。因此从70厘米中减去已知50厘米后，还须补回5厘米，才是另一根的完整长度；不能再减去重合段。\par \[\begin{aligned}&70-50+5\\={}&20+5\\={}&25\text{(厘米)}\end{aligned}\]\end{multicols}\end{snippet}
% p094_s04
\begin{snippet}{505}{11.5}{14}\begin{multicols}{3}\raggedcolumns 拼图要沿外边界逐段检查，只有方向相同、连成一直线的边才能合成一条边。两个完全相同的梯形沿腰可拼平行四边形，特殊的直角梯形可拼长方形，沿合适底边可出现六边形；保留下来的转折不能消去，因此不能拼成题列的三角形。\end{multicols}\end{snippet}
% p094_s05
\begin{snippet}{505}{11.5}{14}\begin{multicols}{3}\raggedcolumns 两组人数不同，不能把140和150直接平均。简便法以140厘米为基准：小刚组4人合计比4份140厘米多40厘米，把这40厘米平均分给全部10人，每人补4厘米。平均数作基准不表示每个人原来都恰好这么高。\par \textbf{方法一：移多补少，简便计算}\end{multicols}\[\begin{aligned}&140+(150-140)\times4\div(6+4)\\={}&140+10\times4\div10\\={}&140+40\div10\\={}&140+4\\={}&144\text{(厘米)}\end{aligned}\]\begin{multicols}{3}\raggedcolumns \textbf{方法二：总身高除以总人数}\end{multicols}\[\begin{aligned}&(140\times6+150\times4)\div(6+4)\\={}&(840+600)\div10\\={}&1440\div10\\={}&144\text{(厘米)}\end{aligned}\]\begin{multicols}{3}\raggedcolumns 选B。\end{multicols}\end{snippet}
% p094_s06
\begin{snippet}{505}{11.5}{14}\begin{multicols}{3}\raggedcolumns \textbf{二、7：}\ 先找一个弧形作为基本图形，比较它在其他位置的朝向，可以看到各图形都围绕共同中心转过相同角度。平移会保持朝向，不能让这些弧形形成绕中心的排列，因此应根据方向的变化判断运动方式。\end{multicols}\end{snippet}
% p094_s07
\begin{snippet}{505}{11.5}{14}\begin{multicols}{3}\raggedcolumns \textbf{二、8：}\ 36元先按每盒4元计算可付费买的盒数，再从这些付费盒中按4盒一组找可获赠的组数。9盒里只有2个完整的4盒组，剩下1盒不能再获赠，最后把买的盒数和赠的盒数相加；赠品不再计入下一组“买4盒”。\end{multicols}\end{snippet}
% p094_s08
\begin{snippet}{505}{11.5}{14}\begin{multicols}{3}\raggedcolumns （1）先算14乘6，再添乘整十数所对应的0；（2）（4）（6）小数加减按数位对齐，整数可补小数0再减。（3）同级乘除从左到右，或把25与除以25配对后仍保留两个乘4，不能把它们抵消。（5）除以1000向左移三位，位数不够要补0，不能写成0.7。\end{multicols}\end{snippet}
% p094_s09
\begin{snippet}{505}{11.5}{14}\begin{multicols}{3}\raggedcolumns 前两题都先对齐小数点并补0，\(8.050-4.357\)\allowbreak 遇到十分位0时，要从个位逐级退位，不能跳过某一数位。（3）304中的3在百位，28乘3所得的积要对着百位写，0占着十位。（4）6270除以30时，落下7后所得27不足30，十位必须商0，再落下0继续除；漏写这个0会把209变成29。\end{multicols}\end{snippet}
% p094_s09_local1
\begin{snippet}{34.25}{10}{14}\SourceSans\setlength{\mathindent}{2bp}\(15.16\)\end{snippet}
% p094_s09_local2
\begin{snippet}{34.25}{10}{14}\SourceSans\setlength{\mathindent}{2bp}\(3.693\)\end{snippet}
% p094_s09_local3
\begin{snippet}{34.25}{10}{14}\SourceSans\setlength{\mathindent}{2bp}\(8512\)\end{snippet}
% p094_s09_local4
\begin{snippet}{34.25}{10}{14}\SourceSans\setlength{\mathindent}{2bp}\(209\)\end{snippet}
% p095_s01
\begin{snippet}{505}{11.5}{14}\begin{multicols}{3}\raggedcolumns \textbf{三、3（1）：}\ 中括号内减去75与67的和，可以改成先减75再减67，使175与75先凑成整百。括号外的40仍乘整个差，先得到33再相乘；去括号时两个加数都变为减数，不能把67改成加数。\end{multicols}\end{snippet}
% p095_s01_local1
\begin{snippet}{172.567}{10.5}{12.8}\SourceSans\setlength{\mathindent}{2bp}\[\begin{aligned}={}&40\times[175-75-67]\\={}&40\times[100-67]\\={}&40\times33\\={}&1320\end{aligned}\]\end{snippet}
% p095_s02
\begin{snippet}{505}{11.5}{14}\begin{multicols}{3}\raggedcolumns \textbf{三、3（2）：}\ 末尾单独的99也是1份99，与前面的39份合成40份，先用乘法分配律合并，得到整十倍数。若从另一方向想，也可把原来的99拆成100少1，将多算的39从整百积中扣掉后再加99；两种思路都必须保留末尾这一份99。\end{multicols}\end{snippet}
% p095_s02_local1
\begin{snippet}{125.56700000000001}{10.5}{12.8}\SourceSans\setlength{\mathindent}{2bp}\[\begin{aligned}={}&(39+1)\times99\\={}&40\times99\\={}&3960\end{aligned}\]\end{snippet}
% p095_s03
\begin{snippet}{505}{11.5}{14}\begin{multicols}{3}\raggedcolumns \textbf{三、3（3）：}\ 两项都含有36，可以把它们看成245份36减145份36。先用乘法分配律合并份数，245减145凑成100，再求100份36；共同的36只保留一次，不能也相减。\end{multicols}\end{snippet}
% p095_s03_local1
\begin{snippet}{192.733}{10.5}{12.8}\SourceSans\setlength{\mathindent}{2bp}\[\begin{aligned}={}&(245-145)\times36\\={}&100\times36\\={}&3600\end{aligned}\]\end{snippet}
% p095_s04
\begin{snippet}{505}{11.5}{14}\begin{multicols}{3}\raggedcolumns \textbf{四、1：}\ 画钝角三角形先保证有一个角大于直角，画等腰直角三角形则用两条相等且互相垂直的格线作直角边。选定一条边为底后，要从与底相对的顶点向底所在直线作高，不能随便找一条竖线。原示例的钝角三角形向上方底边作高，等腰直角三角形从直角顶点向斜边作高，均画虚线并在垂足标直角。\end{multicols}\end{snippet}
% p095_s05
\begin{snippet}{505}{11.5}{14}\begin{multicols}{3}\raggedcolumns \textbf{四、2：}\ 先画长4格的底，再在距底垂直距离2格的平行线上取上底，使两组对边分别平行；斜边长不等于所说的高。以左下角为固定中心，把各顶点逆时针转直角后再连接，原来向右4格的底边会转为向上4格。旋转前后的底长、高和各边长度保持不变，高用虚线并标直角。\end{multicols}\end{snippet}
% p095_s06
\begin{snippet}{505}{11.5}{14}\begin{multicols}{3}\raggedcolumns 运走的量比剩下的多，先从82.6吨减去13.8吨，得到较少的剩余量，再与已运量合并。另法把剩余量先补到与已运量相等，合计两份82.6吨后，再扣掉多补的13.8吨；不能把“运走的多”理解成剩下的还要加差。\par \textbf{方法一：先求剩下质量}\[\begin{aligned} &82.6-13.8+82.6\\={}&68.8+82.6\\={}&151.4\text{(吨)}\end{aligned}\]\textbf{方法二：补成两份运走量}\[\begin{aligned} &82.6+82.6-13.8\\={}&165.2-13.8\\={}&151.4\text{(吨)}\end{aligned}\]答：这批蔬菜共151.4吨。\end{multicols}\end{snippet}
% p095_s07
\begin{snippet}{505}{11.5}{14}\begin{multicols}{3}\raggedcolumns 前3天给的是每天的平均页数，要用21乘3还原总页数，后2天27页已经是总数，只加一次，再除以5天。另法把每天都按21页看，后2天实际比这个标准少15页，把少读的页数平均到全部5天，所以应从21中扣去每天少的部分。\par \textbf{方法一：先求全书页数}\[\begin{aligned} &(21\times3+27)\div5\\={}&(63+27)\div5\\={}&90\div5\\={}&18\text{(页)}\end{aligned}\]\textbf{方法二：与每天21页比较}\[\begin{aligned} &21-(21\times2-27)\div5\\={}&21-(42-27)\div5\\={}&21-15\div5\\={}&21-3\\={}&18\text{(页)}\end{aligned}\]答：小林平均每天读18页。\end{multicols}\end{snippet}
% p096_s01
\begin{snippet}{505}{11.5}{14}\begin{multicols}{3}\raggedcolumns 先用付款减找回的钱求实际支出，再扣掉牛奶总价，余下的才是鸡蛋总价，最后除以鸡蛋每千克的价钱。另法先把84元全部按鸡蛋价格换算成7千克，再扣去牛奶费用相当于的3千克；这只是按同一单价换算费用，并非把牛奶实际称成鸡蛋。\par \textbf{方法一：先扣牛奶费用}\[\begin{aligned} &(100-16-36)\div12\\={}&(84-36)\div12\\={}&48\div12\\={}&4\text{(千克)}\end{aligned}\]\textbf{方法二：先按每千克鸡蛋的价格换算}\[\begin{aligned} &(100-16)\div12-36\div12\\={}&84\div12-3\\={}&7-3\\={}&4\text{(千克)}\end{aligned}\]答：这次买了4千克鸡蛋。\end{multicols}\end{snippet}
% p096_s02
\begin{snippet}{505}{11.5}{14}\begin{multicols}{3}\raggedcolumns 105千瓦时超过第一档上限，前100千瓦时仍按每千瓦时2元，只把超出的5千瓦时按3元计费。另法先全部按3元计，前100千瓦时每千瓦时多算1元，要扣回100份差价；不能把全部105千瓦时都按其中一档收费。\par \textbf{方法一：分两档计算}\[\begin{aligned} &2\times100+3\times(105-100)\\={}&200+3\times5\\={}&200+15\\={}&215\text{(元)}\end{aligned}\]\textbf{方法二：先按高价算再扣多算部分}\[\begin{aligned} &3\times105-(3-2)\times100\\={}&315-1\times100\\={}&315-100\\={}&215\text{(元)}\end{aligned}\]答：10月份应交电费215元。\end{multicols}\end{snippet}
% p096_s03
\begin{snippet}{505}{11.5}{14}\begin{multicols}{3}\raggedcolumns 前100千瓦时收费200元，260元已超过这个金额，才可确定用了第二档电量。方法一用超出200元的费用除以每千瓦时3元，求超出100千瓦时的部分后再加100；方法二把前100千瓦时少收的差价补回，使费用相当于全按3元收取，再求总电量。\par 先核实\(260>2\times100\)\allowbreak ，已用到第二档电量。\textbf{方法一：先扣第一档电费}\[\begin{aligned} &(260-2\times100)\div3+100\\={}&(260-200)\div3+100\\={}&60\div3+100\\={}&20+100\\={}&120\text{(千瓦时)}\end{aligned}\]\textbf{方法二：把费用补成全按高价计费}\[\begin{aligned} &[260+(3-2)\times100]\div3\\={}&[260+1\times100]\div3\\={}&[260+100]\div3\\={}&360\div3\\={}&120\text{(千瓦时)}\end{aligned}\]答：11月份用了120千瓦时电。\end{multicols}\end{snippet}
% p096_s04
\begin{snippet}{505}{11.5}{14}\begin{multicols}{3}\raggedcolumns 先分男生、女生逐人划记，认清四个分段的上下界，每个成绩只能进入一段。126厘米属于104至127厘米段，128和149属于下一段，150则进入最高段，不能把边界邻近数移错。最后分别检查两行人数都为24，再把男女各段合并核对共48人，可发现漏计或重复计数。\par 男生核对\(1+3+9+11=24\)\allowbreak (人)，女生核对\(2+7+11+4=24\)\allowbreak (人)；两组共\(24+24=48\)\allowbreak (人)。\end{multicols}\end{snippet}
% p096_s05
\begin{snippet}{505}{11.5}{14}\begin{multicols}{3}\raggedcolumns 先在男女各行中比较四段人数，男生最高段11人最多，女生128至149厘米段11人最多。两组都是24人，再比较较高分段：150厘米及以上男生11人、女生4人，128厘米及以上男生20人、女生15人，因此男生整体较好；不能只比较“人数最多段”这几个字。\end{multicols}\end{snippet}
% p097_s01
\begin{snippet}{505}{11.5}{14}\begin{multicols}{3}\raggedcolumns 先把原数除以一亿，用带“亿”的数表示：\par\noindent\(\displaystyle 126339000\text{亩}=1.26339\text{亿亩}\)\allowbreak \par 保留一位小数时再看百分位6，向十分位进1。第一空后只印“亩”，填写时必须带“亿”；第二空后已经印“亿亩”，只填近似数的数值，避免数量级丢失或重复。\end{multicols}\end{snippet}
% p097_s02
\begin{snippet}{505}{11.5}{14}\begin{multicols}{3}\raggedcolumns 按数位把3.925分别分成3个一、9个十分之一、2个百分之一和5个千分之一，不能把末位5看成百分之一。画数线标记前，先认清3.9到4.0被分成20小格，每格0.005；3.925比3.9多0.025，应从3.9向右数5小格，让三角形顶点对准刻度。\end{multicols}\end{snippet}
% p097_s03
\begin{snippet}{505}{11.5}{14}\begin{multicols}{3}\raggedcolumns 忽略3.5的小数点，就把这个加数看成35，先比较错误加数与正确加数相差多少。另一个加数保持不变，和增加的量就等于这一差，不能把“变成原来的10倍”误写成“多了10倍”。\end{multicols}\end{snippet}
% p097_s04
\begin{snippet}{505}{11.5}{14}\begin{multicols}{3}\raggedcolumns 左边32乘的是圆与三角形的和，展开时32必须分别与两个加数相乘，两个积之间仍用加号。右边两个空按原来的加数顺序填入相应图形，这一过程是把乘一个和分配到每个加数上，不是只交换因数位置。\end{multicols}\end{snippet}
% p097_s05
\begin{snippet}{505}{11.5}{14}\begin{multicols}{3}\raggedcolumns 直线可以无限延长，要先看延长后是否有交点，不能只看图中短线有没有碰上。①、③、⑤延长后相交，其中③、⑤的交角为直角，所以它们同时属于相交和垂直；②、④方向相同且不相交，再归为平行。\end{multicols}\end{snippet}
% p097_s06
\begin{snippet}{505}{11.5}{14}\begin{multicols}{3}\raggedcolumns \textbf{一、6：}\ 本题按题面“宽1厘米”，梯形高等于长方形宽，填1厘米。\par 重叠梯形的两底分别在长方形两条长边上，它们的垂直距离就是长方形的宽，须按题面1厘米判断，不采用书后4厘米。另看两个角，它们在同一条直线上拼成平角，角2应由180度减去角1的102度；图形斜着画不改变平角关系。\par \[\begin{aligned}&180^\circ-102^\circ\\={}&78^\circ\end{aligned}\]\end{multicols}\end{snippet}
% p097_s07
\begin{snippet}{505}{11.5}{14}\begin{multicols}{3}\raggedcolumns \textbf{一、7（分割图，问题续下一页）：}\ 由上底右端向下底画一条与左腰平行的分割线，配合原来平行的两底，就能保证左侧四边形两组对边分别平行。不能任意连接上下底，否则未必得到平行四边形；三角形护栏的完整求长过程在下一页条件齐全处说明。\end{multicols}\end{snippet}
% p098_s01
\begin{snippet}{505}{11.5}{14}\begin{multicols}{3}\raggedcolumns 沿用前页的分割，先从梯形周长扣除两底长度，再把余下长度平分给两条相等的腰。平行四边形对边相等，分割线与左腰等长，所占下底又与上底等长，因此可依次确定三角形的三条边；护栏围的是三角形四周，须把三边都计算在内。\par 先由等腰梯形周长求腰长：\[\begin{aligned}&(10-(2+4))\div2\\={}&(10-6)\div2\\={}&4\div2\\={}&2\text{(米)}\end{aligned}\]平行四边形在下底上所占线段与上底等长，为\(2\)\allowbreak 米；三角形的底是下底剩余部分，长\(4-2=2\)\allowbreak 米。分割线与左腰等长，三角形的另一条边就是右腰，因此这两条边也各长\(2\)\allowbreak 米。三角形周长：\[\begin{aligned}&2\times3\\={}&6\text{(米)}\end{aligned}\]\end{multicols}\end{snippet}
% p098_s02
\begin{snippet}{505}{11.5}{14}\begin{multicols}{3}\raggedcolumns 循环赛时，每组都要赛7场，按8组统计会把“甲对乙”与“乙对甲”同一场数两次，所以要除以2。列式后利用8能被2整除，可先算8除以2再乘7，减少中间数。淘汰赛每赛一场只淘汰1组，8组剩下1个冠军就要淘汰7组；这两个空采用的赛制不同，不能使用同一种计数方法。\par 循环赛：\[\begin{aligned}&7\times8\div2\\={}&7\times(8\div2)\\={}&7\times4\\={}&28\text{(场)}\end{aligned}\]淘汰赛：\(8-1=7\)\allowbreak (场)。\end{multicols}\end{snippet}
% p098_s03
\begin{snippet}{505}{11.5}{14}\begin{multicols}{3}\raggedcolumns 先从两侧找四舍五入的边界：1.445可以进到1.45，1.455却会进到1.46，不能取到后者。逐个按0.001递增枚举：\(1.445,\allowbreak 1.446,\allowbreak 1.447,\allowbreak 1.448,\allowbreak 1.449,\allowbreak \)\allowbreak \par \(1.450,\allowbreak 1.451,\allowbreak 1.452,\allowbreak 1.453,\allowbreak 1.454\)\allowbreak 。1.450按三位小数的写法应计入，不能因为末尾是0就少算这一个。\end{multicols}\end{snippet}
% p098_s04
\begin{snippet}{505}{11.5}{14}\begin{multicols}{3}\raggedcolumns 本题比较的是现在与原来的变化，只需把接上的4.15米与剪掉的3.8米比较。增加量比减少量多0.35米，所以是净增加；20米原长在前后比较中相同，不必再参与这一差额的计算。\par \[\begin{aligned}&4.15-3.80\\={}&0.35\text{(米)}\end{aligned}\]\end{multicols}\end{snippet}
% p098_s05
\begin{snippet}{505}{11.5}{14}\begin{multicols}{3}\raggedcolumns 三种分割都应先检查是否覆盖整个六边形，且各部分内部不重叠，再把各部分的内角和合并。甲分成4个三角形，乙分成2个四边形，丙分成2个三角形和1个四边形，都得到720度；分割线没有引入需另外扣除的内部周角。\end{multicols}\end{snippet}
% p098_s06
\begin{snippet}{505}{11.5}{14}\begin{multicols}{3}\raggedcolumns 可沿图①的边按顺序看各边长和缺角的位置，再与各选项转动后的排列核对。A、C可通过旋转对应，D保持原朝向可平移对应，B却把缺角翻到了镜像位置；平移、旋转只移动和转向，不会把图形翻面。\end{multicols}\end{snippet}
% p098_s07
\begin{snippet}{505}{11.5}{14}\begin{multicols}{3}\raggedcolumns 原选项没有正确答案。少记32分，补正后全组总分增加32，人数仍为8，平均分增加：\[\begin{aligned}&32\div8\\={}&4\text{(分)}\end{aligned}\]原D清楚写“增加1分”，不能暗改成“增加4分”。\end{multicols}\end{snippet}
% p099_s01
\begin{snippet}{505}{11.5}{14}\begin{multicols}{3}\raggedcolumns 先利用对折把要剪去的相同小角叠在一起，才能一次剪下多处重合部分。展开后要检查所得外轮廓与目标图一致且关于折痕对称；题目数的是剪的次数，对折本身不计作一次剪。\end{multicols}\end{snippet}
% p099_s02
\begin{snippet}{505}{11.5}{14}\begin{multicols}{3}\raggedcolumns （1）至（4）小数加减按相同数位计算，3.4可看成3.40，0.3可看成0.30，不能按末位对齐。（5）2.4是24个十分之一，5倍为120个十分之一；（6）1.25是125个百分之一，8倍为1000个百分之一，按计数单位理解即可。（7）先算除法再加45；（8）先算乘法再减30，不能一律从左到右。\end{multicols}\end{snippet}
% p099_s03
\begin{snippet}{505}{11.5}{14}\begin{multicols}{3}\raggedcolumns \textbf{三、2（1）：}\ 两个减数13.64与6.36的小数部分互补，先求和正好凑成20。利用连续减去两个数等于减去它们的和，再从45.2中一次扣除，既少计算也避免连续退位；括号里应保留加号。\end{multicols}\end{snippet}
% p099_s03_local1
\begin{snippet}{166.9}{10.5}{12.8}\SourceSans\setlength{\mathindent}{2bp}\[\begin{aligned}={}&45.2-(13.64+6.36)\\={}&45.2-20\\={}&25.2\end{aligned}\]\end{snippet}
% p099_s04
\begin{snippet}{505}{11.5}{14}\begin{multicols}{3}\raggedcolumns \textbf{三、2（2）：}\ 两项都是以27为每份数，原来99份减去49份，可先合并为50份。利用乘法分配律反向合并，份数差凑成整十，再与27相乘；共同的27不能再作减法。\end{multicols}\end{snippet}
% p099_s04_local1
\begin{snippet}{163.9}{10.5}{12.8}\SourceSans\setlength{\mathindent}{2bp}\[\begin{aligned}={}&27\times(99-49)\\={}&27\times50\\={}&1350\end{aligned}\]\end{snippet}
% p099_s05
\begin{snippet}{505}{11.5}{14}\begin{multicols}{3}\raggedcolumns \textbf{三、2（3）：}\ 先算最内层的和32加16，得到48，中括号内整体是48除以4的商。除以这个商可转为先除以48再乘4，这里的除以4要变为乘4，才能保持商不变；也可先把中括号算成12再除，检验结果相同。\end{multicols}\end{snippet}
% p099_s05_local1
\begin{snippet}{192.233}{10.5}{12.8}\SourceSans\setlength{\mathindent}{2bp}\[\begin{aligned}={}&960\div[48\div4]\\={}&960\div48\times4\\={}&20\times4\\={}&80\end{aligned}\]\end{snippet}
% p099_s06
\begin{snippet}{505}{11.5}{14}\begin{multicols}{3}\raggedcolumns 最短小路从图书馆所在点出发，向主干道作垂线段并在垂足处标直角，依据点到直线的垂线段最短。北箭头向上，东西方向就沿水平方向，再过A画这条水平直线；不能把A点附近任意一条斜线当作东西方向。\end{multicols}\end{snippet}
% p099_s07
\begin{snippet}{505}{11.5}{14}\begin{multicols}{3}\raggedcolumns \textbf{四、2（1）：}\ 题目未指定对称轴，原示例选右边端点所在的竖直格线为轴，用虚线表示。分别数每个顶点到轴的水平格数，在另一侧同一水平线上取等距点，按轮廓顺序连线；轴上端点保持不动，不能把左半图向右平移来代替对称。\end{multicols}\end{snippet}
% p099_s08
\begin{snippet}{505}{11.5}{14}\begin{multicols}{3}\raggedcolumns \textbf{四、2（2）：}\ 先绕O顺时针转直角，O不动，原来从O向右的3格边转为向下3格，另一顶点也随同一方向转动。画出旋转图后，再以它的三个顶点为起点各向右移动3格，另画平移图；题目要求两幅过程图，不能只留下最后的三角形。\end{multicols}\end{snippet}
% p100_s01
\begin{snippet}{505}{11.5}{14}\begin{multicols}{3}\raggedcolumns 苹果与梨的单价、箱数都不同，原法分别用每箱价钱乘对应箱数，得到两部分收入再合并。另法先把所有27箱按每箱24元计算，只有15箱苹果每箱少算8元，要补15份差价；不能用总箱数去补苹果差价。\par \textbf{方法一：分别求两种收入}\[\begin{aligned} &32\times15+24\times12\\={}&480+288\\={}&768\text{(元)}\end{aligned}\]\textbf{方法二：全部先按梨价再补苹果差价}\end{multicols}\[\begin{aligned} &24\times(15+12)+(32-24)\times15\\={}&24\times27+8\times15\\={}&648+120\\={}&768\text{(元)}\end{aligned}\]\begin{multicols}{3}\raggedcolumns 答：上午共收入768元。\end{multicols}\end{snippet}
% p100_s02
\begin{snippet}{505}{11.5}{14}\begin{multicols}{3}\raggedcolumns 平台服务费只从下午核桃收入中扣除，上午收入不重复扣费。可先求每袋实收，再乘80袋；也可先列下午销售额减总服务费，利用乘法分配律合并共同的80。\par 另一列式可这样化为原位主式：\[\begin{aligned}&768+45\times80-2\times80\\={}&768+(45-2)\times80\end{aligned}\]后续共用原位的完整简便脱式。\end{multicols}\end{snippet}
% p100_s03
\begin{snippet}{505}{11.5}{14}\begin{multicols}{3}\raggedcolumns 把笔杆与笔帽各自全长相加时，套在一起的重叠段被重复计算了一次，减去套好总长就能求出这一段。另法先用总长减笔杆长，求笔帽露出的部分，再从笔帽全长中扣除；最后把米换成厘米，不能把0.04米直接写成0.04厘米。\par \textbf{方法一：两部分全长之和减套好总长}\par\noindent\[\displaystyle \begin{aligned} &0.14+0.05-0.15\\={}&0.19-0.15\\={}&0.04\text{(米)}\end{aligned}\]\par\textbf{方法二：笔帽长减露出的长度}\par\noindent\[\displaystyle \begin{aligned} &0.05-(0.15-0.14)\\={}&0.05-0.01\\={}&0.04\text{(米)}\end{aligned}\]\par \noindent\(\displaystyle 0.04\text{米}=4\text{厘米}\)\allowbreak \par 答：重叠部分长4厘米。\end{multicols}\end{snippet}
% p100_s04
\begin{snippet}{505}{11.5}{14}\begin{multicols}{3}\raggedcolumns 线段图中，720千米由甲车已行、未行240千米和乙车已行三段组成，两车都行驶3小时。方法一先求两车合走的480千米，再除以3求速度和，减甲速；方法二先扣甲车的270千米和未行路程，得到乙车3小时的路程再除以3。\par \textbf{方法一：先求速度和}\[\begin{aligned} &(720-240)\div3-90\\={}&480\div3-90\\={}&160-90\\={}&70\text{(千米/时)}\end{aligned}\]\textbf{方法二：先求乙车已行路程}\[\begin{aligned} &(720-240-90\times3)\div3\\={}&(720-240-270)\div3\\={}&(480-270)\div3\\={}&210\div3\\={}&70\text{(千米/时)}\end{aligned}\]答：乙车每小时约行70千米。\end{multicols}\end{snippet}
% p100_s05
\begin{snippet}{505}{11.5}{14}\begin{multicols}{3}\raggedcolumns 整体较好：36人中有23人达到优秀，占全班一半以上；但仍有3人不合格，应予以帮助。\par （1）先按等级上下界逐个划记，85至100分、70至84分、60至69分和60分以下互不重叠，每人成绩只计一次，再核对总人数36。（2）评价要引用统计数据：优秀23人超过全班一半，优秀与良好合计30人，说明整体较好。同时还有3人不合格，应关注这部分同学，不能只写“很好”而不提供依据。\end{multicols}\end{snippet}
% p101_s01
\begin{snippet}{505}{11.5}{14}\begin{multicols}{3}\raggedcolumns 先把6个一、3个十分之一、5个百分之一放在各自数位，读小数时小数点后的3、5逐位读出。精确到十分位要看它后一位，即百分位的5，需要向十分位进1；不能把“看十分位”误作只照抄原十分位。\end{multicols}\end{snippet}
% p101_s02
\begin{snippet}{505}{11.5}{14}\begin{multicols}{3}\raggedcolumns 先把275替换为它的来源128与147的和，和必须用括号保护，才能先相加再除以25。再把所得的11作为最后乘240的一个因数，写综合算式时保留这一层先后关系；不能写成128加上147除以25。\end{multicols}\end{snippet}
% p101_s03
\begin{snippet}{505}{11.5}{14}\begin{multicols}{3}\raggedcolumns （1）小数末尾的0不改变大小；（2）（3）先统一为元、米再比较：\par\noindent\(\displaystyle 7\text{角}5\text{分}=0.75\text{元}\)\allowbreak \par \noindent\(\displaystyle 7\text{米}2\text{厘米}=7.02\text{米}\)\allowbreak \par（4）\(65\times99\)\allowbreak 是从\(65\times100\)\allowbreak 中减去65，另一边只减1，因此左边较大。（5）连续减去156、244等于减去它们的和，两式相等，不能把括号内写成减法。\end{multicols}\end{snippet}
% p101_s04
\begin{snippet}{505}{11.5}{14}\begin{multicols}{3}\raggedcolumns 已知的是一个底角，另一个底角也为45度，先把这两个角相加再从180度中扣除，求出顶角。顶角恰为90度，按角分类看这个最大角即可；“等腰”是按边分类，不代替本题所问的角的类别。\par \[\begin{aligned}&180^\circ-45^\circ\times2\\={}&180^\circ-90^\circ\\={}&90^\circ\end{aligned}\]\end{multicols}\end{snippet}
% p101_s05
\begin{snippet}{505}{11.5}{14}\begin{multicols}{3}\raggedcolumns “沿相同方向”限制了两次折叠方式，展开后的相应折痕也保持同一方向，彼此不相交。若改成横、竖方向各折一次才可出现垂直折痕，本题不能把另一种折法的结果一并填入。\end{multicols}\end{snippet}
% p101_s06
\begin{snippet}{505}{11.5}{14}\begin{multicols}{3}\raggedcolumns 把撕下的各角顶点重合，按顺序紧挨着摆放，不能重叠或留空，拼出的总角才表示内角和。三角形三个角共180度可拼平角，平行四边形四个角共360度可拼周角；不要只按纸片外轮廓猜角的名称。\end{multicols}\end{snippet}
% p101_s07
\begin{snippet}{505}{11.5}{14}\begin{multicols}{3}\raggedcolumns 可先分别从18人和16人中扣去两项都参加的5人，再把两类只参加者相加。若先加两个社团人数，5名重合者被计两次，而题目要求他们一次也不计，所以要减去两份5；只减一次求到的是至少参加一项的人数。\end{multicols}\end{snippet}
% p102_s01
\begin{snippet}{505}{11.5}{14}\begin{multicols}{3}\raggedcolumns 围成三角形须检查两短边之和大于最长边，B的\(2+3>3\)\allowbreak ，C的\(2.4+2.4>4.2\)\allowbreak ，两项都符合。题目没有给出可以据图量长的条件，不能仅凭衣架外观排除B；A的末长首位又不清，故应保留疑读，原单选至少B、C不唯一。\end{multicols}\end{snippet}
% p102_s02
\begin{snippet}{505}{11.5}{14}\begin{multicols}{3}\raggedcolumns ①千分位5引起连续进位，精确到百分位要保留7.00的两个0；②把37分成30与7，竖式两行积正体现分别相乘再求和的分配律。③即使0.3在两数之间，还有0.21、0.22等小数，继续细分还可找到更多，因此错误。④三角形的形状不易改变，支架采用它是利用稳定性，要把这条正确说法也计入。\end{multicols}\end{snippet}
% p102_s03
\begin{snippet}{505}{11.5}{14}\begin{multicols}{3}\raggedcolumns 125是前三人的平均数，乘3才能还原他们的总次数，再加乐乐145个，并把新总数平均分给4人。也可以125为基准，把乐乐多出的20个平均补给4人；原来三人不一定各跳125个，只是总数与三份125相等，不能直接平均125与145。\par 方法一，先求总次数：\[\begin{aligned}&(125\times3+145)\div4\\={}&(375+145)\div4\\={}&520\div4\\={}&130\text{(个)}\end{aligned}\]方法二，移多补少：\[\begin{aligned}&125+(145-125)\div4\\={}&125+20\div4\\={}&125+5\\={}&130\text{(个)}\end{aligned}\]\end{multicols}\end{snippet}
% p102_s04
\begin{snippet}{505}{11.5}{14}\begin{multicols}{3}\raggedcolumns 同一平面内两条不同直线先分为相交与平行两类，再在相交中找是否成直角。垂直必须有交点，因此应放在相交类内部，既不能放在平行类里面，也不能与相交并列成互不包含的三类。\end{multicols}\end{snippet}
% p102_s05
\begin{snippet}{505}{11.5}{14}\begin{multicols}{3}\raggedcolumns 把下面的线段看作已固定的一边，两端所画圆弧表示另外两边可能伸到的位置。圆弧不交、相切、相交对应能否找到共同的第三个顶点，研究的是三边能否围成三角形；判断重点不在图形有几个角或周长是多少。\end{multicols}\end{snippet}
% p103_s01
\begin{snippet}{505}{11.5}{14}\begin{multicols}{3}\raggedcolumns 短跑比较用时，秒数较少的一方获胜，因此让二组最慢的丙对一组最快的A，先接受这一场失利。再让乙21秒对C23秒、甲16秒对B20秒，二组就能赢这两场，满足三局两胜；不能把秒数较大误当成成绩更好。\end{multicols}\end{snippet}
% p103_s02
\begin{snippet}{505}{11.5}{14}\begin{multicols}{3}\raggedcolumns 连接内点E与四个顶点后，得到4个三角形，先把这4个三角形的内角和合起来。其中E周围还包含一个360度周角，它不属于原四边形的内角，必须扣去；图中编号的小角也没有包含各顶点的全部角，不能只加四个编号角。\par 四个三角形合计\(180^\circ\times4=720^\circ\)\allowbreak ，扣除内点的一周角，\(720^\circ-360^\circ=360^\circ\)\allowbreak 。\end{multicols}\end{snippet}
% p103_s03
\begin{snippet}{505}{11.5}{14}\begin{multicols}{3}\raggedcolumns （1）（2）（3）（5）（7）先对齐数位，0.65与0.35的小数部分可凑成1，8减2.68时8可看成8.00再退位。（4）先算38乘0得到0，再加62，不能把整个和都乘0。（6）除以100向左移两位；（8）乘1000向右移三位，位数不够须补0，分别检查0.006和4320的0所处位置。\end{multicols}\end{snippet}
% p103_s04
\begin{snippet}{505}{11.5}{14}\begin{multicols}{3}\raggedcolumns \textbf{三、2（1）：}\ 把15.9写成15.90，小数点对齐，才能使2个十分之一与9个十分之一相加。十分位满十进1，个位相加时要带上这一进位，不能只在十分位写1后就漏算进来的1。\end{multicols}\end{snippet}
% p103_s04_local1
\begin{snippet}{32.0}{10}{14}\SourceSans\setlength{\mathindent}{2bp}\(25.17\)\end{snippet}
% p103_s05
\begin{snippet}{505}{11.5}{14}\begin{multicols}{3}\raggedcolumns \textbf{三、2（2）：}\ 把两数的小数点对齐，十分位7加4得11，十分位写1并向个位进1。个位5加6还要加进位的1，满十再向十位进1；这道题有连续的两次进位，不能只处理十分位。\end{multicols}\end{snippet}
% p103_s05_local1
\begin{snippet}{32.75}{10}{14}\SourceSans\setlength{\mathindent}{2bp}\(42.1\)\end{snippet}
% p103_s06
\begin{snippet}{505}{11.5}{14}\begin{multicols}{3}\raggedcolumns \textbf{三、2（3）：}\ 10可写成10.00，百分位0不够减3时，要从十位开始逐级退位，不能越过中间的0。退位后十分位剩9，百分位为10，再逐位减去4.03；用所得差加4.03应恢复10，便于检查退位是否漏记。\end{multicols}\end{snippet}
% p103_s06_local1
\begin{snippet}{25.5}{10}{14}\SourceSans\setlength{\mathindent}{2bp}\(5.97\)\end{snippet}
% p103_s07
\begin{snippet}{505}{11.5}{14}\begin{multicols}{3}\raggedcolumns \textbf{三、2（4）：}\ 先将18.2补写为18.20，百分位与十分位都要根据前一步退位后剩下的数继续相减。带星题还须验算，用差11.42加减数6.78，应得到原被减数18.20；验算时要检查结果是否恢复原数，原图采用差加减数恢复被减数的方法。\end{multicols}\end{snippet}
% p103_s07_local1
\begin{snippet}{34.25}{10}{14}\SourceSans\setlength{\mathindent}{2bp}\(11.42\)\end{snippet}
% p103_s08
\begin{snippet}{505}{11.5}{14}\begin{multicols}{3}\raggedcolumns \textbf{三、3（1）：}\ 45.01与0.99的小数部分互补，先配成46；6.1与2.9也能配成整数9。用加法交换律、结合律把这两对分组，再把两个整数相加，比分别按原顺序计算更简便。\end{multicols}\end{snippet}
% p103_s08_local1
\begin{snippet}{208.9}{10.5}{12.8}\SourceSans\setlength{\mathindent}{2bp}\[\begin{aligned}={}&(45.01+0.99)+(6.1+2.9)\\={}&46+9\\={}&55\end{aligned}\]\end{snippet}
% p103_s09
\begin{snippet}{505}{11.5}{14}\begin{multicols}{3}\raggedcolumns \textbf{三、3（2）：}\ 98比100少2，先把45的100倍作为易算的整百积，再扣掉多算的2份45。也可把45拆成50少5再分配，但仍要让括号中的每一部分都与98相乘；本题用接近整百的98来拆更直接。\end{multicols}\end{snippet}
% p103_s09_local1
\begin{snippet}{110.233}{10.5}{12.8}\SourceSans\setlength{\mathindent}{2bp}\[\begin{aligned}={}&45\times(100-2)\\={}&4500-90\\={}&4410\end{aligned}\]\end{snippet}
% p103_s10
\begin{snippet}{505}{11.5}{14}\begin{multicols}{3}\raggedcolumns \textbf{三、3（3）：}\ 末尾减去的101就是一份101，前面的积可看作101份101。利用乘法分配律先把份数101减1凑成100，再乘101；不能把单独的一份误当成只减1。\end{multicols}\end{snippet}
% p103_s10_local1
\begin{snippet}{184.567}{10.5}{12.8}\SourceSans\setlength{\mathindent}{2bp}\[\begin{aligned}={}&101\times(101-1)\\={}&101\times100\\={}&10100\end{aligned}\]\end{snippet}
% p103_s11
\begin{snippet}{505}{11.5}{14}\begin{multicols}{3}\raggedcolumns \textbf{四、（1）：}\ AC是斜放的底，应从与它相对的顶点B出发作垂线，与AC相交于D。用三角尺一条直角边贴住AC、另一条经过B，画虚线线段BD，并在D处标直角；图上原有的AB或BC虽互相垂直，却都不是AC边上所问的高。\end{multicols}\end{snippet}
% p103_s12
\begin{snippet}{505}{11.5}{14}\begin{multicols}{3}\raggedcolumns \textbf{四、（2）：}\ A、B就在对称轴AB上，保持原位；C在轴左侧3格，对应点应在轴右侧同一水平线上3格处。找到这个点再分别与A、B相连，对称轴用虚线表示；对应点连线与轴垂直，不能把三角形整体向右移。\end{multicols}\end{snippet}
% p103_s13
\begin{snippet}{505}{11.5}{14}\begin{multicols}{3}\raggedcolumns \textbf{四、（3）：}\ 先固定中心A，把原来从A向下4格的AB逆时针转直角，变为向右4格，C也按同一方向转动。连接旋转后的三个点后，再各向右数4格画平移图；第二步包括把旋转中心对应点也一起移动，不能仍让它留在原A处。\end{multicols}\end{snippet}
% p104_s01
\begin{snippet}{505}{11.5}{14}\begin{multicols}{3}\raggedcolumns 答：现在有7.1元。\par 按事情顺序看，买笔使钱减少，妈妈给钱使钱增加，依次作减法、加法即可。另法先合并原有的钱和新收入，再扣一次买笔支出，得到的是同一余额；要把收入与支出的方向分清。\par 另法：先合计原有的钱与新收入，再扣支出。\[\begin{aligned} &12.3+2.3-7.5\\={}&14.6-7.5\\={}&7.1\text{(元)}\end{aligned}\]\end{multicols}\end{snippet}
% p104_s01_local1
\begin{snippet}{80}{9.5}{14}\SourceSans\setlength{\mathindent}{2bp}\(\text{(元)}\)\end{snippet}
% p104_s02
\begin{snippet}{505}{11.5}{14}\begin{multicols}{3}\raggedcolumns 买五送一是每6双中付5双的钱，12双正好两组。原法先求付费双数，再用4.68乘10移动小数点，已经简便，予以保留。另法把4.68元换成468分，从12双的总标价中扣2双赠品的标价；计算时用乘法分配律先合并为10份468分，不必先算两个大乘积。\par 另法：\par\noindent\(\displaystyle 4.68\text{元}=468\text{分}\)\allowbreak \par \noindent\end{multicols}\[\displaystyle \begin{aligned}&468\times12-468\times[12\div(5+1)]\\={}&468\times12-468\times[12\div6]\\={}&468\times12-468\times2\\={}&468\times(12-2)\\={}&468\times10\\={}&4680\text{(分)}\end{aligned}\]\begin{multicols}{3}\raggedcolumns \par \noindent\(\displaystyle 4680\text{分}=46.8\text{元}\)\allowbreak \par\end{multicols}\end{snippet}
% p104_s02_local1
\begin{snippet}{242}{10.5}{12.8}\SourceSans\setlength{\mathindent}{2bp}\[\begin{aligned}&4.68\times[5\times(12\div(5+1))]\\={}&4.68\times[5\times(12\div6)]\\={}&4.68\times[5\times2]\\={}&4.68\times10\\={}&46.8\text{(元)}\end{aligned}\]答：12双袜子共花46.8元。\end{snippet}
% p104_s03
\begin{snippet}{505}{11.5}{14}\begin{multicols}{3}\raggedcolumns 足球和篮球各15个，可把各一个配成一套，先将每套单价108与92凑200。分别列两类球总价时，利用乘法分配律合并共同的15，也化为同一简便式。\par 另一列式可这样化为原位主式：\[\begin{aligned}&108\times15+92\times15\\={}&(108+92)\times15\end{aligned}\]后续共用原位的完整简便脱式。\end{multicols}\end{snippet}
% p104_s03_local1
\begin{snippet}{80}{9.5}{14}\SourceSans\setlength{\mathindent}{2bp}\(\text{(元)}\)\end{snippet}
% p104_s04
\begin{snippet}{505}{11.5}{14}\begin{multicols}{3}\raggedcolumns 答：上山每分钟慢12米。\par 按原路返回说明上下山路程相同，先各用1080米除以对应分钟数，求出72和60米/分，再用下山快的速度减上山慢的速度。另法按上山速度计算15分钟走900米，与下山同样15分钟走1080米比较，路程相差180米再除以15，得到每分钟的差；不能直接把时间差当速度差。\par 另法：用原式已求的上山速度60米/分，比较同样15分钟所行路程。\[\begin{aligned} &(1080-60\times15)\div15\\={}&(1080-900)\div15\\={}&180\div15\\={}&12\text{(米/分)}\end{aligned}\]\end{multicols}\end{snippet}
% p104_s04_local1
\begin{snippet}{80}{9.5}{14}\SourceSans\setlength{\mathindent}{2bp}\(\text{(米/分)}\)\end{snippet}
% p104_s05
\begin{snippet}{505}{11.5}{14}\begin{multicols}{3}\raggedcolumns 答：相距1200千米。\par 第一次相距200千米表示尚未相遇，已行路程和要加尚余200千米。可先求两车速度和；分别列两车路程时，利用乘法分配律提取共同的2，将220与280凑500。\par 另一列式可这样化为原位主式：\[\begin{aligned}&220\times2+280\times2+200\\={}&(220+280)\times2+200\end{aligned}\]后续共用原位的完整简便脱式。\end{multicols}\end{snippet}
% p104_s05_local1
\begin{snippet}{80}{9.0}{14}\SourceSans\setlength{\mathindent}{2bp}\(\text{(千米)}\)\end{snippet}
% p105_s01
\begin{snippet}{505}{11.5}{14}\begin{multicols}{3}\raggedcolumns 9.08的小数部分逐位读，十分位0也要读出来；最低位是百分位，先确定一个计数单位是0.01。把9个一换成900个百分之一再加8个百分之一，才能求共有多少个单位。最小两位数10含1000个这样的单位，减去原有908个，就是还需添的个数。\end{multicols}\end{snippet}
% p105_s02
\begin{snippet}{505}{11.5}{14}\begin{multicols}{3}\raggedcolumns （1）千克换克乘1000，6千克先换成6000克，再加原有50克：\par\noindent\(\displaystyle 6\text{千克}50\text{克}=(6000+50)\text{克}=6050\text{克}\)\allowbreak \par（2）元与角的进率是10，小数部分先换成角：\par\noindent\(\displaystyle 0.6\text{元}=6\text{角}\)\allowbreak \par 整数部分仍是15元：\par\noindent\(\displaystyle 15.6\text{元}=15\text{元}6\text{角}\)\allowbreak \par（3）克换千克除以1000：\par\noindent\(\displaystyle 900\text{克}=0.9\text{千克}\)\allowbreak \par（4）千米换米乘1000：\par\noindent\(\displaystyle 0.07\text{千米}=70\text{米}\)\allowbreak \par 与原有整数部分合并：\par\noindent\(\displaystyle 3.07\text{千米}=3\text{千米}70\text{米}\)\allowbreak \par 不能把小数部分的07直接读成7米。\end{multicols}\end{snippet}
% p105_s03
\begin{snippet}{505}{11.5}{14}\begin{multicols}{3}\raggedcolumns 三位小数意味着小数点后恰有三位，四个数字各用一次，整数部分只剩一位。最大时先把最大数字7放在整数部分，余数按从大到小排；最小时整数位取0，余数从小到大排。最大数末位的0虽然不改变数值，也必须保留，才符合用完数字且为三位小数的要求。\end{multicols}\end{snippet}
% p105_s04
\begin{snippet}{505}{11.5}{14}\begin{multicols}{3}\raggedcolumns 先把10千米消耗的631千卡平均分给10份，求每千米消耗量，再用每千米消耗量乘100千米。也可把100千米看成10段10千米，核对总消耗是10份631；“两周”说明完成这段路程的时间，不改变按千米数计算的关系。\end{multicols}\end{snippet}
% p105_s05
\begin{snippet}{505}{11.5}{14}\begin{multicols}{3}\raggedcolumns 底角是42度，另一个底角也相等，先求两个底角共84度，再用180度减去它们求顶角。所得顶角96度超过直角，因此按角分类可确定类别；不能因为两个底角都是锐角，就把整个三角形判成锐角三角形。\par \[\begin{aligned}&180^\circ-42^\circ\times2\\={}&180^\circ-84^\circ\\={}&96^\circ\end{aligned}\]\end{multicols}\end{snippet}
% p105_s06
\begin{snippet}{505}{11.5}{14}\begin{multicols}{3}\raggedcolumns \textbf{一、6：}\ 从4根中取3根，可按小棒长度列出四种组合，逐一检查两短边之和是否大于最长边。\(2+5=7\)\allowbreak 等于最长边不能围成，另两组含2和10时两短边之和还不足10；只有\((5,\allowbreak 7,\allowbreak 10)\)\allowbreak 满足条件，最后把这三根相加求周长。\end{multicols}\end{snippet}
% p105_s07
\begin{snippet}{505}{11.5}{14}\begin{multicols}{3}\raggedcolumns \textbf{一、7：}\ 保留一位小数要看百分位，3.75可以向十分位进1成为3.8，3.74却会舍成3.7。另一端3.84仍可舍成3.8，3.85已进成3.9；由于原数限定两位小数，最大、最小就在这两个允许的端点处，不能再多取一百分之一。\end{multicols}\end{snippet}
% p105_s08
\begin{snippet}{505}{11.5}{14}\begin{multicols}{3}\raggedcolumns 先从全班42人中扣除两种棋都不会的10人，剩下32人是至少会一种棋的人。19人与16人相加为35人，比实际会棋人数多出的3人就是两种都会者，因为这部分被计了两次。求只会象棋，再从19人中扣去这3名重合者，不能从全班人数中直接扣。\par 至少会一种：\(42-10=32\)\allowbreak (人)；两种都会：\(19+16-32=3\)\allowbreak (人)；只会象棋：\(19-3=16\)\allowbreak (人)。\end{multicols}\end{snippet}
% p106_s01
\begin{snippet}{505}{11.5}{14}\begin{multicols}{3}\raggedcolumns 先看各围法有没有把容易变形的四边形用支撑杆分成三角形，B中的斜杆与横杆形成了这样的支撑。三角形三边长度固定后形状不易改变，能加固围栏；不能仅凭杆子多或外观看起来密就判断牢固。\end{multicols}\end{snippet}
% p106_s02
\begin{snippet}{505}{11.5}{14}\begin{multicols}{3}\raggedcolumns 13.8中的8在十分位，表示8个十分之一吨，即0.8吨。进率为：\par\noindent\(\displaystyle 1\text{吨}=1000\text{千克}\)\allowbreak \par 按进率换算：\par\noindent\(\displaystyle 0.8\text{吨}=(0.8\times1000)\text{千克}=800\text{千克}\)\allowbreak \par 不能只看到数字8就直接当成8千克。\end{multicols}\end{snippet}
% p106_s03
\begin{snippet}{505}{11.5}{14}\begin{multicols}{3}\raggedcolumns D把99拆成100少1，两部分分别与48相乘后相减，符合乘法分配律。A右边把加22变成了减22；B拆开44后没有把两部分都正确相乘；C擅自把从左到右的乘除运算改成除以一个积，改变了运算结果。\end{multicols}\end{snippet}
% p106_s04
\begin{snippet}{505}{11.5}{14}\begin{multicols}{3}\raggedcolumns 把整根铁丝按刻度看成12份，每个方案都要先按切点求出三段各有几份，再检查两短段之和是否大于最长段。②④把它分成4、3、5份，\(3+4>5\)\allowbreak 可围成；①②为1、3、8份，②③为4、2、6份，③④为6、1、5份，其余均不满足严格大于的条件。\end{multicols}\end{snippet}
% p106_s05
\begin{snippet}{505}{11.5}{14}\begin{multicols}{3}\raggedcolumns 先认清各图的单位“1”：A百分位4个珠表示4个百分之一，B的100小格中涂4格也表示0.04。D从0到0.1分5份，每份0.02，第二格对应0.04。C涂色是10层中的4整层，表示0.4，不能把立体厚层误作百分格或只数表面小格。\end{multicols}\end{snippet}
% p106_s06
\begin{snippet}{505}{11.5}{14}\begin{multicols}{3}\raggedcolumns 原图①先按每行3个、共6行，列\(3\times6\)\allowbreak ；图②按每列3个、共4列，列\(3\times4\)\allowbreak 。将图①旋转后，才成为每列3个、共6列，与图②并排拼成每列3个、共\(6+4\)\allowbreak 列的长方形。因此两部分的总数也可用\(3\times(6+4)\)\allowbreak 计算，体现乘法分配律；每次列式先写每份数，再写份数，不能直接把原①说成每列3个。\par 将分开计数与拼后计数联系起来：\end{multicols}\[3\times6+3\times4=3\times(6+4)=3\times10=30.\]\end{snippet}
% p107_s01
\begin{snippet}{505}{11.5}{14}\begin{multicols}{3}\raggedcolumns （1）（2）（7）小数加减按数位算，两个加法的小数部分都可凑成整一。（3）（5）（6）分别检查小数点右移两位、左移两位、右移三位，位数不够处要补0。（4）连续除以5、2可合为除以10；（8）先把25与4结合成100，再乘35，利用结构凑整。\end{multicols}\end{snippet}
% p107_s02
\begin{snippet}{505}{11.5}{14}\begin{multicols}{3}\raggedcolumns \textbf{三、2（1）：}\ 把16.8补写成16.80，小数点与2.27对齐，百分位先算0加7。十分位8加2满十，十分位写0，个位计算时还要加进位的1，结果中间的0不能省去。\end{multicols}\end{snippet}
% p107_s02_local1
\begin{snippet}{34.25}{10}{14}\SourceSans\setlength{\mathindent}{2bp}\(19.07\)\end{snippet}
% p107_s03
\begin{snippet}{505}{11.5}{14}\begin{multicols}{3}\raggedcolumns \textbf{三、2（2）：}\ 把12写成12.0，十分位0不够减5，从个位退1换成10个十分之一，再继续相减。验算用差加减数，7.5加4.5应恢复12.0，与原被减数相同；带星号的题不能只写差而略去验算。\end{multicols}\end{snippet}
% p107_s03_local1
\begin{snippet}{34.25}{10}{14}\SourceSans\setlength{\mathindent}{2bp}\(7.5\)\end{snippet}
% p107_s04
\begin{snippet}{505}{11.5}{14}\begin{multicols}{3}\raggedcolumns \textbf{三、2（3）：}\ 先把1.9补成1.90，小数点对齐，10.02的十分位是0，不能把它改看成10.20。十分位不够减9，个位又是0，须先从十位退1到个位，再从个位退1到十分位，逐级记录剩余数字；百分位只算2减0，不跟着随意改变。\end{multicols}\end{snippet}
% p107_s04_local1
\begin{snippet}{34.25}{10}{14}\SourceSans\setlength{\mathindent}{2bp}\(8.12\)\end{snippet}
% p107_s05
\begin{snippet}{505}{11.5}{14}\begin{multicols}{3}\raggedcolumns \textbf{三、3（1）：}\ 先算小括号106加74，正好得到180，于是能发现两个加数都含有180。把单独的180看成1份，与19份合成20份，再乘180；这一简便方法的关键是算出括号后找共同因数，不能先随意分配19。\end{multicols}\end{snippet}
% p107_s05_local1
\begin{snippet}{199.567}{10.5}{12.8}\SourceSans\setlength{\mathindent}{2bp}\[\begin{aligned}={}&180+19\times180\\={}&(1+19)\times180\\={}&20\times180\\={}&3600\end{aligned}\]\end{snippet}
% p107_s06
\begin{snippet}{505}{11.5}{14}\begin{multicols}{3}\raggedcolumns \textbf{三、3（2）：}\ 减去7.38与10.65的和，可以改为连续减去这两个数，并把容易凑整的10.65先减。40.65与10.65的小数部分相同，先减得到30，再扣7.38更方便；去括号后两项都应作减法。\end{multicols}\end{snippet}
% p107_s06_local1
\begin{snippet}{191.233}{10.5}{12.8}\SourceSans\setlength{\mathindent}{2bp}\[\begin{aligned}={}&40.65-10.65-7.38\\={}&30-7.38\\={}&22.62\end{aligned}\]\end{snippet}
% p107_s07
\begin{snippet}{505}{11.5}{14}\begin{multicols}{3}\raggedcolumns \textbf{三、3（3）：}\ 25乘201表示201份25，末尾单独减25就是减去其中1份。先合并为200份25，利用乘法分配律凑成整百份数再计算；不能把末尾的25误看成只把积减1。\end{multicols}\end{snippet}
% p107_s07_local1
\begin{snippet}{200.067}{10.5}{12.8}\SourceSans\setlength{\mathindent}{2bp}\[\begin{aligned}={}&25\times(201-1)\\={}&25\times200\\={}&5000\end{aligned}\]\end{snippet}
% p107_s08
\begin{snippet}{505}{11.5}{14}\begin{multicols}{3}\raggedcolumns \textbf{三、3（4）：}\ 先算最内层81减77得4，再把15与4相乘凑成整十数60，直接用840除以60。这里除以的是积，算出这个易算的积比拆成先除以15再除以4更简便；除数是60，不能误当成15与4的和。\end{multicols}\end{snippet}
% p107_s08_local1
\begin{snippet}{200.567}{10.5}{12.8}\SourceSans\setlength{\mathindent}{2bp}\[\begin{aligned}={}&840\div[15\times4]\\={}&840\div60\\={}&14\end{aligned}\]\end{snippet}
% p107_s09
\begin{snippet}{505}{11.5}{14}\begin{multicols}{3}\raggedcolumns \textbf{四、（1）：}\ 先以小路的一条直边为已知直线，用直尺、三角尺保持方向不变，移动到另一边恰好经过淘气所在点后画路线。小路本身斜放，所画平行线也应同样倾斜，不能按纸面水平方向画；最后检查指定点确实在线上。\end{multicols}\end{snippet}
% p107_s10
\begin{snippet}{505}{11.5}{14}\begin{multicols}{3}\raggedcolumns \textbf{四、（2）：}\ 从笑笑所在的右下角点，向靠近她的小路边界作垂线段，终点落在这条边界上。点到直线的垂线段最短，垂足处直角标记的两边须分别沿小路边界和所画路线；所求是到小路，不要求走到奇思所在的点。\end{multicols}\end{snippet}
% p108_s01
\begin{snippet}{505}{11.5}{14}\begin{multicols}{3}\raggedcolumns 第一问比较两人所得金额，用小芳较多的金额减爸爸的金额，小数点对齐后求差。第二问以妈妈的50元为总量，可以依次扣两人所得，也可以先合计两人所得再一次扣除；两种算法都须扣全，按题面数据核算，不能沿用书后18.65元。\par \textbf{多抢：}\[\begin{aligned} &19.7-15.65\\={}&19.70-15.65\\={}&4.05\text{(元)}\end{aligned}\]答：小芳比爸爸多抢4.05元。\par \textbf{剩余方法一：依次扣除}\[\begin{aligned} &50-15.65-19.7\\={}&34.35-19.7\\={}&14.65\text{(元)}\end{aligned}\]\textbf{剩余方法二：先合计再扣除}\[\begin{aligned} &50-(15.65+19.7)\\={}&50-35.35\\={}&14.65\text{(元)}\end{aligned}\]答：红包还剩14.65元。\end{multicols}\end{snippet}
% p108_s02
\begin{snippet}{505}{11.5}{14}\begin{multicols}{3}\raggedcolumns 先把总页数\((40\times11)\)\allowbreak 括成整体。提前3天，实际8天；按乘积除以一个数的性质，可先算\(40\div8\)\allowbreak ，再乘11。\textbf{方法一：先除再乘}\[\begin{aligned}&(40\times11)\div(11-3)\\={}&(40\times11)\div8\\={}&(40\div8)\times11\\={}&5\times11\\={}&55\text{(页)}\end{aligned}\]\textbf{方法二：把少3天的任务平均补给8天}\[\begin{aligned}&40+(40\times3)\div(11-3)\\={}&40+120\div8\\={}&40+15\\={}&55\text{(页)}\end{aligned}\]答：平均每天要看55页。\end{multicols}\end{snippet}
% p108_s03
\begin{snippet}{505}{11.5}{14}\begin{multicols}{3}\raggedcolumns 每名学生各买一件上衣和一条裤子，可先把每套单价62与38凑100，再乘45套。分别列上衣、裤子总价时，用乘法分配律合并共同的45；人数不能因此变成90。\par \[\begin{aligned}&62\times45+38\times45\\={}&(62+38)\times45\\={}&100\times45\\={}&4500\text{(元)}\end{aligned}\]答：买这批合唱服一共需要4500元。\end{multicols}\end{snippet}
% p108_s04
\begin{snippet}{505}{11.5}{14}\begin{multicols}{3}\raggedcolumns 前4天2500米已是总量，后3天每天平均800米要先乘3求总量，再把7天总路程平均到7天。另法都以每天800米为基准，前4天比这个标准少700米，把少的部分平均扣给全周7天；不能把2500和800这两个不同含义的数直接平均。\par \textbf{方法一：总路程除以总天数}\[\begin{aligned} &(2500+800\times3)\div(4+3)\\={}&(2500+2400)\div7\\={}&4900\div7\\={}&700\text{(米)}\end{aligned}\]\textbf{方法二：以每天800米为基准}\[\begin{aligned} &800-(800\times4-2500)\div7\\={}&800-(3200-2500)\div7\\={}&800-700\div7\\={}&800-100\\={}&700\text{(米)}\end{aligned}\]答：这一周平均每天修路700米。\end{multicols}\end{snippet}
% p108_s05
\begin{snippet}{505}{11.5}{14}\begin{multicols}{3}\raggedcolumns 先按20至29、30至39、40至49、50至59岁逐人划记，每个年龄只进入一个分段，28与40等边界附近数要看清。用\(9+8+6+5=28\)\allowbreak 核对四段人数之和，与原记录人数一致后，再比较人数找最多、最少；相同年龄出现多次代表不同教师，不能去重后再数。\end{multicols}\end{snippet}\end{document}